\documentclass[UTF8,11pt,twoside,openany]{ctexbook}
\usepackage[paperwidth=176mm,paperheight=250mm,inner=21mm,outer=18mm,top=20mm,bottom=20mm,headheight=16pt,headsep=6mm,footskip=10mm]{geometry}
\usepackage{amsmath,amssymb,amsthm,bm,mathtools,booktabs,array,enumitem,float}
\usepackage{tikz,fancyhdr,hyperref}
\usetikzlibrary{arrows.meta,calc,patterns,angles,quotes,decorations.pathmorphing}
\hypersetup{hidelinks,pdftitle={近代物理进阶：物理竞赛二轮讲义},pdfauthor={}}
\setlength{\parindent}{2em}\setlength{\parskip}{3pt}\setlength{\emergencystretch}{2em}
\linespread{1.18}\raggedbottom\let\cleardoublepage\clearpage\setcounter{tocdepth}{1}
\setlength{\intextsep}{8pt plus 1pt minus 1pt}
\ctexset{chapter={name={第,讲},number=\chinese{chapter},format=\LARGE\bfseries\centering,beforeskip=16pt,afterskip=22pt},section={format=\Large\bfseries},subsection={format=\large\bfseries}}
\pagestyle{fancy}\fancyhf{}\fancyhead[LE,RO]{\small\thepage}\fancyhead[LO]{\small 近代物理进阶}\fancyhead[RE]{\small\leftmark}
\renewcommand{\headrulewidth}{0pt}\fancypagestyle{plain}{\fancyhf{}\fancyfoot[C]{\small\thepage}\renewcommand{\headrulewidth}{0pt}}
\theoremstyle{definition}\newtheorem{example}{例}[chapter]\newtheorem{exercise}{习题}[chapter]
\newenvironment{solution}{\par\medskip\noindent\textbf{解}\quad\ignorespaces}{\par\medskip}
\numberwithin{equation}{chapter}
\newcommand{\dd}{\mathop{}\!\mathrm d}\newcommand{\ee}{\mathrm e}\newcommand{\kb}{k_{\mathrm B}}\newcommand{\vct}[1]{\bm{#1}}
\tikzset{axis/.style={->,>=Stealth,line width=.45pt},curve/.style={line width=.7pt},guide/.style={dashed,line width=.35pt},every node/.style={font=\small},particle/.style={circle,fill,inner sep=1.3pt}}
\setlist[enumerate]{itemsep=3pt,topsep=4pt,leftmargin=2em}
\let\LectureContents\tableofcontents\renewcommand{\tableofcontents}{\begingroup\small\linespread{1}\selectfont\LectureContents\endgroup}
\newcommand{\Diagram}[3]{\begin{figure}[H]\centering\begin{tikzpicture}#3\end{tikzpicture}\caption{#2}\label{#1}\end{figure}}
\newcommand{\STAxes}{\draw[axis](-.4,0)--(5,0)node[right]{$x$};\draw[axis](0,-.2)--(0,4.5)node[above]{$ct$};}
\providecommand{\ExpansionDiagram}[3]{\begin{figure}[H]\centering\begin{tikzpicture}#3\end{tikzpicture}\caption{#2}\label{#1}\end{figure}}


% 分册与十编合订本共用的定位页码；局部章、图、题号保持原有对应。
\newcommand{\CoursePartNumber}{9}
\geometry{inner=19.5mm,outer=19.5mm,headheight=16pt}
\AtBeginDocument{%
  \fancyhead{}\fancyhead[C]{\small 近代物理进阶}%
  \fancyfoot[C]{\small\CoursePartNumber-\thepage}%
  \fancypagestyle{plain}{\fancyhf{}\fancyfoot[C]{\small\CoursePartNumber-\thepage}\renewcommand{\headrulewidth}{0pt}}%
}

\begin{document}
\frontmatter
\begin{titlepage}\centering\vspace*{28mm}
{\Large 物理竞赛二轮讲义\quad 第九编\par}\vspace{8mm}
{\Huge\bfseries 近代物理进阶\par}\vspace{12mm}{\Large 物理竞赛二轮讲义\par}\vfill\end{titlepage}
\chapter*{编写说明}
相对论计算首先要选对事件和参考系。同一对端点在什么时候测量，光在什么时候发出，功率按哪一只钟计时，都会改变列式。本编由这些测量问题出发，建立四维量和守恒关系，讨论火箭、延迟成像及电磁场的变换。第四讲转入光量子、热辐射和束缚态，着重说明能量计数与边界条件怎样决定结果。

各讲保留一般参数推导，并用数值情形说明符号和极限。图中的事件、路径及方向标记直接用于解答；例题和习题均附计算过程。涉及半经典、远场或微小过程近似时，适用条件随公式一并给出。
\tableofcontents
\mainmatter
\chapter{事件、四维量与运动光源}
\section{从测量程序确定事件}\label{th:m1events}
惯性系中，一个事件的位置由静止于该系的尺确定，时间由事件所在地已经同步的钟给出。爱因斯坦同步的约定是：从甲钟发出的光在乙处立即反射，若甲钟发射、接收时刻为$t_1,t_2$，就把乙钟在反射事件的读数校准为$(t_1+t_2)/2$。这一定义用到了同一惯性系内的双程光速$c$。

取$S'$相对$S$沿$x$正向以速度$V$运动，两系原点在事件$O$重合并归零。记
\[
 \beta=V/c,\qquad \gamma=(1-\beta^2)^{-1/2}.
\]
从线性变换$x'=a(x-Vt)$出发，要求$x=\pm ct$在另一系仍有$x'=\pm ct'$，可得$t'=a(t-Vx/c^2)$。逆变换与正变换的形式相同，只把$V$改成$-V$；把两式相乘要求恢复原坐标，得到$a^2(1-\beta^2)=1$。取保持时间方向的正根，得到
\begin{equation}\label{eq:m1lorentz}
 ct'=\gamma(ct-\beta x),\qquad x'=\gamma(x-\beta ct),\qquad y'=y,\quad z'=z.
\end{equation}
逆变换则为$ct=\gamma(ct'+\beta x')$、$x=\gamma(x'+\beta ct')$。这些式子同时适用于任意两事件的坐标差，即使两系坐标原点没有这样选取。

\Diagram{fig:m1coordinates}{同一事件在两惯性系中的坐标；虚线为过$P$的两组等时线}{
\STAxes
\draw[axis](0,0)--(2.5,4.1)node[right]{$ct'$};\draw[axis](0,0)--(4.5,2.7)node[right]{$x'$};
\fill(2.2,3.2)circle(1.3pt)node[above left]{$P$};\draw[guide](0,3.2)--(3.8,3.2);\draw[guide](.3,2.06)--(3.6,4.04);\node[below left]at(0,0){$O$};}
图\ref{fig:m1coordinates}中的$x'$轴是$t'=0$，斜率为$ct/x=\beta$；$ct'$轴是$x'=0$，斜率为$ct/x=1/\beta$。图上两轴的欧氏长度不能直接作为另一系的坐标刻度。若某杆端点的世界线为$x_A(t),x_B(t)$，$S$测得的杆长是$x_B(t)-x_A(t)$，两端事件必须位于同一条水平线上。照片则来自若干不同的发光时刻。

\subsection{同时性、钟慢与尺缩的条件}
在$S$中同时发生的两事件满足$\Delta t=0$，但
\[
 \Delta t'=-\gamma V\Delta x/c^2.
\]
只有$\Delta x=0$或$V=0$时仍然同时。对于固定在$S'$的同一只钟，$\Delta x'=0$，逆变换给$\Delta t=\gamma\Delta t'$，这里$\Delta t'$是这只钟经历的固有时。若测$S'$静止杆的长度$l_0=\Delta x'$，必须另取$\Delta t=0$的一对端点事件；正变换给$l_0=\gamma\Delta x$，所以$l=l_0/\gamma$。钟慢沿同一只钟取事件，尺缩沿同一条等时线取端点；使用公式时应保留各自的事件约束。

\begin{example}\label{ex:m1simultaneous}
在$S$中，左右两个信号灯相距$l$，在$t=0$同时闪光。$S'$沿正向以$0.6c$运动，原点在左灯闪光事件重合。求两闪光事件在$S'$中的时间差、位置差，并说明哪一个先发生。
\Diagram{fig:m1lampsq}{同一水平线上的两个闪光事件}{\STAxes\fill(0,0)circle(1.5pt)node[below]{$A$};\fill(3,0)circle(1.5pt)node[below]{$B$};\draw[guide](0,2)--(3,2);\draw[curve](3,0)--(3,3.6);\node[right]at(3,3.4){右灯};\node[below]at(1.5,0){$l$};}
\end{example}
\begin{solution}
图\ref{fig:m1lampsq}给出$A:(0,0)$、$B:(0,l)$，坐标次序取$(ct,x)$。代入式\eqref{eq:m1lorentz}，$A'=(0,0)$，而$B'$为
\[
 ct'_B=-\gamma\beta l=-\frac34l,\qquad x'_B=\gamma l=\frac54l.
\]
故$\Delta t'=t'_B-t'_A=-3l/(4c)$，右灯先闪。这里$5l/4$是两次闪光事件的空间坐标差，不是$S'$测得的两灯间距。
\Diagram{fig:m1lampsa}{在$S'$中取同时事件才能测量灯距，坐标轴按相同比例作图}{\draw[axis](-.3,0)--(4.3,0)node[right]{$x'$};\draw[axis](0,-2.8)--(0,2.2)node[above]{$ct'$};\draw[curve](1.5,-2.5)--(-.9,1.5);\draw[curve](3.9,-2.5)--(1.5,1.5);\fill(0,0)circle(1.4pt)node[left]{$A$};\fill(3.75,-2.25)circle(1.4pt)node[right]{$B$};\draw[guide](0,0)--(2.4,0);\fill(2.4,0)circle(1.4pt)node[above right]{$C$};}
图\ref{fig:m1lampsa}中的$C$取在右灯世界线上，并满足$t'_C=0$。右灯在$S$中$x=l$，由$t'=0$得$t=Vl/c^2$，再代入$x'=\gamma(l-Vt)$，得$x'_C=l/\gamma=4l/5$。这才是$S'$的灯距。
\end{solution}

\begin{example}\label{ex:m1rods}
杆$AB$静止于$S$，杆$A'B'$静止于$S'$，两杆固有长均为$l=ct_0$，$S'$沿正向以$V$运动。$A$与$A'$相遇于$M$，这两只端点钟归零；$B$与$B'$相遇于$N$，那两只钟归零。所有钟此后正常走时。$M$处从$A'$向$B'$发出光，光在$B'$立即反射，向$A'$返回。求光出射时经过$B$、在$B'$反射、返回时经过$B$、到达$A'$及到达$A$这五个事件中，各相应端点钟的读数，并讨论返回光线遇到各端点的次序。
\Diagram{fig:m1rodsq}{两杆及局部归零事件，图示为$S$中的一个同时截面}{\draw[axis](-.4,0)--(6,0)node[right]{$x$};\draw[curve](0,.7)--(3,.7);\draw[curve](.8,1.7)--(3.2,1.7);\foreach\x/\y/\t in{0/.7/A,3/.7/B,.8/1.7/A',3.2/1.7/B'}{\fill(\x,\y)circle(1.3pt)node[above]{$\t$};}\draw[axis](3.7,1.7)--(5,1.7)node[right]{$V$};\draw[<->](0,.3)--(3,.3)node[midway,below]{$l$};\draw[<->](.8,2.2)--(3.2,2.2)node[midway,above]{$l/\gamma$};}
\end{example}
\begin{solution}
以$M$为两坐标系原点。图\ref{fig:m1rodsq}中$B$固定在$x=l$，$B'$的世界线为$x=Vt+l/\gamma$。相遇条件给
\[
 t_N=\frac{l}{V}(1-\gamma^{-1})=\frac{t_0}{\beta}(1-\gamma^{-1}).
\]
因此$B$钟读$t-t_N$，$A$钟读$t$。在$S'$中，$A'$钟读$t'$，$B'$钟读$t'-t'_N$，其中
\begin{align*}
 t'_N&=\gamma(t_N-\beta t_0)
 =\frac{\gamma t_0}{\beta}(1-\gamma^{-1}-\beta^2)\\
 &=\frac{t_0}{\beta}(\gamma^{-1}-1)=-t_N.
\end{align*}
最后一步用$1-\beta^2=\gamma^{-2}$。两个不同的局部归零操作没有把各杆上的两只钟同步。
\Diagram{fig:m1rodsa}{光的全部交点；$P,Q$在$B$世界线上，$R$为反射事件}{
\draw[axis](-.2,0)--(4.9,0)node[right]{$x/l$};\draw[axis](0,-.15)--(0,6.4)node[above]{$ct/l$};
\draw[curve](1,0)--(1,6);\node[above]at(1,6){$B$};\draw[curve](0,0)--(3.6,6);\node[right]at(3.6,6){$A'$};\draw[curve](.8,0)--(4.4,6);\node[right]at(4.4,6){$B'$};
\draw[very thick,->](0,0)--(2,2);\draw[very thick,->](2,2)--(0,4);
\foreach\x/\y/\t in{0/0/M,1/.333/N,1/1/P,2/2/R,1/3/Q,1.5/2.5/T,0/4/U}{\fill(\x,\y)circle(1.4pt)node[above left]{$\t$};}}
图\ref{fig:m1rodsa}按$\beta=0.6$画出比例，计算保留一般$\beta$。出射光$x=ct$，故$P$为$t_P=t_0$。光在$S'$中走过固定长度$l$后到$R$，所以$(ct'_R,x'_R)=(l,l)$，逆变换给
\[
 t_R=\gamma(1+\beta)t_0,\quad x_R=ct_R.
\]
反射光在$S$中的方程是$x=2ct_R-ct$。与$x=l$相交得$t_Q=2t_R-t_0$；与$x=0$相交得$t_U=2t_R$。在$S'$中，光往返的总路程为$2l$，故回到$A'$的事件$T$满足$t'_T=2t_0$。于是所求各钟读数为
\begin{align*}
 B(P):&\quad t_0-t_N,\\
 B'(R):&\quad t_0-t'_N=t_0+t_N,\\
 B(Q):&\quad 2\gamma(1+\beta)t_0-t_0-t_N,\\
 A'(T):&\quad 2t_0,\\
 A(U):&\quad 2\gamma(1+\beta)t_0.
\end{align*}
返回途中，$t_T=2\gamma t_0$，故$t_Q-t_T=(2\gamma\beta-1)t_0$。令此差为零，得$4\beta^2=1-\beta^2$，即$\beta=1/\sqrt5$。因此$0<\beta<1/\sqrt5$时依次遇到$B,A',A$；$\beta>1/\sqrt5$时依次遇到$A',B,A$；临界时$B$与$A'$恰在光路上相遇。若取$\beta=0.6$，则$\gamma=1.25$、$t_N=t_0/3$，五个钟面读数依次为$2t_0/3,4t_0/3,8t_0/3,2t_0,4t_0$，而$T$发生在$Q$之前。事件排序使用同一参考系的坐标时间，不能直接比较未同步的端点钟。
\end{solution}

\section{四维矢量与固有量}\label{th:m1four}
以$x^\mu=(ct,x,y,z)$表示四维位置，取度规$g_{\mu\nu}=\operatorname{diag}(1,-1,-1,-1)$。降指标得到$x_\mu=(ct,-x,-y,-z)$，重复上下指标表示求和。两个四维矢量的内积为
\[
 A\cdot B=A^\mu B_\mu=A^0B^0-A^1B^1-A^2B^2-A^3B^3.
\]
将式\eqref{eq:m1lorentz}代入可逐项验证$\Delta s^2=c^2\Delta t^2-\Delta\vct r^2$不变。矩阵语言是$\Lambda^{\mathsf T}g\Lambda=g$。四维矢量按同一个$\Lambda$变换，其内积因此不变；仅发现某个量的平方不变，还不足以证明任意构造的四个数具有正确变换律。

\subsection{类时、类空与类光间隔}
若$\Delta s^2>0$且$\Delta t>0$，两事件可由亚光速质点依次到达，并存在使两事件发生在同一位置的惯性系。这不要求参与两个事件的不同物体都在该系静止。若$\Delta s^2=0$且事件间隔非零，两事件可由真空光线连接。若$\Delta s^2<0$，存在使两事件同时的惯性系，部分惯性系会交换其时间顺序；它们不能由不超过光速的信号建立因果关系。

对有质量质点的世界线，微小间隔满足
\begin{equation}\label{eq:m1proper}
 c^2\dd\tau^2=c^2\dd t^2-\dd\vct r^2,
 \qquad \dd\tau=\dd t\sqrt{1-v^2/c^2}.
\end{equation}
弯曲世界线的固有时要沿轨迹积分，不能把端点间隔的平方根当作整段固有时。两点间匀速直线世界线的固有时最大：在两端点共同的静止系中，$\int\sqrt{1-v^2/c^2}\dd t\leq\int\dd t$，等号只在全程$v=0$时成立。

\begin{example}\label{ex:m1interval}
两事件在$S$中的坐标差为$c\Delta t=3L$、$\Delta x=5L$。求使它们同时的惯性系及该系中的距离，说明是否有使它们同地的惯性系。
\Diagram{fig:m1intervalq}{类空分离的两个事件位于光锥外}{\STAxes\draw[guide](0,0)--(4.2,4.2);\fill(4,2.4)circle(1.4pt)node[right]{$B$};\node[below left]at(0,0){$A$};\draw[curve](0,0)--(4,2.4);}
\end{example}
\begin{solution}
图\ref{fig:m1intervalq}中两事件的间隔为
\[
 \Delta s^2=(3L)^2-(5L)^2=-16L^2<0.
\]
这是类空间隔，允许改变事件的时间先后。沿$x$作速度$V$的变换，令两事件同时，条件为
\[
 0=c\Delta t'=\gamma(3L-5\beta L),\qquad \beta=V/c.
\]
因此$\beta=3/5$、$\gamma=1/\sqrt{1-9/25}=5/4$。在该系中空间距离为
\[
 \Delta x'=\frac54\left(5-\frac35\times3\right)L=4L.
\]
\Diagram{fig:m1intervala}{同时系中的两事件，间隔绝对值对应空间距离}{\draw[axis](-.3,0)--(5,0)node[right]{$x'$};\draw[axis](0,-.2)--(0,2.4)node[above]{$ct'$};\fill(0,0)circle(1.4pt)node[below]{$A$};\fill(4,0)circle(1.4pt)node[below]{$B$};\draw[<->](0,.6)--(4,.6)node[midway,above]{$4L$};}
图\ref{fig:m1intervala}的$x'$轴通过两事件；它代表$t'=0$的同时线。此时不变量为$-(\Delta x')^2=-16L^2$，与原系相同。若要求两事件同地，则须令
\[
 0=\Delta x'=\gamma(5L-3\beta L),\qquad \beta=5/3.
\]
这个速度不属于惯性系之间的亚光速变换。也可直接从不变量排除：同地时应有$\Delta s^2=c^2\Delta t'^2\geq0$，与题设矛盾。使两事件同时的$x$向速度是唯一的；若还允许横向变换，可在已经同时的系中继续沿垂直于事件分离方向变换，仍保持同时。
\end{solution}

\subsection{体积和密度}
静止小体积$V_0$在运动方向缩短，实验室的同时截面给$V=V_0/\gamma$。若物质粒子的静质量总和为$M_0$，实验室单位体积所含静质量为$\rho_m=\gamma\rho_{m0}$；它与实验室能量密度不是同一个量。无压力冷物质的能量密度为$u=\gamma^2\rho_{m0}c^2$，因为每个粒子的能量还多一因子$\gamma$。有压力或内部应力的介质须使用能量动量张量，不能继续套用这个冷尘结果。

\begin{example}\label{ex:m1density}
一束粒子在共同静止系的数密度为$n_0$，每粒子静质量$m$、电荷$q$。粒子束在实验室沿$x$正向以$v$运动。求数密度、静质量密度、电荷密度、电流密度及能量密度。
\Diagram{fig:m1densityq}{粒子束的静止体积和实验室同时截面}{\draw[curve](0,0)rectangle(2.7,1.4);\draw[curve](4,0)rectangle(6,1.4);\foreach\x in{.4,1.2,2}{\fill(\x,.7)circle(2pt);}\foreach\x in{4.3,4.9,5.5}{\fill(\x,.7)circle(2pt);}\node[above]at(1.35,1.4){$V_0$};\node[above]at(5,1.4){$V_0/\gamma$};\draw[axis](4,-.5)--(5.8,-.5)node[right]{$v$};}
\end{example}
\begin{solution}
按图\ref{fig:m1densityq}，同一批$N=n_0V_0$粒子占据$V_0/\gamma$，所以$n=\gamma n_0$。每粒子的静质量、电荷不因变换改变，故
\[
 \rho_m=\gamma n_0m,\quad \rho_e=\gamma n_0q,\quad j_x=\rho_ev.
\]
实验室每粒子能量$\gamma mc^2$，于是$u=n\gamma mc^2=\gamma^2n_0mc^2$。
\Diagram{fig:m1densitya}{穿过截面的粒子数决定电流，能量通量另乘单粒子能量}{\draw[curve](3,-.3)--(3,1.8);\node[above]at(3,1.8){面积$A$};\draw[curve](0,0)rectangle(3,1.4);\draw[axis](.5,.7)--(2.5,.7)node[midway,above]{$v$};\node[below]at(1.5,0){$v\dd t$};\node[right]at(3.3,.6){$\dd Q=qnAv\dd t$};}
图\ref{fig:m1densitya}中时间$\dd t$内穿过面积$A$的粒子来自长度$v\dd t$的小柱，故$\dd Q/A\dd t=qnv$。同理能流为$uv$。密度和通量必须使用同一参考系的体积、时间与单粒子物理量。
\end{solution}

\section{速度、光行差与方向变换}\label{th:m1velocity}
对式\eqref{eq:m1lorentz}求微分，并除以$\dd t'$，得
\begin{equation}\label{eq:m1vel}
 v'_x=\frac{v_x-V}{1-Vv_x/c^2},\qquad
 v'_{y,z}=\frac{v_{y,z}}{\gamma(1-Vv_x/c^2)}.
\end{equation}
分母来自不同时性的时间变换，不能只对位移做变换。若$v<c$，四速度的时间分量为正，因而分母始终正。由直接代数还可得
\[
 1-\frac{v'^2}{c^2}=\frac{1-v^2/c^2}{\gamma^2(1-Vv_x/c^2)^2}>0,
\]
所以亚光速速度经洛伦兹变换仍为亚光速。

真空光线在$S'$中的方向角为$\theta'$，相对$x'$正向测量。将$v'_x=c\cos\theta'$、$v'_y=c\sin\theta'$代入逆速度变换，得到
\begin{equation}\label{eq:m1aberration}
 \cos\theta=\frac{\cos\theta'+\beta}{1+\beta\cos\theta'},\quad
 \sin\theta=\frac{\sin\theta'}{\gamma(1+\beta\cos\theta')}.
\end{equation}
两式保留了角度象限，优于单独使用正切式。还可写成半角形式
\[
 \tan\frac\theta2=\sqrt{\frac{1-\beta}{1+\beta}}\tan\frac{\theta'}2.
\]
前半球压缩是方向变换的结果，并不意味着在光源静止系中发光机制变得不对称。

\begin{example}\label{ex:m1aberration}
在光源静止系，光子沿与运动方向垂直的方向发出。求实验室中的光线方向；再求静止系前半球在实验室对应的圆锥及其立体角比例。
\Diagram{fig:m1aberrationq}{源静止系的前半球及边缘光线}{\draw[axis](-1.8,0)--(2.6,0)node[right]{$x'$};\draw[axis](0,-1.6)--(0,2)node[above]{$y'$};\draw[curve](0,-1.4)arc(-90:90:1.4);\draw[very thick,->](0,0)--(0,1.5);\draw[axis](2,.8)--(3.2,.8)node[right]{$V$};}
\end{example}
\begin{solution}
取源系光子四动量为$(E'/c,E'\cos\theta'/c,E'\sin\theta'/c,0)$。逆变换给
\[
 E=\gamma E'(1+\beta\cos\theta'),\quad
 cp_x=\gamma E'(\cos\theta'+\beta),\quad cp_y=E'\sin\theta'.
\]
由光子$p=E/c$，图\ref{fig:m1aberrationq}中光线的方向满足
\[
 \cos\theta=\frac{\cos\theta'+\beta}{1+\beta\cos\theta'},\quad
 \sin\theta=\frac{\sin\theta'}{\gamma(1+\beta\cos\theta')}.
\]
令$\theta'=\pi/2$，边缘方向便为$\cos\theta_c=\beta$、$\sin\theta_c=1/\gamma$。
\Diagram{fig:m1aberrationa}{实验室中的前向光锥，边缘满足$\cos\theta_c=\beta$}{\draw[axis](-.5,0)--(4,0)node[right]{$x$};\draw[curve](0,0)--(3,1.5);\draw[curve](0,0)--(3,-1.5);\draw[guide](3,1.5)arc(90:270:.25 and 1.5);\draw[->](.9,0)arc(0:26.565:.9);\node[right]at(.9,.45){$\theta_c$};\draw[very thick,->](0,0)--(2.4,1.2);}
设$\mu'=\cos\theta'$，对方向变换求导，得到
\[
 \frac{\dd\cos\theta}{\dd\mu'}
 =\frac{(1+\beta\mu')-\beta(\mu'+\beta)}{(1+\beta\mu')^2}
 =\frac{1-\beta^2}{(1+\beta\mu')^2}>0.
\]
源系前半球$0\leq\mu'\leq1$因此完整映为图\ref{fig:m1aberrationa}中的$\beta\leq\cos\theta\leq1$。绕$x$轴旋转此角域，其立体角为
\[
 \Omega_c=2\pi\int_0^{\theta_c}\sin\theta\dd\theta
 =2\pi[-\cos\theta]_0^{\theta_c}
 =2\pi(1-\beta),\qquad \frac{\Omega_c}{4\pi}=\frac{1-\beta}{2}.
\]
高速时$1-\beta\simeq1/(2\gamma^2)$，故$\Omega_c\simeq\pi/\gamma^2$、$\theta_c\simeq1/\gamma$。角域缩小描述的是方向变换；同一角域中的光子数和能量份额还须分别按发射分布与频移加权。
\end{solution}

\begin{example}\label{ex:m1tilt}
一根静止在$S'$的细杆，其端点坐标差为$(l_0\cos\alpha',l_0\sin\alpha')$。求$S$测得的杆长和倾角，并与一束沿杆方向传播的光的方向变换比较。
\Diagram{fig:m1tiltq}{杆的方向由两个同时端点确定}{\draw[axis](-.3,0)--(4,0)node[right]{$x'$};\draw[axis](0,-.2)--(0,2.8)node[above]{$y'$};\draw[very thick](0,0)--(3,2);\node[above]at(1.5,1.1){$l_0$};\draw(.8,0)arc(0:33.69:.8);\node at(1.2,.35){$\alpha'$};}
\end{example}
\begin{solution}
图\ref{fig:m1tiltq}中的杆静止，所以其横向分量不变、纵向分量缩短。对$S$中同时的端点事件，$\Delta x'=\gamma\Delta x$，得到
\[
 l_x=\frac{l_0\cos\alpha'}\gamma,\quad l_y=l_0\sin\alpha',\quad
 l=l_0\sqrt{\frac{\cos^2\alpha'}{\gamma^2}+\sin^2\alpha'}.
\]
故$\tan\alpha=\gamma\tan\alpha'$，实际象限由两个分量确定。
\Diagram{fig:m1tilta}{杆的等时方向与沿杆光线的传播方向有不同定义}{\draw[axis](-.3,0)--(4.6,0)node[right]{$x$};\draw[axis](0,-.2)--(0,3)node[above]{$y$};\draw[very thick](0,0)--(2,2);\node[above left]at(2,2){杆};\draw[->,thick](0,0)--(3.7,1.3);\node[right]at(3.7,1.3){光};\draw[guide](2,0)--(2,2);}
光的两个事件是同一光子先后经过两处的事件，须用式\eqref{eq:m1aberration}，并取$\theta'=\alpha'$，得到$\tan\theta=\sin\alpha'/[\gamma(\cos\alpha'+\beta)]$。图\ref{fig:m1tilta}展示了这种差别。发射后杆仍沿$x$方向运动，因此另一系中的光路方向一般不同于杆的同时截面方向。
\end{solution}

\section{四速度、四动量与四力}\label{th:m1dynamics}
固有时是不变量，因此位置四矢量对固有时的导数仍是四矢量：
\[
 u^\mu=\frac{\dd x^\mu}{\dd\tau}=\gamma_v(c,\vct v),
 \quad u\cdot u=c^2.
\]
静质量不变的质点有$p^\mu=mu^\mu=(E/c,\vct p)$，即$E=\gamma_vmc^2$、$\vct p=\gamma_vm\vct v$。它们满足
\begin{equation}\label{eq:m1shell}
 E^2-c^2p^2=m^2c^4.
\end{equation}
静能不是人为增加在动能上的常数：能量与动量须共同按四矢量变换，静能由它们的变换关系决定。低速展开$\gamma_v=1+v^2/(2c^2)+\cdots$后，才恢复$E=mc^2+mv^2/2+\cdots$。

令三力$\vct f=\dd\vct p/\dd t$。对式\eqref{eq:m1shell}求导，$E\dd E=c^2\vct p\cdot\dd\vct p$，并利用$c^2\vct p/E=\vct v$，得
\[
 \frac{\dd E}{\dd t}=\vct f\cdot\vct v=\mathcal P.
\]
因而四力为$F^\mu=\dd p^\mu/\dd\tau=\gamma_v(\mathcal P/c,\vct f)$。四加速度$a^\mu=\dd u^\mu/\dd\tau$与四速度正交：对$u^2=c^2$求导，得$u\cdot a=0$。静质量固定时$F^\mu=ma^\mu$，从而$u\cdot F=0$。若研究吸收能量导致静质量变化的物体，则$F^\mu=m a^\mu+u^\mu\dd m/\dd\tau$，不能略去后一项。

\begin{example}\label{ex:m1restforce}
质点瞬时沿$x$方向运动，实验室三力为$(f_\parallel,f_\perp,0)$。求四力的平方，并求瞬时静止系的三力。
\Diagram{fig:m1restforceq}{速度和三力的平行、垂直分解}{\draw[axis](0,0)--(3.7,0)node[right]{$\vct v$};\draw[axis](0,0)--(2.4,1.8)node[right]{$\vct f$};\draw[guide](2.4,0)--(2.4,1.8)--(0,1.8);\node[below]at(1.2,0){$f_\parallel$};\node[left]at(0,.9){$f_\perp$};}
\end{example}
\begin{solution}
四力定义为$F^\mu=\dd p^\mu/\dd\tau$，而三力为$\vct f=\dd\vct p/\dd t$。图\ref{fig:m1restforceq}中速度沿$x$，功率为$\mathcal P=\vct f\cdot\vct v=f_\parallel v$，故
\[
 F^\mu=\gamma_v(f_\parallel v/c,f_\parallel,f_\perp,0).
\]
取度规$(+,-,-,-)$，逐项求平方：
\begin{align*}
 F^2&=\gamma_v^2[f_\parallel^2(\beta_v^2-1)-f_\perp^2]\\
 &=-f_\parallel^2-\gamma_v^2f_\perp^2.
\end{align*}
这里使用了$\gamma_v^2(1-\beta_v^2)=1$。
\Diagram{fig:m1restforcea}{瞬时静止系中的力；此时四力时间分量为零}{\draw[axis](0,0)--(3.6,0)node[right]{$x'$};\draw[axis](0,0)--(0,3)node[above]{$y'$};\draw[very thick,->](0,0)--(2,2.4)node[right]{$\vct f_0$};\draw[guide](2,0)--(2,2.4)--(0,2.4);\node[below]at(1,0){$f_\parallel$};\node[left]at(0,1.2){$\gamma_vf_\perp$};}
沿$x$以速度$v$变换至图\ref{fig:m1restforcea}的瞬时静止系。时间、纵向和横向分量分别为
\begin{align*}
 F'^0&=\gamma_v(F^0-\beta_vF^x)=0,\\
 F'^x&=\gamma_v(F^x-\beta_vF^0)
 =\gamma_v^2f_\parallel(1-\beta_v^2)=f_\parallel,\\
 F'^y&=F^y=\gamma_v f_\perp.
\end{align*}
该事件处$v'=0$、$\gamma_{v'}=1$，所以四力空间分量就是三力，得到$f_{0\parallel}=f_\parallel$、$f_{0\perp}=\gamma_vf_\perp$，且$f_0^2=-F^2$。恒静质量粒子的四力因而是类空或零矢量。时间分量为零表示该系中瞬时功率为零；粒子静止是充分条件，力与速度垂直时也满足这一条件。
\end{solution}

\section{运动光源的角分布}
\subsection{频率与立体角的雅可比}
光子的四动量为$(h\nu/c,(h\nu/c)\vct n)$。对时间分量作变换，得
\[
 \nu=\gamma\nu'(1+\beta\cos\theta')
 =\frac{\nu'}{\gamma(1-\beta\cos\theta)}=\delta\nu',
\]
其中$\delta$称多普勒因子。例\ref{ex:m1aberration}已给出光行差商式的求导过程；以$\mu'=\cos\theta'$表示，结果为
\[
 \dd\mu=\frac{1-\beta^2}{(1+\beta\mu')^2}\dd\mu',
 \qquad \dd\Omega=\delta^{-2}\dd\Omega'.
\]
方位角不变，因而立体角的变换只需处理这一导数。

\begin{example}\label{ex:m1power}
点光源在自身静止系中各向同性发光，单位固有时、单位立体角的辐射能为$I_0$，总功率$P_0=4\pi I_0$。光源以速度$V$沿$x$运动。求按实验室发射时间定义的角功率$I(\theta)$及总功率，并与远处探测器按接收时间定义的角功率比较。
\Diagram{fig:m1powerq}{同一小束光在两系中占据不同立体角}{\draw[axis](-.5,0)--(4.2,0)node[right]{$V$};\draw[curve](0,0)--(3,1);\draw[curve](0,0)--(3,1.7);\draw[curve](3,1)arc(-90:90:.18 and .35);\node[right]at(3.3,1.4){$\dd\Omega$};\draw[guide](0,0)--(3,1.35);\node at(1.2,.24){$\theta$};}
\end{example}
\begin{solution}
先在图\ref{fig:m1powerq}的一束光中数同一批光子，光子数不变，单光子能量乘$\delta$，故$\dd E=\delta\dd E'$。光源两个发射事件位于同一世界线上，实验室发射时间$\dd t=\gamma\dd\tau$。因此
\begin{equation}\label{eq:m1angular}
 I(\theta)=\frac{\dd E}{\dd t\dd\Omega}
 =I_0\frac\delta\gamma\frac{\dd\Omega'}{\dd\Omega}
 =\frac{I_0}{\gamma^4(1-\beta\cos\theta)^3}.
\end{equation}
积分总功率时使用实验室立体角，令$\mu=\cos\theta$：
\begin{align*}
 P&=\frac{2\pi I_0}{\gamma^4}\int_{-1}^{1}\frac{\dd\mu}{(1-\beta\mu)^3}\\
 &=\frac{2\pi I_0}{\gamma^4}
 \left[\frac{1}{2\beta(1-\beta\mu)^2}\right]_{-1}^{1}\\
 &=\frac{2\pi I_0}{\gamma^4}\frac1{2\beta}
 \frac{(1+\beta)^2-(1-\beta)^2}{(1-\beta^2)^2}
 =4\pi I_0.
\end{align*}
原函数由$z=1-\beta\mu$、$\dd\mu=-\dd z/\beta$及$\int z^{-3}\dd z=-z^{-2}/2$得到。以上写法取$\beta\ne0$；静止情形直接积分$I_0$，也得$P=4\pi I_0$。
这个结果还有整体四动量解释：静止系一次各向同性发光的总动量为零，实验室总能量乘$\gamma$，发射时间也乘$\gamma$。
\Diagram{fig:m1powera}{两次发射到达固定远处探测器的路程差}{\draw[axis](-.3,0)--(4,0)node[right]{$x$};\fill(0,0)circle(1.4pt)node[below]{$S_1$};\fill(1,0)circle(1.4pt)node[below]{$S_2$};\fill(4,2)circle(1.4pt)node[right]{$D$};\draw[curve](0,0)--(4,2);\draw[curve](1,0)--(4,2);\draw[guide](1,0)--(.8,.4);\node[above]at(2,1.5){$\theta$近似不变};}
图\ref{fig:m1powera}中两个发射点相距$V\dd t$，第二束路程缩短$V\dd t\cos\theta$，因此
\[
 \dd t_{\rm arr}=(1-\beta\cos\theta)\dd t,
 \quad I_{\rm arr}(\theta)=I_0\delta^4.
\]
这里探测器足够远，方向在两次发射间近似不变。接收时间的角功率不能再对同一个接收时刻的全天球方向简单积分而解释成同一批发射光的总功率；不同方向对应不同的发射时段。
\end{solution}

\section{坐标测量中的事件选择}
\subsection{从光锥得到变换系数}
惯性系的时空均匀性要求：同一个坐标差采用同一变换，与起点位置、起始时刻无关。坐标关系因而具有仿射形式；将原点对齐后，变换为线性。先把原点对齐，沿$x$的变换可写$x'=a x+bct$、$ct'=d x+ect$。$S'$原点在$S$中满足$x=\beta ct$，因此$b=-a\beta$。

分别代入向右、向左的两条原点光线。$x=ct$要求$x'=ct'$，所以$a+b=d+e$；$x=-ct$要求$x'=-ct'$，所以$-a+b=d-e$。两式相加、相减给$d=b=-a\beta$、$e=a$。至此光锥不变只确定了整体比例$a$，尚未确定其数值。用相对性原理要求逆变换在交换两系后只有$\beta$变号，正逆矩阵相乘便给$a^2(1-\beta^2)=1$。连续到$\beta=0$的恒等变换，选$a=\gamma>0$。

这一推导也说明为什么单独要求光速不变还不能任意省略比例系数。若同时允许改变坐标单位，光锥形状不变但标尺标度也会改变；本书要求两系采用相同物理单位和互易的标准钟尺。

横向坐标不变可由对称性理解：沿$x$的运动没有选出横向某一方向。设横向统一缩放为$k(\beta)$，正逆变换要求$k(\beta)k(-\beta)=1$，空间反射对称给两者相等，连续正根为1。这排除了把横向尺寸也按$\gamma$缩短的做法。

\subsection{同一事件、同一物体与同一时刻}
坐标题有三种常见约束，各自消去不同未知量。两条世界线相遇要求同一事件的$t,x$都相同；跟踪同一物体是在其世界线上取两事件；同时测长度则是在某一系固定$t$后取不同物体的事件。只有把题设约束先转为这些关系，代入变换才有明确意义。

固定一个钟的两个事件满足$x'=\text{常数}$，因而$\Delta x=V\Delta t$。代入时间变换得$\Delta t'=\gamma\Delta t(1-\beta^2)=\Delta t/\gamma$。若改成在$S$同时读取两只钟，便有$\Delta t=0$而不是$\Delta x=V\Delta t$；相同的式子转而给$\Delta t'=-\gamma V\Delta x/c^2$。两种情境都由同一变换式产生，不能把它们当作两个可随意替换的“钟慢规则”。

\Diagram{fig:m1pairings}{沿同一钟的两事件与沿同一等时线的两事件，是不同配对}{\draw[axis](-.2,0)--(4.4,0)node[right]{$x$};\draw[axis](0,0)--(0,4.5)node[above]{$ct$};\draw[curve](0,0)--(2.4,4);\draw[curve](1,0)--(3.4,4);\fill(.6,1)circle(1.5pt)node[left]{$A$};\fill(1.8,3)circle(1.5pt)node[left]{$B$};\fill(2.8,3)circle(1.5pt)node[right]{$C$};\draw[guide](1.8,3)--(2.8,3);}
图\ref{fig:m1pairings}中$A,B$在同一钟世界线上，$B,C$在同一实验室等时线上。$AB$适合求钟经历的时间，$BC$适合求实验室钟距。若问题还给定钟面被设过零，最后须在坐标时间上加钟的设置常数，不能让这个常数反过来改变世界线。

\subsection{四钟模型的完整事件坐标}
例\ref{ex:m1rods}可把所有坐标集中写成无量纲表。定义$g=\gamma(1+\beta)$、$a=(1-1/\gamma)/\beta$，坐标单位分别为$t_0,l$。$A$钟读$t/t_0$，$B$钟读$t/t_0-a$；$A'$钟读$t'/t_0$，$B'$钟读$t'/t_0+a$。表中每行都代表一个事件，列之间只做坐标变换。
\begin{center}\small
\begin{tabular}{ccccc}\toprule
事件&$t/t_0$&$x/l$&$t'/t_0$&$x'/l$\\\midrule
$M$&0&0&0&0\\
$N$&$a$&1&$-a$&1\\
$P$&1&1&$1/g$&$1/g$\\
$R$&$g$&$g$&1&1\\
$Q$&$2g-1$&1&$\gamma(2g-1-\beta)$&$\gamma[1-\beta(2g-1)]$\\
$T$&$2\gamma$&$2\gamma\beta$&2&0\\
$U$&$2g$&0&$2\gamma g$&$-2\gamma\beta g$\\\bottomrule
\end{tabular}
\end{center}
$P$行使用$\gamma(1-\beta)=1/g$；$R$行则由源系最简单的$(1,1)$变来。$N$行的$t'_N=-t_N$是本题两杆等固有长和特定归零事件共同造成的结果，换成不等长两杆后一般不再成立。

表中的$Q,T$两行也重现例\ref{ex:m1rods}已求出的返回次序及临界速度$\beta=1/\sqrt5$。几何上比较$x_T=2\gamma\beta l$与固定$B$位置$l$即可判断：反射光向左走，先遇到更右侧的对象。$\beta=.6$时$A'$在右，故先$T$后$Q$。

这些次序针对有因果联系的返回光线，同一系或另一系都会给出相同先后。局部归零后的两只不同钟却可能读出反向大小，因此“表上较大的读数”不能单独代替事件次序。$\beta\to0$时$a\simeq\beta/2\to0$，公式连续趋于静止往返，但正好$\beta=0$时两杆端点一直重合，所谓第二次相遇归零事件需要另作规定。

\section{四维量的几何和角分布的测度}
\subsection{任意方向的纯推动}
沿单位矢量$\hat{\vct V}$分解空间矢量$\vct r=\vct r_\parallel+\vct r_\perp$。对平行分量使用一维变换，垂直分量保持不变，得到
\[
 ct'=\gamma\left(ct-\frac{\vct V\cdot\vct r}{c}\right),\quad
 \vct r'=\vct r+\left[\frac{\gamma-1}{V^2}(\vct V\cdot\vct r)-\gamma t\right]\vct V.
\]
$V=0$的极限取恒等映射。对任意四矢量，只需把$ct$换成其时间分量、把$\vct r$换成空间分量。四动量的能量变换因而是$E'=\gamma(E-\vct V\cdot\vct p)$。这里$\vct V$是坐标系速度，$\vct p$是被研究粒子的动量，二者没有理由总是平行。

洛伦兹标量在任意允许的惯性变换下都保持数值，如固有时、静质量和两个四矢量的缩并。某一分量在一次特殊变换中恰好不变，还不足以判定它是标量；光子能量就可能在特定方向和推动速度下不变，换一种变换后又改变。

类时四矢量的空间分量可在某个系全为零。非零类光四矢量的闵可夫斯基平方为零，时间与空间分量都不能通过有限惯性变换消去；真空光子因而没有静止系。类空四矢量可以把时间分量变为零，却不能把空间分量全消去。这些分类取决于闵可夫斯基内积；时空图的纸面长度不表示四矢量模。

\subsection{单位角、单位立体角与环带}
轴对称分布中，单位立体角强度$I(\theta)$与极角环带内功率不同。环带固体角$\dd\Omega=2\pi\sin\theta\dd\theta$，故$\dd P/\dd\theta=2\pi I(\theta)\sin\theta$。即使$I(\theta)$在$\theta=0$最大，有限环带功率的密度也因$\sin\theta$在轴上为零而不在那里最大。画角分布前须说明纵轴是哪种密度。

\Diagram{fig:m1solidring}{极角环带宽度相同，立体角仍随$\sin\theta$改变}{\draw[curve](0,0)circle(1.7);\draw[axis](0,0)--(2.7,0)node[right]{轴};\draw[curve](0,0)--(1.472,.85);\draw[curve](0,0)--(1.302,1.093);\draw[thick](1.472,.85)arc(30:40:1.7);\node[above]at(.4,2){$\dd\Omega=2\pi\sin\theta\dd\theta$};}
图\ref{fig:m1solidring}以截面短弧表示球面环带。极坐标曲线的半径表示所选的角密度，积分时仍须乘立体角雅可比。运动源的功率份额可先选定一个系中的角域，再变换边界计算，往往能省去完整角分布的积分。

更一般地，源系辐射不必各向同性。设单位固有时发出的总能量为$P_0$，总$x$向动量率为$\dot p'_x$。整个发射小过程的四动量变换给$\dd E=\gamma(\dd E'+\beta c\dd p'_x)$，而发射时间$\dd t=\gamma\dd\tau$，于是实验室总发射功率为
\[
 P=P_0+\beta c\dot p'_x.
\]
各向同性源有$\dot p'_x=0$，所以总功率不变；向后单束光有$\dot p'_x=-P_0/c$，所以$P=P_0(1-\beta)$。两种结果都由同一四动量式得到。总功率不变需要源系净辐射动量为零，各向同性发射满足这一条件。

\subsection{密度变换中的超平面与通量}
粒子数守恒并不表示不同参考系中的数密度相等。计数所用的体积必须处在该系的同时面上；同一批粒子在两个同时面上被截取的位置不同。对共同四速度为$U^\mu$的稀疏粒子流，定义数流$N^\mu=n_0U^\mu$。在实验室中，
\[
 N^\mu=(cn,\vct j_N)=n_0\gamma_v(c,\vct v),\qquad
 n=\gamma_vn_0,\quad \vct j_N=n\vct v.
\]
穿过时空小面元的粒子数为$\dd\mathcal N=N^\mu\dd\Sigma_\mu/c$。实验室等时小体积有$\dd\Sigma_\mu=(\dd V,0,0,0)$，给出$n\dd V$；跨越固定空间面的面元则给出粒子通量。密度与流量是同一四维对象的不同分量，变换时不能仅替换长度而保持时间率不变。

无热运动的粒子束还具有能量动量流$T^{\mu\nu}=n_0mU^\mu U^\nu$。其分量为
\[
 T^{00}=n_0m\gamma_v^2c^2=u,\quad
 cT^{0x}=uv_x,\quad T^{xx}=n_0m\gamma_v^2v_x^2.
\]
最后一项是单位面积、单位时间跨面的$x$动量；即使束在自身系没有压力，运动系也测得定向动量通量。对存在随机热运动的气体，须对每个粒子的速度分布求和，不能再把整个系统视为只有一个共同$U^\mu$的冷束。
\Diagram{fig:m1currentcut}{同一粒子流由不同同时面截取；斜线不代表粒子轨迹的偏折}{\draw[axis](0,0)--(4.5,0)node[right]{$x$};\draw[axis](0,0)--(0,3.7)node[above]{$ct$};\foreach\a in{.4,1.2,2,2.8}{\draw[curve](\a,0)--({\a+.7},3.2);}\draw[guide](.1,1.1)--(3.8,1.1)node[right]{$t=\text{常量}$};\draw[densely dotted](.1,1.6)--(3.8,2.7)node[right]{$t'=\text{常量}$};}
图\ref{fig:m1currentcut}的两条截线各自连接一组不同事件。密度变换应同时指明计数的同时面和跨面通量，两者由同一四流联系。静质量密度$n m$、总能量密度$u$和$u/c^2$分别包含不同的洛伦兹因子，不能互作同义量。

\section{有限装置中的同时事件与雷达测距}
相对论题中的“同时打开”“一起读数”必须绑定某一组静止同步钟。装置的连接条件还可能受信号传播限制。几何上可任意选择一条等时线与若干世界线的交点；物理上是否能用同一触发信号控制这些交点，要另看它们的因果关系。下列模型把这两个层次分开。

\begin{example}\label{ex:m1tunnel}
隧道静止长度$L$，列车固有长也是$L$，列车以$\beta c$进入。地面系在列车中点经过隧道中点时，同时短暂关闭两门并立即打开。证明地面看来列车能在两门之间，并求车系两门关闭的时间差和位置差，解释车系隧道更短却无矛盾。
\Diagram{fig:m1tunnelq}{地面同时关门，列车等时长度为$L/\gamma$}{\draw[axis](-.5,0)--(6,0)node[right]{$x$};\draw[very thick](0,-.8)--(0,1.3);\draw[very thick](5,-.8)--(5,1.3);\draw[curve](.6,.2)rectangle(4.4,.8);\draw[axis](2.3,1.6)--(3.7,1.6)node[right]{$V$};\node[below]at(2.5,-.8){两门距离$L$};}
\end{example}
\begin{solution}
取关门中点事件为$(ct,x)=(0,0)$，两门关闭为$D_L=(0,-L/2)$、$D_R=(0,L/2)$。图\ref{fig:m1tunnelq}的列车两端在该时刻位于$\pm L/(2\gamma)$，严格处于两门之间。变换两关闭事件，得
\[
 ct'_L=\frac{\gamma\beta L}2,\quad x'_L=-\frac{\gamma L}2;
 \qquad ct'_R=-\frac{\gamma\beta L}2,\quad x'_R=\frac{\gamma L}2.
\]
所以右门先关，时间差$t'_R-t'_L=-\gamma\beta L/c$。两次关门的事件位置差是$\gamma L$，并非车系同时隧道长度$L/\gamma$。
\Diagram{fig:m1tunnela}{车系中两关闭事件不同时；虚线为车系等时线}{\draw[axis](-2.6,0)--(3,0)node[right]{$x'$};\draw[axis](0,-2)--(0,2.4)node[above]{$ct'$};\draw[curve](-2,1.2)--(0,-2.1);\draw[curve](0,2.1)--(2,-1.2);\fill(-2,1.2)circle(1.4pt)node[left]{$D_L$};\fill(2,-1.2)circle(1.4pt)node[right]{$D_R$};\draw[guide](-2,1.2)--(1.5,1.2);\draw[guide](-1.5,-1.2)--(2,-1.2);}
图\ref{fig:m1tunnela}中前门关闭时，列车尾部尚在隧道外；后门稍后关闭时，前门早已打开，车头已可通过。车系从未声称两门同时夹住整列车。若要求门持续关闭至列车撞门并停止，则题目变为受力和信号传播问题，不能继续把列车视为始终匀速的刚体。
\end{solution}

\begin{example}\label{ex:m1radar}
观察者固定在$x=0$，匀速飞行的反射镜满足$x=L+Vt$，$L>0,V>0$。观察者在$t=0$发光，光在镜面反射后返回。求反射、返回时刻，以及观察者按雷达约定赋给反射事件的时间和距离。
\Diagram{fig:m1radarq}{出射光追赶向外运动的镜面}{\STAxes\draw[curve](1,0)--(3,4);\node[right]at(3,4){镜};\draw[very thick,->](0,0)--(2,2);\node[below]at(1,0){$L$};}
\end{example}
\begin{solution}
图\ref{fig:m1radarq}的交点满足$ct_R=L+Vt_R$，所以$t_R=L/(c-V)$、$x_R=cL/(c-V)$。反射光以$c$返回原点，返回时刻$t_A=t_R+x_R/c=2L/(c-V)$。
\Diagram{fig:m1radara}{反射事件的雷达时间是发射、返回时刻的平均}{\draw[axis](-.2,0)--(4.6,0)node[right]{$x$};\draw[axis](0,-.2)--(0,5)node[above]{$ct$};\draw[very thick,->](0,0)--(2,2);\draw[very thick,->](2,2)--(0,4);\draw[guide](0,2)--(2,2);\node[left]at(0,2){$ct_{\rm rad}$};\node[right]at(2,2){$R$};\node[left]at(0,4){$ct_A$};}
雷达定义$t_{\rm rad}=(0+t_A)/2=L/(c-V)$，$r_{\rm rad}=ct_A/2=cL/(c-V)$，与图\ref{fig:m1radara}中反射事件的惯性坐标$t_R,x_R$完全一致。雷达同步把距离赋给反射事件，因此这里应与反射时刻的镜距比较。发射时的$L$和接收时的$L+Vt_A$属于另外两个事件。
\end{solution}

\begin{example}\label{ex:m1meeting}
两质点在实验室$t=0$位于$x=0,L$，速度分别为$+u,-w$，其中$0<u,w<c$。求相遇事件及在第一质点静止系中的相遇时间，说明为什么该系的初始两点距离不能直接取$L/\gamma_u$再任意相除。
\Diagram{fig:m1meetingq}{两条匀速世界线的交点}{\draw[axis](-.3,0)--(4.7,0)node[right]{$x$};\draw[axis](0,-.2)--(0,4.5)node[above]{$ct$};\draw[curve](0,0)--(2,3.5);\draw[curve](4,0)--(2,3.5);\node[below]at(4,0){$L$};\node[right]at(2,3.5){$C$};}
\end{example}
\begin{solution}
图\ref{fig:m1meetingq}满足$ut=L-wt$，所以$t_C=L/(u+w)$、$x_C=uL/(u+w)$。第一质点系以$u$运动，且把其初始事件归零，因此
\[
 t'_C=\gamma_u(t_C-ux_C/c^2)=\frac{L}{\gamma_u(u+w)}.
\]
这是第一质点从出发到相遇的固有时。若要在其系中用距离除以速度，必须重新求$t'=0$与第二质点世界线的交点。
\Diagram{fig:m1meetinga}{在第一质点系中，第二质点初始位置取在$t'=0$线上}{\draw[axis](-.2,0)--(4.3,0)node[right]{$x'$};\draw[axis](0,-.2)--(0,4.4)node[above]{$ct'$};\draw[curve](3.4,0)--(0,3.7);\draw[very thick](0,0)--(0,3.7);\node[below]at(3.4,0){$D$};\node[left]at(0,3.7){$C$};}
令第二质点$x=L-wt$满足$t-ux/c^2=0$，得到$t_D=uL/(c^2+uw)$，$x_D=Lc^2/(c^2+uw)$，进而$x'_D=L/[\gamma_u(1+uw/c^2)]$。第二质点相对速度为$-(u+w)/(1+uw/c^2)$。图\ref{fig:m1meetinga}距离除以相对速大小，恰得上式$t'_C$。$L$是两运动物体间的某系等时距离，不是一个共同静止杆的固有长。
\end{solution}

\section{光锥、波面和速度空间}
\begin{example}\label{ex:m1flash}
点光源在事件$O$瞬时闪光。证明实验室固定时刻$t$的光前沿是以$O$空间位置为中心的球面，并说明与运动光源角分布的关系。
\Diagram{fig:m1flashq}{光锥的水平截面是球面，不随之后的光源位置移动}{\draw[axis](-2.4,0)--(2.6,0)node[right]{$x$};\draw[axis](0,-.2)--(0,4)node[above]{$ct$};\draw[curve](0,0)--(2,3.2);\draw[curve](0,0)--(-2,3.2);\draw[curve](0,3.2)ellipse(2 and .4);\draw[guide](0,0)--(.9,3.2);\node[right]at(.9,2.6){源世界线};}
\end{example}
\begin{solution}
每个光前沿事件与$O$的间隔为零。在源静止系$\Delta s'^2=0$，变换不改间隔，故$c^2t^2-x^2-y^2-z^2=0$，取未来支$t>0$便得半径$ct$球面。图\ref{fig:m1flashq}中的源在闪光后已移动，其当前位置不是这一次闪光球面的中心。
\Diagram{fig:m1flasha}{球面形状与能量分布分属几何和光子统计两个层面}{\draw[curve](0,0)circle(1.6);\foreach\a in{0,20,40,60,90,140,210,270,300,330,345}{\draw[->](0,0)--({1.6*cos(\a)},{1.6*sin(\a)});}\draw[axis](0,0)--(2.8,0)node[right]{$V$};\node[below]at(0,-2){球面上能量可不均匀};}
图\ref{fig:m1flasha}用较密箭头示意前向辐射更集中。球面由传播速度和发射事件决定，角能量密度由光行差、单光子频移决定；各向同性波速并不要求各向同性光子分布。对持续运动源，每一次发射有各自的球面，所有球面叠加的稳态场又是另一问题。
\end{solution}

\begin{example}\label{ex:m1lightellipse}
在源系，一次闪光到光子发射后固定源系时刻$t'_0$的全部前沿点组成圆。将这一组事件变到实验室，求它们在$x,y$坐标投影中的曲线，并说明为什么不违背上一例的球面结论。
\Diagram{fig:m1lightellipseq}{本题固定源系时间，所选事件不在实验室同时面上}{\draw[axis](-2.2,0)--(2.5,0)node[right]{$x'$};\draw[axis](0,-2)--(0,2.3)node[above]{$y'$};\draw[curve](0,0)circle(1.5);\node[above]at(1.5,1.5){$r'=ct'_0$};}
\end{example}
\begin{solution}
按图\ref{fig:m1lightellipseq}，令$r_0=ct'_0$，各事件$x'=r_0\cos\theta'$、$y'=r_0\sin\theta'$。逆变换给
\[
 x=\gamma r_0(\cos\theta'+\beta),\quad y=r_0\sin\theta',\quad
 ct=\gamma r_0(1+\beta\cos\theta').
\]
\Diagram{fig:m1lightellipsea}{异时事件的空间投影是平移的椭圆}{\draw[axis](-1.5,0)--(3.8,0)node[right]{$x$};\draw[axis](0,-2)--(0,2.2)node[above]{$y$};\draw[curve](1.322876,0)ellipse(2 and 1.5);\draw[guide](1.322876,-1.5)--(1.322876,1.5);\node[above]at(1.322876,1.6){中心$\gamma\beta r_0$};}
消去$\theta'$，图\ref{fig:m1lightellipsea}的投影满足$(x-\gamma\beta r_0)^2/(\gamma^2r_0^2)+y^2/r_0^2=1$。但这些点的实验室时间随角度变化，范围为$\gamma t'_0(1\pm\beta)$。上一例固定的是同一实验室时间，两题选择的事件集合不同，因此投影为椭圆并不表示实验室真空光速各向异性。
\end{solution}

\begin{example}\label{ex:m1relativeangle}
在$S'$中，两粒子速度大小同为$u$，方向分别与$x'$正向成$\pm\alpha$。求实验室中的夹角，并讨论当$u=c$和$u\ll c$时有何差别。
\Diagram{fig:m1relativeangleq}{关于运动轴对称的两条发射速度}{\draw[axis](0,0)--(3.7,0)node[right]{$V$};\draw[axis](0,0)--(2.5,1.6)node[right]{$u$};\draw[axis](0,0)--(2.5,-1.6)node[right]{$u$};\node at(1,.35){$\alpha$};}
\end{example}
\begin{solution}
图\ref{fig:m1relativeangleq}两粒子都有$u'_x=u\cos\alpha$，横向分量为$\pm u\sin\alpha$。逆速度变换给
\[
 v_x=\frac{u\cos\alpha+V}{1+Vu\cos\alpha/c^2},\quad
 v_y=\frac{\pm u\sin\alpha}{\gamma(1+Vu\cos\alpha/c^2)}.
\]
当变换后的速度非零时，每条速度与$x$的有向角可用$\vartheta=\operatorname{atan2}(u\sin\alpha,\gamma(u\cos\alpha+V))$表示，两线夹角取$2\vartheta$的较小几何夹角。若$u\sin\alpha=0$且$u\cos\alpha+V=0$，粒子在实验室静止，方向及夹角均无定义，不能调用$\operatorname{atan2}(0,0)$。
\Diagram{fig:m1relativeanglea}{同向变换压缩前向发射角，但须保留纵向分量的符号}{\draw[axis](0,0)--(4,0)node[right]{$x$};\draw[axis](0,0)--(3.2,.8);\draw[axis](0,0)--(3.2,-.8);\draw[guide](0,0)--(2,1.5);\draw[guide](0,0)--(2,-1.5);\node[right]at(3.2,.8){实验室};\node[right]at(2,1.5){源系方向示意};}
若$u=c$，此式化为光行差，所有粒子的速率仍为$c$。若$u,V\ll c$，$\gamma\simeq1$且分母近1，恢复伽利略速度矢量平移。图\ref{fig:m1relativeanglea}取$\cos\alpha>0$；若发射近后方使$u\cos\alpha+V<0$，两粒子仍可能向后运动，单独使用正切会误把它们画到前半平面。
\end{solution}

\section{角分布积分与四维微分关系}
\begin{example}\label{ex:m1hemipower}
各向同性静止光源以速度$V$运动，按实验室发射时间定义总功率。求实验室前半球$0\leq\theta\leq\pi/2$收到的功率份额，并与源系前半球光子集合的份额区分。
\Diagram{fig:m1hemipowerq}{实验室前半球是按实验室方向筛选的集合}{\draw[axis](-1.8,0)--(2.6,0)node[right]{$x$};\draw[curve](0,-1.6)arc(-90:90:1.6);\draw[guide](0,-1.8)--(0,1.8);\draw[axis](0,0)--(1.4,.8);}
\end{example}
\begin{solution}
使用式\eqref{eq:m1angular}，图\ref{fig:m1hemipowerq}的前半球积分为
\[
 P_F=\frac{2\pi I_0}{\gamma^4}\int_0^1\frac{\dd\mu}{(1-\beta\mu)^3}
 =\frac{2\pi I_0}{\gamma^4}
 \left[\frac{(1-\beta\mu)^{-2}}{2\beta}\right]_0^1
 =\frac{\pi I_0}{\beta\gamma^4}[(1-\beta)^{-2}-1].
\]
除以总功率$4\pi I_0$，再用$\gamma^{-4}=(1-\beta)^2(1+\beta)^2$，得到
\[
 \frac{P_F}P=\frac{(2-\beta)(1+\beta)^2}{4}
 =\frac{2+3\beta-\beta^3}{4}.
\]
\Diagram{fig:m1hemipowera}{前半球功率份额从一半连续趋向全部}{\draw[axis](0,0)--(4.2,0)node[right]{$\beta$};\draw[axis](0,0)--(0,2.8)node[above]{$P_F/P$};\draw[curve,domain=0:1,samples=70]plot({3.5*\x},{2.3*(2+3*\x-\x^3)/4});\node[left]at(0,1.15){$1/2$};\node[left]at(0,2.3){$1$};\node[below]at(3.5,0){$1$};}
图\ref{fig:m1hemipowera}在$\beta\to0$给$1/2$，在$\beta\to1$趋1。这里边界是实验室$\cos\theta=0$，对应源系$\cos\theta'=-\beta$，包含源系后半球的一部分；不能替换成例\ref{ex:m1coneenergy}中源系前半球所占的$1/2+\beta/4$能量份额。
\end{solution}

\begin{example}\label{ex:m1curvedproper}
粒子在实验室作半径$R$、速率$v$的匀速圆周运动，转一周后回到同一位置。求粒子固有时、四加速度及其不变量，说明速度方向变化是否影响钟慢因子。
\Diagram{fig:m1curvedq}{匀速圆周速度切向、加速度径向}{\draw[curve](0,0)circle(1.5);\fill(1.5,0)circle(2pt);\draw[axis](1.5,0)--(1.5,1.2)node[right]{$v$};\draw[axis](1.5,0)--(.3,0)node[above]{$v^2/R$};}
\end{example}
\begin{solution}
图\ref{fig:m1curvedq}中实验室周期$T=2\pi R/v$。速率恒定，故$\tau=\int_0^T\sqrt{1-v^2/c^2}\dd t=[t/\gamma_v]_0^T=T/\gamma_v$。方向变化不改变局部钟慢因子，固有时仍沿粒子的实际世界线累积。
\Diagram{fig:m1curveda}{圆周世界线在时空中成螺旋，固有时沿世界线累积}{\draw[axis](0,0)--(0,4.2)node[above]{$ct$};\draw[curve,domain=0:720,samples=150]plot({1.5*cos(\x)},{.005*\x+.3*sin(\x)});\node[right]at(1.7,3.6){空间圆周投影示意};}
先对$\gamma_v=(1-v^2/c^2)^{-1/2}$用链式法则求导：$\dot\gamma_v=\gamma_v^3\vct v\cdot\vct a/c^2=0$。再对$u^\mu=(\gamma_vc,\gamma_v\vct v)$逐项求导，
\[
 a^\mu=\gamma_v\frac{\dd u^\mu}{\dd t}
 =\gamma_v(c\dot\gamma_v,\dot\gamma_v\vct v+\gamma_v\vct a)
 =(0,\gamma_v^2\vct a).
\]
所以$a^2=-\gamma_v^4v^4/R^2$，固有加速度$\alpha=\gamma_v^2v^2/R$。图\ref{fig:m1curveda}是螺旋的二维投影，投影交叉点未必是同一时空事件。留在起点的实验室钟走过$T$，粒子钟少走$T(1-1/\gamma_v)$。
\end{solution}

\section{相对速率的标量表达与辐射的局部能量}
\subsection{任意方向的相对速率}
两粒子的实验室速度分别为$\vct u,\vct v$，四速度为$U=\gamma_u(c,\vct u)$、$V=\gamma_v(c,\vct v)$。在第二粒子静止系，$V'=(c,0,0,0)$，故$U'\cdot V'=\gamma_{\rm rel}c^2$。内积不变立即给
\[
 \gamma_{\rm rel}=\gamma_u\gamma_v
 \left(1-\frac{\vct u\cdot\vct v}{c^2}\right),\qquad
 w=c\sqrt{1-\gamma_{\rm rel}^{-2}}.
\]
这里$w$指任一粒子静止系中测得的另一粒子速率，交换两粒子后数值不变。在第三个参考系中作速度矢量差，得到的是该系等时位置差的变化率；它与上述相对速率采用了不同的测量程序。

若两速度共线同向，公式化为$w=|u-v|/(1-uv/c^2)$；共线反向则分子变为$u+v$、分母为$1+uv/c^2$。非共线时，先算夹角余弦形成的点积通常比逐个变换三个分量更直接。若还需要相对运动方向，才补做速度变换；一个标量速率不足以确定方向。

\begin{example}正交运动的两粒子
实验室中$A$沿$x$以$0.8c$运动，$B$沿$y$以$0.6c$运动，两者在原点相遇。求它们的相对速率，并在$B$的静止系中用分量变换核对。
\ExpansionDiagram{fig:dep-m1-relative-q}{相遇事件处的两条实验室速度；虚线表示坐标速度差}{
\draw[->](0,0)--(4,0)node[right]{$x$};\draw[->](0,0)--(0,3)node[above]{$y$};\draw[->,thick](0,0)--(3.2,0)node[midway,below]{$0.8c$};\draw[->,thick](0,0)--(0,2.4)node[midway,left]{$0.6c$};\draw[dashed,->](0,2.4)--(3.2,0);\node[below left]at(0,0){$O$};}
\end{example}
\begin{solution}
图\ref{fig:dep-m1-relative-q}中两速度点积为零，$\gamma_A=5/3$、$\gamma_B=5/4$，所以
\[
 \gamma_{\rm rel}=\frac{25}{12},\qquad
 w=\frac{\sqrt{481}}{25}c\simeq0.8773c.
\]
变到沿$+y$以$0.6c$运动的系，$A$原来的$u_y=0$，速度变换的公共分母为$1-u_yV/c^2=1$。纵向分量变为$u_y'=-0.6c$，横向为$u_x'=u_x/\gamma_B=0.64c$。
\ExpansionDiagram{fig:dep-m1-relative-a}{在$B$静止系中，$A$的横向分量也发生变化}{
\draw[->](0,0)--(3.8,0)node[right]{$x'$};\draw[->](0,-2.8)--(0,1)node[above]{$y'$};\draw[->,thick](0,0)--(2.56,-2.4)node[right]{$\vct u'_A$};\draw[dashed](2.56,0)--(2.56,-2.4)--(0,-2.4);\node[above]at(1.3,0){$0.64c$};\node[left]at(0,-1.4){$-0.6c$};}
由图\ref{fig:dep-m1-relative-a}，$w^2/c^2=0.64^2+0.6^2=481/625$，与内积法一致。若直接使用原图两速度差，会得到$c$；错误来自把实验室的同时性和时间间隔原封不动地带入$B$系。
\end{solution}

\subsection{局部能量密度为什么还与压力有关}
光子束的粒子数四流与单个光子的四动量分别变换。设某源系方向余弦为$\mu$的一小束光在源系能量密度为$\dd u_*$；源系相对实验室沿$+x$运动，速度为$\beta c$。光子能量乘$\gamma(1+\beta\mu)$，同一方向粒子束的数密度也乘这个因子，故
\[
 \dd u=\gamma^2(1+\beta\mu)^2\dd u_*.
\]
这里仍以源系方向为积分标签，尚未把积分测度换为实验室立体角。用同一标签积分可避免重复引入角度雅可比。该束在实验室的$x$方向余弦为$(\mu+\beta)/(1+\beta\mu)$，于是其能流为$c\dd u$乘这个余弦。

对源系各向同性辐射，角平均的权重是$\dd\mu/2$。所需角度矩为
\[
 \langle\mu\rangle=\frac12\int_{-1}^1\mu\dd\mu
 =\left[\frac{\mu^2}{4}\right]_{-1}^1=0,
 \qquad
 \langle\mu^2\rangle=\frac12\int_{-1}^1\mu^2\dd\mu
 =\left[\frac{\mu^3}{6}\right]_{-1}^1=\frac13.
\]
将能量因子平方展开为$1+2\beta\mu+\beta^2\mu^2$，能流的角度因子展开为$(1+\beta\mu)(\mu+\beta)$，逐项取上述平均，得到
\[
 u=\gamma^2u_*(1+\beta^2/3),\qquad
 S_x=\frac43\gamma^2\beta c u_*.
\]
再以方向余弦平方求法向动量通量，有$p_x=\gamma^2u_*(1/3+\beta^2)$；横向压力仍为$u_*/3$。辐射运动后不再各向同性，不能继续令所有方向压力都等于实验室$u/3$。

\begin{example}移动辐射腔中的能量与能流
长方形腔在自身静止系装有各向同性平衡辐射，能量密度$u_*$、体积$\mathcal V_*$。腔相对实验室沿$x$以$0.6c$运动。求辐射在实验室的局部能量密度、能流、纵向压力和等时截面总能量。解释为何最后一个量不等于单独把$u_*\mathcal V_*$乘$\gamma$。
\ExpansionDiagram{fig:dep-m1-radiation-q}{腔系辐射各向同性；实验室中的等时截面另有方向权重}{
\draw(0,0)rectangle(4,2);\foreach\a in{0,45,90,135,180,225,270,315}{\draw[->](2,1)--++(\a:.7);}\draw[->,thick](4.4,1)--(6,1)node[right]{$V$};\node[below]at(2,0){腔系体积$\mathcal V_*$};}
\end{example}
\begin{solution}
图\ref{fig:dep-m1-radiation-q}中的源系角平均给$\gamma=5/4$。代入局部公式，
\[
 u=\frac74u_*,\qquad S_x=\frac54cu_*,\qquad
 p_x=\frac{13}{12}u_*.
\]
横向压力为$u_*/3$。能流虽非零，但$S_x/(cu)=5/7<1$；它是不同方向光子贡献的合成，不是说所有光子共同以某个亚光速运动。
\ExpansionDiagram{fig:dep-m1-radiation-a}{压力和能流共同参与运动辐射的局部描述}{
\draw(0,0)rectangle(3.2,1.8);\draw[->,thick](1,.9)--(4.8,.9)node[right]{$S_x$};\draw[->](3.2,1.4)--(4.2,1.4)node[right]{$p_x$};\draw[->](1.6,1.8)--(1.6,2.6)node[above]{$u_*/3$};\node[below]at(1.6,0){$u=7u_*/4$};}
实验室等时体积为$\mathcal V_*/\gamma=4\mathcal V_*/5$，所以辐射总能量$E_{\rm rad}=7u_*\mathcal V_*/5$，而$\gamma u_*\mathcal V_*=5u_*\mathcal V_*/4$。

两者不同并不违反四动量变换。图\ref{fig:dep-m1-radiation-a}中的腔壁持续承受辐射压力，辐射子系统通过壁面交换动量，并非孤立封闭系统。将腔壁的能量、动量和应力一并纳入，封闭整体的总四动量才按四矢量变换。不能把一团受约束辐射的局部密度积分，无条件当作自由粒子的质量能公式。
\end{solution}

\section{习题与解答}
\begin{exercise}\label{ex:m1anisotropicsource}
光源静止系单位固有时、单位立体角辐射能为$I'(\theta')=I_0(1+a\cos\theta')$，$|a|\leq1$。源沿对称轴以$V$运动。求实验室发射时间角功率和总功率，讨论何时总功率大于静止系值。
\Diagram{fig:m1anisotropicq}{源系辐射含沿运动轴的一阶不对称项}{\draw[axis](-2,0)--(3.1,0)node[right]{$x',V$};\draw[curve,domain=0:360,samples=100]plot({1.2*(1+.6*cos(\x))*cos(\x)},{1.2*(1+.6*cos(\x))*sin(\x)});\node[above]at(0,1.8){图取$a>0$};}
\end{exercise}
\begin{solution}
图\ref{fig:m1anisotropicq}总静止功率仍为$P_0=4\pi I_0$，因为$\cos\theta'$全空间积分为零。对同一批光子使用能量、时间与立体角三因子，
\[
 I(\theta)=\frac{\delta^3}{\gamma}I_0
 \left[1+a\frac{\cos\theta-\beta}{1-\beta\cos\theta}\right].
\]
要积分总功率，改在源系按$\mu'=\cos\theta'$积分更直接。每束的实验室能量乘$\gamma(1+\beta\mu')$，除发射时间$\gamma\dd\tau$，所以
\begin{align*}
 P&=2\pi I_0\int_{-1}^1(1+a\mu')(1+\beta\mu')\dd\mu'\\
 &=2\pi I_0\int_{-1}^1[1+(a+\beta)\mu'+a\beta\mu'^2]\dd\mu'\\
 &=2\pi I_0\left[\mu'+\frac{a+\beta}{2}\mu'^2+
 \frac{a\beta}{3}\mu'^3\right]_{-1}^{1}
 =4\pi I_0(1+a\beta/3).
\end{align*}
\Diagram{fig:m1anisotropica}{总功率的额外项来自源系净辐射动量}{\draw[curve](0,0)rectangle(2.7,1);\node at(1.35,.5){$P_0$};\draw[->](2.9,.5)--(3.7,.5);\node[right]at(3.9,.5){$P=P_0+\beta c\dot p'_x$};\node[below]at(2,-.7){$\dot p'_x=aP_0/(3c)$};}
图\ref{fig:m1anisotropica}中，总功率的修正项由源系净辐射动量决定。$a\beta>0$时实验室功率更大，反号时更小，$a=0$才恢复各向同性源的总功率不变。给定发射机制若长期维持不对称辐射，光源会有反冲；本题假定外部装置维持其匀速运动，或只研究一个瞬时发射小过程。
\end{solution}

\begin{exercise}[选做]\label{ex:m1antitelephone}
假设存在在任意发射者静止系均可沿任意方向以同一速率$u>c$发送的信号。地球$A$固定在$x=0$，飞船$B$满足$x=L+Vt$。$A$在$t=0$发信号，$B$收到后立即在自身系以$-u$回复。求回复到达地球时刻，讨论因果矛盾的条件。
\Diagram{fig:m1antiq}{假设的超光速信号是类空路径；此图不表示真实光线}{\draw[axis](-.3,0)--(3.7,0)node[right]{$x/L$};\draw[axis](0,-1.2)--(0,3.4)node[above]{$ct/L$};\draw[curve](1,0)--(3.7,3);\draw[->,very thick](0,0)--(1.818182,.909091);\draw[->,very thick](1.818182,.909091)--(0,-.413223);\node[right]at(1.818182,.909091){$R$};\node[left]at(0,-.413223){$C$};}
\end{exercise}
\begin{solution}
出发信号$x=ut$与飞船相交，图\ref{fig:m1antiq}的接收事件为$t_R=L/(u-V)$、$x_R=uL/(u-V)$。以飞船系中从回复发出起经过的坐标时间$s>0$作参数，回复有$\Delta t'=s$、$\Delta x'=-us$。逆变换逐项给
\[
 \Delta t=\gamma(1-Vu/c^2)s,\qquad
 \Delta x=\gamma(V-u)s.
\]
回复到达地球要求$x_R+\Delta x=0$，所以$s=x_R/[\gamma(u-V)]>0$。代入时间增量，
\[
 t_C=t_R+\frac{(1-Vu/c^2)x_R}{u-V}
 =\frac{L[2u-V(1+u^2/c^2)]}{(u-V)^2}.
\]
当$Vu=c^2$，回复路径在实验室的时间增量为零，仍由以上参数式处理；直接以$\Delta x/\Delta t$表示斜率则会出现零分母。当$Vu>c^2$，沿回复路径的实验室时间增量为负，不能再用面向未来的普通速度图像解释其方向。
\Diagram{fig:m1antia}{若回复早于出发，地球可在发送之前收到由该发送触发的回信}{\draw[axis](0,-1.4)--(0,2)node[above]{地球钟};\fill(0,0)circle(1.5pt)node[right]{$A:t=0$发送};\fill(0,-.9)circle(1.5pt)node[right]{$C:t_C<0$收到回复};\draw[->,thick](-.5,-.8)--(-.5,-.1);}
因此$V>2u/(1+u^2/c^2)$时$t_C<0$。对$u>c$，这个门槛小于$c$，可由亚光速飞船达到。图\ref{fig:m1antiq}取$u=2c,V=.9c$作示意，回复位于地球出发事件的过去。图\ref{fig:m1antia}所示矛盾依赖“每个惯性系都能双向按同一超光速规则发信号”的假设；视超光速图样不具有这种信号能力，不能用它构造本装置。
\end{solution}

\begin{exercise}\label{ex:m1clocktrain}
列车固有长$L$，以速度$V$运动。车头和车尾各有一只在车系同步的钟。地面在$t=0$同时读取两钟，求车头减车尾的读数差，并求地面收到这两次读数的先后是否必然相同。
\Diagram{fig:m1trainq}{同时读取运动列车的两端钟}{\draw[axis](-.5,0)--(5,0)node[right]{$x$};\draw[curve](0,.5)rectangle(3.5,1.3);\node[above]at(0,1.3){尾};\node[above]at(3.5,1.3){头};\draw[axis](3.8,.9)--(4.8,.9)node[above]{$V$};\draw[guide](0,2)--(3.5,2);\node[above]at(1.75,2){$t=0$};}
\end{exercise}
\begin{solution}
图\ref{fig:m1trainq}中的同时截面长为$L/\gamma$。用$\Delta t'=\gamma(\Delta t-V\Delta x/c^2)$，得到头钟减尾钟
\[
 \Delta t'=-\gamma\frac{V}{c^2}\frac{L}{\gamma}=-\frac{VL}{c^2}.
\]
头钟在地面同时截面上落后。读数事件之间是类空间隔，其光信号到达同一观察者的时间还要加各自传播时间。
\Diagram{fig:m1traina}{地面接收者的位置影响读数信号的到达次序}{\draw[curve](0,0)--(4,0);\fill(0,0)circle(1.4pt)node[below]{尾};\fill(4,0)circle(1.4pt)node[below]{头};\fill(1,1.5)circle(1.4pt)node[above]{$D$};\draw[->](0,0)--(1,1.5);\draw[->](4,0)--(1,1.5);\node[below]at(2,-.5){$\Delta t_{\rm arr}=(r_{\text{头}}-r_{\text{尾}})/c$};}
图\ref{fig:m1traina}中的接收者若更靠近车尾，先收到尾部信号；靠近车头则相反；在垂直平分线上同时收到。钟面读数的差与信号到达次序没有一一对应关系。
\end{solution}

\begin{exercise}\label{ex:m1lifetime}
不稳定粒子的平均固有寿命为$\tau_0$，在实验室以$v$匀速飞行。求在距离$L$处仍未衰变的比例。再在粒子系说明为什么得到同一结果。
\Diagram{fig:m1lifeq}{粒子从产生点飞向距离为$L$的探测平面}{\draw[axis](-.3,0)--(5,0)node[right]{$x$};\fill(0,0)circle(2pt)node[below]{$O$};\draw[curve](4,-.7)--(4,1.2)node[above]{探测面};\draw[->,thick](.4,.3)--(2,.3)node[above]{$v$};\draw[<->](0,-.5)--(4,-.5)node[midway,below]{$L$};}
\end{exercise}
\begin{solution}
图\ref{fig:m1lifeq}中实验室飞行时间为$t=L/v$，粒子经历的固有时为$\tau=t/\gamma_v$。恒定衰变概率对应$\dd N/\dd\tau=-N/\tau_0$。分离变量并用$N(0)=N_0$，
\[
 \int_{N_0}^{N(\tau)}\frac{\dd n}{n}
 =-\frac1{\tau_0}\int_0^\tau\dd s,
 \qquad \ln\frac{N(\tau)}{N_0}=-\frac{\tau}{\tau_0}.
\]
取指数并代入飞行固有时，得
\[
 \frac N{N_0}=\exp\left[-\frac{L}{\gamma_vv\tau_0}\right].
\]
\Diagram{fig:m1lifea}{粒子系中探测面向静止粒子运动}{\fill(0,0)circle(2pt)node[above]{粒子};\draw[curve](3,-.7)--(3,1.2);\draw[axis](2.7,.5)--(1.5,.5)node[midway,above]{$v$};\draw[<->](0,-.5)--(3,-.5)node[midway,below]{$L/\gamma_v$};}
在图\ref{fig:m1lifea}的粒子系，产生点到探测面在同一时刻的距离缩短为$L/\gamma_v$，探测面以$v$靠近，所需时间$L/(\gamma_vv)$，恰为上面的固有时间。这里“平均寿命”不代表所有粒子都在同一时刻衰变，结果是统计存活比例。
\end{solution}

\begin{exercise}\label{ex:m1crossvelocity}
飞船$S'$以$V$相对实验室运动，在船上沿$y'$方向发射速度为$u$的粒子。求实验室速度、洛伦兹因子及粒子运动到船系坐标$y'=H$时的实验室位置。
\Diagram{fig:m1crossq}{船内的横向发射}{\draw[curve](0,0)rectangle(3,2.2);\fill(1,0)circle(1.5pt);\draw[axis](1,0)--(1,1.8)node[left]{$u$};\node[right]at(1,2){$H$};\draw[axis](3.4,1)--(4.6,1)node[right]{$V$};}
\end{exercise}
\begin{solution}
图\ref{fig:m1crossq}中$v'_x=0,v'_y=u$。逆速度变换给$v_x=V,v_y=u/\gamma$。于是
\[
 1-v^2/c^2=(1-\beta^2)(1-u^2/c^2),\quad \gamma_v=\gamma\gamma_u.
\]
到达上边界的事件在船系为$t'=H/u,x'=0,y'=H$。逆坐标变换得$t=\gamma H/u$，$x=\gamma VH/u$，$y=H$。
\Diagram{fig:m1crossa}{实验室中粒子轨迹斜向前方，与船保持相同纵向速度}{\draw[axis](-.3,0)--(4.8,0)node[right]{$x$};\draw[axis](0,-.2)--(0,2.5)node[above]{$y$};\draw[very thick,->](0,0)--(3.7,2);\draw[guide](0,2)--(3.7,2)--(3.7,0);\node[left]at(0,2){$H$};\node[below]at(3.7,0){$\gamma VH/u$};}
图\ref{fig:m1crossa}的斜率$H/(\gamma VH/u)=u/(\gamma V)$也等于$v_y/v_x$。若直接使用伽利略合成$v_y=u$，会使坐标变换与轨迹速度互相矛盾。
\end{solution}

\begin{exercise}\label{ex:m1coneenergy}
静止系各向同性的单次闪光总能量为$E_0$，各方向的光谱相同。实验室观察前向锥$\theta\leq\arccos\beta$。求锥内光子数与能量各占整次闪光的比例。
\Diagram{fig:m1coneq}{前向锥恰好由静止系前半球变换而来}{\draw[axis](-.3,0)--(4,0)node[right]{$x$};\draw[curve](0,0)--(3,1.3);\draw[curve](0,0)--(3,-1.3);\draw[guide](3,-1.3)--(3,1.3);\node[right]at(3,1.3){$\cos\theta_c=\beta$};}
\end{exercise}
\begin{solution}
由例\ref{ex:m1aberration}，图\ref{fig:m1coneq}中的光锥来自$0\leq\mu'\leq1$。静止系光子数各向同性，变换不改变光子数，所以数目占比为$1/2$。
\Diagram{fig:m1conea}{在静止系积分能量权重，权重为$\gamma(1+\beta\mu')$}{\draw[axis](0,0)--(3.8,0)node[right]{$\mu'$};\draw[axis](0,0)--(0,2.8)node[above]{能量权重};\draw[curve](0,1)--(3,2.3);\draw[guide](3,0)--(3,2.3);\node[below]at(3,0){$1$};\node[left]at(0,1){$\gamma$};}
能量按图\ref{fig:m1conea}的权重变换，锥内能量为
\[
 E_c=\frac{\gamma E_0}{2}\int_0^1(1+\beta\mu')\dd\mu'
 =\frac{\gamma E_0}{2}\left[\mu'+\frac\beta2\mu'^2\right]_0^1
 =\gamma E_0\left(\frac12+\frac\beta4\right).
\]
整个闪光的能量为$\gamma E_0$，故能量份额为$1/2+\beta/4$。即使高相对论速度，静止系前半球所含光子的实验室能量份额也只趋向$3/4$；这与小立体角内较大的角能量密度并不冲突。
\end{solution}

\begin{exercise}\label{ex:m1transverse}
运动光源的静止频率为$\nu_0$。分别求实验室观察方向与速度垂直、以及光源系发射方向与速度垂直时的实验室频率。
\Diagram{fig:m1transq}{“横向”可指两个不同惯性系中的方向}{\draw[axis](0,0)--(3.4,0)node[right]{$V$};\draw[axis](0,0)--(0,2)node[above]{$\theta=\pi/2$};\draw[->,thick](0,0)--(2.6,1.6)node[right]{$\theta'=\pi/2$的光};}
\end{exercise}
\begin{solution}
图\ref{fig:m1transq}中竖直的实验室光线有$\cos\theta=0$，所以$\nu=\nu_0/\gamma$，发生红移。若取源系$\cos\theta'=0$，则由逆能量变换$\nu=\gamma\nu_0$，发生蓝移。
\Diagram{fig:m1transa}{两个横向条件对应不同实验室角度}{\draw[axis](0,0)--(3.5,0)node[right]{$x$};\draw[curve](0,0)--(0,2);\draw[curve](0,0)--(2.6,1.6);\node[left]at(0,1.7){$\nu_0/\gamma$};\node[right]at(2.6,1.6){$\gamma\nu_0$};\draw[->](.9,0)arc(0:31.6:.9);\node[above]at(1.3,1.1){$\cos\theta=\beta$};}
第二种光线在实验室满足$\cos\theta=\beta$，如图\ref{fig:m1transa}。将其代入$\nu=\nu_0/[\gamma(1-\beta\cos\theta)]$，分母成为$\gamma(1-\beta^2)=1/\gamma$，与直接变换吻合。叙述横向多普勒效应时必须说明横向角在哪个系定义。
\end{solution}

\begin{exercise}\label{ex:m1fourwork}
静质量$m$的粒子由静止开始受恒定三力$f$沿$x$加速。求其动量、能量、速度和位移随实验室时间的变化，以及经历的固有时。
\Diagram{fig:m1workq}{实验室恒定三力的直线运动}{\draw[axis](-.3,0)--(4.7,0)node[right]{$x$};\fill(0,0)circle(2pt)node[below]{$m$};\draw[very thick,->](.2,.5)--(2,.5)node[midway,above]{$f$};}
\end{exercise}
\begin{solution}
以图\ref{fig:m1workq}的出发事件为$t=x=0$。恒三力满足$\dd p/\dd t=f$；用$p(0)=0$，有$\int_0^{p(t)}\dd p_1=f\int_0^t\dd t_1$，即$p=ft$。质壳关系和$v=pc^2/E$给
\[
 E=mc^2\sqrt{1+s^2},\qquad
 v=\frac{cs}{\sqrt{1+s^2}},\qquad s=\frac{ft}{mc}.
\]
位移的定积分为
\begin{align*}
 x(t)&=\int_0^t v(t_1)\dd t_1
 =\frac{mc^2}f\int_0^s\frac{q\dd q}{\sqrt{1+q^2}}\\
 &=\frac{mc^2}f[\sqrt{1+q^2}]_0^s
 =\frac{mc^2}f[\sqrt{1+s^2}-1].
\end{align*}
于是$fx=E-mc^2$，外力做功恰等于动能增加。
\Diagram{fig:m1worka}{恒力世界线渐近于光线，虚线斜率为1}{\draw[axis](-.2,0)--(4.5,0)node[right]{$fx/(mc^2)$};\draw[axis](0,-.2)--(0,4.5)node[above]{$ft/(mc)$};\draw[curve,domain=0:4,samples=80]plot({sqrt(1+\x*\x)-1},\x);\draw[guide](0,1)--(3.4,4.4);}
固有时间应沿图\ref{fig:m1worka}的加速世界线累计。由$\gamma_v=\sqrt{1+s^2}$，
\[
 \tau=\frac{mc}f\int_0^s\frac{\dd q}{\sqrt{1+q^2}}
 =\frac{mc}f\operatorname{arsinh}s.
\]
最后一步令$q=\sinh\eta$，则$\dd q=\cosh\eta\dd\eta$、$\sqrt{1+q^2}=\cosh\eta$，上下限为$0,\operatorname{arsinh}s$，积分为$[\eta]_0^{\operatorname{arsinh}s}$。反解$s=\sinh(f\tau/mc)$，代回得到
\[
 ct=\frac{mc^2}f\sinh\frac{f\tau}{mc},\qquad
 x=\frac{mc^2}f\left(\cosh\frac{f\tau}{mc}-1\right).
\]
从$v(t)=cs(1+s^2)^{-1/2}$用乘积与链式法则求导，
\[
 \frac{\dd v}{\dd t}=c\frac{\dd s}{\dd t}
 \left[(1+s^2)^{-1/2}-s^2(1+s^2)^{-3/2}\right]
 =\frac fm(1+s^2)^{-3/2}.
\]
纵向固有加速度为$\gamma_v^3\dd v/\dd t=f/m$。恒三力对应恒纵向固有加速度，实验室加速度则不断减小。$s\ll1$时$x\simeq ft^2/(2m)$、$\tau\simeq t$，任何有限$t$都有$v<c$。
\end{solution}
\begin{exercise}\label{ex:m1twoemitters}
静止系中两灯位于$\pm L/2$，同时向中点发光。求两束光相遇事件在运动系的坐标，以及两发射事件在运动系是否同时。
\Diagram{fig:m1emitq}{两灯同时发射相向光线}{\draw[axis](-2.3,0)--(2.5,0);\fill(-2,0)circle(1.5pt)node[below]{$A$};\fill(2,0)circle(1.5pt)node[below]{$B$};\draw[->](-1.8,.5)--(-.2,.5);\draw[->](1.8,.5)--(.2,.5);}
\end{exercise}
\begin{solution}
在图\ref{fig:m1emitq}的$S$系中，左右发射事件分别为$A:(0,-L/2)$、$B:(0,L/2)$，坐标次序为$(ct,x)$。两束光到中点均走$L/2$，故相遇于$C:(L/2,0)$。变换至沿$+x$以$V$运动的$S'$，得到
\begin{align*}
 A'&:(\gamma\beta L/2,-\gamma L/2),\\
 B'&:(-\gamma\beta L/2,\gamma L/2),\\
 C'&:(\gamma L/2,-\gamma\beta L/2).
\end{align*}
所以右灯先发光，$t'_B-t'_A=-\gamma\beta L/c$。
\Diagram{fig:m1emita}{运动系中先发射的光走得更久，相遇点仍唯一}{\draw[axis](-2.7,0)--(2.5,0)node[right]{$x'$};\draw[axis](0,-1.7)--(0,3)node[above]{$ct'$};\draw[very thick](-2,1.2)--(-1.2,2)--(2,-1.2);\node[left]at(-2,1.2){$A$};\node[right]at(2,-1.2){$B$};\node[above]at(-1.2,2){$C$};}
图\ref{fig:m1emita}中，左行和右行光仍具有光速。由$A'$到$C'$，
\[
 c(t'_C-t'_A)=\frac{\gamma L}2(1-\beta)
 =x'_C-x'_A;
\]
由$B'$到$C'$，
\[
 c(t'_C-t'_B)=\frac{\gamma L}2(1+\beta)
 =x'_B-x'_C.
\]
两条光路的传播时间不同，正好补偿了两次发射的不同步。相遇是同一时空事件，不会因参考系改变而消失。$x'_B-x'_A=\gamma L$是两个不同时发射事件的位置差；两灯在$S'$的同时距离仍为$L/\gamma$。
\end{solution}

\begin{exercise}\label{ex:m1clockoffset}
两个静止于$S'$的钟相距$L$，钟面相对坐标时间偏移分别为$b_A,b_B$。实验室在$t=0$同时读取，求读数差；什么偏移能使这次读数相同？
\Diagram{fig:m1offsetq}{钟面读数含可人为设置的常数偏移}{\draw[curve](0,0)circle(.5);\draw[curve](4,0)circle(.5);\node at(0,0){$A$};\node at(4,0){$B$};\node[above]at(0,.7){$t'+b_A$};\node[above]at(4,.7){$t'+b_B$};\draw[axis](1,-1)--(2.5,-1)node[right]{$V$};}
\end{exercise}
\begin{solution}
在图\ref{fig:m1offsetq}的实验室同时截面，两钟坐标时间差$\Delta t'=-VL/c^2$，因此钟面差$T_B-T_A=b_B-b_A-VL/c^2$。使本次读数相同需$b_B-b_A=VL/c^2$。
\Diagram{fig:m1offseta}{两钟速率相同，偏移只改变读数的常数差}{\draw[axis](0,0)--(4.5,0)node[right]{$t$};\draw[axis](0,0)--(0,3.6)node[above]{读数};\draw[curve](0,.4)--(3.8,2.8);\draw[guide](0,1)--(3.8,3.4);\node[right]at(3.8,2.8){固定偏移};}
图\ref{fig:m1offseta}中两钟在实验室的速率均为$1/\gamma$，所以在每个实验室同时截面上，该读数差都相同；这个偏移相当于按实验室同时性而非车系爱因斯坦同时性设置钟。钟的正常走时速率与一组钟的同步约定是不同问题。
\end{solution}

\begin{exercise}\label{ex:m1propersegments}
粒子先以$v_1$运动实验室时间$T_1$，再以$v_2$运动$T_2$，方向任意，变速时间忽略。求总固有时；若已知两端点坐标差，何时可以用端点间隔代替？
\Diagram{fig:m1segmentsq}{折线世界线的两段固有时须相加}{\STAxes\draw[very thick](0,0)--(2,2.5)--(1,4);\draw[guide](0,0)--(1,4);\node[right]at(2,2.5){$B$};\node[right]at(1,4){$C$};}
\end{exercise}
\begin{solution}
图\ref{fig:m1segmentsq}每段匀速，故$\tau=T_1\sqrt{1-v_1^2/c^2}+T_2\sqrt{1-v_2^2/c^2}$。端点间隔给的是直接连接两端的匀速世界线固有时$\tau_{AC}=\sqrt{(T_1+T_2)^2-|\vct v_1T_1+\vct v_2T_2|^2/c^2}$。
\Diagram{fig:m1segmentsa}{在端点共同静止系中，直线停留的钟积累时间最多}{\draw[axis](-1.5,0)--(2,0)node[right]{$x_*$};\draw[axis](0,0)--(0,4.3)node[above]{$ct_*$};\draw[very thick](0,0)--(1.3,2.1)--(0,3.8);\draw[guide](0,0)--(0,3.8);}
在图\ref{fig:m1segmentsa}的端点静止系，实际轨迹$\int\sqrt{1-v_*^2/c^2}\dd t_*\leq\int\dd t_*=\tau_{AC}$，等号要求全程$v_*=0$，即两段本来就是同一匀速速度矢量。仅两段速率相同而方向改变，仍不能用端点间隔代替曲线积分。
\end{solution}

\begin{exercise}\label{ex:m1inversevel}
实验室测得粒子$v_x=.8c,v_y=.3c$。求使其纵向速度为零的$x$向惯性系，以及该系中的横向速度和洛伦兹因子。
\Diagram{fig:m1inverseq}{选择变换速度消去纵向分量}{\draw[axis](0,0)--(4,0)node[right]{$x$};\draw[axis](0,0)--(0,2.3)node[above]{$y$};\draw[->,thick](0,0)--(3.2,1.2);\draw[guide](3.2,0)--(3.2,1.2);}
\end{exercise}
\begin{solution}
图\ref{fig:m1inverseq}中，设所求参考系沿$x$以$V$运动。速度变换的分母$1-Vv_x/c^2>0$，所以$v'_x=0$等价于$V=v_x=.8c$，其洛伦兹因子为$\gamma=5/3$。横向速度为
\[
 v'_y=\frac{v_y}{\gamma(1-Vv_x/c^2)}
 =\frac{.3c}{(5/3)(1-.64)}=.5c.
\]
因此新系中粒子仅沿$y'$运动，$\gamma_{v'}=(1-.25)^{-1/2}=2/\sqrt3$。
\Diagram{fig:m1inversea}{新系中轨迹纯横向，速率仍小于$c$}{\draw[axis](0,0)--(3.7,0)node[right]{$x'$};\draw[axis](0,0)--(0,2.7)node[above]{$y'$};\draw[very thick,->](0,0)--(0,2.1)node[right]{$.5c$};}
横向速度增大的原因可由微小事件差看出：$\dd y'=\dd y$，而
\[
 \dd t'=\gamma(1-Vv_x/c^2)\dd t=.6\dd t.
\]
相同横向位移除以更短的新系坐标时间，得到图\ref{fig:m1inversea}的$.5c$，并非横向位移本身被放大。四速度提供独立核对：原速率平方$.73c^2$，故
\[
 \gamma_v=\frac1{\sqrt{.27}},\qquad
 \gamma_{v'}=\gamma\gamma_v(1-.64)
 =\frac{.6}{\sqrt{.27}}=\frac2{\sqrt3}.
\]
只消去纵向速度仍未使粒子静止。若再在$S'$中沿$y'$以$.5c$变换，才到达粒子的静止系；两次非共线变换的空间轴关系须另行确定。
\end{solution}

\begin{exercise}\label{ex:m1numberangular}
各向同性光源在固有时内每秒发射$\dot N_0$个光子。求实验室按发射时间定义的角光子率，积分求总光子率，并与总功率不变比较。
\Diagram{fig:m1numberq}{同一批光子的数目不变，角密度与时间率改变}{\draw[axis](0,0)--(3.8,0)node[right]{$V$};\draw[curve](0,0)--(3,1);\draw[curve](0,0)--(3,1.5);\node[right]at(3,1.5){$\dd\Omega$};}
\end{exercise}
\begin{solution}
图\ref{fig:m1numberq}中，源系在$\dd\tau$内向$\dd\Omega'$发出的光子数为$\dd N=\dot N_0\dd\tau\dd\Omega'/(4\pi)$。令$\delta=[\gamma(1-\beta\cos\theta)]^{-1}$。由光行差，
\[
 \frac{\dd\cos\theta}{\dd\cos\theta'}
 =\frac1{\gamma^2(1+\beta\cos\theta')^2}=\delta^{-2},
\]
方位角不变，所以$\dd\Omega=\delta^{-2}\dd\Omega'$。再用实验室发射时间$\dd t=\gamma\dd\tau$，得到
\[
 \frac{\dd\dot N}{\dd\Omega}=\frac{\dot N_0}{4\pi\gamma}\delta^2.
\]
\Diagram{fig:m1numbera}{单次光子数不变，时间膨胀降低实验室总发射率}{\draw[axis](0,0)--(4.2,0)node[right]{$\beta$};\draw[axis](0,0)--(0,2.7)node[above]{$\dot N/\dot N_0$};\draw[curve,domain=0:1,samples=70]plot({3.5*\x},{2.3*sqrt(1-\x*\x)});\node[left]at(0,2.3){$1$};}
全角积分可以直接计算。以$\mu=\cos\theta$，
\begin{align*}
 \int\delta^2\dd\Omega
 &=\frac{2\pi}{\gamma^2}\int_{-1}^1\frac{\dd\mu}{(1-\beta\mu)^2}\\
 &=\frac{2\pi}{\gamma^2}
 \left[\frac{1}{\beta(1-\beta\mu)}\right]_{-1}^1\\
 &=\frac{2\pi}{\beta\gamma^2}
 \left[\frac1{1-\beta}-\frac1{1+\beta}\right]=4\pi.
\end{align*}
在$\beta=0$处取连续极限。因此图\ref{fig:m1numbera}给出总光子率$\dot N=\dot N_0/\gamma$。若源系每光子能量为$E_0$，同一批各向同性光子的总动量为零，整体能量变换便给平均能量$\langle E\rangle=\gamma E_0$。于是$P=\dot N\langle E\rangle=\dot N_0E_0$。按实验室发射时间计，光子数率因时间膨胀减小，平均单光子能量增加，两者相抵。接收时间率还须计入传播延迟。
\end{solution}

\begin{exercise}\label{ex:m1redblue}
单色运动源静止频率$\nu_0$。求实验室观察到无频移的方向，并确定蓝移、红移的角域。
\Diagram{fig:m1redq}{多普勒因子随实验室观察角变化}{\draw[axis](0,0)--(4,0)node[right]{$V$};\draw[curve](0,0)--(2.3,1.6);\draw[curve](0,0)--(-1,1.6);\draw[guide](0,0)--(0,2.2);\node at(1,.35){$\theta$};}
\end{exercise}
\begin{solution}
图\ref{fig:m1redq}中，$0<\beta<1$时$\nu/\nu_0=[\gamma(1-\beta\cos\theta)]^{-1}$。令此比为1，得到
\[
 \cos\theta_* =\frac{1-1/\gamma}{\beta}
 =\frac{\gamma-1}{\gamma\beta}.
\]
利用$(\gamma-1)(\gamma+1)=\gamma^2\beta^2$，可改写为$\cos\theta_*=\gamma\beta/(\gamma+1)$。该值介于0和1之间，对应唯一的前向无频移圆锥。
\Diagram{fig:m1reda}{无频移锥内蓝移，锥外红移}{\draw[axis](-1,0)--(4,0)node[right]{$x$};\draw[curve](0,0)--(1.8,2);\draw[curve](0,0)--(1.8,-2);\node[right]at(2,.5){蓝移};\node[left]at(-.2,1.5){红移};\node[above]at(1.8,2){$\theta_*$};}
图\ref{fig:m1reda}的锥内满足$\cos\theta>\cos\theta_*$，多普勒分母小于1，因而蓝移；锥外红移。将边界方向反变换到源系，有
\[
 \cos\theta'_* =\frac{\cos\theta_*-\beta}{1-\beta\cos\theta_*}
 =-\frac{\gamma\beta}{\gamma+1}=-\cos\theta_*.
\]
源系略向后的发射方向也可能在实验室无频移。低速展开$\gamma=1+\beta^2/2+\cdots$给$\cos\theta_*\simeq\beta/2$，边界略小于$\pi/2$；高速时$\cos\theta_*\to1$。当$\beta=0$时所有方向都无频移，不能把$\beta\to0$所得的边界锥误认为静止光源只有一个无频移方向。
\end{solution}



\chapter{时空过程、相对论火箭与延迟成像}
\section{坐标原点与仿射变换}\label{th:m2affine}
洛伦兹变换作用于事件差。若两系并不把同一个事件取为零点，坐标关系应写为$x'^\mu=\Lambda^\mu{}_\nu x^\nu+C^\mu$。常数$C^\mu$由一件已知事件的坐标确定。平移不改变速度、微分间隔和固有时间，却会改变某个事件的坐标读数。

例如一只钟在事件$B=(ct_B,x_B)$换为速度$-V$匀速运动，并规定它在$B$显示$\tau_B$。以该钟为$x_2=0$原点的返程系可以定义为
\begin{align}\label{eq:m2return}
 ct_2&=c\tau_B+\gamma[(ct-ct_B)+\beta(x-x_B)],\\
 x_2&=\gamma[(x-x_B)+\beta(ct-ct_B)].\notag
\end{align}
把$B$代入，得到$(ct_2,x_2)=(c\tau_B,0)$；令$x_2=0$，又得到$x-x_B=-V(t-t_B)$，所以原点世界线确实沿返程运动。这两个代入可以避免把平移常数的符号写反。

\begin{example}\label{ex:m2twins}
地球与恒星在$S$中相距$l=10$光年。旅行者由地球以$V=0.8c$飞往恒星，立即换向并以同速返回。忽略转向所占时间。求地球和旅行者经过的时间；建立去程、返程坐标，并求转向前后旅行者判定与转向同时的地球时刻。
\Diagram{fig:m2twinsq}{地球与旅行者的世界线及转向事件}{\draw[axis](-.4,0)--(4,0)node[right]{$x$};\draw[axis](0,-.2)--(0,5.8)node[above]{$ct$};\draw[curve](0,0)--(0,5);\draw[curve](0,0)--(2,2.5)--(0,5);\draw[guide](2,0)--(2,5.2);\node[below]at(2,0){$l$};\fill(2,2.5)circle(1.4pt)node[right]{$B$};\node[left]at(0,5){$C$};\node[below left]at(0,0){$A$};}
\end{example}
\begin{solution}
图\ref{fig:m2twinsq}中$B$的坐标为$(cl/V,l)$，$C$的坐标为$(2cl/V,0)$。地球时间$T=2l/V=25$年。每段旅行者时间为$l/(\gamma V)$，$\gamma=5/3$，故总固有时间15年，去返各7.5年。

去程系原点由$A$出发，使用$ct_1=\gamma(ct-\beta x)$、$x_1=\gamma(x-\beta ct)$。在$B$，$ct_{1B}=l/(\gamma\beta)$。返程原点必须经过$B$且沿$-V$运动，钟面连续，故式\eqref{eq:m2return}取$ct_B=l/\beta$、$x_B=l$、$c\tau_B=l/(\gamma\beta)$，整理得
\[
 ct_2=\gamma(ct+\beta x)-2\gamma\beta l,
 \qquad x_2=\gamma(x+\beta ct)-2\gamma l.
\]
逆变换为
\[
 ct=\gamma(ct_2-\beta x_2),\qquad
 x=\gamma(x_2-\beta ct_2)+2l.
\]
\Diagram{fig:m2twinsa}{转向事件处的两条同时线与地球世界线相交于$D_-,D_+$}{\draw[axis](-.2,0)--(4,0)node[right]{$x$};\draw[axis](0,-.2)--(0,5.7)node[above]{$ct$};\draw[curve](0,0)--(2,2.5)--(0,5);\draw[guide](0,.9)--(3.1,3.38);\draw[guide](0,4.1)--(3.1,1.62);\fill(2,2.5)circle(1.4pt)node[right]{$B$};\node[left]at(0,.9){$D_-$};\node[left]at(0,4.1){$D_+$};}
去程同时线满足$ct-\beta x=ct_B-\beta l$，与地球$x=0$相交得
\[
 t_- =\frac{l}{\beta c}-\frac{\beta l}c=4.5\text{年}.
\]
返程同时线满足$ct+\beta x=ct_B+\beta l$，给$t_+=l/(\beta c)+\beta l/c=20.5$年。图\ref{fig:m2twinsa}中两交点相差$2\beta l/c=16$年。这是改变同时性定义后对远处地球时刻的重新赋值，地球钟本身连续走时，旅行者在转向处也没有瞬间看到地球老16年。
\end{solution}

\begin{example}\label{ex:m2fixed}
在例\ref{ex:m2twins}的坐标约定下，求同时满足$ct_2=ct$、$x_2=x$的事件。说明它是否位于旅行者转向处。
\Diagram{fig:m2fixedq}{返程系的坐标网原点世界线经过$B$}{\draw[axis](-.3,0)--(4,0)node[right]{$x$};\draw[axis](0,-.2)--(0,5.6)node[above]{$ct$};\draw[curve](0,0)--(2,2.5)--(0,5);\node[right]at(2,2.5){$B$};\draw[guide](2,0)--(2,5.2);}
\end{example}
\begin{solution}
将相同坐标条件代入逆变换，得到
\[
 (\gamma-1)ct=\gamma\beta x,
 \qquad(\gamma-1)x-\gamma\beta ct=-2l.
\]
第一式给$ct=\gamma\beta x/(\gamma-1)$。代入第二式后，$x$的系数为
\[
 (\gamma-1)-\frac{\gamma^2\beta^2}{\gamma-1}
 =(\gamma-1)-(\gamma+1)=-2,
\]
所以$x=l$，$ct=l(\gamma+1)/(\gamma\beta)$。取本题参数得到$ct=2l$，即$t=20$年。
\Diagram{fig:m2fixeda}{坐标不动事件$F$与转向事件$B$不同}{\draw[axis](-.2,0)--(3.8,0)node[right]{$x/l$};\draw[axis](0,-.2)--(0,4.6)node[above]{$ct/l$};\draw[curve](0,0)--(1.5,1.875)--(0,3.75);\fill(1.5,1.875)circle(1.4pt)node[right]{$B:(1.25,1)$};\fill(1.5,3)circle(1.4pt)node[right]{$F:(2,1)$};\draw[guide](1.5,0)--(1.5,3.8);}
图\ref{fig:m2fixeda}中$B$的地球时刻12.5年，$F$是恒星处更晚的事件。不同坐标系给某件事相同的数字，与该事件处是否有两物体相遇无关；此结果依赖既定的坐标平移约定。
\end{solution}

\section{共线变换与快度}\label{th:m2rapidity}
沿同一方向的洛伦兹矩阵在$(ct,x)$子空间内为
\[
 \Lambda(\beta)=\begin{pmatrix}\gamma&-\gamma\beta\\-\gamma\beta&\gamma\end{pmatrix}.
\]
将两矩阵相乘，其非对角元与对角元之比给
\[
 \beta_{12}=\frac{\beta_1+\beta_2}{1+\beta_1\beta_2},\quad
 \gamma_{12}=\gamma_1\gamma_2(1+\beta_1\beta_2).
\]
定义快度$\chi$满足$\beta=\tanh\chi$，于是$\gamma=\cosh\chi$、$\gamma\beta=\sinh\chi$。双曲函数加法公式说明共线快度直接相加。若使用虚数时间$ict$，形式上可把双曲变换写成虚角转动；这里的角度只是代数类比，不表示在普通空间中转动物体。

对一维运动，$u^\mu=c(\cosh\chi,\sinh\chi)$，故
\[
 a^\mu=c\frac{\dd\chi}{\dd\tau}(\sinh\chi,\cosh\chi),\quad
 \alpha=\sqrt{-a^2}=c\left|\frac{\dd\chi}{\dd\tau}\right|.
\]
恒定固有加速度对应快度随固有时线性增加。实验室速度则按双曲正切变化。

\begin{example}\label{ex:m2boosts}
飞船相对地球以$0.6c$前进，在船上又沿前方发射相对船速$0.6c$的探测器。求探测器相对地球的速度、快度和洛伦兹因子。
\Diagram{fig:m2boostsq}{两次同向速度变换}{\draw[axis](-.3,0)--(6,0)node[right]{$x$};\draw[curve](0,.3)rectangle(2,1);\node[above]at(1,1){飞船};\draw[axis](2.2,.65)--(3.4,.65)node[above]{$0.6c$};\fill(4.3,.65)circle(2pt);\draw[axis](4.4,.65)--(5.6,.65)node[above]{$0.6c$船系};}
\end{example}
\begin{solution}
图\ref{fig:m2boostsq}的两次速度属于不同参考系。令飞船和探测器在各自上一级参考系中的速度为$u_1=u_2=.6c$，实验室速度为
\[
 \beta=\frac{.6+.6}{1+.6^2}=\frac{15}{17}.
\]
每次变换的快度满足$\ee^{2\chi_1}=(1+.6)/(1-.6)=4$，故$\chi_1=\ln2$，两次同轴变换后$\chi=2\ln2=\ln4$。
\Diagram{fig:m2boostsa}{能量—动量双曲线上的快度相加}{\draw[axis](0,0)--(4.7,0)node[right]{$p/(mc)$};\draw[axis](0,0)--(0,4.8)node[above]{$E/(mc^2)$};\draw[curve,domain=0:3.8,samples=70]plot(\x,{sqrt(1+\x*\x)});\fill(.75,1.25)circle(1.5pt)node[left]{$\chi_1$};\fill(1.875,2.125)circle(1.5pt)node[above left]{$2\chi_1$};}
图\ref{fig:m2boostsa}采用$(p/(mc),E/(mc^2))$坐标。由双曲函数定义，
\[
 \gamma=\cosh\chi=\frac{4+1/4}{2}=\frac{17}{8},\qquad
 \gamma\beta=\sinh\chi=\frac{4-1/4}{2}=\frac{15}{8}.
\]
也可在第一次变换后的四速度上再作一次变换：$(\gamma_1,\gamma_1\beta_1)=(5/4,3/4)$，因此
\[
 \gamma=\frac54\frac54+\frac34\frac34=\frac{17}{8},\quad
 \gamma\beta=2\frac54\frac34=\frac{15}{8}.
\]
两项满足$\gamma^2-(\gamma\beta)^2=1$，商为$15/17$。快度相加以两次共线为前提；非共线速度不能把速度大小各自换成快度后直接相加。
\end{solution}

\section{从守恒定律推导火箭方程}\label{th:m2rocket}
火箭喷气是变质量系统，必须先选封闭的微小过程：燃烧前火箭，与燃烧后的火箭加刚排出的燃气。取燃烧前火箭的瞬时静止系，火箭静质量为$m$；燃烧后为$m+\dd m$，其中$\dd m<0$，获得微小速度$\dd v_*$。排出物沿后方以固定速度$u$运动，其静质量为$\dd\mu>0$。

在一阶精度下，火箭动能是$(\dd v_*)^2$阶，可在能量式中略去；燃气自身的相对论因子$\gamma_u$却不能略去。因此
\begin{align*}
 mc^2&=(m+\dd m)c^2+\gamma_u\dd\mu c^2,\\
 0&=m\dd v_*-\gamma_u u\dd\mu.
\end{align*}
第一式给$\gamma_u\dd\mu=-\dd m$，代入第二式得$m\dd v_*=-u\dd m$。火箭在相邻瞬时静止系间的快度变化为$\dd\chi=\dd v_*/c$，所以$\dd\chi=-(u/c)\dd m/m$。用初值$m=M_0,\chi=0$积分，
\[
 \int_0^\chi\dd q=-\frac uc\int_{M_0}^{m}\frac{\dd\mu}{\mu}
 =-\frac uc[\ln\mu]_{M_0}^{m}.
\]
由此得到
\begin{equation}\label{eq:m2rocket}
 \dd\chi=-\frac uc\frac{\dd m}m,\qquad
 \chi=\frac uc\ln\frac{M_0}m,\qquad
 v=c\tanh\left(\frac uc\ln\frac{M_0}m\right).
\end{equation}
式中$M_0$包含全部初始燃料和飞船。忽略外力、喷流方向共线、$u$在瞬时火箭系恒定，是此式的适用条件。把排出物静质量误写成$-\dd m$，会漏掉释放为喷流动能的质量能。

\Diagram{fig:m2rocketbalance}{瞬时静止系中，喷流和剩余火箭构成守恒系统}{\draw[curve](2,0)rectangle(4,1);\node at(3,.5){$m+\dd m$};\draw[axis](4.2,.5)--(5.5,.5)node[above]{$\dd v_*$};\fill(1,.5)circle(2pt);\draw[axis](.8,.5)--(-.4,.5)node[above]{$u$};\node[below]at(1,.2){$\dd\mu$};}

\begin{example}\label{ex:m2exhaust}
排出物在火箭系速度$u=0.8c$。若飞船静质量减少$1\,\mathrm{kg}$，求排出物的静质量、动能和总能量，并求燃烧前质量为$1000\,\mathrm{kg}$的飞船获得的近似快度增量。
\Diagram{fig:m2exhaustq}{质量能在喷流静能和动能之间分配}{\draw[curve](0,0)rectangle(2.2,1.1);\node at(1.1,.55){$-\dd m=1$ kg};\draw[->](2.4,.55)--(3.5,.55);\draw[curve](3.7,0)rectangle(6.2,1.1);\node at(4.95,.55){喷流总能};}
\end{example}
\begin{solution}
题中的$1\,\mathrm{kg}$损失相对于$1000\,\mathrm{kg}$很小。先按连续喷流的瞬时船系计算。设船的静质量改变$\dd m<0$，排出静质量$\dd\mu>0$。图\ref{fig:m2exhaustq}的一阶能量守恒为
\[
 -\dd m=\gamma_u\dd\mu,\qquad \gamma_u=\frac1{\sqrt{1-.8^2}}=\frac53.
\]
因此整个燃烧过程中排出静质量的和为$.6\,\mathrm{kg}$；把每一微元在其发射瞬间船系测得的能量相加，总能量为$1\,\mathrm{kg}\,c^2$，其中静能$.6\,\mathrm{kg}\,c^2$、动能$.4\,\mathrm{kg}\,c^2$。
\Diagram{fig:m2exhausta}{喷流能量分成静能与动能；两部分之和由质量损失决定}{\draw[curve](0,0)rectangle(5,1);\draw[curve](3,0)--(3,1);\node at(1.5,.5){$0.6c^2$ kg};\node at(4,.5){$0.4c^2$ kg};\node[above]at(1.5,1){静能};\node[above]at(4,1){动能};}
图\ref{fig:m2exhausta}的能量分割采用上述局部计量。喷流在瞬时船系中向后运动，动量大小为$\gamma_u u\dd\mu=-u\dd m$；飞船获得向前的动量$mc\dd\chi$，从而
\begin{align*}
 \dd\chi&=-\frac uc\frac{\dd m}{m},\\
 \Delta\chi&=-.8\int_{1000}^{999}\frac{\dd m}{m}
 =-.8[\ln m]_{1000}^{999}\\
 &=.8\ln\frac{1000}{999}=8.0040\times10^{-4}.
\end{align*}
若只保留损失比例的一阶，便为$8.0\times10^{-4}$。有限过程在初始实验室系的总喷流能量略有不同：末船能量为$999\,\mathrm{kg}\,c^2\cosh\Delta\chi$，所以
\[
 E_{\rm jet}^{\rm lab}
 =[1000-999\cosh\Delta\chi]\,\mathrm{kg}\,c^2
 \simeq.999680\,\mathrm{kg}\,c^2.
\]
约$.000320\,\mathrm{kg}\,c^2$成为船的反冲动能。各次发射的瞬时船系并不相同，局部能量相加后仍须区别于固定实验室系能量。喷流静质量的和为$.6\,\mathrm{kg}$。
\end{solution}

\begin{example}\label{ex:m2fuel}
比较两种火箭由静止达到$0.01c$所需的初末静质量比：一种喷气速度$3000\,\mathrm{m/s}$，另一种把排出能量全部变为向后的光。取$c=3.0\times10^8\,\mathrm{m/s}$。
\Diagram{fig:m2fuelq}{相同末速度下，喷气和光子排气的动量携带能力不同}{\draw[curve](0,0)rectangle(2,1);\draw[axis](-.1,.5)--(-1.3,.5)node[above]{$u$};\draw[curve](4,0)rectangle(6,1);\draw[decorate,decoration={snake,amplitude=1.2pt,segment length=5pt}](3.9,.5)--(2.7,.5);\node[above]at(3.3,.7){光};}
\end{example}
\begin{solution}
所需快度$\chi=\operatorname{artanh}(.01)=\tfrac12\ln(1.01/.99)\simeq0.01000033$。图\ref{fig:m2fuelq}左侧喷气$u/c=10^{-5}$，所以
\[
 \ln\frac{M_0}{M_f}=\frac c u\chi\simeq1000.033,
 \quad \log_{10}\frac{M_0}{M_f}\simeq434.31.
\]
质量比约为$2.0\times10^{434}$，这个模型中的困难来自低喷气速度。不能把$0.01c$称为足够低速就忽略连续变质量过程中质量比的指数。
\Diagram{fig:m2fuela}{光子火箭初末质量和快度的关系}{\draw[axis](0,0)--(4.7,0)node[right]{$\chi$};\draw[axis](0,0)--(0,3.4)node[above]{$M_0/M_f$};\draw[curve,domain=0:1.05,samples=50]plot({3*\x},{exp(\x)});\node[left]at(0,1){$1$};\node[right]at(3.1,2.8){$\ee^\chi$};}
光子火箭的排出物没有静质量，应直接以能量$\dd E_\gamma=-c^2\dd m$、动量$\dd E_\gamma/c$写守恒，得到式\eqref{eq:m2rocket}的$u=c$极限。图\ref{fig:m2fuela}给
\[
 \frac{M_0}{M_f}=\ee^\chi=\sqrt{\frac{1+\beta}{1-\beta}}\simeq1.01005.
\]
这是假定完全定向、完全质量能转化的理想模型，不包含推进装置的质量和效率限制。
\end{solution}

\begin{example}\label{ex:m2burn}
理想光子火箭从实验室原点由静止出发，固有质量损失率恒定为$-\dd m/\dd\tau=a>0$。求$m(\tau)$、$v(\tau)$、世界线和固有加速度。初质量为$M_0$。
\Diagram{fig:m2burnq}{质量在固有时间内线性减少}{\draw[axis](0,0)--(4.5,0)node[right]{$\tau$};\draw[axis](0,0)--(0,3)node[above]{$m$};\draw[curve](0,2.5)--(3.7,0);\node[left]at(0,2.5){$M_0$};\node[below]at(3.7,0){$M_0/a$};}
\end{example}
\begin{solution}
图\ref{fig:m2burnq}中$m=M_0-a\tau$，定义$s=m/M_0=1-a\tau/M_0$。物理区间为$0<s\leq1$，实际还须在干质量处停机。由光子火箭式，$\chi=\ln(M_0/m)=-\ln s$，所以
\[
 \gamma=\tfrac12(s^{-1}+s),\quad \gamma\beta=\tfrac12(s^{-1}-s),\quad
 \beta=\frac{1-s^2}{1+s^2}.
\]
世界线对固有时的导数为$\dd(ct)/\dd\tau=c\gamma$、$\dd x/\dd\tau=c\gamma\beta$。分别从0积分至$\tau$，利用
\[
 \int_0^\tau\frac{\dd q}{1-aq/M_0}
 =\left[-\frac{M_0}a\ln(1-aq/M_0)\right]_0^\tau
 =-\frac{M_0}a\ln s,
\]
\[
 \int_0^\tau(1-aq/M_0)\dd q
 =\left[q-\frac{aq^2}{2M_0}\right]_0^\tau
 =\tau-\frac{a\tau^2}{2M_0},
\]
得到
\begin{align*}
 ct&=\frac c2\left[-\frac{M_0}a\ln s+\tau-\frac{a\tau^2}{2M_0}\right],\\
 x&=\frac c2\left[-\frac{M_0}a\ln s-\tau+\frac{a\tau^2}{2M_0}\right].
\end{align*}
\Diagram{fig:m2burna}{世界线趋近光速；接近质量耗尽时固有加速度增大}{\draw[axis](-.2,0)--(4.8,0)node[right]{$ax/(cM_0)$};\draw[axis](0,-.2)--(0,4.8)node[above]{$at/M_0$};\draw[curve,domain=.02:1,samples=90]plot({(-ln(\x)-(1-\x*\x)/2)/2},{(-ln(\x)+(1-\x*\x)/2)/2});\draw[guide](0,.5)--(3.7,4.2);}
图\ref{fig:m2burna}的两个无量纲坐标之差为$(1-s^2)/2$，趋向有限值$1/2$，而二者各自对数发散。固有加速度
\[
 \alpha=c\frac{\dd\chi}{\dd\tau}
 =-\frac c s\frac{\dd s}{\dd\tau}=\frac{ca}{m},\qquad
 \frac{\dd v}{\dd t}=\frac{\alpha}{\gamma^3}.
\]
初始$\alpha=ca/M_0$，与喷光动量通量除以初质量吻合。若把$\ee^\chi$写成$1+a\tau/M_0$，就不满足质量线性减少的条件。
\end{solution}

\section{延迟时间与视运动}\label{th:m2retarded}
固定观察者在$\vct r_D$，物体在发光时刻$t_e$位于$\vct r(t_e)$，接收时刻满足
\begin{equation}\label{eq:m2retarded}
 t_o=t_e+\frac{|\vct r_D-\vct r(t_e)|}{c}.
\end{equation}
求照片要固定$t_o$解每个发光点的$t_e$，再把光线投影到成像面。求物体长度要固定$t_e$，两种操作不可互换。若观察距离远大于物体尺度，记朝向观察者的单位矢量为$\vct n$，则$\dd t_o=(1-\vct n\cdot\vct v/c)\dd t_e$。这里光程压缩产生的快慢属于接收过程，不改变局部运动速度。

\begin{example}\label{ex:m2super}
发光斑点以$\beta c$匀速运动，与朝向远处观察者的视线夹角为$\theta$。求视横向速度及其最大值；当$\theta=45^\circ$时，求超过$c$的条件。
\Diagram{fig:m2superq}{相邻发光点的横向位移和沿视线位移}{\draw[axis](0,0)--(5.3,0)node[right]{朝观察者};\draw[very thick,->](0,0)--(3,2);\draw[guide](3,0)--(3,2);\node[right]at(3,1){$v\dd t_e\sin\theta$};\node[below]at(1.5,0){$v\dd t_e\cos\theta$};\draw(.8,0)arc(0:33.69:.8);\node at(1.2,.35){$\theta$};}
\end{example}
\begin{solution}
图\ref{fig:m2superq}中第二次发光更靠近观察者，接收时间差为$\dd t_o=(1-\beta\cos\theta)\dd t_e$。将横向位移除以它，得到
\[
 v_{\rm app}=\frac{\beta c\sin\theta}{1-\beta\cos\theta}.
\]
在$45^\circ$时，条件$\beta/\sqrt2>1-\beta/\sqrt2$给$\beta>1/\sqrt2$。
\Diagram{fig:m2supera}{最大视速度的几何条件$\cos\theta=\beta$}{\draw[axis](0,0)--(4.3,0)node[right]{$\cos\theta$};\draw[axis](0,0)--(0,3)node[above]{$v_{\rm app}/c$};\draw[curve,domain=0:1,samples=90]plot({3.5*\x},{.8*sqrt(1-\x*\x)/(1-.8*\x)});\draw[guide](2.8,0)--(2.8,1.333);\node[below]at(2.8,0){$\beta$};\node[above]at(2.8,1.333){$\beta\gamma$};}
求导得
\[
 \frac1c\frac{\dd v_{\rm app}}{\dd\theta}
 =\frac{\beta\cos\theta(1-\beta\cos\theta)-\beta^2\sin^2\theta}
 {(1-\beta\cos\theta)^2}
 =\frac{\beta(\cos\theta-\beta)}{(1-\beta\cos\theta)^2}.
\]
极大满足$\cos\theta=\beta$，代回得$v_{\rm app,max}=\beta\gamma c$，如图\ref{fig:m2supera}。信号本身仍沿光锥传播；视速度是由两次不同位置、不同发光时刻构成的接收投影速度，不能用于发送超光速信号。
\end{solution}

\begin{example}\label{ex:m2photo}
半径$r$的圆截面在自身静止系位于$x'y'$平面，以速度$V$沿$x$运动。远处观察者位于$(0,-b,0)$，$b\gg r$。取快门接收时刻$t_o=b/c$。用参数$\varphi$标记圆周，求同时截面、进入照片的发光位置和横向投影，并解释旋转角。
\Diagram{fig:m2photoq}{运动圆截面及远处沿$-y$方向观察的光线}{\draw[axis](-2,0)--(3,0)node[right]{$x$};\draw[axis](0,-2.7)--(0,2.1)node[above]{$y$};\draw[curve](0,0)ellipse(1.2 and 1.5);\draw[axis](1.5,1.2)--(2.8,1.2)node[above]{$V$};\fill(0,-2.2)circle(2pt)node[right]{$D,\ b\gg r$};\draw[guide](-1,-.8)--(0,-2.2);\draw[guide](1,.8)--(0,-2.2);}
\end{example}
\begin{solution}
圆周在$S'$中是$x'=r\cos\varphi,y'=r\sin\varphi$。在实验室固定$t$，尺缩给$x=Vt+(r/\gamma)\cos\varphi$、$y=r\sin\varphi$，所以同时截面为椭圆。图\ref{fig:m2photoq}中的观察者距发光点约为$b+y$，由式\eqref{eq:m2retarded}及固定$t_o=b/c$得$t_e=-y/c$，从而
\[
 x_e=\frac r\gamma\cos\varphi-\beta r\sin\varphi,
 \qquad y_e=r\sin\varphi.
\]
定义$\cos\alpha=1/\gamma$、$\sin\alpha=\beta$，横向投影为
\[
 x_{\rm photo}=x_e=r\cos(\varphi+\alpha),\qquad \alpha=\arcsin\beta.
\]
\Diagram{fig:m2photoa}{照片横向坐标是两项投影之和；发光点并不同时}{\draw[axis](-2.4,0)--(2.9,0)node[right]{$x_{\rm photo}$};\draw[axis](0,-1.8)--(0,2.2)node[above]{$y$};\draw[curve,domain=0:360,samples=100]plot({1.5*(.8*cos(\x)-.6*sin(\x))},{1.5*sin(\x)});\draw[guide](0,0)--(1.2,0);\node[below]at(1.2,0){$\varphi=0$};\fill(-.9,1.5)circle(1.4pt)node[above left]{$\varphi=\pi/2$};\node[right]at(1.8,1.4){发光位置集合};}
图\ref{fig:m2photoa}画的是各个延迟发光点在实验室的坐标集合，不是照片上的二维圆形。沿$-y$观察时，$y$是视线深度，成像面只保留$x$和$z$。对有标记的圆柱侧面，$x_{\rm photo}$的参数平移等效于把圆柱旋转$\alpha$后投影；它不表示物体真的转动，也不表示任何有限距离、任意形状物体的照片都恰为刚体旋转。忽略的距离二阶项为物体尺度平方除以$b$。
\end{solution}
\section{连续转向与加速参考系}
\subsection{瞬时惯性系的适用范围}
光滑类时世界线在每个事件都有瞬时静止惯性系，可用来定义固有加速度和局部推进功率。若要由这些局部惯性系建立加速坐标，还须指定空间坐标和同步方式。双生子题中赋给远处地球的同时读数可随转向而跃变，旅行者的本地钟却连续走时。

对任意有限加速过程，旅行者固有时始终为$\int\dd t\sqrt{1-v^2/c^2}$，只需已知实验室轨迹。瞬时转向模型是把加速段的持续时间和路程同时趋零的极限，保留两段惯性速度的差。若加速持续时间不可略，去返总时间、最远距离以及固有时都应包括该段积分。

在一维运动中，$\chi=\operatorname{artanh}(v/c)$。链式法则给$\dd\chi/\dd t=\dot v/[c(1-v^2/c^2)]=\gamma_v^2\dot v/c$；再乘$\dd t/\dd\tau=\gamma_v$，得到有向固有加速度$a_*=c\dd\chi/\dd\tau=\gamma_v^3\dot v$，其大小为$\alpha=|a_*|$。制动时$a_*$可为负，而$\alpha$仍非负。随着$\gamma_v$增大，有限的固有加速度可对应越来越小的实验室加速度；题设须说明控制的是哪一个量。

\subsection{恒固有加速度的坐标网}
令参考飞船的固有加速度为$\alpha>0$，设$L_\alpha=c^2/\alpha$。为与飞船保持固定局部距离，可引入坐标$(\tau,\xi)$：
\begin{align*}
 ct&=(L_\alpha+\xi)\sinh(\alpha\tau/c),\\
 x&=(L_\alpha+\xi)\cosh(\alpha\tau/c)-L_\alpha.
\end{align*}
参考飞船在$\xi=0$，此处$\tau$正是其固有时。记$\chi=\alpha\tau/c$，乘积法则给
\begin{align*}
 \dd(ct)&=\sinh\chi\dd\xi
 +(L_\alpha+\xi)\cosh\chi\frac\alpha c\dd\tau,\\
 \dd x&=\cosh\chi\dd\xi
 +(L_\alpha+\xi)\sinh\chi\frac\alpha c\dd\tau.
\end{align*}
相减平方时，混合项相等而抵消；$\dd\xi^2$的系数为$\sinh^2\chi-\cosh^2\chi=-1$，$\dd\tau^2$的系数为$(L_\alpha+\xi)^2\alpha^2/c^2$，得到
\[
 \dd s^2=c^2(1+\alpha\xi/c^2)^2\dd\tau^2-\dd\xi^2.
\]
该坐标网只取$\xi>-L_\alpha$以保持同一时间方向。固定$\xi$的其他钟在同一坐标时间内经历$\dd\tau_\xi=(1+\alpha\xi/c^2)\dd\tau$，其固有加速度为$\alpha_\xi=\alpha/(1+\alpha\xi/c^2)$。要让不同位置的点保持这种固定距离，它们不能全取相同固有加速度。

\Diagram{fig:m2rindler}{固定加速坐标的世界线是共用光状渐近线的双曲线}{\draw[axis](-1.5,0)--(4.3,0)node[right]{$x/L_\alpha$};\draw[axis](0,0)--(0,4.5)node[above]{$ct/L_\alpha$};\foreach\a in{.7,1,1.5}{\draw[curve,domain=0:1.6,samples=70]plot({\a*cosh(\x)-1},{\a*sinh(\x)});}\draw[guide](-1,0)--(3.2,4.2);\node[below]at(-1,0){$-1$};}
图\ref{fig:m2rindler}中越靠前的固定坐标点，所需固有加速度越小。若一根杆的各处在实验室同一时刻开始按完全相同速度程序加速，实验室端点距离可保持不变，但随共同速度增大，其瞬时静止系端点距离会增大。这与保持局部固有距离不变是不同的加速安排，物体内部会承受应力，不能假设任意加速都保留一个绝对刚性形状。

\subsection{加速系光信号的有限传播}
沿$\xi$方向的光满足$\dd s^2=0$，所以坐标速度为$\dd\xi/\dd\tau=\pm c(1+\alpha\xi/c^2)$。局部观察者用自身钟和局部尺测得的光速仍为$c$。取$-L_\alpha<\xi_1<\xi_2$及正号，分离变量：
\[
 \Delta\tau=\frac1c\int_{\xi_1}^{\xi_2}
 \frac{\dd\xi}{1+\alpha\xi/c^2}.
\]
令$z=1+\alpha\xi/c^2$，则$\dd\xi=(c^2/\alpha)\dd z$，上下限为$z_i=1+\alpha\xi_i/c^2>0$，故
\[
 \Delta\tau=\frac c\alpha\int_{z_1}^{z_2}\frac{\dd z}z
 =\frac c\alpha[\ln z]_{z_1}^{z_2}
 =\frac c\alpha\ln\frac{1+\alpha\xi_2/c^2}{1+\alpha\xi_1/c^2}.
\]
若两站间距$\Delta\xi=\xi_2-\xi_1$很小，展开对数得到$\Delta\tau\simeq\Delta\xi/[c(1+\alpha\xi_1/c^2)]$。只有两站还位于$|\alpha\xi/c^2|\ll1$的原点附近时，才进一步近似为$\Delta\xi/c$。靠近$\xi=-L_\alpha$时坐标传播时间发散，对应图中的光状边界；它限制持续加速观察者的信号联系范围。

\section{推进方程的能量和质量含义}
\subsection{喷流速度在哪里定义}
火箭方程中的$u$是喷流在每次燃烧瞬间的飞船静止系测量值。火箭持续变速后，不同时刻喷出的两团燃气在初始实验室中的速度不同；它们无需以某一固定实验室速度向后运动。若题设给的是实验室排气速度，必须先用速度变换换成瞬时$u$，一般会得到随$m$或$\chi$变化的微分方程。

微分能量式$-\dd m=\gamma_u\dd\mu$表示飞船静质量损失中包括喷流静能和喷流动能。这里$m$是整个剩余飞船连同未用燃料的静质量，不等于某个发动机标称干质量。燃料中可转化的化学能很小，实际喷出的静质量大部分只是原料质量的搬运；光子火箭则取极端理想情况，排出物没有静质量。

在燃烧前瞬时静止系，残余飞船获得$\dd v_*$。它的反冲动能约$m\dd v_*^2/2$，相对一阶喷流能是二阶小量，所以微分推导可以忽略。有限脉冲反冲则不能忽略；例\ref{ex:m2finitepulse}保留它之后仍给相同积分质量比，原因是发射光子的实验室能量随着火箭快度连续变化，而非每一步都在同一个静止系中计量。

\subsection{质量比、推进效率与任务边界}
若真实系统只有部分失去的质量能用于定向喷流，可把静止系单位质量损失携带的后向动量记为$\eta c(-\dd m)$，$0\leq\eta\leq1$。四动量守恒给$\dd\chi=-\eta\dd m/m$。固定排气速度的理想模型有$\eta=u/c$，向后纯光束有$\eta=1$，等量前后光束有$\eta=0$。这个定义把发动机效率、废热和喷流分布压缩为一个明确动量效率，但不能仅由“多少热量释放”确定$\eta$。

对每个$\eta>0$的推进段，有$-\dd m/m=|\dd\chi|/\eta$，积分后该段的对数质量比为$\ln(m_i/m_f)=\int|\dd\chi|/\eta$。若$\eta=0$却仍排放质量，$\dd\chi=0$不能反推出质量损失，须按实际消耗另计。若巡航段不损失质量，它只增加旅行时间和距离。制动、返航各有相应推进消耗；抛弃空壳时，还须把抛弃后的剩余质量作为下一段初质量。比较燃料消耗应采用相同末态和载荷条件。

\section{延迟成像的投影几何}
照片中每一个像素对应一条到达观察者的光线。远场近似把这些光线看作平行，成像面取垂直视线的平面。设朝向观察者的单位矢量$\vct n$固定，接收时间可去掉一个共同常数写成$t_o=t_e-\vct n\cdot\vct r(t_e)/c$。固定$t_o$首先确定发射事件，随后投影
\[
 \vct R_{\rm image}=(I-\vct n\vct n^{\mathsf T})\vct r(t_e).
\]
这个投影矩阵删去沿视线深度分量，保留横向两分量。照片只保留横向投影；不同$t_e$的发射点集合也不属于某一实验室同时截面。

\Diagram{fig:m2projection}{先沿光锥选择发射事件，再向垂直视线的平面投影}{\draw[axis](0,0)--(4.5,0)node[right]{$\vct n$};\draw[axis](0,0)--(0,3)node[above]{像面横向};\fill(3,2)circle(1.5pt)node[right]{$\vct r(t_e)$};\draw[guide](3,2)--(0,2);\fill(0,2)circle(1.5pt)node[left]{投影};\draw[->](0,0)--(3,2);}
图\ref{fig:m2projection}中的横向位移除以接收时间间隔，得到视横向速度；实际位移除以发射时间间隔，得到物体的局部运动速度。具体地，$\vct v_{\rm app,\perp}=\vct v_\perp/(1-\vct n\cdot\vct v/c)$。对亚光速源，分母为正，接收顺序与发射顺序一致；视横向速度却可超过$c$，如例\ref{ex:m2super}所示。

圆柱照片中的参数转角也是投影意义上的结果。源系圆周参数$\varphi$只表示表面标记，远场横向坐标中出现$\cos(\varphi+\alpha)$，使某些原本侧面的标记进入可见投影。要判断它是否真的可见，还需考虑不透明物体的自遮挡、前后表面和发光特性；透明发光圆环与不透明圆柱不能共用同一套可见点集合。

有限观察距离时，光线不再平行，距离展开中$R_\perp^2/(2b)$项会改变延迟时刻，透视比例也随深度变化。若只保留一阶远场项，就应把“等效旋转”限定在该精度和相应物形，不把它推广为任意物体照片的普遍刚体旋转定理。

\subsection{变质量推进的局部守恒与全局能量}
火箭方程中的$m$是船体、载荷和全部剩余燃料的静质量。前文已在燃烧前的瞬时船系证明$-\dd m=\gamma_u\dd\mu$、$mc\dd\chi=\gamma_u u\dd\mu$。以下用这两个局部守恒式核对固定实验室中的能量和动量；其中$\dd\mu$是喷流静质量，$u$是其局部后向速度大小。

局部一阶守恒可积成有限的准确火箭方程。计算固定实验室中的总喷流能量时，须先把各瞬时船系的微元变换到同一实验室系：
\[
 \dd E_{\rm jet}=\gamma_v\gamma_u\dd\mu(c^2-vu),\qquad
 \dd p_{\rm jet}=\gamma_v\gamma_u\dd\mu(v-u).
\]
对船的$E=mc^2\cosh\chi$、$p=mc\sinh\chi$用乘积法则，代入$\dd m=-\gamma_u\dd\mu$及$\dd\chi=\gamma_u u\dd\mu/(mc)$，逐项得到
\begin{align*}
 \dd E&=c^2\cosh\chi\,\dd m+mc^2\sinh\chi\,\dd\chi\\
 &=-\gamma_v\gamma_u c^2\dd\mu
   +\gamma_v\gamma_u vu\dd\mu=-\dd E_{\rm jet},\\
 \dd p&=c\sinh\chi\,\dd m+mc\cosh\chi\,\dd\chi\\
 &=-\gamma_v\gamma_u v\dd\mu
   +\gamma_v\gamma_u u\dd\mu=-\dd p_{\rm jet}.
\end{align*}
这里用了$\cosh\chi=\gamma_v$、$c\sinh\chi=\gamma_vv$；两个守恒式中的喷流量与船体量均在同一实验室系计量。
\Diagram{fig:m2jetframes}{各次喷流的局部静止系不同；总能量需统一到同一参考系}{\draw[axis](0,0)--(6.1,0)node[right]{$x$};\foreach\x/\v in{.7/v_1,2.7/v_2,4.7/v_3}{\draw[curve](\x,.3)rectangle({\x+.8},.9);\draw[->]({\x+.9},.6)--({\x+1.35},.6)node[above]{$\v$};\draw[->](\x,.6)--({\x-.5},.6);}\node[below]at(3,0){同一实验室系中的连续发射位置};}
图\ref{fig:m2jetframes}只示意沿途发射，不把几个船体画成同时存在。对有限过程，先求末态船四动量，再从初态总四动量中扣除它，常比逐一积分全部喷流更省算式；两种方法必须给出同一结果。

\section{往返过程中的钟和坐标网}
\begin{example}\label{ex:m2generalturn}
旅行者去程速率$V_1$、返程速率$V_2$，在地球系距离$l$处瞬时换向。求地球和旅行者的总时间，并求换向前后同时地球时刻的跃变。
\Diagram{fig:m2generalq}{不同去返速率使世界线两段斜率不同}{\draw[axis](-.2,0)--(3.8,0)node[right]{$x$};\draw[axis](0,-.2)--(0,5.5)node[above]{$ct$};\draw[curve](0,0)--(2,2.6)--(0,4.7);\node[right]at(2,2.6){$B$};\node[below]at(2,0){$l$};}
\end{example}
\begin{solution}
在图\ref{fig:m2generalq}中取出发事件为$A:(0,0)$，转向为$B:(t_B,l)$。两段均匀速，故
\[
 t_B=\frac l{V_1},\quad T=\frac l{V_1}+\frac l{V_2},\quad
 \tau=\frac l{\gamma_1V_1}+\frac l{\gamma_2V_2}.
\]
瞬时转向不积累有限固有时。求旅行者在$B$赋给地球的同时刻时，须分别采用转向前、后的惯性系。
\Diagram{fig:m2generala}{两段同时线在转向事件处改变倾斜方向}{\draw[axis](-.2,0)--(3.7,0);\draw[axis](0,0)--(0,5.6)node[above]{$ct$};\draw[curve](0,0)--(2,2.6)--(0,4.7);\draw[guide](0,1.061538)--(3,3.369231);\draw[guide](0,4.504762)--(3,1.647619);\node[left]at(0,1.061538){$D_-$};\node[left]at(0,4.504762){$D_+$};}
去程系中与$B$同时的事件满足$\Delta t-(V_1/c^2)\Delta x=0$。令地球事件为$(t_-,0)$，则
\[
 t_--t_B=\frac{V_1}{c^2}(0-l),\qquad
 t_-=t_B-\frac{V_1l}{c^2}.
\]
返程速度为$-V_2$，同理$t_+=t_B+V_2l/c^2$。所以图\ref{fig:m2generala}中$D_+$与$D_-$的时间差为$(V_1+V_2)l/c^2$。可把旅行者赋予地球的时间分成三部分：
\[
 \frac{l}{\gamma_1^2V_1},\qquad
 \frac{(V_1+V_2)l}{c^2},\qquad
 \frac{l}{\gamma_2^2V_2}.
\]
第一项是$A$到$D_-$，末项是$D_+$到重逢；它们与中间的同时性改变量相加，利用$1/\gamma_i^2=1-V_i^2/c^2$恰恢复$T$。地球钟没有跳走，变化的是旅行者用不同同时线选取的远处事件。
\end{solution}

\begin{example}\label{ex:m2outbackopt}
规定地球往返总时间为$T$，旅行者在中途换向一次，并全程以相同速率$v$运动。求最远距离和旅行者时间；若把时间改为两段不同速率，在相同端点条件下哪类路径能积累最多固有时？
\Diagram{fig:m2optq}{固定重逢端点、一次换向的对称路径}{\draw[axis](-.2,0)--(4,0)node[right]{$x$};\draw[axis](0,0)--(0,5.4)node[above]{$ct$};\draw[curve](0,0)--(2,2.5)--(0,5);\draw[guide](0,2.5)--(2,2.5);\node[left]at(0,5){$cT$};}
\end{example}
\begin{solution}
图\ref{fig:m2optq}对称点发生在地球时间$T/2$，所以$l=vT/2$，旅行者时间$\tau=T\sqrt{1-v^2/c^2}$。若只固定出发和重逢事件，完全不离开地球的惯性直线能积累$T$，是最大值。
\Diagram{fig:m2opta}{固定端点的世界线越偏离直线，其固有时越小}{\draw[axis](-.3,0)--(3.7,0);\draw[axis](0,0)--(0,5.4)node[above]{$ct$};\draw[guide](0,0)--(1,2.5)--(0,5);\draw[curve](0,0)--(2,2.5)--(0,5);\draw[very thick](0,0)--(0,5);}
若再规定必须经过远处某个事件$B$，则$A$到$B$、$B$到$C$各自的匀速直线分别使固有时最大，总和也最大；这由每段的积分不等式逐段得到。图\ref{fig:m2opta}中的宽折线和窄折线若转向事件不同，固有时比较应保留相同端点和途经事件的约束；它取决于整条世界线。
\end{solution}

\section{火箭的有限点火与燃料分配}
\begin{example}\label{ex:m2finitepulse}
初始静止、质量$M$的理想飞船一次向后发出能量为$E_\gamma$的光脉冲。$E_\gamma$按初始飞船系测量。发射后飞船静质量$m$。求脉冲能量、末速度，并与连续光子火箭质量比比较。
\Diagram{fig:m2pulseq}{一次有限脉冲过程须保留剩余飞船的反冲能}{\draw[curve](2,0)rectangle(4,1);\node at(3,.5){初始$M$};\draw[->](1.7,.5)--(.1,.5)node[above]{$E_\gamma$};\draw[->](4.2,.5)--(5.5,.5)node[above]{$v$};}
\end{example}
\begin{solution}
在图\ref{fig:m2pulseq}的初始船系中，向后光脉冲的动量为$-E_\gamma/c$。发射后飞船向前运动，守恒式为
\[
 Mc^2=\gamma_vmc^2+E_\gamma,\qquad
 \gamma_vmv=E_\gamma/c.
\]
用前式表示末船能量，用后式表示其动量，再作质壳平方：
\begin{align*}
 m^2c^4&=(Mc^2-E_\gamma)^2-E_\gamma^2\\
 &=M^2c^4-2Mc^2E_\gamma.
\end{align*}
故$E_\gamma=(M^2-m^2)c^2/(2M)$，末船能量为$(M^2+m^2)c^2/(2M)$。
\Diagram{fig:m2pulsea}{初始系能量分为脉冲和反冲后飞船能量}{\draw[curve](0,0)rectangle(5,1);\draw[curve](2,0)--(2,1);\node at(1,.5){$E_\gamma$};\node at(3.5,.5){$\gamma_vmc^2$};\node[above]at(2.5,1){$Mc^2$};}
图\ref{fig:m2pulsea}中，末船速率和反冲动能分别为
\[
 \beta=\frac{pc}{E}=\frac{M^2-m^2}{M^2+m^2},\qquad
 K=\frac{(M-m)^2}{2M}c^2.
\]
直接相加$E_\gamma+K=(M-m)c^2$，质量损失的能量去向因而完整闭合。再由速度式计算
\[
 \frac{1+\beta}{1-\beta}=\frac{M^2}{m^2},\qquad
 \chi=\tfrac12\ln\frac{1+\beta}{1-\beta}=\ln\frac Mm.
\]
它与连续定向光子火箭的质量比相同，因为所有排出光子的四动量共线，总和仍为类光。若向多个方向发光，辐射总四动量通常不再类光，这一有限脉冲关系就需要重算。
\end{solution}

\begin{example}\label{ex:m2twosided}
飞船在瞬时静止系发出总光功率$P$，其中向后比例为$f$、向前比例为$1-f$。保持$f$不变，求质量与快度的关系。若各向同性发光，飞船能否加速？
\Diagram{fig:m2twosidedq}{相反方向光功率的差产生净推力}{\draw[curve](2,0)rectangle(4,1);\draw[->](1.8,.5)--(.1,.5)node[above]{$fP$};\draw[->](4.2,.5)--(5.9,.5)node[above]{$(1-f)P$};}
\end{example}
\begin{solution}
图\ref{fig:m2twosidedq}总能量损失给$\dd m/\dd\tau=-P/c^2$。向后减向前的光动量率为$(2f-1)P/c$，所以瞬时推力$F_0=(2f-1)P/c$，快度率$\dd\chi/\dd\tau=F_0/(mc)$。
\Diagram{fig:m2twosideda}{同样质量损失，只有光束的不对称部分产生快度变化}{\draw[axis](0,0)--(4.3,0)node[right]{$f$};\draw[axis](0,-2)--(0,2.2)node[above]{$\chi/\ln(M_0/m)$};\draw[curve](0,-1.6)--(3.6,1.6);\node[below]at(1.8,0){$1/2$};\node[below]at(3.6,0){$1$};}
消去$\tau$后$\dd\chi=-(2f-1)\dd m/m$。以初始静止、质量$M_0$为初值，
\[
 \int_0^\chi\dd q=-(2f-1)\int_{M_0}^{m}\frac{\dd\mu}{\mu}
 =-(2f-1)[\ln\mu]_{M_0}^{m},
 \quad \chi=(2f-1)\ln\frac{M_0}{m}.
\]
这给出图\ref{fig:m2twosideda}的直线。$f=1/2$时两束相抵，质量下降而速度不变；各向同性辐射也没有静止系净动量。$f<1/2$则向后加速，符号由喷光方向确定。
\end{solution}

\begin{example}\label{ex:m2varyexhaust}
某理想火箭排气速度随剩余质量变化，$u(m)=u_0m/M_0$，$u_0<c$。由初始静止加速至剩余质量$M_f$，求末快度。与把排气速度粗略取末值的结果比较。
\Diagram{fig:m2varyq}{排气速度随质量减少而下降}{\draw[axis](0,0)--(4.3,0)node[right]{$m/M_0$};\draw[axis](0,0)--(0,2.8)node[above]{$u/u_0$};\draw[curve](0,0)--(3.5,2.3);\node[below]at(3.5,0){$1$};\node[left]at(0,2.3){$1$};}
\end{example}
\begin{solution}
式\eqref{eq:m2rocket}仍可逐微元使用。图\ref{fig:m2varyq}中的$u$随质量变化，积分为
\[
 \chi_f=\int_{M_f}^{M_0}\frac{u(m)}c\frac{\dd m}m
 =\frac{u_0}{cM_0}\int_{M_f}^{M_0}\dd m
 =\frac{u_0}{cM_0}[m]_{M_f}^{M_0}
 =\frac{u_0}c\left(1-\frac{M_f}{M_0}\right).
\]
因此$\beta_f=\tanh\chi_f$。即使把$M_f$趋向0，此控制模型的快度仍只趋$u_0/c$，不发散。
\Diagram{fig:m2varya}{真实快度是对数质量坐标下排气速度的积分}{\draw[axis](0,0)--(3.4,0)node[right]{$\ln(M_0/m)$};\draw[axis](0,0)--(0,2.8)node[above]{$u/c$};\draw[curve,domain=0:2,samples=60]plot(\x,{2.3*exp(-\x)});\draw[guide](0,{2.3*exp(-2)})--(2,{2.3*exp(-2)});\draw[guide](2,0)--(2,{2.3*exp(-2)});\node[above right]at(2,{2.3*exp(-2)}){末值近似};}
若固定取末速度$u_f=u_0M_f/M_0$，会得到$\chi_{\rm app}=(u_0/c)(M_f/M_0)\ln(M_0/M_f)$。图\ref{fig:m2varya}中末值矩形低于整条曲线，所以该近似低估快度；直接用算术平均速度也无依据，正确权重是$\dd m/m$。
\end{solution}

\begin{example}\label{ex:m2fueloptimal}
火箭的全部点火满足同一固定排气速度$u$。需先加速到巡航速度，再减速到静止，初末质量比$R=M_0/M_f$已知。求可达巡航速度；若两次点火各使用相等的绝对燃料质量，是否能恰好停止？
\Diagram{fig:m2optimalq}{两次相反点火，中间质量为$M_1$}{\draw[curve](0,0)rectangle(1.5,1);\node at(.75,.5){$M_0$};\draw[->](1.7,.5)--(2.5,.5);\draw[curve](2.7,0)rectangle(4.2,1);\node at(3.45,.5){$M_1$};\draw[->](4.4,.5)--(5.2,.5);\node[right]at(5.3,.5){$M_f$};}
\end{example}
\begin{solution}
图\ref{fig:m2optimalq}中令两次燃烧之间的质量为$M_1$。加速和制动方向相反，要求的快度变化大小相同：
\[
 \chi=\frac uc\ln\frac{M_0}{M_1}
 =\frac uc\ln\frac{M_1}{M_f}.
\]
两式相等给$M_1^2=M_0M_f$；相加给$2\chi=(u/c)\ln R$。所以$M_1=\sqrt{M_0M_f}$，巡航速度为$c\tanh[(u/2c)\ln R]$。
\Diagram{fig:m2optimala}{两次相等的是质量比；后段在更小质量下产生相同快度变化}{\draw[axis](0,0)--(5.2,0)node[right]{过程};\draw[axis](0,0)--(0,3)node[above]{$\ln m$};\draw[curve](0,2.6)--(2.3,1.5)--(4.6,.4);\draw[guide](0,1.5)--(2.3,1.5);\node[right]at(2.3,1.5){$\ln\sqrt{M_0M_f}$};}
图\ref{fig:m2optimala}按$\ln m$标出两段等长。若改为各消耗相同绝对燃料质量，则$M_1=(M_0+M_f)/2$，其最终快度为
\begin{align*}
 \chi_f&=\frac uc\left(\ln\frac{M_0}{M_1}-\ln\frac{M_1}{M_f}\right)\\
 &=\frac uc\ln\frac{4M_0M_f}{(M_0+M_f)^2}<0
 \quad(M_0>M_f).
\end{align*}
后段制动超过所需，船最终反向运动。由于每消耗一单位质量的快度增量与当前质量成反比，同样千克数的燃料在后段作用更大。满足停车条件时，两段消耗量为$M_0-\sqrt{M_0M_f}$和$\sqrt{M_0M_f}-M_f$，前段较多，其比值为$\sqrt R$。
\end{solution}

\section{近场接收与远场视速度}
\begin{example}\label{ex:m2nearretarded}
观察者在原点，发光点沿$x$轴以$v$运动，轨迹为$(vt,b)$，$b>0$。推导接收时间对发射时间的导数，说明为什么严格近场中没有单个固定观察角。
\Diagram{fig:m2nearq}{有限冲量参数$b$的直线运动}{\draw[axis](-2.5,0)--(3,0)node[right]{$x$};\draw[axis](0,0)--(0,3)node[above]{$y$};\draw[curve](-2.2,1.5)--(2.7,1.5);\fill(1.8,1.5)circle(1.5pt)node[above]{$S$};\draw[guide](0,0)--(1.8,1.5);\node[left]at(0,1.5){$b$};}
\end{example}
\begin{solution}
图\ref{fig:m2nearq}的源距$R(t_e)=\sqrt{v^2t_e^2+b^2}$，接收时间$t_o=t_e+R/c$。求导得
\[
 \frac{\dd t_o}{\dd t_e}
 =1+\frac1{2c}(v^2t_e^2+b^2)^{-1/2}(2v^2t_e)
 =1+\frac{v^2t_e}{c\sqrt{v^2t_e^2+b^2}}.
\]
若定义从源指向观察者的方向$\vct n=-(vt_e,b)/R$，就恢复$1-\vct n\cdot\vct v/c$。该导数在$1-\beta$与$1+\beta$之间，始终为正。
\Diagram{fig:m2neara}{最近点前信号间隔压缩，经过后拉长}{\draw[axis](-2.7,0)--(3,0)node[right]{$vt_e/b$};\draw[axis](0,0)--(0,2.8)node[above]{$\dd t_o/\dd t_e$};\draw[curve,domain=-2.5:2.5,samples=80]plot(\x,{1.4*(1+.8*\x/sqrt(1+\x*\x))});\draw[guide](-2.5,1.4)--(2.5,1.4);}
图\ref{fig:m2neara}在$t_e=0$导数为1，但对应接收时刻为$b/c$。由单调性，亚光速源每个接收时刻只有一个延迟发射时刻；不能把物体“现在的位置”代入传播距离后再解释照片。远场短时间近似才允许把$\vct n$视为常量。
\end{solution}

\begin{example}\label{ex:m2timecompression}
向观察者运动的光源在固有时$\Delta\tau$内完成一次亮度变化，远场观察角为$\theta$。求接收变化持续时间；求光源同一段运动对应的视横向长度。
\Diagram{fig:m2compressq}{固有持续时间、发射时间和到达时间的三层关系}{\draw[curve](0,0)rectangle(1.6,1);\node at(.8,.5){$\Delta\tau$};\draw[->](1.8,.5)--(2.6,.5);\node at(3.1,.5){$\Delta t_e$};\draw[->](3.7,.5)--(4.5,.5);\node[right]at(4.6,.5){$\Delta t_o$};}
\end{example}
\begin{solution}
令$\vct n$从源指向远处观察者，图\ref{fig:m2compressq}中的夹角为$\vct n$与$\vct v$的夹角。源连续两次发射间隔$\Delta t_e$，第二次更靠近观察者$\vct n\cdot\vct v\Delta t_e$，故在方向近似不变的条件下
\[
 \Delta t_o=\Delta t_e-\frac{\vct n\cdot\vct v}{c}\Delta t_e
 =(1-\beta\cos\theta)\gamma\Delta\tau
 =\frac{\Delta\tau}{\delta}.
\]
这里先使用时间膨胀把源的固有时换成发射时间，再计算光程的缩短，两步各有不同来源。
\Diagram{fig:m2compressa}{横向投影长度用发射期间的真实位移，视速度用接收时间相除}{\draw[axis](0,0)--(4.2,0)node[right]{视线};\draw[->,thick](0,0)--(2.7,1.8);\draw[guide](2.7,0)--(2.7,1.8);\node[right]at(2.7,.9){$v\gamma\Delta\tau\sin\theta$};}
图\ref{fig:m2compressa}的同段横向位移为$\Delta X_\perp=v\gamma\Delta\tau\sin\theta$，因此
\[
 v_{\rm app}=\frac{\Delta X_\perp}{\Delta t_o}
 =\frac{v\sin\theta}{1-\beta\cos\theta}.
\]
角度函数与例\ref{ex:m2super}相同，其求导过程已在那里给出：最大视横速出现在$\cos\theta=\beta$，为$\gamma v$。时间压缩因子$\delta$则在$\theta=0$最大，为$\gamma(1+\beta)$；此时横向位移为零。最快亮度变化与最大视横速对应不同观察角。若近距离使视线在变化期间显著转动，应沿实际运动积分微分延迟关系。
\end{solution}

\section{有限加速装置的距离与信号}
\subsection{固定固有距离、等实验室距离与雷达距离}
本讲的加速坐标已给出$\dd s^2=c^2(1+\alpha\xi/c^2)^2\dd\tau^2-\dd\xi^2$。在同一坐标时间的局部静止切片上，沿杆的长度元是$\dd\xi$。所以前后端固定在$\xi=0,L$时，杆的固有长度保持$L$；前端固有加速度却为$\alpha/(1+\alpha L/c^2)$。这是有限装置保持固有长度的一种安排，需要不同位置的动力控制配合。

另一种安排是两个飞船从实验室同时启动，始终使用相同实验室运动程序，它们的实验室间距保持原值。这不保持加速过程中的共动距离。第三种距离是某一站按发射、接收光的往返时间定义的雷达距离，它还包含信号传播期间的加速影响。题目中的“距离保持不变”必须先说明属于哪一种。

\begin{example}加速杆上的钟与后端雷达
一根保持固有长度$L$的杆沿$x$方向加速，后端$A$的固有加速度恒为$\alpha$，前端$B$固定在加速坐标$\xi=L$。后端发出光，经前端理想反射后返回。求后端测得的往返固有时间和雷达距离，并求前端向后端持续发光时的接收频率比。
\ExpansionDiagram{fig:dep-m2-radar-q}{固定加速坐标中的两站及往返信号，图示$\alpha L/c^2=1$}{
\draw[->](0,0)--(6,0)node[right]{$\xi$};\draw[thick](1,-.4)--(1,2.5)node[above]{$A:0$};\draw[thick](5,-.4)--(5,2.5)node[above]{$B:L$};\draw[->,domain=1:5,samples=70]plot(\x,{.3+ln(1+(\x-1)/4)});\draw[->,domain=5:1,samples=70]plot(\x,{.3+2*ln(2)-ln(1+(\x-1)/4)});\node[left]at(.9,.3){发射};\node[left]at(.9,1.6863){返回};\node[left]at(.9,2.3){$\tau$};}
\end{example}
\begin{solution}
在图\ref{fig:dep-m2-radar-q}的坐标中，前行光满足$\dd\xi/\dd\tau=c(1+\alpha\xi/c^2)$。分离变量，从0积分到$L$，单程所需后端坐标时间为
\[
 t_+=\int_0^L\frac{\dd\xi}{c(1+\alpha\xi/c^2)}
 =\left[\frac c\alpha\ln(1+\alpha\xi/c^2)\right]_0^L
 =\frac c\alpha\ln(1+\alpha L/c^2).
\]
反向光的坐标速度反号，但区间次序也反向，故$t_-=t_+$。后端钟的固有时就是该坐标时间，得到
\[
 \Delta\tau_A=\frac{2c}{\alpha}\ln(1+\alpha L/c^2),\qquad
 D_A=\frac{c\Delta\tau_A}{2}
 =\frac{c^2}{\alpha}\ln(1+\alpha L/c^2).
\]
\ExpansionDiagram{fig:dep-m2-radar-a}{雷达距离与固定固有长度的比较}{
\draw[->](0,0)--(4,0)node[right]{$\alpha L/c^2$};\draw[->](0,0)--(0,3.4)node[above]{$\alpha D_A/c^2$};\draw[dashed](0,0)--(3.1,3.1);\draw[thick,domain=0:3.1,samples=80]plot(\x,{ln(1+\x)});\node[right]at(2.9,1.2){$\ln(1+x)$};}
图\ref{fig:dep-m2-radar-a}中$D_A<L$。弱加速展开为$D_A=L-\alpha L^2/(2c^2)+\cdots$，这与局部测尺得到$L$不同，是两种测量程序的差别。

前端相邻两个波峰的发射固有时间间隔为$\Delta\tau_B$。固定两站之间的坐标传播时间不随发射时刻变化，因此相邻波峰的接收坐标时间间隔等于发射坐标间隔。又$\Delta\tau_B=(1+\alpha L/c^2)\Delta\tau$、$\Delta\tau_A=\Delta\tau$，所以
\[
 \frac{\nu_A}{\nu_B}=1+\frac{\alpha L}{c^2}.
\]
这里$\nu_B$按前端本地钟定义。若只讨论后端发出再由固定前端反射回来的光，两个方向的频率因子相乘抵消，返回后端的频率等于原发射频率；单程两站比较与同站往返比较不可混用。
\end{solution}

\subsection{相同加速程序造成的有限长度效应}
两端距离很短时，瞬时长度收缩关系可以估算共动长度。有限间距下，同一瞬时惯性系选出的两个端点事件往往处于不同加速阶段，须先求同时面与两条世界线的交点，再计算其空间距离。

\begin{example}始终保持实验室间距的两飞船
两飞船初始静止，实验室间距$d>0$。从$t=0$起，两者均沿$+x$保持相同固有加速度$\alpha$，并持续加速。后船$A$的快度为$\chi$时，在其瞬时静止惯性系中，与该事件同时的前船$B$事件具有何种快度？求该系中的两船距离。
\ExpansionDiagram{fig:dep-m2-bell-q}{两条世界线在实验室等时切片上的间距恒为$d$}{
\draw[->](-.3,0)--(4,0)node[right]{$x$};\draw[->](0,0)--(0,3.4)node[above]{$ct$};\draw[thick,domain=0:1.5,samples=60]plot({cosh(\x)-1},{sinh(\x)});\draw[thick,domain=0:1.5,samples=60]plot({cosh(\x)-.4},{sinh(\x)});\draw[dashed](.543,1.175)--(1.143,1.175)node[midway,above]{$d$};\node[below]at(0,0){$A$};\node[below]at(.6,0){$B$};}
\end{example}
\begin{solution}
设$\ell_\alpha=c^2/\alpha$。图\ref{fig:dep-m2-bell-q}的两条世界线参数为
\[
 ct_A=\ell_\alpha\sinh\chi,\quad
 x_A=\ell_\alpha(\cosh\chi-1),
\]
\[
 ct_B=\ell_\alpha\sinh\eta,\quad
 x_B=d+\ell_\alpha(\cosh\eta-1).
\]
在$A$的瞬时系中同时，要求$\Delta(ct)'=\cosh\chi\,\Delta(ct)-\sinh\chi\,\Delta x=0$。将两条世界线代入，使用双曲正弦差角公式得到
\[
 \ell_\alpha\sinh(\eta-\chi)=d\sinh\chi,
 \qquad \eta=\chi+\operatorname{arsinh}
 \left(\frac d{\ell_\alpha}\sinh\chi\right).
\]
这表明前船上所选的同时事件晚于实验室等时所选事件，且在持续加速期间已经具有更大的快度。
\ExpansionDiagram{fig:dep-m2-bell-a}{后船瞬时系的同时线与前船世界线相交于更晚事件}{
\draw[->](-.3,0)--(4,0)node[right]{$x$};\draw[->](0,0)--(0,3.4)node[above]{$ct$};\draw[thick,domain=0:1.8,samples=70]plot({cosh(\x)-1},{sinh(\x)});\draw[thick,domain=0:1.8,samples=70]plot({cosh(\x)-.4},{sinh(\x)});\draw[dashed](.543,1.175)--(3.5,1.175);\draw[->](.543,1.175)--(2.9,2.97)node[right]{同时线};\fill(.543,1.175)circle(1.5pt)node[left]{$A$};\fill(2.324,2.533)circle(1.5pt)node[right]{$B$};}
距离须变换空间分量：
\begin{align*}
 D_A&=\cosh\chi\,\Delta x-\sinh\chi\,\Delta(ct)\\
 &=d\cosh\chi+\ell_\alpha[\cosh(\eta-\chi)-1]\\
 &=d\cosh\chi+\ell_\alpha\left[
 \sqrt{1+(d/\ell_\alpha)^2\sinh^2\chi}-1\right].
\end{align*}
图\ref{fig:dep-m2-bell-a}中的同时线不是实验室水平切片，所以有限间距的精确结果比$d\cosh\chi$还大。当$(d/\ell_\alpha)|\sinh\chi|\ll1$时，第二项约为$d^2\sinh^2\chi/(2\ell_\alpha)$，此时才可用$D_A\simeq d\gamma$近似。若两船在某实验室时刻停止加速，则前船上所取事件可能已经进入匀速段，必须分段接续世界线，不能继续沿用此式。
\end{solution}

\section{同一实验室时刻怎样读取加速杆两端的钟}
比较加速杆两端的钟，需要先指定读取事件。加速坐标的同时面和实验室的同时面选取不同的事件对，钟面读数也随之改变。以下令杆的固有长度保持不变，比较实验室在同一时刻读取的两端固有时。

设杆后端固有加速度为$\alpha$，固有长度$L$，记$s=1+\alpha L/c^2>1$。本讲加速坐标给后端$ct=(c^2/\alpha)\sinh(\alpha\tau/c)$；前端则为$ct=s(c^2/\alpha)\sinh(\alpha\tau/c)$。在两端都取实验室$t=0$读数为零后，前端固有时还要乘局部钟率因子$s$。

\begin{example}同时拍摄两只加速钟的坐标读数
理想记录装置已经消除光传播延迟，得到实验室同一时刻$t$的后端钟读数$T_A$和前端钟读数$T_B$。两钟按上述保持固有长度的方案从静止开始持续加速。求两读数，并判断哪只钟读数更大。这里讨论坐标同时读数，不讨论未经延迟修正的照片。
\ExpansionDiagram{fig:dep-m2-clock-q}{实验室水平同时线选取两条加速世界线上的不同加速坐标时刻}{
\draw[->](0,0)--(4,0)node[right]{$x$};\draw[->](0,0)--(0,3.8)node[above]{$ct$};\draw[thick,domain=0:1.8,samples=70]plot({cosh(\x)-1},{sinh(\x)});\draw[thick,domain=0:1.6,samples=70]plot({1.5*cosh(\x)-1},{1.5*sinh(\x)});\draw[dashed](0,2)--(3.3,2)node[right]{同一$t$};\node[below]at(0,0){$A$};\node[below]at(.5,0){$B$};}
\end{example}
\begin{solution}
图\ref{fig:dep-m2-clock-q}要求先分别反解世界线时间。记$s=1+\alpha L/c^2>1$。后端固有时就是其加速坐标时间，前端固有时为局部钟率因子$s$乘其坐标时间，所以
\[
 T_A=\frac c\alpha\operatorname{arsinh}\frac{\alpha t}{c},\qquad
 T_B=\frac{sc}{\alpha}\operatorname{arsinh}\frac{\alpha t}{sc}.
\]
不能把$T_B$直接写为$sT_A$，因为那样比较的是相同加速坐标时间，不是图中的实验室水平切片。
\ExpansionDiagram{fig:dep-m2-clock-a}{初始读数相同，随后前端钟在实验室相同时间内积累更多固有时}{
\draw[->](0,0)--(5,0)node[right]{$\alpha t/c$};\draw[->](0,0)--(0,3)node[above]{$\alpha T/c$};\draw[thick,domain=0:4.5,samples=90]plot(\x,{ln(\x+sqrt(1+\x^2))});\draw[dashed,domain=0:4.5,samples=90]plot(\x,{1.5*ln(\x/1.5+sqrt(1+(\x/1.5)^2))});\node[right]at(3.5,2.7){$T_B$};\node[right]at(3.5,1.7){$T_A$};}
两式求导，
\[
 \frac{\dd T_A}{\dd t}
 =\frac c\alpha\frac{\alpha/c}{\sqrt{1+(\alpha t/c)^2}}
 =\frac1{\sqrt{1+(\alpha t/c)^2}},
\]
\[
 \frac{\dd T_B}{\dd t}
 =\frac{sc}\alpha\frac{\alpha/(sc)}{\sqrt{1+(\alpha t/(sc))^2}}
 =\frac1{\sqrt{1+(\alpha t/(sc))^2}}.
\]
因$s>1$，对$t>0$有$\dot T_B>\dot T_A$，结合初始两读数为零，得到$T_B>T_A$。图\ref{fig:dep-m2-clock-a}画$s=1.5$的曲线。初期展开可得$T_B-T_A\simeq\alpha^2t^3(1-s^{-2})/(6c^2)$，差值不是从启动瞬间就线性增长。

物理上，保持固有长度要求前端加速度较小。在同一实验室时刻，前端速率较低，钟慢效应较弱，所以固有时积累更多。这一解释依赖本题的具体加速安排；若两端使用相同固有加速度，世界线和所比较距离都改变，不能照搬这些读数。
\end{solution}

\section{习题与解答}
\begin{exercise}\label{ex:m2horizon}
飞船从地球由静止出发，持续保持正向固有加速度$\alpha$。地球在$t_e>0$发出追赶飞船的光。求飞船接收该信号的固有时，并判断地球所有未来信号是否都能追上。
\Diagram{fig:m2horizonq}{恒固有加速世界线及追赶它的光线}{\draw[axis](-.2,0)--(4.5,0)node[right]{$\alpha x/c^2$};\draw[axis](0,0)--(0,4.6)node[above]{$\alpha t/c$};\draw[curve,domain=0:2,samples=90]plot({cosh(\x)-1},{sinh(\x)});\draw[->](0,.5)--(.25,.75);\draw[guide](0,1)--(3.4,4.4);\node[left]at(0,1){$1$};}
\end{exercise}
\begin{solution}
飞船世界线$ct=L_\alpha\sinh\chi$、$x=L_\alpha(\cosh\chi-1)$，$\chi=\alpha\tau/c$。图\ref{fig:m2horizonq}的信号满足$x=c(t-t_e)$，代入后
\[
 \frac{\alpha t_e}c=\sinh\chi-\cosh\chi+1=1-\ee^{-\chi}.
\]
所以仅当$0<t_e<c/\alpha$有有限接收固有时
\[
 \tau_R=-\frac c\alpha\ln(1-\alpha t_e/c).
\]
\Diagram{fig:m2horizona}{发射时刻逼近$c/\alpha$时，接收固有时发散}{\draw[axis](0,0)--(4.5,0)node[right]{$\alpha t_e/c$};\draw[axis](0,0)--(0,3.8)node[above]{$\alpha\tau_R/c$};\draw[curve,domain=0:.96,samples=90]plot({3.5*\x},{-ln(1-\x)});\draw[guide](3.5,0)--(3.5,3.5);\node[below]at(3.5,0){$1$};}
图\ref{fig:m2horizona}在边界发散。$t_e\geq c/\alpha$的信号永不追上持续加速的飞船，虽然飞船在每个有限时刻速度都小于$c$。这来自其未来世界线不断趋近一条光状渐近线。若飞船后来停机或转向，未来轨迹改变，不能继续使用永久加速的信号边界。
\end{solution}

\begin{exercise}\label{ex:m2labconstant}
光子火箭从静止出发，按初始实验室时间计算的后向光功率恒为$P$。求火箭剩余静质量和速度随实验室时间的变化，并与恒固有功率模型区分。初始质量$M_0$。
\Diagram{fig:m2labconstantq}{功率固定在同一实验室中定义}{\draw[curve](2,0)rectangle(4,1);\draw[->](1.8,.5)--(.1,.5)node[above]{$P$实验室};\draw[->](4.2,.5)--(5.7,.5)node[above]{$v(t)$};}
\end{exercise}
\begin{solution}
图\ref{fig:m2labconstantq}截至时间$t$发出的总光能$Pt$、总后向动量$-Pt/c$。全系统守恒给飞船总能量$E=M_0c^2-Pt$、动量$p=Pt/c$。令$s=Pt/(M_0c^2)$，质壳关系得到
\[
 m=M_0\sqrt{1-2s},\qquad \beta=\frac{pc}E=\frac s{1-s},\quad 0\leq s<1/2.
\]
实际还须在$m$达到干质量时停机，不能把$m=0$极限当有质量飞船状态。
\Diagram{fig:m2labconstanta}{实验室定功率与固有定功率对应不同质量时间关系}{\draw[axis](0,0)--(4.4,0)node[right]{$s$};\draw[axis](0,0)--(0,2.8)node[above]{$m/M_0$};\draw[curve,domain=0:.499,samples=90]plot({7*\x},{2.3*sqrt(1-2*\x)});\node[below]at(3.5,0){$1/2$};\node[left]at(0,2.3){$1$};}
图\ref{fig:m2labconstanta}的质量不是实验室时间的线性函数。由$P=P_*(1-\beta)$，所需固有功率$P_*=P(1-s)/(1-2s)$会增大。固有时间满足$\dd\tau=\sqrt{1-2s}\dd t/(1-s)$。令$K=M_0c^2/P$、$w=\sqrt{1-2s}$，则$\dd t=K\dd s=-Kw\dd w$、$1-s=(1+w^2)/2$，所以
\[
 \tau=-2K\int_1^w\frac{q^2}{1+q^2}\dd q
 =-2K[q-\arctan q]_1^w.
\]
代入上下限，得到
\[
 \tau=\frac{2M_0c^2}P\left[1-w+\arctan w-\frac\pi4\right].
\]
在$t=0$有$w=1,\tau=0$。可见恒功率条件随参考系而变。固有功率恒定时采用$m=M_0-P_*\tau/c^2$，对应另一种推进程序。
\end{solution}

\begin{exercise}\label{ex:m2midburn}
光子火箭需从静止到达距离$L$处并停止，前半程保持固有加速度$\alpha$，后半程保持同样大小的固有减速度。求地球时间、船上时间和初末质量比。
\Diagram{fig:m2midq}{中点反转推力，速度连续}{\draw[axis](-.3,0)--(5.8,0)node[right]{$x$};\fill(0,0)circle(1.5pt);\fill(5,0)circle(1.5pt);\draw[guide](2.5,-.7)--(2.5,1.3);\draw[->](.5,.6)--(1.8,.6)node[above]{$\alpha$};\draw[->](4.5,.6)--(3.2,.6)node[above]{$-\alpha$};}
\end{exercise}
\begin{solution}
对图\ref{fig:m2midq}的前半程，用恒固有加速度世界线，令$q=\alpha L/(2c^2)$，得到中点快度$\chi_m=\operatorname{arcosh}(1+q)$。该段地球时间$c\sinh\chi_m/\alpha$、船上时间$c\chi_m/\alpha$。
\Diagram{fig:m2mida}{固有时中的快度先线性增加再线性减少}{\draw[axis](0,0)--(5.2,0)node[right]{$\tau$};\draw[axis](0,0)--(0,3)node[above]{$\chi$};\draw[curve](0,0)--(2.3,2.5)--(4.6,0);\node[above]at(2.3,2.5){$\chi_m$};}
后半程是时间反向的减速过程，图\ref{fig:m2mida}两段持续时间相同。因此$T=2c\sqrt{q(2+q)}/\alpha$、$\tau=2c\operatorname{arcosh}(1+q)/\alpha$。光子火箭消耗与总绝对快度$2\chi_m$对应，$M_0/M_f=\ee^{2\chi_m}=[1+q+\sqrt{q(2+q)}]^2$。减速阶段仍耗质量，不能让快度符号变化抵消燃料消耗。
\end{solution}

\begin{exercise}\label{ex:m2staging}
固定排气速度$u$的火箭初质量100，燃烧至60时丢弃质量20的空壳，再燃烧至最终质量20。求总快度，并与同样初质量、最终携带40质量但不丢壳的单级模型比较。所有质量使用同一单位。
\Diagram{fig:m2stageq}{无相对弹射的丢壳瞬间，飞船速度不变}{\node at(0,0){$100$};\draw[->](.6,0)--(1.5,0);\node at(2,0){$60$};\draw[->](2.6,0)--(3.5,0);\node at(4,0){$40$};\draw[->](4.6,0)--(5.5,0);\node at(6,0){$20$};\node[below]at(3,-.6){丢壳20};}
\end{exercise}
\begin{solution}
图\ref{fig:m2stageq}第一段快度$(u/c)\ln(100/60)$。丢壳时空壳与飞船相对静止，四动量按共同四速度分开，没有快度突变；第二段快度$(u/c)\ln(40/20)$。总快度$\chi=(u/c)\ln(10/3)$。
\Diagram{fig:m2stagea}{对数质量图中竖直丢壳不贡献推进快度}{\draw[axis](0,0)--(5.2,0)node[right]{阶段};\draw[axis](0,0)--(0,3.5)node[above]{$\ln m$};\draw[curve](0,3)--(1.8,2.3)--(1.8,1.7)--(4.2,.7);\node[right]at(1.8,2){丢壳};}
不丢壳的对照模型结束时含壳和有效载荷总质量40。其快度为
\[
 \chi_1=\frac uc\ln\frac{100}{40},\qquad
 \chi-\chi_1=\frac uc\ln\frac43>0.
\]
图\ref{fig:m2stagea}对应的分级增益为后一式。若把无推进丢掉的20也按火箭式积分，会错误地多算一段快度；比较模型的有效载荷、燃料和空壳质量必须逐项一致。
\end{solution}

\begin{exercise}\label{ex:m2labpower}
光子火箭沿$+x$运动，瞬时静止系向后辐射功率为$P_*$。求按实验室发射时间计算的光功率，并由它和飞船动能变化核对质量能守恒。
\Diagram{fig:m2labpowerq}{后向光子在实验室红移}{\draw[curve](2,0)rectangle(4,1);\draw[axis](4.2,.5)--(5.6,.5)node[above]{$V$};\draw[axis](1.8,.5)--(.2,.5)node[above]{$P_*$};}
\end{exercise}
\begin{solution}
图\ref{fig:m2labpowerq}后向光子能量因子为$\gamma(1-\beta)$，发射时间因子为$\gamma$，所以$P_\gamma=P_*(1-\beta)$。同时$\dd m/\dd t=-P_*/(\gamma c^2)$，飞船快度率$\dd\chi/\dd t=P_*/(\gamma mc^2)$。
\Diagram{fig:m2labpowera}{飞船总能量减小率等于后向辐射实验室功率}{\draw[curve](0,0)rectangle(2.2,1);\node at(1.1,.5){$\gamma mc^2$};\draw[->](2.4,.5)--(3.5,.5);\node[right]at(3.7,.5){$P_*(1-\beta)$};}
对图\ref{fig:m2labpowera}的飞船总能量求导，利用$\dd\gamma/\dd\chi=\gamma\beta$，有
\[
 \frac{\dd(\gamma mc^2)}{\dd t}
 =c^2\left(\gamma\frac{\dd m}{\dd t}
 +m\frac{\dd\gamma}{\dd\chi}\frac{\dd\chi}{\dd t}\right)
 =-P_*+mc^2\gamma\beta\frac{P_*}{\gamma mc^2}
 =-P_*(1-\beta).
\]
飞船虽加速，总能量仍因辐射流失而下降；静质量损失的一部分转为更大的运动因子。不能把“速度增加”直接等同于飞船总能量必增。
\end{solution}

\begin{exercise}\label{ex:m2relay}
沿用地球与恒星相距10光年、速度$0.8c$的过程，但旅行者不转向。另一只返航钟在恒星处与其相遇，把读数设置为7.5年后飞回地球。比较重逢读数与原双生子结果，解释是否必须由同一只钟经历加速才能得到差值。
\Diagram{fig:m2relayq}{两只匀速钟在恒星处交接读数}{\draw[axis](-.2,0)--(4.2,0);\draw[axis](0,0)--(0,5.4)node[above]{$ct$};\draw[curve](0,0)--(2,2.5)--(3.7,4.625);\draw[curve](3.6,.5)--(0,5);\fill(2,2.5)circle(1.4pt)node[right]{$B$交接};}
\end{exercise}
\begin{solution}
图\ref{fig:m2relayq}去程钟从出发到$B$经历7.5年，返航钟从$B$到地球也经历7.5年。它在$B$被设置为7.5年，所以到地球读15年，地球仍为25年。
\Diagram{fig:m2relaya}{交接后的读数拼接两段不同世界线的固有时间}{\draw[axis](0,0)--(5.1,0)node[right]{地球时间};\draw[axis](0,0)--(0,3)node[above]{拼接读数};\draw[curve](0,0)--(4.6,2.4);\draw[guide](2.3,0)--(2.3,1.2);\node[below]at(2.3,0){12.5年};\node[right]at(4.6,2.4){15年};}
图\ref{fig:m2relaya}的拼接读数沿两段组成的折线累积，几何上与瞬时转向世界线相同。没有单只返航钟需要在$B$改变速度，但必须实际校准读数；两只钟若早在各自惯性系随意同步，初始偏移未必满足本题。差值由选择的时空路径决定，加速是连接路径的一种实现方式。
\end{solution}

\begin{exercise}\label{ex:m2retardedzero}
发光点在实验室沿$(vt,b)$运动，观察者在原点。观察者在$t_o=0$拍照，而该坐标时刻物体恰在最近点$x=0$。求照片中光的发射时刻和发射位置。
\Diagram{fig:m2zeroq}{同时最近点与照片发光点不是同一点}{\draw[axis](-3,0)--(2.5,0)node[right]{$x$};\draw[axis](0,0)--(0,2.8)node[above]{$y$};\draw[curve](-2.6,1.6)--(2.2,1.6);\fill(0,1.6)circle(1.5pt)node[above right]{同时位置};}
\end{exercise}
\begin{solution}
图\ref{fig:m2zeroq}的延迟条件为$0=t_e+\sqrt{v^2t_e^2+b^2}/c$，要求$t_e<0$。移项平方后$c^2t_e^2=v^2t_e^2+b^2$，所以$t_e=-\gamma b/c$、$x_e=-\beta\gamma b$。
\Diagram{fig:m2zeroa}{入眼光线来自最近点之前，路程为$\gamma b$}{\draw[axis](-3.2,0)--(2.3,0)node[right]{$x$};\draw[axis](0,0)--(0,2.7)node[above]{$y$};\fill(-2,1.6)circle(1.5pt)node[above]{$x_e=-\beta\gamma b$};\draw[->](-2,1.6)--(0,0);\draw[guide](0,0)--(0,1.6);\node[right]at(0,.8){$b$};}
代回原式，光程$R_e=\sqrt{\beta^2\gamma^2b^2+b^2}=\gamma b$，恰在$\gamma b/c$后到达，验证选取了正确负根。图\ref{fig:m2zeroa}中的源虽在坐标时间0已到最近点，照片仍显示此前位置；照片记录的是经过传播后到达的旧信号。
\end{solution}

\begin{exercise}\label{ex:m2fuelbudget}
理想光子火箭总初末质量比为$R$，需把给定质量的有效载荷送到远方并相对地球停下。不规定距离，忽略巡航损失。求最大巡航洛伦兹因子以及质量比达到100时的巡航速度。
\Diagram{fig:m2budgetq}{燃料预算分成加速和减速两个相等快度段}{\draw[axis](0,0)--(5.2,0);\draw[curve](0,.3)--(2.3,1.8)--(4.6,.3);\node[above]at(2.3,1.8){巡航};\node[below]at(0,0){静止};\node[below]at(4.6,0){静止};}
\end{exercise}
\begin{solution}
图\ref{fig:m2budgetq}总绝对快度为$\ln R$，要最终停止，两段均为$\chi=\tfrac12\ln R$。所以$\gamma_{\max}=\cosh(\tfrac12\ln R)=(\sqrt R+1/\sqrt R)/2$，$\beta=(R-1)/(R+1)$。
\Diagram{fig:m2budgeta}{保留制动燃料会显著降低同一初末质量比的巡航速度}{\draw[axis](0,0)--(4.8,0)node[right]{$R$};\draw[axis](0,0)--(0,2.8)node[above]{$\beta$};\draw[curve,domain=1:100,samples=80]plot({.04*\x},{2.3*(\x-1)/(\x+1)});\node[below]at(4,0){100};}
取$R=100$，$\gamma=5.05$、$\beta=99/101\simeq.980198$，如图\ref{fig:m2budgeta}。若把全部燃料只用来加速，速度可达$(R^2-1)/(R^2+1)$，但剩余质量不足以再完成制动。任务若要求到达后停下，质量比还要包含制动段的消耗。
\end{solution}

\begin{exercise}\label{ex:m2signals}
例\ref{ex:m2twins}中地球每年发出一束时间标记光信号。求旅行者去程和返程接收到的频率，并据此还原地球25年的连续钟信号。
\Diagram{fig:m2signalsq}{地球连续发光，旅行者先远离后靠近}{\draw[axis](-.2,0)--(3.5,0)node[right]{$x$};\draw[axis](0,-.2)--(0,5.5)node[above]{$ct$};\draw[curve](0,0)--(2,2.5)--(0,5);\foreach\y in{.1,.3,.5}{\draw[guide](0,\y)--({4*\y},{5*\y});}\foreach\y in{1.3,2.1,2.9,3.7,4.5}{\draw[guide](0,\y)--({(4-.8*\y)/1.8},{(4+\y)/1.8});}}
\end{exercise}
\begin{solution}
去程接收者远离，频率比为$\sqrt{(1-\beta)/(1+\beta)}=1/3$；返程靠近，频率比为3。旅行者两段固有时间各7.5年，所以去程接收相当于地球2.5年内发出的信号，返程接收其余22.5年。
\Diagram{fig:m2signalsa}{转向瞬间接收到的光来自地球2.5年时刻}{\draw[axis](-.2,0)--(3.5,0)node[right]{$x$};\draw[axis](0,-.2)--(0,5.5)node[above]{$ct$};\draw[curve](0,0)--(2,2.5)--(0,5);\draw[very thick,->](0,.5)--(2,2.5);\node[left]at(0,.5){2.5年};\node[right]at(2,2.5){$B$};}
图\ref{fig:m2signalsa}中由光速连线也可算发射时刻$t_e=t_B-l/c=12.5-10=2.5$年。与旅行者去程、返程赋予地球的同时刻4.5年、20.5年都不同。接收信号连续，只有接收频率在理想瞬时换向模型中突变。
\end{solution}

\begin{exercise}\label{ex:m2fuelreturn}
光子火箭需由地球静止起飞，加速至$\beta c$，到远处减速至相对地球静止，随后原路返回并最终停下。忽略巡航耗能，求包括全部返航燃料在内的初末质量比。
\Diagram{fig:m2returnq}{四段速度改变与两段巡航}{\draw[axis](-.3,0)--(6,0)node[right]{$x$};\fill(0,0)circle(1.5pt)node[below]{地球};\fill(5,0)circle(1.5pt)node[below]{远端};\draw[->](.3,.8)--(4.7,.8)node[midway,above]{$+\beta c$};\draw[->](4.7,-.8)--(.3,-.8)node[midway,below]{$-\beta c$};}
\end{exercise}
\begin{solution}
令$\chi=\operatorname{artanh}\beta>0$。图\ref{fig:m2returnq}中每次由静止到巡航或由巡航到静止，都需消耗绝对快度变化$\chi$。喷流方向在减速时反转，火箭质量仍减少，故应对快度变化取绝对值。
\Diagram{fig:m2returna}{快度变化总长度为$4\chi$}{\draw[axis](-.3,0)--(6,0)node[right]{过程};\draw[axis](0,-1.5)--(0,1.8)node[above]{$\chi_{\rm ship}$};\draw[very thick](0,0)--(1,1)--(2,1)--(3,0)--(4,-1)--(5,-1)--(6,0);\node[left]at(0,1){$\chi$};\node[left]at(0,-1){$-\chi$};}
图\ref{fig:m2returna}中四次点火的质量比逐次相乘，故$M_0/M_f=\ee^{4\chi}=[(1+\beta)/(1-\beta)]^2$。若$\beta=.8$，质量比为81。不能先算单程消耗，再以最终干质量为基准简单加倍，因为起飞时必须推动全部后续燃料。
\end{solution}

\begin{exercise}\label{ex:m2constalpha}
飞船从静止出发，沿$x$保持固有加速度$\alpha$，直到到达地球距离$L$处。求地球时间、船上时间和末速度。
\Diagram{fig:m2alphaq}{固定终点距离的持续加速}{\draw[axis](-.3,0)--(5,0)node[right]{$x$};\fill(0,0)circle(1.5pt);\draw[curve](4,-.5)--(4,1.2);\draw[<->](0,-.5)--(4,-.5)node[midway,below]{$L$};\draw[axis](.5,.5)--(2,.5)node[above]{$\alpha$};}
\end{exercise}
\begin{solution}
图\ref{fig:m2alphaq}取飞船由原点静止出发。固有加速度定义给$\dd\chi/\dd\tau=\alpha/c$，所以$\chi=\alpha\tau/c$。再由四速度的时间和空间分量，
\[
 t=\int_0^\tau\cosh(\alpha s/c)\dd s
 =\left[\frac c\alpha\sinh(\alpha s/c)\right]_0^\tau
 =\frac c\alpha\sinh\chi,
\]
\[
 x=\int_0^\tau c\sinh(\alpha s/c)\dd s
 =\left[\frac{c^2}\alpha\cosh(\alpha s/c)\right]_0^\tau
 =\frac{c^2}\alpha(\cosh\chi-1).
\]
两个积分常数已由$t=x=0$确定。
\Diagram{fig:m2alphaa}{终点落在恒固有加速度双曲线上}{\draw[axis](-.2,0)--(4.6,0)node[right]{$\alpha x/c^2$};\draw[axis](0,-.2)--(0,4.6)node[above]{$\alpha t/c$};\draw[curve,domain=0:2,samples=70]plot({cosh(\x)-1},{sinh(\x)});\fill(2,2.828)circle(1.5pt)node[right]{$x=L$};\draw[guide](2,0)--(2,2.828);}
在图\ref{fig:m2alphaa}中令$q=\alpha L/c^2$。到达时$\cosh\chi_f=1+q$，正向运动选择$\sinh\chi_f=+\sqrt{(1+q)^2-1}=\sqrt{q(2+q)}$，因此
\[
 t_f=\frac c\alpha\sqrt{q(2+q)},\quad
 \tau_f=\frac c\alpha\operatorname{arcosh}(1+q),\quad
 \beta_f=\frac{\sqrt{q(2+q)}}{1+q}.
\]
平方后$1-\beta_f^2=(1+q)^{-2}>0$，也给$\gamma_f=1+q$。小$q$时$\operatorname{arcosh}(1+q)\simeq\sqrt{2q}$，两种时间同趋$\sqrt{2L/\alpha}$、末速趋$\sqrt{2\alpha L}$。大$q$时$t_f\simeq L/c+c/\alpha$，而$\tau_f\simeq(c/\alpha)\ln(2q)$，远距离的船上时间仅呈对数增长；这依赖全程持续提供恒固有加速度。
\end{solution}

\begin{exercise}\label{ex:m2powerengine}
理想光子火箭在自身瞬时静止系的辐射功率恒为$P_*$。将它与恒固有加速度光子火箭比较：分别求质量随固有时间的变化及固有功率随时间的条件。
\Diagram{fig:m2engineq}{定向光束的固有功率与推力}{\draw[curve](2,0)rectangle(4.3,1);\node at(3.15,.5){火箭$m$};\draw[axis](1.8,.5)--(.2,.5)node[above]{$P_*/c$};\draw[axis](4.5,.5)--(5.7,.5)node[above]{$\alpha$};}
\end{exercise}
\begin{solution}
图\ref{fig:m2engineq}中，瞬时船系向后发光的能量率为$P_*$，动量率大小为$P_*/c$。一阶能量与动量守恒分别给
\[
 \frac{\dd m}{\dd\tau}=-\frac{P_*}{c^2},\quad
 mc\frac{\dd\chi}{\dd\tau}=\frac{P_*}{c},\quad
 \alpha=c\frac{\dd\chi}{\dd\tau}=\frac{P_*}{mc}.
\]
若$P_*$恒定，以$m(0)=M_0$积分：$\int_{M_0}^{m}\dd\mu=-(P_*/c^2)\int_0^\tau\dd s$，得到$m=M_0-P_*\tau/c^2$。再相除两个微分方程，有$\dd\chi=-\dd m/m$，从初始静止积分得$\chi=-[\ln\mu]_{M_0}^{m}=\ln(M_0/m)$。推力恒定而惯性质量减小，故固有加速度不断增大。
\Diagram{fig:m2enginea}{恒加速度需要功率与剩余质量同步下降}{\draw[axis](0,0)--(4.3,0)node[right]{$\alpha\tau/c$};\draw[axis](0,0)--(0,2.8)node[above]{$m/M_0,\ P_*/P_{*0}$};\draw[curve,domain=0:3.5,samples=70]plot(\x,{2.3*exp(-\x)});\node[left]at(0,2.3){$1$};}
若要求图\ref{fig:m2enginea}中的$\alpha$恒定，必须令功率随质量改变：$P_*=mc\alpha$。质量方程于是变为
\[
 \frac{\dd m}{m}=-\frac\alpha c\dd\tau,
 \quad \int_{M_0}^{m}\frac{\dd\mu}{\mu}
 =-\frac\alpha c\int_0^\tau\dd s,
 \quad \ln\frac m{M_0}=-\frac{\alpha\tau}c.
\]
取指数并代回功率条件，
\[
 m=M_0\ee^{-\alpha\tau/c},\qquad
 P_*=M_0c\alpha\ee^{-\alpha\tau/c}.
\]
此时$\chi=\alpha\tau/c$，与恒固有加速度世界线一致。若干质量为$M_d$，两种程序都在$m=M_d$时停机：恒功率的可燃固有时间为$(M_0-M_d)c^2/P_*$，恒加速度则为$(c/\alpha)\ln(M_0/M_d)$。在相同初末质量比下，两种程序给出相同末快度；燃烧历时和沿途位置由功率程序决定。
\end{solution}

\begin{exercise}\label{ex:m2oppositejets}
一对同时由同一点射出的对称光斑在源系以同速$\beta c$沿相反方向运动。视线与靠近光斑速度夹角为$\theta$。测得两侧视横向速度分别为$u_+,u_-$。求$\beta\cos\theta$及$\beta\sin\theta$。
\Diagram{fig:m2jetsq}{靠近和远离光斑的投影}{\draw[axis](-.3,0)--(4.5,0)node[right]{视线};\draw[very thick,->](0,0)--(2.7,1.7)node[right]{$+$};\draw[very thick,->](0,0)--(-1.9,-1.2)node[left]{$-$};\node[below right]at(0,0){源};}
\end{exercise}
\begin{solution}
图\ref{fig:m2jetsq}的两侧横向位移大小相同，但接收时间分别缩短和拉长，故
\[
 \frac{u_+}c=\frac{\beta\sin\theta}{1-\beta\cos\theta},\quad
 \frac{u_-}c=\frac{\beta\sin\theta}{1+\beta\cos\theta}.
\]
当$u_++u_->0$时，两式相除后解得$A=\beta\cos\theta=(u_+-u_-)/(u_++u_-)$。再乘回任一分母，得$B=\beta\sin\theta=2u_+u_-/[c(u_++u_-)]$。若两侧视横向速度都为零，可能是$\beta=0$，也可能是运动轴沿视线；此时横向数据不足以唯一反演，不能使用上述零分母公式。
\Diagram{fig:m2jetsa}{两个可测投影组合给出速度和视角}{\draw[axis](0,0)--(3.8,0)node[right]{$A$};\draw[axis](0,0)--(0,2.8)node[above]{$B$};\draw[very thick,->](0,0)--(2.7,1.8)node[right]{$\beta$};\draw[guide](2.7,0)--(2.7,1.8)--(0,1.8);\node at(.9,.3){$\theta$};}
图\ref{fig:m2jetsa}给$\beta=\sqrt{A^2+B^2}$、$\theta=\operatorname{atan2}(B,A)$。若由观测反演得到$\beta\geq1$，则对称、同速、远场或距离标定中至少一项不满足模型假设，不能把不相容的数据解释成真实超光速。
\end{solution}

\begin{exercise}\label{ex:m2imagebar}
长度$L_0$的细杆沿$x$运动，杆与速度平行。观察者在远处，视线单位矢量$\vct n$与$x$夹角$\theta$。求同一照片中杆端横向投影的长度，并与等时长度比较。
\Diagram{fig:m2barq}{照片由两端不同的发光事件组成}{\draw[curve](0,0)--(3.5,0);\fill(0,0)circle(1.5pt)node[below]{$A$};\fill(3.5,0)circle(1.5pt)node[below]{$B$};\draw[axis](1.7,.5)--(3,.5)node[above]{$V$};\draw[guide](0,0)--(1.5,2.2);\draw[guide](3.5,0)--(5,2.2);\node[right]at(5,2.2){$\vct n$};}
\end{exercise}
\begin{solution}
图\ref{fig:m2barq}中，实验室同时杆长为$L=L_0/\gamma$。取尾端、头端世界线$x_A=Vt_A$、$x_B=Vt_B+L$。远处观察者的方向$\vct n$由物体指向相机，源点到相机的距离差近似为$-\vct n\cdot\Delta\vct r_e=-\cos\theta\Delta x_e$。同一照片要求
\[
 0=\Delta t_o=t_B-t_A-\frac{\cos\theta}{c}\Delta x_e.
\]
再由两端世界线相减，$\Delta x_e=V(t_B-t_A)+L$，消去发射时间差得
\[
 \Delta x_e=\frac{L}{1-\beta\cos\theta}.
\]
\Diagram{fig:m2bara}{发光事件空间差向成像面的投影}{\draw[curve](0,0)--(4,0);\draw[guide](0,0)--(2,2.5);\draw[guide](4,0)--(5.56,1.95);\draw[very thick](2,2.5)--(4.44,.55);\draw[guide](4,0)--(4.44,.55);\node[above right]at(3.4,1.6){$L_{\rm photo}$};\node[below]at(2,0){$\Delta x_e$};}
在图\ref{fig:m2bara}的像平面内，端点分离矢量须扣去沿视线分量，故
\begin{align*}
 L_{\rm photo}^2
 &=|\Delta\vct r_e|^2-(\vct n\cdot\Delta\vct r_e)^2\\
 &=(\Delta x_e)^2\sin^2\theta,
\end{align*}
因此$L_{\rm photo}=L_0\sin\theta/[\gamma(1-\beta\cos\theta)]$。在$\theta=\pi/2$时恢复$L_0/\gamma$。其角度函数与例\ref{ex:m2super}相同，沿用该例已经展开的商式求导，最大值在$\cos\theta=\beta$取得；代入$\sin\theta=1/\gamma$，得$L_{\rm photo,max}=L_0$。式中的$L$、$\Delta x_e$、$L_{\rm photo}$分别表示同时长度、延迟发射位置差和像面投影长，三者须按各自定义使用。
\end{solution}


\chapter{相对论动力学与电磁场}
\section{相位与波四矢量}\label{th:m3wave}
单色平面波的复表示为$\mathcal E=\mathcal E_0\exp[i(\vct k\cdot\vct r-\omega t)]$。同一波面的事件在不同坐标系中仍具有相同相位，所以$\omega t-\vct k\cdot\vct r$是洛伦兹标量。将逆坐标变换代入相位，按$t',x',y',z'$收集系数，得到
\begin{align}\label{eq:m3wave}
 \omega'&=\gamma(\omega-Vk_x),& k_x'&=\gamma(k_x-V\omega/c^2),\\
 k_y'&=k_y,&k_z'&=k_z.\notag
\end{align}
因此$k^\mu=(\omega/c,\vct k)$是波四矢量。真空中$\omega=c|\vct k|$，故$k^2=0$，任何惯性系中仍有真空色散关系。若在介质中使用波矢变换，介质的运动也要一起变换，不能同时假定两系介质都静止且折射率相同。

由$k_x=(\omega/c)\cos\theta$可得$\omega'=\gamma\omega(1-\beta\cos\theta)$。再以$k_x'c/\omega'$求方向，便重现光行差公式。频率和方向应作为同一四矢量的分量一起求出。

\begin{example}\label{ex:m3movingmirror}
一束频率$\nu$的光沿$x$正向照射以速度$V>0$远离的理想平面镜，镜面垂直$x$。忽略镜子反冲。求反射光在实验室中的频率。
\Diagram{fig:m3mirrorq}{运动镜面与入射、反射光}{\draw[curve](4,-1)--(4,1.4);\draw[axis](4.2,1)--(5.5,1)node[above]{$V$};\draw[very thick,->](0,.6)--(3.7,.6)node[midway,above]{$\nu$};\draw[very thick,->](3.7,-.3)--(.2,-.3)node[midway,below]{$\nu_r$};}
\end{example}
\begin{solution}
先把图\ref{fig:m3mirrorq}的入射光变到镜系。入射$k_x=\omega/c$，故$\omega'_i=\gamma\omega(1-\beta)$。镜系中镜面静止，反射保持频率并使$k_x'$反号：$\omega'_r=\omega'_i$，$k'_{xr}=-\omega'_i/c$。
\Diagram{fig:m3mirrora}{镜系反射保持频率，再将反向波变回实验室}{\draw[curve](3,-1)--(3,1.3);\draw[->](0,.5)--(2.8,.5)node[midway,above]{$\omega'_i$};\draw[->](2.8,-.4)--(0,-.4)node[midway,below]{$\omega'_r=\omega'_i$};\node[right]at(3.2,.2){镜系};}
逆变换$\omega_r=\gamma(\omega'_r+Vk'_{xr})$给$\omega_r=\gamma(1-\beta)\omega'_i$。因此图\ref{fig:m3mirrora}对应的两次频率变换合成
\[
 \frac{\nu_r}\nu=\gamma^2(1-\beta)^2=\frac{1-\beta}{1+\beta}.
\]
若镜子靠近，只需令$\beta<0$，频率升高。低速时$\nu_r/\nu\simeq1-2\beta$，因频移发生在入射与返回两段，不是一次多普勒因子。
\end{solution}

\section{总四动量、质心系与阈值}\label{th:m3momentum}
封闭系统的总四动量$P^\mu=\sum_i p_i^\mu$守恒。若$P^2>0$，存在总空间动量为零的质心系，系统不变质量为
\[
 M^2c^4=E_{\rm tot}^2-c^2\vct P^2.
\]
其速度为$\vct V_{\rm cm}=c^2\vct P/E_{\rm tot}$。在该系总能量等于$Mc^2$，其中包含各粒子的相对运动能、内部能及相互作用能；一般$M\ne\sum_i m_i$。讨论散射时取入射和出射粒子相隔很远的状态，可把相互作用能归入粒子静质量及可分辨辐射，直接使用粒子四动量求和。

\subsection{功—能关系与一维加速}
单粒子满足$E^2=p^2c^2+m^2c^4$，故$\dd E=c^2\vct p\cdot\dd\vct p/E=\vct v\cdot\dd\vct p$。由此$\dd E=\vct f\cdot\dd\vct r$，无须先假定牛顿动能。对直线运动，$\dd p/\dd v=m\gamma_v^3$；对瞬时垂直方向的速度微扰，$\dd\vct p_\perp=m\gamma_v\dd\vct v_\perp$。相同三力在两方向产生的实验室加速度不同。

\begin{example}\label{ex:m3photonpair}
两个光子能量为$E_1,E_2$，传播方向夹角$\theta$。求它们组成系统的不变质量；讨论平行、反向两种极限及能否产生静质量$m$的正反粒子对。
\Diagram{fig:m3pairq}{两光子四动量的空间夹角}{\draw[axis](0,0)--(3.6,0)node[right]{$\vct p_1$};\draw[axis](0,0)--(1.8,2.5)node[right]{$\vct p_2$};\draw(.8,0)arc(0:54.25:.8);\node at(1,.5){$\theta$};}
\end{example}
\begin{solution}
图\ref{fig:m3pairq}中$|\vct p_i|=E_i/c$，所以
\begin{align*}
 M^2c^4&=(E_1+E_2)^2-c^2(p_1^2+p_2^2+2p_1p_2\cos\theta)\\
 &=2E_1E_2(1-\cos\theta).
\end{align*}
平行时$M=0$，不存在质心静止系；反向时$Mc^2=2\sqrt{E_1E_2}$。若反向且$E_1=E_2$，实验室就是质心系。
\Diagram{fig:m3paira}{质心系中的粒子对阈值：末态相对动量趋于零}{\draw[axis](-2.7,0)--(2.7,0);\draw[->,thick](-2.5,.6)--(-.3,.6)node[midway,above]{$E'_1$};\draw[->,thick](2.5,.6)--(.3,.6)node[midway,above]{$E'_2$};\fill(-.5,-.5)circle(2pt)node[below]{$m$};\fill(.5,-.5)circle(2pt)node[below]{$m$};}
图\ref{fig:m3paira}所示的最小末态质心能量为$2mc^2$，因而阈值条件为$E_1E_2(1-\cos\theta)\geq2m^2c^4$。这只是能量动量允许的条件，还须满足电荷等其他守恒律和实际相互作用机制。单个自由光子不可能在真空中产生一对有质量粒子，因为其总不变质量为零。
\end{solution}

\begin{example}\label{ex:m3inelastic}
静质量$m$的粒子以动能$K$撞击同样质量的静止粒子，两者完全粘合且不向外辐射。求粘合体的静质量、速度及其内部增加的能量。
\Diagram{fig:m3inelasticq}{实验室中的完全非弹性碰撞}{\fill(0,0)circle(3pt)node[below]{$m$};\draw[axis](.2,.4)--(1.8,.4)node[above]{$K$};\fill(3,0)circle(3pt)node[below]{$m$静止};\draw[->](3.5,0)--(4.5,0);\draw[curve](5,0)circle(.35);\node[below]at(5,-.35){$M$};}
\end{example}
\begin{solution}
图\ref{fig:m3inelasticq}中，入射粒子的总能量为$E_1=mc^2+K$，靶粒子能量为$mc^2$。入射动量由质壳关系给出：
\[
 p^2c^2=E_1^2-m^2c^4=K^2+2mc^2K.
\]
粘合体携带全部四动量，所以$E_T=2mc^2+K$、$P_T=p$，其静质量满足
\begin{align*}
 M^2c^4&=(2mc^2+K)^2-p^2c^2\\
 &=4m^2c^4+2mc^2K.
\end{align*}
故$M=2m\sqrt{1+K/(2mc^2)}$，速度为$V=Pc^2/E_T=c\sqrt{K^2+2mc^2K}/(2mc^2+K)$。
\Diagram{fig:m3inelastica}{质心系内相向运动的动能变为粘合体内部能}{\draw[->](-2,.4)--(-.5,.4);\draw[->](2,.4)--(.5,.4);\fill(-1.5,0)circle(2.5pt);\fill(1.5,0)circle(2.5pt);\draw[curve](0,-1)circle(.35);\node[right]at(.5,-1){$\Delta U=(M-2m)c^2$};}
图\ref{fig:m3inelastica}是在总动量为零的质心系，粘合前的相向动能全部转为内部能，因此
\[
 \Delta U=(M-2m)c^2,
 \qquad K_{\rm bulk}=E_T-Mc^2=K-\Delta U.
\]
两项之和为初始实验室动能$K$。其中$\Delta U$是粘合体自身系的内能增加，$K_{\rm bulk}$是实验室中的整体动能。令$k=K/(2mc^2)$，展开$\sqrt{1+k}=1+k/2-k^2/8+\cdots$，得到
\[
 \Delta U=\frac K2-\frac{K^2}{16mc^2}+\cdots,\qquad
 K_{\rm bulk}=\frac K2+\frac{K^2}{16mc^2}+\cdots.
\]
牛顿极限中两者各占一半。相对论情形不能把所有入射动能都算成内部能，因为总动量守恒要求粘合体继续平动。
\end{solution}

\section{四力、加速度与瞬时静止系}\label{th:m3force}
设实验室粒子速度为$\vct u$，使用$\gamma_u$区别于参考系相对运动的$\gamma$。三力$\vct f=\dd\vct p/\dd t$对应
\[
 F^\mu=\gamma_u(\vct f\cdot\vct u/c,f_x,f_y,f_z).
\]
沿$x$变换后，粒子的新洛伦兹因子$\gamma_{u'}=\gamma\gamma_u(1-Vu_x/c^2)$。以$F'^i/\gamma_{u'}$求三力，得
\begin{align}\label{eq:m3forcetrans}
 f'_x&=\frac{f_x-V(\vct f\cdot\vct u)/c^2}{1-Vu_x/c^2},\\
 f'_{y,z}&=\frac{f_{y,z}}{\gamma(1-Vu_x/c^2)}.\notag
\end{align}
在瞬时静止系中$u'=0$，四力时间分量为零，$F^2=-f_0^2$。这里$f_0$专指瞬时静止系三力的大小，固定惯性系中的三力须另作变换。

四加速度需同时求导$\gamma_u$与$\vct u$：
\[
 \dot\gamma_u=-\frac12(1-u^2/c^2)^{-3/2}
 \left(-\frac{2\vct u\cdot\vct a}{c^2}\right)
 =\gamma_u^3\frac{\vct u\cdot\vct a}{c^2}.
\]
再用$\dd/\dd\tau=\gamma_u\dd/\dd t$对$u^\mu=\gamma_u(c,\vct u)$逐项求导，得
\[
 a^0=\gamma_u^4\frac{\vct u\cdot\vct a}{c},\quad
 \vct a_{(4)}=\gamma_u^2\vct a+\gamma_u^4\frac{\vct u\cdot\vct a}{c^2}\vct u.
\]
沿瞬时速度方向分解，四加速度的纵、横空间分量分别为$\gamma_u^4a_\parallel$、$\gamma_u^2a_\perp$，时间分量为$\gamma_u^4ua_\parallel/c$。因此
\[
 \alpha^2=-a_\mu a^\mu
 =\gamma_u^8(1-u^2/c^2)a_\parallel^2+\gamma_u^4a_\perp^2
 =\gamma_u^6a_\parallel^2+\gamma_u^4a_\perp^2.
\]
恒静质量时$F^\mu=ma^\mu$，这也给出四力模与固有加速度的关系。

\begin{example}\label{ex:m3accel}
粒子在实验室瞬时以速度$u$沿$x$运动，其瞬时静止系加速度分量为$(\alpha_x,\alpha_y,0)$。从四加速度变换推导实验室两分量。
\Diagram{fig:m3accelq}{瞬时静止系中的加速度分解}{\draw[axis](0,0)--(3.8,0)node[right]{$x_0$};\draw[axis](0,0)--(0,2.8)node[above]{$y_0$};\draw[->,thick](0,0)--(2.4,1.8);\draw[guide](2.4,0)--(2.4,1.8)--(0,1.8);\node[below]at(1.2,0){$\alpha_x$};\node[left]at(0,.9){$\alpha_y$};}
\end{example}
\begin{solution}
把四加速度记为$A^\mu=\dd U^\mu/\dd\tau$，实验室三加速度记为$\vct a$，以免同一字母混淆。图\ref{fig:m3accelq}的瞬时静止系中，$A_*^\mu=(0,\alpha_x,\alpha_y,0)$。沿$x$逆变换得到
\[
 A^0=\gamma_u\beta_u\alpha_x,\quad
 A^x=\gamma_u\alpha_x,\quad A^y=\alpha_y.
\]
这里变换的是同一事件处的四矢量，不把该瞬时静止系延伸为整个加速过程的静止系。
\Diagram{fig:m3accela}{实验室加速度的纵、横分量分别多出不同的时间因子}{\draw[axis](0,0)--(3.8,0)node[right]{$x,\ u$};\draw[axis](0,0)--(0,2.8)node[above]{$y$};\draw[->,thick](0,0)--(1.3,1.8);\draw[guide](1.3,0)--(1.3,1.8)--(0,1.8);\node[below]at(1.3,0){$\alpha_x/\gamma_u^3$};\node[left]at(0,1){$\alpha_y/\gamma_u^2$};}
为将图\ref{fig:m3accela}中的$A^i$还原为三加速度，先对$\gamma_u$求导：
\[
 \frac{\dd\gamma_u}{\dd t}=\gamma_u^3\frac{\vct u\cdot\vct a}{c^2}
 =\gamma_u^3\frac{ua_x}{c^2}.
\]
因$\dd/\dd\tau=\gamma_u\dd/\dd t$，该事件处$u_y=0$，所以
\begin{align*}
 A^x&=\gamma_u\frac{\dd(\gamma_u u_x)}{\dd t}
 =\gamma_u^2a_x+\gamma_u^4\frac{u^2a_x}{c^2}
 =\gamma_u^4a_x,\\
 A^y&=\gamma_u\frac{\dd(\gamma_u u_y)}{\dd t}=\gamma_u^2a_y.
\end{align*}
与变换所得分量比较，得到$a_x=\alpha_x/\gamma_u^3$、$a_y=\alpha_y/\gamma_u^2$。时间分量$A^0=\gamma_u^4ua_x/c$也与上式一致。纵向多出的一个因子来自速度变化引起的$\gamma_u$变化，不能把横向时间膨胀的推算直接套到纵向。
\end{solution}

\begin{example}\label{ex:m3longfield}
均匀电场沿$x$，带电粒子瞬时也沿$x$运动。用瞬时静止系三力和加速度关系说明为何平行电场不变；再求实验室加速度。
\Diagram{fig:m3longq}{平行电场中的带电粒子}{\foreach\y in{0,.6,1.2}{\draw[axis](0,\y)--(4,\y);}\node[right]at(4,.6){$E_x$};\fill(1.5,.6)circle(3pt)node[above]{$q,m$};\draw[axis](1.7,1.6)--(3,1.6)node[above]{$u$};}
\end{example}
\begin{solution}
图\ref{fig:m3longq}中实验室$f_x=qE_x$。把式\eqref{eq:m3forcetrans}取$V=u$且$\vct f\cdot\vct u=f_xu$，得$f_{0x}=f_x$。粒子在瞬时静止系受到电力$qE_{0x}$，故$qE_{0x}=qE_x$，即$E_{0x}=E_x$。
\Diagram{fig:m3longa}{相同平行电力对应不同的坐标加速度}{\draw[axis](0,0)--(4.3,0)node[right]{$u/c$};\draw[axis](0,0)--(0,2.8)node[above]{$a_x/(qE_x/m)$};\draw[curve,domain=0:1,samples=70]plot({3.5*\x},{2.3*(1-\x*\x)^1.5});\node[left]at(0,2.3){$1$};\node[below]at(3.5,0){$1$};}
瞬时静止系$\alpha_x=qE_x/m$，由例\ref{ex:m3accel}，图\ref{fig:m3longa}的实验室加速度$a_x=qE_x/(m\gamma_u^3)$。随着速度增加，电力保持不变，加速度减小；能量输入率仍为$qE_xu$，并不因加速度变小而趋于零。
\end{solution}

\section{四电流与电荷守恒}\label{th:m3current}
对于共同运动的带电粒子群，静止系密度为$\rho_0$，实验室电荷密度$\rho=\gamma_u\rho_0$，电流$\vct j=\rho\vct u$，因此
\[
 J^\mu=(c\rho,\vct j)=\rho_0u^\mu,\quad J^2=c^2\rho_0^2.
\]
多种不同速度的带电成分应分别构造四电流再求和；总四电流未必类时，也未必存在把所有电荷同时静止的惯性系。连续性方程
\[
 \frac{\partial\rho}{\partial t}+\nabla\cdot\vct j=0
\]
可以写为$\partial_\mu J^\mu=0$，其积分形式是某固定体积内电荷减少等于流出边界的电荷。沿$x$的变换给$\rho'=\gamma(\rho-Vj_x/c^2)$、$j'_x=\gamma(j_x-V\rho)$，横向电流分量不变。

\begin{example}\label{ex:m3wire}
无限直导线沿$x$，实验室中正离子静止，线电荷密度$+\lambda$；电子速度为$-u$，电子实验室线密度为$-\lambda$。求沿$+V$运动的观察系中的总线电荷密度和电流。规定常规电流沿正向。
\Diagram{fig:m3wireq}{实验室中电中性的载流导线}{\draw[curve](-.3,0)--(5.8,0);\foreach\x in{0,1,2,3,4,5}{\node at(\x,.4){$+$};\node at(\x,-.4){$-$};}\draw[axis](3,-.9)--(1.5,-.9)node[midway,below]{$u$};\draw[axis](1,1)--(2.5,1)node[midway,above]{$I=\lambda u$};}
\end{example}
\begin{solution}
图\ref{fig:m3wireq}的总$\lambda_T=0$，$I=\lambda u$。对单位横截面的积分四电流同样变换，得到
\[
 \lambda_T'=-\gamma\frac{VI}{c^2},\qquad I'=\gamma I.
\]
也可逐成分求解：离子$\lambda'_+=\gamma\lambda$；电子$\lambda'_- =\gamma[-\lambda-V(\lambda u)/c^2]$。两者相加即得上式。
\Diagram{fig:m3wirea}{运动观察系中负电荷密度的绝对值大于正电荷密度}{\draw[curve](0,0)--(5.8,0);\foreach\x in{.3,1.3,2.3,3.3,4.3,5.3}{\node at(\x,.4){$+$};}\foreach\x in{0,.7,1.4,2.1,2.8,3.5,4.2,4.9,5.6}{\node at(\x,-.4){$-$};}\node[below]at(3,-1){$\lambda'_T<0$，示意密度};}
图\ref{fig:m3wirea}说明电中性是与参考系有关的密度条件，整段材料上包含多少总电荷则还涉及边界和同时截面。若直接把离子和电子都缩短相同倍数，便错误地忽略了电子在实验室已经运动以及密度变换中的电流项。
\end{solution}

\section{四势、规范与场方程}\label{th:m3potential}
由$\nabla\cdot\vct B=0$，局部可写$\vct B=\nabla\times\vct A$。法拉第定律给$\nabla\times(\vct E+\partial_t\vct A)=0$，所以$\vct E=-\nabla\phi-\partial_t\vct A$。势不唯一：$\vct A\to\vct A+\nabla\zeta$、$\phi\to\phi-\partial_t\zeta$保持场不变。

把势代入高斯定律，得到$-\nabla^2\phi-\partial_t\nabla\cdot\vct A=\rho/\varepsilon_0$。取洛伦兹规范
\[
 \nabla\cdot\vct A+\frac1{c^2}\partial_t\phi=0,
\]
就有$(c^{-2}\partial_t^2-\nabla^2)\phi=\rho/\varepsilon_0$。再把势代入安培—麦克斯韦定律，利用$\nabla\times\nabla\times\vct A=\nabla(\nabla\cdot\vct A)-\nabla^2\vct A$，梯度项由规范条件抵消，得到
\[
 \left(\frac1{c^2}\partial_t^2-\nabla^2\right)\vct A=\mu_0\vct j.
\]
令$A^\mu=(\phi/c,\vct A)$、$\Box=c^{-2}\partial_t^2-\nabla^2$，两式合并为$\Box A^\mu=\mu_0J^\mu$。在保持此协变规范的选择下，四势按四矢量变换。若另行改用其他规范，还须补相应规范变换，不能把某个规范下的分量形式当作场本身。

\subsection{导数变换}
由链式法则，$\partial_{ct'}=\gamma(\partial_{ct}+\beta\partial_x)$、$\partial_{x'}=\gamma(\partial_x+\beta\partial_{ct})$，$\partial_{y',z'}=\partial_{y,z}$。特别注意导数的正负号与坐标正变换不同，因为这里是对逆坐标式求导。将两导数组合平方，交叉项抵消，就可验证$\Box'=\Box$。

\section{电磁场的完整变换}\label{th:m3fields}
对四势作沿$x$的变换，$\phi'=\gamma(\phi-VA_x)$、$A_x'=\gamma(A_x-V\phi/c^2)$、$A_y'=A_y$、$A_z'=A_z$。例如
\begin{align*}
 E_y'&=-\partial_y\phi'-\partial_{t'}A_y'\\
 &=-\gamma\partial_y\phi+\gamma V\partial_y A_x
 -\gamma\partial_tA_y-\gamma V\partial_xA_y\\
 &=\gamma(E_y-VB_z).
\end{align*}
这里$B_z=\partial_xA_y-\partial_yA_x$，负号由叉乘方向决定。同样将$E_x'$完整展开，时间与空间交叉导数抵消，剩下$\gamma^2(1-\beta^2)E_x=E_x$。对$\vct B'=\nabla'\times\vct A'$逐分量计算，得到
\begin{equation}\label{eq:m3field}
 \begin{aligned}
 E_x'&=E_x,&B_x'&=B_x,\\
 E_y'&=\gamma(E_y-VB_z),&B_y'&=\gamma(B_y+VE_z/c^2),\\
 E_z'&=\gamma(E_z+VB_y),&B_z'&=\gamma(B_z-VE_y/c^2).
 \end{aligned}
\end{equation}
相对任意$\vct V$的矢量形式是
\[
 \vct E'_\parallel=\vct E_\parallel,\quad
 \vct E'_\perp=\gamma(\vct E+\vct V\times\vct B)_\perp,
\]
\[
 \vct B'_\parallel=\vct B_\parallel,\quad
 \vct B'_\perp=\gamma(\vct B-\vct V\times\vct E/c^2)_\perp.
\]
从分量式可直接验证$E^2-c^2B^2$和$\vct E\cdot\vct B$不变。两个不变量可以排除某些变换目标，但具体参考系速度仍须代入变换式求出。

\begin{example}\label{ex:m3crossed}
实验室中$\vct E=E\hat{\vct y}$、$\vct B=B\hat{\vct z}$，$E,B>0$。求能使电场或磁场消失的沿$x$运动惯性系，并讨论$E=cB$。
\Diagram{fig:m3crossq}{正交电磁场与可选的变换方向}{\draw[axis](0,0)--(3.6,0)node[right]{$x$};\draw[axis](0,0)--(0,2.5)node[above]{$E\hat{\vct y}$};\draw(1.5,1)circle(.18);\fill(1.5,1)circle(1.2pt);\node[right]at(1.8,1){$B\hat{\vct z}$};}
\end{example}
\begin{solution}
图\ref{fig:m3crossq}中$E_y'=\gamma(E-VB)$、$B_z'=\gamma(B-VE/c^2)$。若$E<cB$，取$V=E/B<c$，可令$E'=0$。此时$B'=\sqrt{B^2-E^2/c^2}$，仍沿$z$正向。
\Diagram{fig:m3crossa}{由不变量符号区分磁占优、电占优和光状场}{\draw[axis](0,0)--(5,0)node[right]{$E/(cB)$};\draw[curve](2.5,-.6)--(2.5,.8);\node[below]at(2.5,-.6){$1$};\node[above]at(1,.2){可令$E'=0$};\node[above]at(4,.2){可令$B'=0$};}
若$E>cB$，取$V=c^2B/E<c$，得到$B'=0$和$E'=\sqrt{E^2-c^2B^2}$。图\ref{fig:m3crossa}的分界$E=cB$要求$V=c$才能消场，不是合法惯性系。此时两个不变量均为零，任何有限速度系中仍有非零正交场且$E'=cB'$，真空平面波正属于这类情况。
\end{solution}

\begin{example}\label{ex:m3wavefield}
真空平面波沿$x$传播，电场沿$y$，磁场沿$z$。沿波方向以速度$V$运动的观察者测得何种频率、振幅与光强？
\Diagram{fig:m3waveq}{沿波传播方向追赶平面波}{\draw[axis](-.2,0)--(5.5,0)node[right]{$\vct k,V$};\draw[curve,domain=0:5,samples=100]plot(\x,{.7*sin(180*\x)});\node[above]at(2,1){$E_y$};\draw[axis](3,-1)--(4.5,-1);}
\end{example}
\begin{solution}
图\ref{fig:m3waveq}中$B_z=E_y/c$。式\eqref{eq:m3field}给$E_y'=\gamma(1-\beta)E_y$、$B_z'=E_y'/c$。频率同样有$\omega'=\gamma(1-\beta)\omega$。
\Diagram{fig:m3wavea}{振幅和频率同比例红移，波长相应伸长}{\draw[axis](0,0)--(5.6,0)node[right]{$x'$};\draw[curve,domain=0:5,samples=90]plot(\x,{.45*sin(100*\x)});\draw[<->](.9,-.9)--(4.5,-.9)node[midway,below]{$\lambda'=\lambda/[\gamma(1-\beta)]$};}
时间平均光强$I=\varepsilon_0cE_0^2/2$，所以图\ref{fig:m3wavea}中的$I'/I=\gamma^2(1-\beta)^2=(1-\beta)/(1+\beta)$。频率比和振幅比相同，而光强比是其平方。此结果对整个无限平面波的局部能流成立，与点源角功率的立体角雅可比属于不同问题。
\end{solution}
\section{守恒方程中的未知量与参考系}
\subsection{质心系的速度和能量}
总四动量为$(E/c,\vct P)$，若$E^2-c^2P^2>0$，令参考系速度沿$\vct P$，大小$V=c^2P/E$。变换后的平行动量为$P'=\gamma(P-VE/c^2)=0$。相应$\gamma=E/\sqrt{E^2-c^2P^2}$，新系能量
\[
 E'=\gamma(E-VP)=\sqrt{E^2-c^2P^2}=Mc^2.
\]
总四动量为类时时，质心系速度由它唯一确定。同向光子束的总四动量为类光，形式上需要$V=c$，因而没有质心惯性系。固定靶实验的实验室一般有非零总动量，阈值时末态仍须携带整体动能。

\subsection{两体末态的通用解}
在质心系中，两末态粒子的动量必为$\vct p_*,-\vct p_*$。设总能量为$W$，质量$m_1,m_2$，则$E_{1*}+E_{2*}=W$，并有$E_{1*}^2-E_{2*}^2=(m_1^2-m_2^2)c^4$。将左侧因式分解，得到
\[
 E_{1*}=\frac{W^2+(m_1^2-m_2^2)c^4}{2W},\quad
 E_{2*}=\frac{W^2+(m_2^2-m_1^2)c^4}{2W}.
\]
再由$E_{1*}^2-p_*^2c^2=m_1^2c^4$，得到
\[
 p_*c=\frac{\sqrt{[W^2-(m_1+m_2)^2c^4]
 [W^2-(m_1-m_2)^2c^4]}}{2W}.
\]
末态粒子各有正能量$E_{i*}=\sqrt{m_i^2c^4+p_*^2c^2}\geq m_ic^2$，所以物理分支必须满足$W\geq(m_1+m_2)c^2$，这才是两体阈值。仅要求根号为实还可能保留$W\leq|m_1-m_2|c^2$的代数分支，它不满足上述两个正能量条件。在阈值极限$p_*=0$，末态没有相对运动，但整体在实验室仍可运动。内部激发应计入$m_i$，不能在公式外再次添加同一激发能。

\Diagram{fig:m3twobody}{两体质心系中动量等大反向，能量未必相等}{\draw[axis](-3,0)--(3.2,0);\fill(-1,0)circle(2pt)node[below]{$m_1$};\fill(1,0)circle(2pt)node[below]{$m_2$};\draw[->,thick](-1,.5)--(-2.7,.5)node[above]{$p_*$};\draw[->,thick](1,.5)--(2.7,.5)node[above]{$p_*$};}
图\ref{fig:m3twobody}给出反向动量，不要求两个粒子以相同速率运动。若$m_1\ne m_2$，能量也不同。得到$E_{i*},p_*$后，再按题目给定出射角变回实验室，通常比直接在实验室同时联立角度、速度、能量更清楚。

\subsection{相对速度与入射能量不是同一输入}
在两个相同粒子的对撞中，某系观察到每束速率为$v$，彼此相对速率是$2v/(1+v^2/c^2)$。质心能量却是$2\gamma_vmc^2$，不能把相对速率直接代入“每束能量”公式。相对速率属于一粒子的静止系，束能量属于实验室，所选参考系必须在整个守恒式中一致。

同样，“完全非弹性”只规定末态相对运动消失，不规定静质量守恒。动能变为热或激发后，总四动量仍守恒，而复合物静质量增加。若相互作用放出辐射，须把辐射四动量留在末态，复合体的质量和速度会相应改变；不能在不封闭的子系统中同时要求能量、动量各自不变。

\section{四势推导中的全部横向分量}
\subsection{磁场变换的一项完整展开}
以$B_z'=\partial_{x'}A_y'-\partial_yA_x'$为例，代入导数和势的变换：
\begin{align*}
 B_z'&=\gamma(\partial_x+V\partial_t/c^2)A_y
 -\gamma\partial_y(A_x-V\phi/c^2)\\
 &=\gamma(\partial_xA_y-\partial_yA_x)
 +\frac{\gamma V}{c^2}(\partial_tA_y+\partial_y\phi)\\
 &=\gamma(B_z-VE_y/c^2).
\end{align*}
每个括号都是原系已定义的场，不需要假设场静态。对$B_y'=\partial_zA_x'-\partial_{x'}A_z'$，同样给$\gamma(B_y+VE_z/c^2)$；符号变化来自旋度的分量次序。$B_x'=\partial_yA_z'-\partial_zA_y'=B_x$，因为这两势分量和这两个导数都未变。

纵向电场中抵消较多，完整写成
\begin{align*}
 E_x'&=-\gamma^2(\partial_x+V\partial_t/c^2)(\phi-VA_x)\\
 &\quad-\gamma^2(\partial_t+V\partial_x)(A_x-V\phi/c^2).
\end{align*}
展开后，$V\partial_xA_x$和$V\partial_t\phi/c^2$的交叉项分别抵消，余下$-\gamma^2(1-\beta^2)(\partial_x\phi+\partial_tA_x)=E_x$。若漏掉$\partial_{t'}$中的空间导数项，就会错误地产生纵向场也乘$\gamma$的结果。

\subsection{场不变量的代数核对}
沿$x$推动时，纵向$E_x^2-c^2B_x^2$不变。将横向组合为$E_y'^2-c^2B_z'^2$，代入式\eqref{eq:m3field}：
\begin{align*}
 E_y'^2-c^2B_z'^2
 &=\gamma^2[(E_y-VB_z)^2-c^2(B_z-VE_y/c^2)^2]\\
 &=\gamma^2(1-\beta^2)(E_y^2-c^2B_z^2)\\
 &=E_y^2-c^2B_z^2.
\end{align*}
另一对$(E_z,B_y)$相同，三项相加就得$E'^2-c^2B'^2=E^2-c^2B^2$。点积中的横向交叉项也成对抵消，得到$\vct E'\cdot\vct B'=\vct E\cdot\vct B$。

非零真空平面波的两个场不变量均为零，其能量密度和坡印廷矢量仍非零。有限洛伦兹变换可逆，所以全零场只能由全零场变来；非零场至多在满足条件时消去其中一种场。

\section{麦克斯韦方程与电荷连续性}
安培—麦克斯韦方程为$\nabla\times\vct B=\mu_0\vct j+c^{-2}\partial_t\vct E$。取散度，左侧旋度的散度恒为零，右侧用$\nabla\cdot\vct E=\rho/\varepsilon_0$，得$0=\mu_0\nabla\cdot\vct j+\mu_0\partial_t\rho$。因此连续性方程是场方程相容的要求，位移电流项对变化电荷必不可少。

采用$\mathcal F^{0i}=-E_i/c$、$\mathcal F^{i0}=E_i/c$、$\mathcal F^{ij}=-\epsilon_{ijk}B_k$，其中$i,j,k=1,2,3$，完整矩阵及洛伦兹力的分量核对见例\ref{ex:m3fourtensor}。含源方程为$\partial_\nu\mathcal F^{\nu\mu}=\mu_0J^\mu$。取$\mu=0$，空间散度给$\nabla\cdot\vct E/c=\mu_0c\rho$，就是高斯定律；取空间分量，得到$\nabla\times\vct B-c^{-2}\partial_t\vct E=\mu_0\vct j$。再缩并$\partial_\mu$，由于$\partial_\mu\partial_\nu$对称而场张量反对称，左侧恒零，故$\partial_\mu J^\mu=0$。指标次序须与场张量定义一致。

\subsection{规范自由度与实际场}
洛伦兹规范下$\Box A^\mu=\mu_0J^\mu$，规范函数若满足$\Box\zeta=0$，改变势后仍保持同一规范。因此即使固定洛伦兹规范，势也未完全唯一。可观测场由势的导数组合决定，某个分量为零不能单独作为“没有电场”或“没有磁场”的证据。

对局域源，场方程的物理解还需选择边界条件，例如不从无限远无缘无故入射的辐射，以及对过去源历史的延迟响应。仅验证某个势满足波动方程，并不能说明它已满足实际初始、边界条件。匀速电荷场、加速电荷辐射场、外加平面波虽然都遵循同一方程，其边界数据不同。

\subsection{场的能量与参考系}
真空电磁能量密度为$u_{\rm em}=\varepsilon_0E^2/2+B^2/(2\mu_0)$，能流$\vct S=\vct E\times\vct B/\mu_0$。纯静电场的能流可以为零，变到运动系出现磁场后能流一般非零，这与能量和动量共同变换一致。不能只变换能量密度而保留原系“动量为零”的判断。

\Diagram{fig:m3energyflow}{运动电荷场的电磁能流与场的角分布相联系}{\draw[axis](-2.3,0)--(2.7,0)node[right]{$V$};\fill(0,0)circle(2pt);\draw[axis](0,1.5)--(0,2.5)node[above]{$E_y$};\draw(.5,1.6)circle(.14);\fill(.5,1.6)circle(1pt);\draw[axis](.6,2.3)--(2,2.3)node[right]{$\vct S$};}
图\ref{fig:m3energyflow}在电荷横向上方给$E_y$、$B_z$，所以能流沿正$x$。局部能量守恒为$\partial_tu_{\rm em}+\nabla\cdot\vct S=-\vct j\cdot\vct E$。其中$\vct j\cdot\vct E$是场对物质做功的功率密度，右侧的负号表示这部分能量从场中流失；磁力不直接做功。不同参考系可有不同的电力、磁力分量，总能量动量守恒仍成立。

\subsection{电磁场不变量怎样限制可选参考系}
若某系中电场为零，则$\vct E\cdot\vct B=0$且$c^2B^2-E^2>0$；若磁场为零，则第一不变量仍为零，而第二不变量小于零。这些是必要条件，也给出正交场中的构造方法。设$\vct E\perp\vct B$，沿$\vct E\times\vct B$方向取速度，场变换中平行速度的分量均为零，横向分量为
\[
 \vct E'=\gamma(\vct E+\vct V\times\vct B),\quad
 \vct B'=\gamma(\vct B-\vct V\times\vct E/c^2).
\]
磁占优时取$\vct V=\vct E\times\vct B/B^2$。矢量恒等式$(\vct E\times\vct B)\times\vct B=-B^2\vct E$使$\vct E'=0$，而$V=E/B<c$。电占优时取$\vct V=c^2\vct E\times\vct B/E^2$，得到$\vct B'=0$，且$V=c^2B/E<c$。

\Diagram{fig:m3fieldclasses}{正交场沿电磁能流方向选择变换；等强边界不能由有限惯性变换静止化}{\draw[axis](0,0)--(4,0)node[right]{$\vct E\times\vct B$};\draw[->,thick](0,0)--(0,2.5)node[above]{$\vct E$};\draw(1.5,1.6)circle(.16);\fill(1.5,1.6)circle(1.5pt);\node[right]at(1.8,1.6){$\vct B$朝外};\draw[->](.5,-.6)--(3,-.6)node[midway,below]{$\vct V$};}
图\ref{fig:m3fieldclasses}的方向构造还可从能流理解：正交场的坡印廷矢量沿$\vct E\times\vct B$。但$E=cB\ne0$时，两种尝试均要求$V=c$，没有可消去单个场的惯性系；平面电磁波属于这一类零不变量场。两个不变量都为零并不表示场本身为零。

当$\vct E\cdot\vct B\ne0$时，任一场在任何惯性系中都非零。在电、磁场平行的系中，两场的有向大小仍须同时满足两个不变量；确定场型后，再求相应的亚光速变换速度。

\section{从三力到轨迹}
对$\vct p=m\gamma_u\vct u$求导，可把三力写成
\[
 \vct f=m\gamma_u\vct a+m\gamma_u^3
 \frac{\vct u\cdot\vct a}{c^2}\vct u.
\]
平行于瞬时速度的投影给$f_\parallel=m\gamma_u^3a_\parallel$，垂直投影给$f_\perp=m\gamma_ua_\perp$。再反解为
\[
 \vct a=\frac1{m\gamma_u}
 \left[\vct f-\frac{\vct f\cdot\vct u}{c^2}\vct u\right].
\]
这个矢量式比引入方向依赖的“质量”更直接：静质量$m$是同一不变量，方向差异出自动量与速度的非线性关系。

\begin{example}\label{ex:m3transconstant}
粒子初始动量为$p_0\hat{\vct x}$，从$t=0$起受固定方向恒定三力$f\hat{\vct y}$。求速度和轨迹，说明为何它一般不是匀速圆周运动。
\Diagram{fig:m3transq}{固定方向力只在初始瞬间垂直速度}{\draw[axis](0,0)--(3.7,0)node[right]{$p_0$};\draw[axis](0,0)--(0,2.6)node[above]{$f$};\fill(0,0)circle(2pt);}
\end{example}
\begin{solution}
图\ref{fig:m3transq}中，三力固定沿$y$，故$\dd p_x/\dd t=0$、$\dd p_y/\dd t=f$，初值给$p_x=p_0$、$p_y=ft$。记$E_a=\sqrt{m^2c^4+p_0^2c^2}$，由质壳关系
\[
 E=\sqrt{E_a^2+f^2c^2t^2},\qquad
 u_x=\frac{c^2p_0}{E},\quad u_y=\frac{c^2ft}{E}.
\]
令$s=fct/E_a$，位移的两个定积分分别为
\begin{align*}
 x&=\frac{cp_0}f\int_0^s\frac{\dd q}{\sqrt{1+q^2}}
 =\frac{cp_0}f[\operatorname{arsinh}q]_0^s
 =\frac{cp_0}f\operatorname{arsinh}s,\\
 y&=\frac{E_a}f\int_0^s\frac{q\dd q}{\sqrt{1+q^2}}
 =\frac{E_a}f[\sqrt{1+q^2}]_0^s
 =\frac{E_a}f(\sqrt{1+s^2}-1).
\end{align*}
\Diagram{fig:m3transa}{动量沿直线增长，空间轨迹逐渐转向力的方向}{\draw[axis](0,0)--(4,0)node[right]{$x$};\draw[axis](0,0)--(0,4.1)node[above]{$y$};\draw[curve,domain=0:3.5,samples=80]plot({ln(\x+sqrt(1+\x*\x))},{sqrt(1+\x*\x)-1});\draw[->](1.8,2)--(1.95,2.6);}
当$p_0\ne0$时消去$s$，空间轨迹为
\[
 y=\frac{E_a}f\left[\cosh\left(\frac{fx}{cp_0}\right)-1\right].
\]
图\ref{fig:m3transa}中的曲线逐渐转向力的方向。由上式$E-E_a=fy$，可直接核对恒力做功；只在初始瞬间力与速度正交，随后$fu_y>0$，速率增大，因此不是匀速圆周运动。短时展开给$x\simeq c^2p_0t/E_a$、$y\simeq fc^2t^2/(2E_a)$，初始横向加速度为$f/(m\gamma_{u0})$。若$p_0=0$，粒子始终沿$y$直线运动，含$1/p_0$的消参式不适用，应使用原来的参数解。
\end{solution}

\begin{example}\label{ex:m3massgain}
质量$M$的静止物体完全吸收一个能量$E$的光子，不向外辐射。求吸收后物体静质量与速度；为什么不能把其静质量简单写成$M+E/c^2$？
\Diagram{fig:m3gainq}{吸收一个光子同时带入能量和动量}{\draw[->,thick](-2,.5)--(0,.5)node[midway,above]{$E$};\draw[curve](.3,0)rectangle(2,1);\node at(1.15,.5){$M$};}
\end{example}
\begin{solution}
图\ref{fig:m3gainq}中，入射光子动量为$E/c$，吸收后物体必须反冲。全系统总能量、动量分别为$E_T=Mc^2+E$、$P_T=E/c$，故
\[
 M_f^2c^4=E_T^2-c^2P_T^2=M^2c^4+2Mc^2E.
\]
令$\epsilon=E/(Mc^2)$，得到$M_f=M\sqrt{1+2\epsilon}$、$V=c\epsilon/(1+\epsilon)$。后式由$V=P_Tc^2/E_T$求出，不能把总动量除以末静质量当作速度。
\Diagram{fig:m3gaina}{光能一部分成为物体静能，一部分成为整体反冲动能}{\draw[curve](0,0)rectangle(5,1);\draw[curve](3.4,0)--(3.4,1);\node at(1.7,.5){$(M_f-M)c^2$};\node at(4.2,.5){$K_f$};\node[above]at(2.5,1){总输入$E$};}
图\ref{fig:m3gaina}的两部分分别为
\[
 \Delta U=Mc^2(\sqrt{1+2\epsilon}-1),\quad
 K_f=Mc^2(1+\epsilon-\sqrt{1+2\epsilon}).
\]
后一括号可有理化为$\epsilon^2/[1+\epsilon+\sqrt{1+2\epsilon}]$，因此正且在低能时为$\epsilon^2/2$。于是$\Delta U\simeq E-E^2/(2Mc^2)$、$K_f\simeq E^2/(2Mc^2)$，总和严格为$E$。

若错误地取$M_f=M+E/c^2$，便已把全部光能算入静能，却又保留非零末动量，末态总能量会超过初态。两束等能反向光子同时吸收时总动量为零，才可把全部输入能量计为静能。对吸收所致的变质量过程，$u\cdot F=c^2\dd M/\dd\tau$一般非零，与静质量不变点电荷的洛伦兹四力应区别。
\end{solution}

\section{电荷、电流与场的同一变换}
\begin{example}\label{ex:m3movingcharge}
点电荷$q$沿$x$以$V$匀速运动，在$t=0$经过原点。由静止库仑场求实验室时刻$t$、位置$(x,y,z)$的电磁场，描述纵向与横向方向的差别。
\Diagram{fig:m3chargeq}{相对电荷同时位置的矢量$\vct R$}{\draw[axis](-1,0)--(4.5,0)node[right]{$x$};\fill(1,0)circle(2pt)node[below]{$q,\ Vt$};\fill(3,1.8)circle(1.5pt)node[right]{$P$};\draw[->](1,0)--(3,1.8)node[midway,above]{$\vct R$};\draw[axis](1,.4)--(2,.4)node[above]{$V$};}
\end{example}
\begin{solution}
取图\ref{fig:m3chargeq}的$R_x=x-Vt$、$\vct R_\perp=(y,z)$。变到电荷静止系，场点相对电荷坐标为$(\gamma R_x,y,z)$，距离$r'=\sqrt{\gamma^2R_x^2+R_\perp^2}$，库仑场$\vct E'=q(\gamma R_x,y,z)/(4\pi\varepsilon_0r'^3)$、$\vct B'=0$。逆场变换保持平行分量，横向分量乘$\gamma$，所以
\[
 \vct E=\frac{q\gamma\vct R}{4\pi\varepsilon_0(\gamma^2R_x^2+R_\perp^2)^{3/2}},\quad
 \vct B=\frac{\vct V\times\vct E}{c^2}.
\]
\Diagram{fig:m3chargea}{等距球面上横向电场增强，纵向减弱}{\draw[axis](-2.4,0)--(2.7,0)node[right]{$V$};\draw[guide](0,0)circle(1.5);\draw[axis](0,1.5)--(0,2.7)node[above]{$\gamma E_C$};\draw[axis](1.5,0)--(2,0);\node[below]at(1.8,-.3){$E_C/\gamma^2$};\fill(0,0)circle(2pt);}
沿纵向$R_\perp=0$，$E=E_C/\gamma^2$；纯横向$R_x=0$，$E=\gamma E_C$，其中$E_C=q/(4\pi\varepsilon_0R^2)$。图\ref{fig:m3chargea}说明角分布被压向横向。场沿同时位置的$\vct R$指向并不意味着作用瞬时传播；这是匀速轨迹的完整延迟场化简结果，电荷加速时不能继续照用。
\end{solution}

\begin{example}\label{ex:m3wireforce}
实验室电中性无限直导线沿$x$，电流$I>0$。测试电荷$q$位于$y=r>0$并以$V\hat{\vct x}$运动。求实验室和电荷瞬时静止系中的横向力，用导线的变换电荷密度解释后一结果。
\Diagram{fig:m3wireforceq}{导线上方的磁场沿$+z$，正电荷受力向下}{\draw[axis](-2,0)--(3,0)node[right]{$I$};\fill(0,1.7)circle(2pt)node[left]{$q$};\draw[axis](0,1.7)--(1.6,1.7)node[above]{$V$};\draw(.5,.8)circle(.12);\fill(.5,.8)circle(.8pt);\node[right]at(.8,.8){$B_z$};}
\end{example}
\begin{solution}
图\ref{fig:m3wireforceq}中$B_z=\mu_0I/(2\pi r)$、$E_y=0$，实验室$f_y=-qVB_z$。变到电荷系，$E_y'=-\gamma VB_z$，电荷静止只受电力，故$f_y'=qE_y'=-\gamma qVB_z=\gamma f_y$，与横向三力变换一致。
\Diagram{fig:m3wireforcea}{电荷系导线带负线密度，产生向导线的电场}{\draw[curve](-2,0)--(3,0);\foreach\x in{-1.8,-1.1,-.4,.3,1,1.7,2.4}{\node at(\x,.25){$-$};}\fill(0,1.7)circle(2pt);\draw[axis](0,1.5)--(0,.5)node[right]{$E_y'<0$};}
由四电流变换，$\lambda'=-\gamma VI/c^2$。无限线电荷在图\ref{fig:m3wireforcea}处电场为
\[
 E_y'=\lambda'/(2\pi\varepsilon_0r)=-\gamma V\mu_0I/(2\pi r),
\]
与场变换完全相同。这个模型把实验室的纯磁力解释为电荷系的电力，但不意味着可以在所有空间点、所有电荷速度下选同一个系消去全部磁场。
\end{solution}

\begin{example}\label{ex:m3drift}
正交均匀场$\vct E=E\hat{\vct y}$、$\vct B=B\hat{\vct z}$满足$E<cB$。求可匀速直线通过而不偏转的带电粒子速度，并从消电场参考系解释。
\Diagram{fig:m3driftq}{速度选择器中电力与磁力相反}{\draw[axis](0,0)--(3.7,0)node[right]{$v_d$};\draw[axis](1,0)--(1,1.5)node[right]{$qE$};\draw[axis](1,0)--(1,-1.5)node[right]{$qv_dB$};\draw(2,1)circle(.15);\fill(2,1)circle(1pt);\node[right]at(2.3,1){$B$};}
\end{example}
\begin{solution}
图\ref{fig:m3driftq}中令一般速度为$(v_x,v_y,v_z)$。叉乘逐分量为$\vct v\times\vct B=(Bv_y,-Bv_x,0)$，零力条件给
\[
 Bv_y=0,\qquad E-Bv_x=0.
\]
因此全部无偏转速度构成一族：
\[
 \vct v=\frac EB\hat{\vct x}+v_z\hat{\vct z},\qquad
 |v_z|<\sqrt{c^2-\frac{E^2}{B^2}}.
\]
沿磁场的速度分量不产生磁力，因此保留为自由参数。只有额外规定$v_z=0$时，才得到唯一的二维选择速度$E/B$。
\Diagram{fig:m3drifta}{消电场系中的无偏转速度沿磁场；实验室还具有共同横向漂移}{\begin{scope}\draw[axis](0,0)--(2.6,0)node[right]{$x$};\draw[axis](0,-1.4)--(0,1.8)node[above]{$z$};\draw[->,thick](0,0)--(1.7,1.1)node[right]{$\vct v$};\draw[guide](0,1.1)--(1.7,1.1)--(1.7,0);\node[below]at(1.4,0){$E/B$};\node[left]at(0,.9){$v_z$};\end{scope}\begin{scope}[xshift=4cm]\draw[axis](0,0)--(2,0)node[right]{$x\prime$};\draw[axis](0,-1.4)--(0,1.8)node[above]{$z\prime$};\draw[->,thick](0,0)--(0,1.2)node[right]{$v_z\prime$};\draw[->](1.1,-.6)--(1.1,1)node[above]{$\vct B\prime$};\node[below]at(.7,-1){$\vct E\prime=0$};\end{scope}}
取消电场系沿$+x$以$V_d=E/B$运动，记$\gamma_d=(1-V_d^2/c^2)^{-1/2}$。场变换给
\[
 E'_y=\gamma_d(E-V_dB)=0,\qquad
 B'_z=\gamma_d(B-V_dE/c^2)=B/\gamma_d.
\]
对上述速度族，$v'_x=v'_y=0$，而
\[
 v'_z=\frac{v_z}{\gamma_d(1-V_d^2/c^2)}=\gamma_dv_z.
\]
图\ref{fig:m3drifta}中粒子在新系沿$\vct B'$直行，仍不受力；$|v'_z|<c$正好给出实验室的速度界限。若新系另有垂直于$\vct B'$的速度分量，粒子就会回旋，那是有偏转运动，已不属于题目所求的无偏转速度族。$E\geq cB$时不存在有质量粒子的这一亚光速速度族。
\end{solution}

\section{散射与衰变中的不变质量}
\begin{example}\label{ex:m3decay}
静质量$M$的粒子在自身静止系衰变为两个光子。实验室母粒子以速度$V$运动。求两光子实验室能量的范围；若源系两光子都垂直运动轴，求实验室夹角。
\Diagram{fig:m3decayq}{母粒子系两光子能量相等、方向相反}{\draw[->,thick](0,0)--(2,1.3)node[right]{$Mc^2/2$};\draw[->,thick](0,0)--(-2,-1.3)node[left]{$Mc^2/2$};\draw[axis](0,0)--(3,0)node[right]{$V$};}
\end{example}
\begin{solution}
图\ref{fig:m3decayq}静止系动量守恒给两光子能量均为$E_*=Mc^2/2$。若第一光子源系方向余弦为$\mu_*$，则$E_1=\gamma E_*(1+\beta\mu_*)$、$E_2=\gamma E_*(1-\beta\mu_*)$，两者之和为$\gamma Mc^2$。单个能量范围为$[\gamma E_*(1-\beta),\gamma E_*(1+\beta)]$。
\Diagram{fig:m3decaya}{源系横向光子在实验室形成前向对称双束}{\draw[axis](0,0)--(4,0)node[right]{$x$};\draw[->,thick](0,0)--(3,1.6);\draw[->,thick](0,0)--(3,-1.6);\node[right]at(3,1.6){$E_1=\gamma E_*$};\node[right]at(3,-1.6){$E_2=\gamma E_*$};}
若$\mu_*=0$，每个光子在实验室有$\cos\theta=\beta$，夹角为$2\arccos\beta$，如图\ref{fig:m3decaya}。再由两光子不变质量$2E_1E_2(1-\cos2\theta)=M^2c^4$核对：$1-\cos2\theta=2/\gamma^2$，正好恢复母粒子质量。
\end{solution}

\begin{example}\label{ex:m3nuclearpair}
光子撞击静止核，产生电子正电子对，核最终仍为原静质量$M$，电子质量$m$。忽略束缚能，求光子能量阈值，解释核反冲不能完全忽略的原因。
\Diagram{fig:m3nuclearq}{核提供反冲动量，末态总静质量至少$M+2m$}{\draw[->,thick](-2,0)--(0,0)node[midway,above]{$E$};\fill(.5,0)circle(4pt)node[below]{$M$};\draw[->](1,0)--(2,0);\foreach\x/\t in{2.7/M,3.7/{e^-},4.7/{e^+}}{\fill(\x,0)circle(2pt)node[below]{$\t$};}}
\end{example}
\begin{solution}
图\ref{fig:m3nuclearq}初态不变量$s=M^2c^4+2Mc^2E$。阈值时末态三粒子在质心系无相对动量，总质心能量$(M+2m)c^2$。因此
\[
 M^2c^4+2Mc^2E_{\rm th}=(M+2m)^2c^4,
\quad E_{\rm th}=2mc^2\left(1+\frac mM\right).
\]
\Diagram{fig:m3nucleara}{阈值在质心系三粒子同速，实验室仍须整体反冲}{\draw[curve](0,-.5)rectangle(4,.7);\foreach\x/\t in{.7/M,2/{e^-},3.3/{e^+}}{\fill(\x,0)circle(2pt)node[above]{$\t$};}\draw[axis](4.2,.1)--(5.5,.1)node[above]{$V_{\rm cm}$};}
大核质量极限才趋$2mc^2$。图\ref{fig:m3nucleara}中阈值末态在实验室具有共同质心速度，额外项提供反冲能量。若题设改为电子作靶且末态多一个电子，必须重新计总静质量，不能只把$M=m$代入后忽略相同粒子与末态计数。
\end{solution}

\begin{example}\label{ex:m3fourtensor}
在本书度规约定下，构造使$F^\mu=q\mathcal F^{\mu\nu}u_\nu$再现洛伦兹三力的反对称场张量，并说明反对称性怎样保证静质量不变。
\Diagram{fig:m3tensorq}{同一洛伦兹力的时间分量是功率，空间分量是三力}{\draw[axis](0,0)--(3.2,0)node[right]{$\vct u$};\draw[axis](0,0)--(1.3,2.1)node[above]{$\vct E$};\draw(2,1)circle(.15);\fill(2,1)circle(1pt);\node[right]at(2.3,1){$\vct B$};}
\end{example}
\begin{solution}
按$(ct,x,y,z)$排序、度规取$(+,-,-,-)$，图\ref{fig:m3tensorq}所用场张量为
\[
 \mathcal F^{\mu\nu}=\begin{pmatrix}
 0&-E_x/c&-E_y/c&-E_z/c\\
 E_x/c&0&-B_z&B_y\\
 E_y/c&B_z&0&-B_x\\
 E_z/c&-B_y&B_x&0
 \end{pmatrix}.
\]
粒子四速度降指标后为$u_\nu=\gamma_u(c,-u_x,-u_y,-u_z)$。矩阵第0行给$F^0=q\gamma_u\vct E\cdot\vct u/c$；第$x$行给
\[
 F^x=q\gamma_u(E_x+u_yB_z-u_zB_y).
\]
其余两行循环置换指标，得到$\vct F=q\gamma_u(\vct E+\vct u\times\vct B)$。再除以$\gamma_u$便还原洛伦兹三力。
\Diagram{fig:m3tensora}{反对称张量与对称速度乘积缩并为零}{\draw[curve](0,0)rectangle(2.6,1.1);\node at(1.3,.55){$\mathcal F^{\mu\nu}=-\mathcal F^{\nu\mu}$};\draw[->](2.8,.55)--(3.7,.55);\node[right]at(3.8,.55){$u_\mu F^\mu=0$};}
图\ref{fig:m3tensora}中的正交关系来自反对称性。设$Q=u_\mu\mathcal F^{\mu\nu}u_\nu$，交换两个求和哑指标得
\[
 Q=u_\nu\mathcal F^{\nu\mu}u_\mu
 =-u_\mu\mathcal F^{\mu\nu}u_\nu=-Q,
\]
故$Q=0$。为说明它与静质量的关系，暂不预设$m$恒定，将$p^\mu=mu^\mu$代入：
\[
 u_\mu\frac{\dd p^\mu}{\dd\tau}
 =c^2\frac{\dd m}{\dd\tau}
 +\frac m2\frac{\dd(u_\mu u^\mu)}{\dd\tau}
 =c^2\frac{\dd m}{\dd\tau}.
\]
左端由洛伦兹四力为零，所以$\dd m/\dd\tau=0$。电磁场可改变点电荷的运动能量和动量，但在此无内部激发的模型中不改变其静质量。
\end{solution}

\section{弹性碰撞的能量转移与交叉场轨迹}
\subsection{固定靶碰撞中的最大反冲能量}
设入射粒子静质量$M$、实验室能量$E=\gamma Mc^2$、动量$p$，静止靶粒子质量$m$，碰撞前后两粒子的静质量不变。用$s$记总四动量平方乘$c^2$，则
\[
 s=M^2c^4+m^2c^4+2mc^2E,
 \qquad V_* =\frac{pc^2}{E+mc^2},\quad
 \gamma_* =\frac{E+mc^2}{\sqrt s}.
\]
质心系中两粒子动量大小为$p_*$。把两体公式的根号因式展开，利用$E^2-p^2c^2=M^2c^4$，得到
\[
 p_*^2c^2=\frac{m^2c^4(E^2-M^2c^4)}s,
 \qquad p_* =\frac{mc^2p}{\sqrt s}.
\]
靶粒子在质心系中的能量固定，实验室能量随其质心系出射角$\vartheta_*$变化：$E_m=\gamma_*(E_{m*}+V_*p_*\cos\vartheta_*)$。初态靶在该系向后运动，实验室能量为$mc^2=\gamma_*(E_{m*}-V_*p_*)$。两式相减，直接得到反冲动能
\[
 T=\gamma_*V_*p_*(1+\cos\vartheta_*).
\]
最大值对应靶沿质心系正前方射出，而不是任意把实验室两粒子都假定为向前且等速。代入各参数有
\[
 T_{\max}=\frac{2m c^2(\gamma^2-1)}
 {1+2\gamma m/M+(m/M)^2}.
\]
这是纯运动学上界；相互作用的微分截面可能使达到上界的事件很少，但不改变该上界的位置。

\begin{example}重粒子对轻靶的最大能量转移
入射粒子质量为$M=100m$，洛伦兹因子为2，靶粒子初始静止。求最大反冲能量占入射动能的比例；再检查同质量碰撞和非相对论极限。
\ExpansionDiagram{fig:dep-m3-transfer-q}{实验室固定靶与质心系反冲端点}{
\draw[->,thick](0,0)--(2,0)node[midway,above]{$p$};\fill(2.4,0)circle(2pt)node[below]{$M$};\fill(4.5,0)circle(2pt)node[below]{$m$静止};\draw[dashed](5.6,-.6)--(5.6,1.2);\draw[->](6.2,0)--(8,0)node[midway,above]{$p_*$};\node[below]at(7,-.4){质心系靶正向};}
\end{example}
\begin{solution}
由图\ref{fig:dep-m3-transfer-q}的正向质心系端点，
\[
 T_{\max}=\frac{6mc^2}{1+0.04+0.0001}
 \simeq5.7687mc^2.
\]
入射动能$K=(\gamma-1)Mc^2=100mc^2$，最大转移比例约为$5.77\%$。尽管入射动能很大，质量相差悬殊仍使反冲受限；能量和动量必须由同一组末态同时满足。
\ExpansionDiagram{fig:dep-m3-transfer-a}{反冲能量对质心系角度的依赖}{
\draw[->](0,0)--(5,0)node[right]{$\cos\vartheta_*$};\draw[->](0,0)--(0,2.6)node[above]{$T/T_{\max}$};\draw[thick](.5,0)--(4.5,2);\node[below]at(.5,0){$-1$};\node[below]at(2.5,0){$0$};\node[below]at(4.5,0){1};\node[left]at(0,2){1};}
图\ref{fig:dep-m3-transfer-a}的直线关系为$T/T_{\max}=(1+\cos\vartheta_*)/2$，只描述允许能量与角度的对应，不表示各角度等概率。

若$M=m$，分母为$2(1+\gamma)$，所以$T_{\max}=mc^2(\gamma-1)=K$，同质量的一维弹性碰撞可以交换运动状态。若入射速度很小，$\gamma^2-1\simeq v^2/c^2$，得到$T_{\max}\simeq2mM^2v^2/(M+m)^2$；再除以$Mv^2/2$，恢复经典比例$4mM/(M+m)^2$。这两个极限分别检查相同质量和低速度的退化情形。
\end{solution}

\subsection{消去电场后仍须变换整条世界线}
磁占优的正交场可以找到纯磁参考系，但这不意味着实验室中粒子必作匀速直线运动。只有粒子在那个参考系恰好静止时，才对应无偏转速度。一般初态仍有回旋运动，变回实验室后既有漂移又有速率变化。

\begin{example}由静止释放的粒子在正交场中运动
实验室有$\vct E=E\hat{\vct y}$、$\vct B=B\hat{\vct z}$，其中$0<E<cB$。质量$m$、电量$q>0$的粒子在原点由静止释放。忽略辐射，求参数轨迹、平均漂移速度和最大动能。
\ExpansionDiagram{fig:dep-m3-cross-q}{磁场朝纸外，电场向上，无偏转漂移方向向右}{
\draw[->](0,0)--(5,0)node[right]{$x$};\draw[->](0,0)--(0,2.7)node[above]{$y$};\draw[->,thick](1,1)--(1,2)node[above]{$E$};\draw(3,1.5)circle(.18);\fill(3,1.5)circle(1.4pt)node[right=7pt]{$B$};\fill(0,0)circle(2pt)node[below left]{$q$};\draw[->](1,-.6)--(3,-.6)node[midway,below]{$V=E/B$};}
\end{example}
\begin{solution}
按图\ref{fig:dep-m3-cross-q}取$S'$相对实验室以$V=E/B$沿$+x$运动，记$\beta=V/c$、$\Gamma=(1-\beta^2)^{-1/2}$。场变换给$E'=0$、$B'=B/\Gamma$。粒子初始在$S'$的速度为$-V\hat{\vct x}$，所以其恒定洛伦兹因子是$\Gamma$，回旋角频率为
\[
 \Omega=\frac{qB'}{\Gamma m}=\frac{qB}{\Gamma^2m}.
\]
洛伦兹力分量为$\dot v'_x=\Omega v'_y$、$\dot v'_y=-\Omega v'_x$。将第一式再求导并代入第二式，得$\ddot v'_x+\Omega^2v'_x=0$，通解为$A\cos\Omega t'+B\sin\Omega t'$。初值$v'_x(0)=-V$、$\dot v'_x(0)=\Omega v'_y(0)=0$分别给$A=-V,B=0$；再由$v'_y=\dot v'_x/\Omega$得$v'_y=V\sin\Omega t'$。令$s=\Omega t'$，从原点积分：
\begin{align*}
 x'&=\int_0^{t'}-V\cos(\Omega q)\dd q
 =\left[-\frac V\Omega\sin(\Omega q)\right]_0^{t'},\\
 y'&=\int_0^{t'}V\sin(\Omega q)\dd q
 =\left[-\frac V\Omega\cos(\Omega q)\right]_0^{t'}.
\end{align*}
代入两端后，
\[
 x'=-\frac V\Omega\sin s,\quad
 y'=\frac V\Omega(1-\cos s),\quad t'=s/\Omega.
\]
逆变换必须同时作用于时间和位置，得到
\[
 x=\frac{\Gamma V}{\Omega}(s-\sin s),\qquad
 y=\frac V\Omega(1-\cos s),\qquad
 t=\frac\Gamma\Omega(s-\beta^2\sin s).
\]
\ExpansionDiagram{fig:dep-m3-cross-a}{实验室轨迹的无量纲摆线；每个尖点处粒子瞬时静止}{
\draw[->](0,0)--(7,0)node[right]{$x\Omega/(\Gamma V)$};\draw[->](0,0)--(0,2.7)node[above]{$y\Omega/V$};\draw[thick,domain=0:6.283185,samples=180]plot({\x-sin(\x r)},{1-cos(\x r)});\fill(0,0)circle(1.4pt);\fill(6.283185,0)circle(1.4pt);\node[below]at(6.283185,0){$2\pi$};}
图\ref{fig:dep-m3-cross-a}中横纵轴使用不同长度单位，轨迹形状对应上述参数式。因为$\dd t/\dd s=\Gamma(1-\beta^2\cos s)/\Omega>0$，参数没有倒转时间。在一周$s:0\to2\pi$内，$\Delta x=2\pi\Gamma V/\Omega$、$\Delta t=2\pi\Gamma/\Omega$，所以平均漂移为$V=E/B$。

磁力不做功，实验室动能等于$qEy$，于是
\[
 K(s)=\frac{qEV}{\Omega}(1-\cos s)
 =\Gamma^2mV^2(1-\cos s),\quad
 K_{\max}=2\Gamma^2mV^2.
\]
也可变换四动量得粒子能量$\Gamma^2mc^2(1-\beta^2\cos s)$，减去$mc^2$得到同一结果。一周后高度回零，动能也回零，但位置已漂移；这与“在纯磁系能量恒定”相容，两个系的能量定义不同。
\end{solution}

\section{两体末态的能谱与连续光压推进}
\subsection{各向同性衰变怎样形成均匀能谱}
母粒子静质量$M$，在自身静止系衰变为两个同质量$m$的粒子。每个子粒子的$E_*=Mc^2/2$、$p_*c=\sqrt{M^2/4-m^2}\,c^2$。若衰变在该系各向同性，单个已标记子粒子的$\mu_*=\cos\theta_*$在$[-1,1]$上均匀。

母粒子在实验室沿$x$以$\beta c$运动，洛伦兹因子为$\Gamma$。子粒子能量为
\[
 E=\Gamma(E_*+\beta cp_*\mu_*).
\]
在$\beta p_*>0$时，这是$\mu_*$的严格递增线性函数，由$\dd P=\dd\mu_*/2$及$\dd E=\Gamma\beta cp_*\dd\mu_*$，全角接收的能量概率密度为
\[
 \frac{\dd P}{\dd E}=\frac1{2\Gamma\beta cp_*},\qquad
 E_-\leq E\leq E_+,\quad
 E_\pm=\Gamma(E_*\pm\beta cp_*).
\]
若$\beta=0$或$p_*=0$，两端点合并为$E=\Gamma E_*$，成为单能分布，不能再使用上式的零分母密度。实验室角分布另含方向变换的雅可比；有限角域探测须先反解源系接受范围，再截取能谱。

\begin{example}由能谱端点恢复子粒子质量
已知母粒子实验室洛伦兹因子为2，两体等质量衰变的单子粒子全角能谱范围为$4\mathcal E\leq E\leq16\mathcal E$，$\mathcal E$为能量单位。假定母粒子静止系中衰变方向各向同性。求母、子静质量，求$E>12\mathcal E$的概率，并说明同一次衰变中两子粒子能量的相关性。
\ExpansionDiagram{fig:dep-m3-spectrum-q}{全角接收的两体末态形成有两个运动学端点的能谱}{
\draw[->](0,0)--(7,0)node[right]{$E/\mathcal E$};\draw[->](0,0)--(0,2.5)node[above]{概率密度};\draw[thick](1.5,0)--(1.5,1.5)--(6,1.5)--(6,0);\node[below]at(1.5,0){4};\node[below]at(6,0){16};\node[above]at(3.8,1.5){均匀};}
\end{example}
\begin{solution}
图\ref{fig:dep-m3-spectrum-q}的中点为$E_c=10\mathcal E=\Gamma Mc^2/2$，所以$Mc^2=10\mathcal E$。半宽$\Delta=6\mathcal E=\Gamma\beta cp_*$，而$\Gamma\beta=\sqrt3$，得$cp_*=2\sqrt3\mathcal E$。由$E_*^2=m^2c^4+p_*^2c^2$，
\[
 m^2c^4=(5\mathcal E)^2-(2\sqrt3\mathcal E)^2
 =13\mathcal E^2,\qquad mc^2=\sqrt{13}\mathcal E.
\]
矩形分布中，高于$12\mathcal E$的区间占长度$4\mathcal E$，总宽为$12\mathcal E$，所以概率$1/3$。
\ExpansionDiagram{fig:dep-m3-spectrum-a}{同一次衰变的能量对被限制在总能量直线上}{
\draw[->](0,0)--(5,0)node[right]{$E_1$};\draw[->](0,0)--(0,4.5)node[above]{$E_2$};\draw[thick](1,4)--(4,1);\fill(1,4)circle(1.4pt);\fill(4,1)circle(1.4pt);\node[right]at(3.3,3.4){$E_1+E_2=20\mathcal E$};}
图\ref{fig:dep-m3-spectrum-a}中$E_2=20\mathcal E-E_1$，两者完全反相关。令$\xi=E_1-E_c$，则$\dd E_1=\dd\xi$，积分限为$-\Delta,\Delta$，密度为$1/(2\Delta)$。对称性给$\langle\xi\rangle=0$，方差为
\[
 \operatorname{Var}E_1=\frac1{2\Delta}\int_{-\Delta}^{\Delta}\xi^2\dd\xi
 =\left[\frac{\xi^3}{6\Delta}\right]_{-\Delta}^{\Delta}
 =\frac{\Delta^2}{3}=12\mathcal E^2.
\]
又有$E_2-E_c=-(E_1-E_c)$，故协方差为$-\langle\xi^2\rangle=-12\mathcal E^2$。若把两个均匀谱独立抽样，会产生不守恒的能量对。若$\Gamma$也未知，两个端点只约束$\Gamma M$及$\beta\sqrt{1-4m^2/M^2}$，通常不足以唯一确定三个未知量。
\end{solution}

\subsection{移动镜面的光压还要计算追上镜面的光子率}
平面镜沿光束方向后退，实验室固定截面上的入射功率为$P$。若镜速为$\beta c$，单位实验室时间追上镜面的入射能量是$P(1-\beta)$，反射单光子能量比为$(1-\beta)/(1+\beta)$。每个光子反向弹回，镜获得的动量是入射与反射光子动量模之和，因此
\[
 F=\frac{P(1-\beta)}c\left(1+\frac{1-\beta}{1+\beta}\right)
 =\frac{2P}{c}\frac{1-\beta}{1+\beta}.
\]
这里忽略单光子的有限反冲，采用大量弱光子的连续光压极限。若遗漏追赶率$1-\beta$，会高估高速镜所受的力。

\begin{example}恒定入射功率推动理想光帆
质量$M$的理想反射光帆初始静止，实验室光源提供固定截面功率$P$的连续准直光束，光帆全部截获该束并沿束方向运动，静质量不变。忽略衍射、束宽变化及辐射吸收。求速度和位置随时间的参数式，并核对光子能量损失与光帆做功。
\ExpansionDiagram{fig:dep-m3-sail-q}{光帆后退使光子到达率降低，反射光同时红移}{
\draw[->,thick](0,.7)--(4,.7)node[midway,above]{$P$};\draw[very thick](4,-.5)--(4,1.7);\draw[->](4,1.9)--(5.5,1.9)node[right]{$\beta c$};\draw[->](4,-.1)--(1,-.1)node[midway,below]{红移反射};\node[below]at(4,-.6){光帆$M$};}
\end{example}
\begin{solution}
图\ref{fig:dep-m3-sail-q}中用快度$\chi$表示$\beta=\tanh\chi$，有$(1-\beta)/(1+\beta)=\ee^{-2\chi}$。动量$Mc\sinh\chi$的导数满足
\[
 Mc\cosh\chi\frac{\dd\chi}{\dd t}=\frac{2P}c\ee^{-2\chi}.
\]
因此$\dd t/\dd\chi=Mc^2\ee^{2\chi}\cosh\chi/(2P)$。展开双曲函数，并用初值$t(0)=0$，
\[
 t=\frac{Mc^2}{4P}\int_0^\chi(\ee^{3q}+\ee^q)\dd q
 =\frac{Mc^2}{4P}\left[\frac{\ee^{3q}}3+\ee^q\right]_0^\chi.
\]
令$u=\ee^\chi\geq1$，得
\[
 t=\frac{Mc^2}{12P}(u^3+3u-4),\qquad
 \beta=\frac{u^2-1}{u^2+1}.
\]
位置由$\dd x/\dd\chi=\beta c\dd t/\dd\chi$积分。用$\beta=\tanh\chi$及$x(0)=0$，
\[
 x=\frac{Mc^3}{4P}\int_0^\chi(\ee^{3q}-\ee^q)\dd q
 =\frac{Mc^3}{4P}\left[\frac{\ee^{3q}}3-\ee^q\right]_0^\chi,
\]
得到
\[
 x=\frac{Mc^3}{12P}(u^3-3u+2).
\]
\ExpansionDiagram{fig:dep-m3-sail-a}{连续推进的速度单调趋近光速，所需时间随快度迅速增长}{
\draw[->](0,0)--(5,0)node[right]{$u=\ee^\chi$};\draw[->](0,0)--(0,2.5)node[above]{$v/c$};\draw[thick,domain=1:4.5,samples=80]plot({\x-1},{2*(\x^2-1)/(\x^2+1)});\draw[dashed](0,2)--(4.5,2);\node[left]at(0,2){1};\node[below]at(0,0){1};}
图\ref{fig:dep-m3-sail-a}的$u=1$为初始时刻。小时间下$F\simeq2P/c$，恢复初始加速度$2P/(Mc)$；高速时$t$按$u^3$增长，而$1-\beta\simeq2/u^2$，始终不能在有限时间达到$c$。

每单位实验室时间反射离开的光能为$P_{\rm refl}=P(1-\beta)^2/(1+\beta)$。真正撞到光帆的入射能量率为$P_{\rm hit}=P(1-\beta)$，两者之差
\[
 P_{\rm hit}-P_{\rm refl}=\frac{2P\beta(1-\beta)}{1+\beta}=Fv,
\]
与光帆动能增长率一致。固定截面功率$P$和镜面实际截获功率不同；未追上光帆的光还在传播途中，不应直接算作已交给光帆的能量。
\end{solution}

\section{习题与解答}
\begin{exercise}\label{ex:m3electricdominant}
正交均匀场$\vct E=E\hat{\vct y}$、$\vct B=B\hat{\vct z}$满足$E>cB$。带正电粒子初始实验室速度为$V\hat{\vct x}$，其中$V=c^2B/E$。求其轨迹，忽略辐射和边界。
\Diagram{fig:m3edominantq}{电占优场可变为纯电场，初速度选为相应推动速度}{\draw[axis](0,0)--(3.7,0)node[right]{$V$};\draw[axis](0,0)--(0,2.6)node[above]{$E$};\draw(1.8,1)circle(.15);\fill(1.8,1)circle(1pt);\node[right]at(2.1,1){$B$};}
\end{exercise}
\begin{solution}
图\ref{fig:m3edominantq}沿$V$推动后$B'=0$，$E'=\sqrt{E^2-c^2B^2}$，粒子初始静止。设$\alpha=qE'/m$，沿用习题\ref{ex:m1fourwork}已逐项积分的恒力解，并将加速轴改为$y'$，得到
\[
 ct'=\frac{c^2}\alpha\sinh(\alpha\tau/c),\quad
 y'=\frac{c^2}\alpha[\cosh(\alpha\tau/c)-1],\quad x'=0.
\]
逆推动给$t=\gamma t'$、$x=Vt$、$y=y'$，消去固有时得到
\[
 y(t)=\frac{c^2}\alpha\left[\sqrt{1+(\alpha t/(\gamma c))^2}-1\right].
\]
\Diagram{fig:m3edominanta}{轨迹按$X=\alpha x/(\gamma Vc)$、$Y=\alpha y/c^2$归一化，虚线为渐近线}{\draw[axis](0,0)--(4.5,0)node[right]{$X$};\draw[axis](0,0)--(0,3.5)node[above]{$Y$};\draw[curve,domain=0:4,samples=80]plot(\x,{sqrt(1+\x*\x)-1});\draw[guide](1,0)--(4.2,3.2);\node[right]at(3.7,2.4){$Y=X-1$};}
图\ref{fig:m3edominanta}的长时轨迹趋于直线斜向传播。横向速度为$v_y=(c/\gamma)\tanh(\alpha\tau/c)$，所以$v^2=V^2+c^2\tanh^2(\alpha\tau/c)/\gamma^2<c^2$。$v_x$恒定不代表$x$向动量恒定，因为总$\gamma_v$仍在增加；磁力的$x$分量正好提供这部分动量变化。
\end{solution}

\begin{exercise}\label{ex:m3decaycone}
静质量$M>2m$的粒子衰变为两个相同质量$m$的粒子，母粒子实验室速度为$V$。求子粒子在质心系的速率$u_*$；当$V>u_*$时，求所有实验室出射方向被限制的最大半角。
\Diagram{fig:m3decayconeq}{母粒子系两子粒子动量等大反向}{\draw[axis](-2.5,0)--(3.5,0)node[right]{$V$};\draw[->,thick](0,0)--(1.4,1.8)node[right]{$m$};\draw[->,thick](0,0)--(-1.4,-1.8)node[left]{$m$};}
\end{exercise}
\begin{solution}
图\ref{fig:m3decayconeq}中，母粒子静止系总动量为零，两子粒子能量相同，故$E_*=Mc^2/2$。质壳关系给
\[
 p_*c=\sqrt{M^2c^4/4-m^2c^4},\qquad
 u_*=\frac{p_*c^2}{E_*}=c\sqrt{1-4m^2/M^2}.
\]
令$z=\cos\theta_*$，母粒子沿$x$以$V$运动。速度变换的共同分母为$1+Vu_*z/c^2>0$，相除后
\[
 \tan\theta=\frac{u_*\sqrt{1-z^2}}{\gamma(V+u_*z)}.
\]
在$V>u_*$时分母始终正，所有非零速度都向前。
\Diagram{fig:m3decayconea}{亚光速子粒子可全部集中于一个前向圆锥}{\draw[axis](0,0)--(4.4,0)node[right]{$V$};\draw[curve](0,0)--(3,1.4);\draw[curve](0,0)--(3,-1.4);\draw[guide](3,-1.4)--(3,1.4);\node[above]at(2.2,1.2){$\theta_{\max}$};}
对$-1<z<1$求对数导数，最大角满足
\[
 -\frac z{1-z^2}-\frac{u_*}{V+u_*z}=0
 \quad\Longrightarrow\quad Vz+u_*=0.
\]
取$z=-u_*/V$，则$\sqrt{1-z^2}=\sqrt{V^2-u_*^2}/V$、$V+u_*z=(V^2-u_*^2)/V$，代入得到图\ref{fig:m3decayconea}的
\[
 \tan\theta_{\max}=\frac{u_*}{\gamma\sqrt{V^2-u_*^2}}.
\]
两端$z=\pm1$均有$\theta=0$，内点驻值因而为全局最大。速度相等时则出现退化：$V=u_*$时，正后方发射的子粒子在实验室静止，方向未定义；其余均向前，出射角可趋近但达不到$\pi/2$。$V<u_*$时取$z=-1$便有负的纵向速度，出现后向出射，且可有横向出射。光子$u_*=c$不可能满足$V>u_*$，因而不能用这个有限前向圆锥代替全角光行差。
\end{solution}

\begin{exercise}\label{ex:m3fizeau}
透明介质在自身静止系的非色散折射率为$n$，光沿介质运动方向传播，介质相对实验室速度$V$。求实验室相速度，并展开至$V/c$一阶。
\Diagram{fig:m3fizeauq}{介质系光速为$c/n$，介质整体沿同向运动}{\draw[curve](0,-.6)rectangle(4,1);\draw[axis](.4,.2)--(2.8,.2)node[above]{$c/n$};\draw[axis](4.2,.2)--(5.5,.2)node[above]{$V$};}
\end{exercise}
\begin{solution}
图\ref{fig:m3fizeauq}的波面速度沿$x$，由相位变换或速度合成都有$v_{\rm ph}=(c/n+V)/(1+V/(nc))$。低速把分母倒数展开为$1-V/(nc)+O(V^2/c^2)$，得$v_{\rm ph}=c/n+V(1-1/n^2)+O(V^2/c)$。
\Diagram{fig:m3fizeaua}{拖曳系数在$n=1$时为零}{\draw[axis](0,0)--(4.5,0)node[right]{$n$};\draw[axis](0,0)--(0,2.8)node[above]{$1-1/n^2$};\draw[curve,domain=1:4,samples=70]plot(\x,{2.4*(1-1/(\x*\x))});\node[below]at(1,0){$1$};}
图\ref{fig:m3fizeaua}在$n=1$时恢复真空光速$c$。本题采用非色散近似；真实色散介质中，变换也改变频率，应在相应介质静止频率处取$n$，一阶展开还会出现色散修正。
\end{solution}

\begin{exercise}\label{ex:m3currentrest}
某区域总电荷密度为零、电流密度$\vct j=j\hat{\vct x}\ne0$。是否存在使总电流为零的惯性系？解释为什么四电流不总能写为某个实固有密度乘四速度。
\Diagram{fig:m3currentq}{等量正负电荷相反运动，密度相消而电流相加}{\foreach\x in{0,1,2,3}{\node at(\x,.5){$+$};\node at(\x,-.5){$-$};}\draw[axis](.5,1)--(2,1);\draw[axis](2,-1)--(.5,-1);}
\end{exercise}
\begin{solution}
图\ref{fig:m3currentq}中，总四电流为$J^\mu=(0,j,0,0)$，不变量$J^2=-j^2<0$。假设某惯性系有$\vct j'=0$，则$J'^2=c^2\rho'^2\geq0$，与不变量符号矛盾，所以任何方向的惯性变换都不能消去总电流。

若仅沿$x$变换，可逐项看到$\rho'=-\gamma Vj/c^2$、$j'_x=\gamma j$。电中性不保持，但电流仍不为零。
\Diagram{fig:m3currenta}{类空四电流无法通过亚光速变换转到纯时间方向}{\draw[axis](-2,0)--(2.6,0)node[right]{$j_x$};\draw[axis](0,-1.8)--(0,2.4)node[above]{$c\rho$};\draw[guide](-1.8,-1.8)--(2,2);\draw[guide](-1.8,1.8)--(1.8,-1.8);\draw[->,thick](0,0)--(2,0);}
图\ref{fig:m3currenta}的四电流位于光锥外。单一带电成分共同运动时，$J^\mu=\rho_0U^\mu$且$J^2=\rho_0^2c^2\geq0$；中性电流却可能由正负两组四流相加。例如实验室密度为$+\rho$的正电荷静止，密度为$-\rho$的负电荷以$-u\hat{\vct x}$运动，总电荷为零、总电流为$+\rho u\hat{\vct x}$。各组分别有自己的静止系，总和没有共同电流静止系。

更一般地，只有$J^2>0$的非零四电流才可变到纯时间方向；$J^2<0$可变到纯空间方向，$J^2=0$且非零时两种静止化都不可行。总四电流按带符号电荷加权，与物质总能量动量不同；其类空性质不要求任何载体超光速。
\end{solution}

\begin{exercise}\label{ex:m3obliquewave}
平面波实验室传播方向与$x$成60度，观察系以$0.6c$沿$x$运动。求频率比和新方向角，核对横向波矢不变。
\Diagram{fig:m3obliquewaveq}{斜向波矢的纵横分解}{\draw[axis](0,0)--(3.8,0)node[right]{$x$};\draw[axis](0,0)--(1.5,2.598)node[right]{$\vct k$};\draw[guide](1.5,0)--(1.5,2.598);\node at(.9,.4){60度};}
\end{exercise}
\begin{solution}
图\ref{fig:m3obliquewaveq}有$\gamma=1.25$、$\cos\theta=.5$，频率比$\omega'/\omega=1.25(1-.6\times.5)=.875$。方向余弦$\cos\theta'=(.5-.6)/(1-.3)=-1/7$，所以$\theta'\simeq98.21^\circ$，略向后。
\Diagram{fig:m3obliquewavea}{两系采用相同波矢比例，横向分量的高度相同}{
\draw[axis](-1.8,0)--(2.2,0)node[right]{$k'_x$};\draw[axis](0,0)--(0,3.2)node[above]{$k'_y$};
\draw[->,thick](0,0)--(-.375,2.598)node[above left]{$\vct k'$};
\draw[guide](-.375,0)--(-.375,2.598)--(1.5,2.598)--(1.5,0);
\draw[->,dashed](0,0)--(1.5,2.598)node[above right]{$\vct k$};
\node[right]at(1.6,1.7){$k'_y=k_y$};}
新系$\sin\theta'=\sqrt{48}/7=4\sqrt3/7$，因此$k'_y=(\omega'/c)\sin\theta'=(7/8)(\omega/c)(4\sqrt3/7)=\sqrt3\omega/(2c)=k_y$。图\ref{fig:m3obliquewavea}的角度必须用负余弦判断象限，不能把正弦反求成81.79度。
\end{solution}

\begin{exercise}\label{ex:m3recoilphoton}
静止质量$M$的理想镜被能量$E$的单个光子正向撞击，光子反向弹回，镜的静质量不变。求反射光能量和镜反冲动能。
\Diagram{fig:m3recoilq}{有限质量镜不能视为预先给定速度的无限重反射面}{\draw[very thick](2,-1)--(2,1.2);\draw[->](0,.5)--(1.9,.5)node[midway,above]{$E$};\draw[->](1.9,-.3)--(0,-.3)node[midway,below]{$E_r$};\node[right]at(2.3,0){$M$};}
\end{exercise}
\begin{solution}
图\ref{fig:m3recoilq}镜末能量$Mc^2+E-E_r$，末动量$(E+E_r)/c$。代入镜的质壳关系，
\[
 (Mc^2+E-E_r)^2-(E+E_r)^2=M^2c^4.
\]
展开得$2Mc^2(E-E_r)=4EE_r$，所以$E_r=E/[1+2E/(Mc^2)]$。
\Diagram{fig:m3recoila}{镜的反冲动能等于光子损失能量}{\draw[axis](0,0)--(4.3,0)node[right]{$E/(Mc^2)$};\draw[axis](0,0)--(0,2.6)node[above]{$E_r/E$};\draw[curve,domain=0:3.5,samples=70]plot(\x,{2.2/(1+2*\x)});}
反冲动能$K_M=E-E_r=2E^2/(Mc^2+2E)$。图\ref{fig:m3recoila}在$M\to\infty$趋无频移；有限质量时必红移。这里不能先选某个“镜速”代入两次多普勒式，因为速度在单次碰撞中改变，需用全程守恒确定。
\end{solution}

\begin{exercise}\label{ex:m3invariantfields}
在某系有$\vct E=E_0\hat{\vct y}$、$\vct B=2E_0\hat{\vct z}/c$。求消电场的速度、剩余磁场和该系中粒子静止时对应实验室的速度。
\Diagram{fig:m3invariantq}{磁占优正交场}{\draw[axis](0,0)--(3.5,0)node[right]{$x$};\draw[axis](0,0)--(0,2.1)node[above]{$E_0$};\draw(1.6,1)circle(.15);\fill(1.6,1)circle(1pt);\node[right]at(1.9,1){$2E_0/c$};}
\end{exercise}
\begin{solution}
图\ref{fig:m3invariantq}给出$E=E_0$、$B=2E_0/c$，且$\vct E\perp\vct B$，磁占优。沿$+x$变换，有$E_y'=\gamma(E_0-V\,2E_0/c)$。令其为零，得到$V=c/2$、$\gamma=2/\sqrt3$。再求剩余磁场，
\[
 B_z'=\frac2{\sqrt3}\left(\frac{2E_0}c-\frac{E_0}{2c}\right)
 =\frac{\sqrt3E_0}c.
\]
也可由$B'^2=B^2-E^2/c^2=3E_0^2/c^2$独立核对其大小，方向由具体变换保留为$+z$。
\Diagram{fig:m3invarianta}{消场系中静止对应实验室匀速漂移}{\fill(0,0)circle(2pt);\draw[axis](0,.5)--(2.5,.5)node[right]{$V=c/2$};\node[below]at(1,-.4){$\vct E+\vct V\times\vct B=0$};}
图\ref{fig:m3invarianta}中的粒子在新系静止，逆速度变换给实验室$\vct v=(c/2)\hat{\vct x}$。其磁力方向为$\hat{\vct x}\times\hat{\vct z}=-\hat{\vct y}$，大小$q(c/2)(2E_0/c)=qE_0$，正好与电力相抵。

若只要求新系粒子沿$z'$运动而非静止，则逆变换仍有$v_x=c/2$，还允许$v_z=v_z'/\gamma$，范围$|v_z|<\sqrt3c/2$。这进一步说明消电场不等于消去所有粒子的力：沿$\vct B'$运动仍零力，垂直$\vct B'$运动却受磁力并回旋。参考系选择由场决定，具体粒子是否无偏转还要代入其速度。
\end{solution}

\begin{exercise}\label{ex:m3properpower}
粒子静质量不变，某实验室三力始终平行瞬时速度。证明四力模等于三力模，并说明这并不意味着四力时间分量为零。
\Diagram{fig:m3properq}{三力与速度平行而功率非零}{\draw[axis](0,0)--(3.5,0)node[right]{$\vct v$};\draw[axis](0,.6)--(2.5,.6)node[right]{$\vct f$};}
\end{exercise}
\begin{solution}
图\ref{fig:m3properq}中，取瞬时速度沿$x$正向，平行三力的有向分量为$f$。四力为
\[
 F^\mu=\gamma_v(\beta_v f,f,0,0),\qquad
 F^2=\gamma_v^2f^2(\beta_v^2-1)=-f^2.
\]
故固有三力大小$f_0=\sqrt{-F^2}=|f|$。这里所谓四力模是类空四力的$\sqrt{-F^2}$，不是四个分量平方和的欧氏模。
\Diagram{fig:m3propera}{类空四力可有非零时间分量；零功率参考系中的时间分量为零}{\draw[axis](0,0)--(3.5,0)node[right]{$F^x$};\draw[axis](0,0)--(0,2.6)node[above]{$F^0$};\draw[->,thick](0,0)--(2.8,1.4);\draw[guide](0,0)--(2.4,2.4);\node[right]at(2.8,1.4){$F^0/F^x=v/c$};}
图\ref{fig:m3propera}中$F^0=\gamma_vfv/c$，只要$fv\ne0$便非零。变到粒子瞬时静止系后，
\[
 F_*^0=\gamma_v(F^0-\beta_vF^x)=0,\qquad F_*^x=f.
\]
但零时间分量并非瞬时静止系独有。以该静止系为起点，再沿$y$以任意$w<c$变换：原四力$(0,f,0,0)$的$0,y$分量均为零，变换后仍是$(0,f,0,0)$；粒子却以$-w\hat{\vct y}$运动。此系三力为$(f/\gamma_w,0,0)$，与粒子速度垂直，功率仍为零。

一般判据是$F^0=\gamma\vct f\cdot\vct v/c=0$，即该系中瞬时功率为零。四维正交$u\cdot F=0$则在所有参考系成立，它不要求每个参考系中的时间分量都为零。
\end{solution}

\begin{exercise}\label{ex:m3threshold}
静质量$m$的入射粒子撞击同样静质量的静止靶粒子，末态需产生四个质量均为$m$的粒子。忽略其他限制，求入射动能阈值，并与两束等能反向对撞比较。
\Diagram{fig:m3thresholdq}{固定靶与对撞的入射配置}{\draw[->](0,.5)--(2,.5);\fill(2.8,.5)circle(2pt);\node[above]at(1,.5){$E$};\node[right]at(3,.5){固定靶};\draw[->](0,-1)--(1.7,-1);\draw[->](3.6,-1)--(1.9,-1);\node[right]at(4,-1){对撞};}
\end{exercise}
\begin{solution}
在图\ref{fig:m3thresholdq}的固定靶实验室系，入射总能量记为$E$，动量记为$p$，有$E^2-p^2c^2=m^2c^4$。初态总能量$E+mc^2$、总动量$p$，因此
\begin{align*}
 s&=(E+mc^2)^2-p^2c^2\\
 &=E^2-p^2c^2+2mc^2E+m^2c^4\\
 &=2m^2c^4+2mc^2E.
\end{align*}
阈值时四个末态粒子在质心系相对静止，最小总能量为$4mc^2$，故$s=16m^2c^4$，解得$E=7mc^2$、$K=6mc^2$。
\Diagram{fig:m3thresholda}{阈值末态在质心系中无相对运动}{\foreach\x in{0,1,2,3}{\fill(\x,0)circle(3pt)node[below]{$m$};}\draw[curve](-.4,-.5)rectangle(3.4,.6);\node[above]at(1.5,.6){总静能$4mc^2$};}
图\ref{fig:m3thresholda}的末态静止是指质心系。实验室的初态总能量为$8mc^2$，总动量为$\sqrt{7^2-1}\,mc=4\sqrt3\,mc$，因而质心速度$V=Pc^2/E_T=\sqrt3c/2$、$\gamma_V=2$。末态整体动能为
\[
 K_{\rm bulk}=(\gamma_V-1)\,4mc^2=4mc^2.
\]
六份入射动能中，$2mc^2$用来增加末态静能，$4mc^2$用于整体运动。若改成两束等能反向对撞，总动量一开始为零，阈值只需每束总能量$2mc^2$，每束动能$mc^2$，两束共$2mc^2$。两种装置的末态最小静能相同，固定靶所需额外能量来自总动量约束。
\end{solution}

\begin{exercise}\label{ex:m3compton}
能量$E$的光子撞击静止电子，散射光子与入射方向成$\theta$角。用四动量守恒推导波长改变量。
\Diagram{fig:m3comptonq}{散射光子与反冲电子}{\draw[->,thick](-2,0)--(0,0)node[midway,above]{$E$};\draw[->,thick](0,0)--(2,1.7)node[right]{$E'$};\draw[->](0,0)--(2,-1.1)node[right]{电子};\draw[guide](0,0)--(2.5,0);\node at(.8,.35){$\theta$};}
\end{exercise}
\begin{solution}
图\ref{fig:m3comptonq}中，初始电子四动量为$P=(mc,\vct0)$，初末光子四动量记为$k,k'$，其能量分别为$E,E'$。末电子$P'=P+k-k'$，电子静质量不变，故
\[
 0=P'^2-P^2=2P\cdot k-2P\cdot k'-2k\cdot k'.
\]
光子各自平方为零，两个光子的空间夹角为$\theta$，因此
\[
 P\cdot k=mE,\quad P\cdot k'=mE',\quad
 k\cdot k'=\frac{EE'}{c^2}(1-\cos\theta).
\]
代入并除以$EE'$，得$1/E'-1/E=(1-\cos\theta)/(mc^2)$。
\Diagram{fig:m3comptona}{散射角越大，光子能量损失越大}{\draw[axis](0,0)--(4.1,0)node[right]{$\theta$};\draw[axis](0,0)--(0,2.8)node[above]{$\lambda'-\lambda$};\draw[curve,domain=0:180,samples=70]plot({\x/50},{1-cos(\x)});\node[below]at(3.6,0){$\pi$};\node[left]at(0,2){$2h/(mc)$};}
以$E=hc/\lambda$表示，即为图\ref{fig:m3comptona}的康普顿关系
\[
 \lambda'-\lambda=\frac h{mc}(1-\cos\theta),\qquad
 E'=\frac E{1+(E/mc^2)(1-\cos\theta)}.
\]
电子反冲不能在守恒中省掉。取散射光子在$x,y$平面且$y$分量为正，末电子分量为
\[
 p_{ex}=\frac{E-E'\cos\theta}c,\qquad
 p_{ey}=-\frac{E'\sin\theta}c,
 \quad K_e=E-E'.
\]
这些分量应与题图中的反冲方向一致。由它们平方得到$p_e^2c^2=E^2+E'^2-2EE'\cos\theta$；结合$K_e$又有$p_e^2c^2=K_e^2+2mc^2K_e$，两式相等恰恢复上面的散射能量式。正向散射无能量转移，反向散射波长改变量最大为$2h/(mc)$。
\end{solution}

\begin{exercise}\label{ex:m3circular}
相对论带电粒子在均匀磁场中作圆周运动，速度垂直磁场。忽略辐射。求半径、实验室角频率、固有加速度及瞬时静止系电场。
\Diagram{fig:m3circleq}{垂直磁场中的圆周运动，取正电荷}{\draw[curve](0,0)circle(1.5);\fill(1.5,0)circle(2pt);\draw[axis](1.5,0)--(1.5,1.2)node[right]{$v$};\draw[axis](1.5,0)--(.3,0)node[above]{$qvB$};\draw(0,0)circle(.15);\draw(-.1,-.1)--(.1,.1);\draw(-.1,.1)--(.1,-.1);\node[below]at(0,-.3){$B$入纸面};}
\end{exercise}
\begin{solution}
磁力始终垂直速度，$\dd E/\dd t=0$，故$\gamma_v$恒定。图\ref{fig:m3circleq}中动量转向率$p v/R=|q|vB$，因此$R=p/(|q|B)=\gamma_vmv/(|q|B)$，$\omega=v/R=|q|B/(\gamma_vm)$。
\Diagram{fig:m3circlea}{局部瞬时静止系中的电场指向圆心}{\fill(0,0)circle(2pt)node[below]{粒子};\draw[axis](0,0)--(-2,0)node[above]{$\vct E_0=\gamma_v\vct v\times\vct B$};\draw[guide](0,-1)--(0,1.2);\node[right]at(.2,.7){$\vct v$切向};}
实验室向心加速度$a=v^2/R$属垂直分量，所以固有加速度$\alpha=\gamma_v^2v^2/R=\gamma_v|q|vB/m$。瞬时静止系场变换给图\ref{fig:m3circlea}的$\vct E_0=\gamma_v\vct v\times\vct B$；静止粒子所受电力大小$|q|\gamma_vvB=m\alpha$，与加速度关系相符。该系电场可做瞬时加速，却不意味着实验室磁场做功。
\end{solution}

\begin{exercise}\label{ex:m3sheet}
无限带电平面静止系面密度为$\sigma_0$，平面为$xy$面。观察系沿$x$以$V$运动。求平面两侧的电磁场，并用变换后的面电流核对磁场方向。
\Diagram{fig:m3sheetq}{静止带电平面两侧的电场}{\draw[curve](-2,0)--(2,0);\foreach\x in{-1.5,-.5,.5,1.5}{\node at(\x,0){$+$};\draw[axis](\x,.2)--(\x,1.2);\draw[axis](\x,-.2)--(\x,-1.2);}\node[right]at(2,0){$xy$面};\node[right]at(1.5,1.2){$z>0$};}
\end{exercise}
\begin{solution}
图\ref{fig:m3sheetq}中$E_z=\operatorname{sgn}(z)\sigma_0/(2\varepsilon_0)$、$B=0$。沿$x$变换得$E_z'=\gamma E_z$、$B_y'=\gamma VE_z/c^2$。面密度$\sigma'=\gamma\sigma_0$，平面在新系沿$-x$运动，面电流$K_x'=-\gamma\sigma_0V$。
\Diagram{fig:m3sheeta}{新系面电流沿$-x$，上侧磁场沿$+y$}{\draw[axis](-2,0)--(2,0)node[right]{$x'$};\draw[axis](0,-1.7)--(0,1.8)node[above]{$z'$};\draw[very thick,->](1,0)--(-1,0)node[below]{$\vct K'$};\draw(.8,1)circle(.15);\draw(.7,.9)--(.9,1.1);\draw(.7,1.1)--(.9,.9);\node[right]at(1.1,1){$B_y'$入纸面};}
按图\ref{fig:m3sheeta}的右手定则，无限面电流上侧$B_y'=-\mu_0K_x'/2=\mu_0\gamma\sigma_0V/2$，与$\gamma VE_z/c^2$完全一致，因为$\varepsilon_0\mu_0c^2=1$。上下侧场都反号。这里$x'$向右、$z'$向上时，$+y'$朝纸内，基底方向需在图中明确。
\end{solution}

\begin{exercise}\label{ex:m3gauge}
真空平面波的四势为$A^\mu=\epsilon^\mu\ee^{-ik_\nu x^\nu}$，满足洛伦兹规范。求极化四矢量的约束，并说明沿波矢的四势分量为何不改变场。
\Diagram{fig:m3gaugeq}{横向极化与传播方向}{\draw[axis](0,0)--(4,0)node[right]{$\vct k$};\draw[axis](1,0)--(1,1.8)node[above]{$\vct\epsilon_\perp$};\draw[guide](0,1)--(4,1);\node[below]at(2,-.4){等相位波面垂直$\vct k$};}
\end{exercise}
\begin{solution}
规范条件$\partial_\mu A^\mu=0$给$k_\mu\epsilon^\mu=0$。若再选无标势的横向规范$\epsilon^0=0$，则$\vct k\cdot\vct\epsilon=0$，得到图\ref{fig:m3gaugeq}的两个独立横向极化方向。
\Diagram{fig:m3gaugea}{规范变化沿四维波矢方向，横向物理场不变}{\draw[axis](0,0)--(4.3,0)node[right]{$k^\mu$方向};\draw[axis](0,0)--(0,2.2)node[above]{横向};\draw[->,thick](0,0)--(1.5,1.6)node[above]{$\epsilon^\mu$};\draw[->,thick](0,0)--(3.2,1.6)node[right]{$\epsilon^\mu+ak^\mu$};\draw[guide](1.5,1.6)--(3.2,1.6);}
场张量由反对称组合$k^\mu\epsilon^\nu-k^\nu\epsilon^\mu$给出。把图\ref{fig:m3gaugea}中$\epsilon^\mu$改为$\epsilon^\mu+ak^\mu$，新增项$a(k^\mu k^\nu-k^\nu k^\mu)=0$，所以场不变。并且$k^2=0$使$k\cdot\epsilon$约束仍成立。图为四维分量关系示意，不是普通空间中把电场沿传播方向倾斜。
\end{solution}

\begin{exercise}\label{ex:m3parallelfields}
某处电磁场平行且均非零。能否通过惯性变换同时消去其中一个场？若两场本来不平行，两个场不变量可以提供哪些必要限制？
\Diagram{fig:m3parallelq}{平行电磁场的点积不为零}{\draw[axis](0,0)--(0,1.65)node[above]{$\vct E$};\draw[axis](1.2,0)--(1.2,1.35)node[above]{$\vct B$};\node[below]at(.6,0){同向或反向均可};}
\end{exercise}
\begin{solution}
图\ref{fig:m3parallelq}中$\vct E\cdot\vct B=\pm EB\ne0$，它是洛伦兹不变量。若在某系消去任意一个场，点积就为零，矛盾。因此不存在这种惯性系。
\Diagram{fig:m3parallela}{消场必须先满足两个不变量的限制}{\draw[axis](0,0)--(5,0)node[right]{$E^2-c^2B^2$};\draw[curve](2.5,-.6)--(2.5,.8);\node[below]at(2.5,-.6){$0$};\node[above]at(1,.2){纯磁场可能};\node[above]at(4,.2){纯电场可能};\node[below]at(2.5,-1.1){还须$\vct E\cdot\vct B=0$};}
一般场若可变为纯电场，必要有$\vct E\cdot\vct B=0$且$E^2-c^2B^2>0$；若可变为纯磁场，则后一不变量小于0。图\ref{fig:m3parallela}的零边界需要另行处理：两个不变量都零的非零场不能被消去，因为若某系变成零场，逆变换会使原系也为零。对正交场，例\ref{ex:m3crossed}已给出满足这些条件的具体速度。
\end{solution}


\chapter{光量子、热辐射与半经典量子模型}
\section{光电效应与表面势垒}\label{th:m4photo}
电子经电压加速，撞击靶材可以把动能转化为X射线。若单个电子最多把全部动能交给一个光子，截止波长满足$hc/\lambda_{\min}=eU$。光电效应研究相反的能量转移：入射光使材料中的电子逸出。单光子模型中，光子能量为$h\nu$，电子逸出需克服功函数$W$，最大动能
\[
 K_{\max}=h\nu-W,\qquad eU_s=K_{\max}.
\]
其中$e>0$，$U_s$为使最快光电子也到不了收集极的反向电压大小。只有$\nu\geq W/h$才可逸出。提高入射功率通常增加单位时间可吸收的光子数，不会把每个光子的能量改为更大；截止频率和遏止电压随频率的变化，因而由单光子能量决定。

真实金属还存在能带占据、散射、表面状态及逸出方向分布。下面把这些细节简化为明确的势垒模型，用以讨论方向筛选；其逃逸比例不应直接当作任意材料的实测量子效率。

\begin{example}\label{ex:m4escape}
平面表面法线沿$z$正向，材料位于$z<0$。每个被吸收的光子使一个初始动能可忽略的电子获得总动能$h\nu$，速度方向在整个$4\pi$内各向同性。表面势垒只改变法向动量，势垒高度为$W$。求逃逸概率，并求收集极施加反向电压$U\geq0$时的电流。吸收光功率为$P$，忽略空间电荷和二次散射。
\Diagram{fig:m4escapeq}{电子方向角相对外法线量取；只有法向能量用于越过势垒}{
\draw[curve](-2,0)--(2,0);\node[left]at(-2,-.5){材料};
\draw[axis](0,-.7)--(0,2.5)node[above]{$z$};
\draw[->,thick](0,0)--(1.3,1.8)node[right]{$\vct v$};
\draw[guide](0,0)--(-1.3,1.8);
\draw(0,.72)arc(90:54.16:.72);\node at(.31,.92){$\theta$};
\node[right]at(2,.2){$W$};
}
\end{example}
\begin{solution}
图\ref{fig:m4escapeq}的电子法向初动能为$h\nu\cos^2\theta$。平行表面的动量守恒，不能用于克服法向势垒；逸出还必须$\cos\theta>0$。若$h\nu<W$，概率为零；否则允许锥满足$\cos\theta\geq\sqrt{W/(h\nu)}$，所以
\[
 \mathcal P_0=\frac1{4\pi}\int_0^{2\pi}\dd\varphi
 \int_0^{\theta_c}\sin\theta\dd\theta
 =\frac12[-\cos\theta]_0^{\theta_c}
 =\frac12\left(1-\sqrt{\frac W{h\nu}}\right).
\]
\Diagram{fig:m4escapea}{反向电压缩小可被收集的方向锥}{\draw[axis](0,0)--(0,2.8)node[above]{法线};\draw[curve](0,0)--(1.5,2);\draw[curve](0,0)--(-1.5,2);\draw[very thick](0,0)--(.7,2);\draw[very thick](0,0)--(-.7,2);\node[right]at(1.5,2){$U=0$};\node[right]at(.7,2.6){$U>0$};}
反向电场进一步消耗法向能量$eU$。图\ref{fig:m4escapea}中的窄锥满足$h\nu\cos^2\theta\geq W+eU$，由相同积分得
\[
 I(U)=\frac{eP}{2h\nu}
 \left[1-\sqrt{\frac{W+eU}{h\nu}}\right],\quad
 0\leq eU\leq h\nu-W.
\]
超过该区间时$I=0$。即使$h\nu\gg W$，此模型零电压的逃逸比例也只趋向$1/2$，因为另半球的电子朝材料内部。阈值电压仍为$U_s=(h\nu-W)/e$，与简单能量关系一致。
\end{solution}

\begin{example}\label{ex:m4stopping}
同一材料先后受频率$\nu_1,\nu_2$照射，测得遏止电压$U_1,U_2$，其中$\nu_2>\nu_1$。推导由测量求$h/e$和功函数的方法，并解释光电流加倍是否改变结果。
\Diagram{fig:m4stoppingq}{光电管的反向电压测量}{\draw[curve](0,0)rectangle(4,2);\draw[very thick](.8,.4)--(.8,1.6);\draw[very thick](3.2,.4)--(3.2,1.6);\draw[->](-1,1.4)--(.7,1.4)node[midway,above]{$\nu$};\node[below]at(.8,.3){阴极};\node[below]at(3.2,.3){收集极};\node[below]at(2,-.4){反向电压$U$};}
\end{example}
\begin{solution}
图\ref{fig:m4stoppingq}对每种频率分别满足$eU_i=h\nu_i-W$。两式相减，得到$h/e=(U_2-U_1)/(\nu_2-\nu_1)$，再代回$W/e=(h/e)\nu_1-U_1$。
\Diagram{fig:m4stoppinga}{遏止电压对频率的直线斜率与截距}{\draw[axis](0,0)--(4.6,0)node[right]{$\nu$};\draw[axis](0,-.7)--(0,3)node[above]{$U_s$};\draw[curve](.8,0)--(4,2.6);\fill(1.7,.731)circle(1.5pt);\fill(3.2,1.95)circle(1.5pt);\node[below]at(.8,0){$W/h$};\node[right]at(3.3,2.6){斜率$h/e$};}
图\ref{fig:m4stoppinga}的横截距给截止频率。在线性单光子区且材料表面未改变时，增大光功率只增加电子数，直线斜率和截距不变。若电流很大产生空间电荷，或进入多光子吸收区，就超出了此模型。
\end{solution}

\section{腔辐射的模态与能量分布}\label{th:m4modes}
黑体完全吸收照到它的辐射，热平衡时其辐射谱只由温度决定。小孔腔体可近似黑体：外光多次反射后被吸收，孔口射出的光取样于腔内平衡辐射。记$u_\nu(T)\dd\nu$为单位体积在频率区间内的能量，$u=\int u_\nu\dd\nu$。

\subsection{从腔内能量到孔口通量}
各向同性辐射的单位立体角能量密度为$u/(4\pi)$。在$\dd t$内，从方向$\theta$穿过面积$A$的光来自体积$Ac\dd t\cos\theta\dd\Omega$，只积分向外半球。因此出射功率面密度
\[
 M=\frac{cu}{4\pi}\int_0^{2\pi}\dd\varphi\int_0^{\pi/2}\cos\theta\sin\theta\dd\theta=\frac{cu}4.
\]
不能写成$cu/2$，因为斜向光的法向投影使平均值再多一因子$1/2$。

\subsection{三维电磁驻波计数}
取体积$V=L^3$的立方腔，先使用周期边界，波矢各分量间隔$2\pi/L$。每个波矢有两个横向极化态，半径$k$的球壳中模态数
\[
 \dd N=2\frac{V}{(2\pi)^3}4\pi k^2\dd k.
\]
真空$k=2\pi\nu/c$，所以单位体积、单位频率的态密度
\begin{equation}\label{eq:m4density}
 g(\nu)=\frac1V\frac{\dd N}{\dd\nu}=\frac{8\pi\nu^2}{c^3}.
\end{equation}
若改用理想导体驻波边界，正整数波数间隔为$\pi/L$，只取球壳八分之一，再乘两极化，宏观极限结果相同。边界附近少数特殊模态只影响低阶有限尺寸修正。计数时不可把周期边界的全空间与驻波边界的八分球重复叠加。

\Diagram{fig:m4modes}{周期边界波矢格点与球壳；每点有两种横向极化}{\draw[axis](0,0)--(4,0)node[right]{$k_x$};\draw[axis](0,0)--(0,3.6)node[above]{$k_y$};\foreach\x in{.5,1,1.5,2,2.5,3}{\foreach\y in{.5,1,1.5,2,2.5,3}{\fill(\x,\y)circle(.7pt);}}\draw[curve](3,0)arc(0:90:3);\draw[curve](3.3,0)arc(0:90:3.3);\node[right]at(2.8,2){$k\to k+\dd k$};}

\section{从经典配分到普朗克谱}\label{th:m4planck}
经典谐振模态有两个二次型能量项，平均各为$\kb T/2$，所以每模平均总能为$\kb T$。也可把模态的相空间按能量积分：谐振子的等能椭圆所围面积正比能量，能量态密度为常数，故
\[
\langle E\rangle=\frac{\int_0^\infty E\ee^{-E/(\kb T)}\dd E}{\int_0^\infty\ee^{-E/(\kb T)}\dd E}.
\]
令$b=\kb T$。分母为$[-b\ee^{-E/b}]_0^\infty=b$；分子分部积分为$[-bE\ee^{-E/b}]_0^\infty+b\int_0^\infty\ee^{-E/b}\dd E=b^2$，故$\langle E\rangle=b=\kb T$。
若只积分一个动量变量的高斯分布，得到的是一个二次项的$\kb T/2$，不能把它当成整模能量。经典谱$u_\nu=8\pi\nu^2\kb T/c^3$在高频积分发散，这称为紫外灾难。

取单模可交换能量为$E_n=nh\nu$，$n=0,1,\ldots$。令$q=\ee^{-h\nu/(\kb T)}$，归一化配分和为$Z=\sum q^n=1/(1-q)$。利用$q\dd Z/\dd q=\sum nq^n=q/(1-q)^2$，得
\[
 \langle E\rangle=h\nu\frac q{1-q}=\frac{h\nu}{\ee^{h\nu/(\kb T)}-1}.
\]
乘态密度得到普朗克频率谱
\begin{equation}\label{eq:m4planck}
 u_\nu=\frac{8\pi h\nu^3}{c^3}\frac1{\ee^{h\nu/(\kb T)}-1}.
\end{equation}
低频$h\nu\ll\kb T$时展开指数，恢复经典谱；高频端指数压低，积分收敛。谐振子总能量还可含$h\nu/2$零点项，但这一与温度无关的基态能不作为黑体热辐射可交换能量计入此处。

\begin{example}\label{ex:m4stefan}
由普朗克谱推导$M=\sigma T^4$，并给出$\sigma$。积分中可使用$\sum_{n=1}^\infty n^{-4}=\pi^4/90$，但需说明它怎样进入结果。
\Diagram{fig:m4stefanq}{腔内平衡辐射经小孔向外半球射出}{\draw[curve](0,0)--(0,3)--(3,3)--(3,1.8);\draw[curve](3,1.2)--(3,0)--(0,0);\draw[->](3,1.5)--(5,2.4);\draw[->](3,1.5)--(5,1.5);\draw[->](3,1.5)--(5,.6);\node at(1.5,1.5){$u(T)$};\node[right]at(5,1.5){$M$};}
\end{example}
\begin{solution}
令$x=h\nu/(\kb T)$，由式\eqref{eq:m4planck}，
\[
 u=\frac{8\pi\kb^4T^4}{h^3c^3}\int_0^\infty\frac{x^3\dd x}{\ee^x-1}.
\]
对$x>0$，$(\ee^x-1)^{-1}=\sum_{n=1}^\infty\ee^{-nx}$。各项非负，允许交换求和积分。令$y=nx$，则$\dd x=\dd y/n$、上下限仍为$0,\infty$，每一项为$n^{-4}\int_0^\infty y^3\ee^{-y}\dd y$。分部积分取$\dd v=\ee^{-y}\dd y$、$v=-\ee^{-y}$，依次得
\begin{align*}
 \int_0^\infty y^3\ee^{-y}\dd y
 &=[-y^3\ee^{-y}]_0^\infty+3\int_0^\infty y^2\ee^{-y}\dd y,\\
 \int_0^\infty y^2\ee^{-y}\dd y
 &=[-y^2\ee^{-y}]_0^\infty+2\int_0^\infty y\ee^{-y}\dd y,\\
 \int_0^\infty y\ee^{-y}\dd y
 &=[-y\ee^{-y}]_0^\infty+\int_0^\infty\ee^{-y}\dd y=1.
\end{align*}
三个边界项都为零，最后一个积分为$[-\ee^{-y}]_0^\infty=1$，所以三个积分依次为$6,2,1$。因此
\[
 \int_0^\infty\frac{x^3\dd x}{\ee^x-1}=6\sum_{n=1}^\infty n^{-4}=\frac{\pi^4}{15}.
\]
\Diagram{fig:m4stefana}{角度投影积分把腔内$cu$变为孔口$cu/4$}{\draw[curve](-1.8,0)--(1.8,0);\draw[axis](0,0)--(0,2.8)node[above]{法线};\draw[axis](0,0)--(1.5,2)node[right]{$c\dd t$};\draw[guide](1.5,0)--(1.5,2)--(0,2);\node[left]at(0,1){$c\dd t\cos\theta$};}
图\ref{fig:m4stefanq}的孔口通量按图\ref{fig:m4stefana}作半球投影积分：
\[
 M=\frac{cu}{4\pi}\int_0^{2\pi}\dd\varphi
 \int_0^{\pi/2}\cos\theta\sin\theta\dd\theta
 =\frac{cu}{2}\left[\frac{\sin^2\theta}{2}\right]_0^{\pi/2}
 =\frac{cu}{4}.
\]
代入能量密度，得到
\[
 \sigma=\frac{2\pi^5\kb^4}{15h^3c^2}.
\]
因$u$是体积能量密度，$cu$的量纲为功率每面积，与$M$一致。温度的四次方来自三次频率因子和频率积分的一次因子。
\end{solution}

\begin{example}\label{ex:m4wien}
分别求频率谱$u_\nu$和波长谱$u_\lambda$峰值对应的无量纲方程，说明为何$\lambda_{\max}\nu_{\max}\ne c$。
\Diagram{fig:m4wienq}{相同能量小区间在频率轴和波长轴上宽度不同}{\draw[axis](0,1.5)--(5,1.5)node[right]{$\nu$};\draw[axis](0,0)--(5,0)node[right]{$\lambda$};\draw[very thick](1,1.5)--(2,1.5);\draw[very thick](2,0)--(4,0);\draw[guide](1,1.5)--(4,0);\draw[guide](2,1.5)--(2,0);}
\end{example}
\begin{solution}
图\ref{fig:m4wienq}中的频率区间与波长区间方向相反。能量相等应写$u_\lambda\dd\lambda=-u_\nu\dd\nu$，因$\nu=c/\lambda$，
\[
 u_\lambda=u_\nu(c/\lambda)\left|\frac{\dd\nu}{\dd\lambda}\right|
 =\frac{8\pi hc}{\lambda^5[\ee^{hc/(\lambda\kb T)}-1]}.
\]
频率谱中令$x=h\nu/(\kb T)$，只需极大化$f_3(x)=x^3/(\ee^x-1)$。对一般$f_j(q)=q^j/(\ee^q-1)$使用商法则，
\[
 f_j'(q)=\frac{jq^{j-1}(\ee^q-1)-q^j\ee^q}{(\ee^q-1)^2}
 =\frac{q^{j-1}\ee^q}{(\ee^q-1)^2}
 [j(1-\ee^{-q})-q].
\]
令$j=3$，非零峰便满足$3(1-\ee^{-x})=x$。
\Diagram{fig:m4wiena}{频率谱和波长谱的峰值各由自己的密度函数决定}{\draw[axis](0,0)--(5,0)node[right]{无量纲变量};\draw[axis](0,0)--(0,2.6)node[above]{归一化谱};\draw[curve,domain=.1:7,samples=90]plot({.6*\x},{1.4*\x^3/(exp(\x)-1)});\draw[guide](1.693,0)--(1.693,2);\node[below]at(1.693,0){$2.8214$};\node[right]at(4.3,1.8){$u_\nu$示意};}
图\ref{fig:m4wiena}的频率峰根为$x_\nu\simeq2.82144$。对波长谱令$y=hc/(\lambda\kb T)$，它正比$f_5(y)$。链式法则给$\dd u_\lambda/\dd\lambda\propto f_5'(y)\dd y/\dd\lambda$，其中$\dd y/\dd\lambda=-hc/(\kb T\lambda^2)\ne0$。在上式取$j=5$，峰条件为$5(1-\ee^{-y})=y$，非零正根$y_\lambda\simeq4.96511$。因此
\[
 \nu_{\max}=\frac{x_\nu\kb T}h,\quad
 \lambda_{\max}=\frac{hc}{y_\lambda\kb T},\quad
 \lambda_{\max}\nu_{\max}=c\frac{x_\nu}{y_\lambda}\simeq.56825c.
\]
两个方程都有形式根0，但谱在原点或相应无穷远边界趋零，它不是所求内部峰。两谱各自只有一个非零正峰；其位置不同来自密度换元的雅可比$c/\lambda^2$，不表示同一个光子不满足$\lambda\nu=c$。
\end{solution}

\section{玻尔氢样原子与相对论圆轨道}\label{th:m4bohr}
固定核电荷$+Ze$，电子质量$m$，忽略核反冲。非相对论圆轨道满足$m v^2/r=Ze^2/(4\pi\varepsilon_0r^2)$、$mvr=n\hbar$。消去$v$得到
\[
 r_n=\frac{a_0n^2}Z,\quad v_n=\frac{Z\alpha c}n,\quad
 E_n=-\frac{mc^2(Z\alpha)^2}{2n^2},
\]
其中$\alpha=e^2/(4\pi\varepsilon_0\hbar c)$、$a_0=\hbar/(mc\alpha)$，$n=1,2,\ldots$。总能量等于动能加库仑势能，利用向心关系势能为动能的$-2$倍，所以能量为负。长度$a_0$、动量$m\alpha c$、能量$m\alpha^2c^2$组成一组方便的原子单位；氢基态束缚能为这一能量单位的一半。

光谱跃迁频率由$h\nu=E_{n_i}-E_{n_f}>0$确定。圆轨道的转动频率只是大量子数极限下相邻能级间隔的近似，不可对任意$n_i,n_f$直接取轨道频率作为谱线频率。

\begin{example}\label{ex:m4relbohr}
保留圆轨道角动量量子化$pr=n\hbar$，把动力学改为相对论，库仑势保持$-Ze^2/(4\pi\varepsilon_0r)$。求$v_n,r_n$、包含静能的总能量及其低$Z\alpha/n$展开。
\Diagram{fig:m4relbohrq}{圆轨道中动量切向、库仑力径向}{\draw[curve](0,0)circle(1.5);\fill(0,0)circle(2pt)node[below]{核};\fill(1.5,0)circle(2pt);\draw[axis](1.5,0)--(1.5,1.3)node[right]{$p$};\draw[axis](1.5,0)--(.3,0)node[above]{$F$};\node[below]at(.75,0){$r$};}
\end{example}
\begin{solution}
图\ref{fig:m4relbohrq}中动量大小不变但方向转动，径向力为$pv/r$，故$pv=Ze^2/(4\pi\varepsilon_0r)$。代入$pr=n\hbar$，得到$v/c=Z\alpha/n\equiv\eta$。于是$p=\gamma_vmv$，半径为
\[
 r=\frac{n\hbar}{\gamma_vmv}=\frac{a_0n^2}Z\sqrt{1-\eta^2}.
\]
包含静能的总能量$\mathcal E=\gamma_vmc^2-Ze^2/(4\pi\varepsilon_0r)$。利用势能绝对值$pv=\gamma_vmv^2$，得
\[
 \mathcal E=\gamma_vmc^2(1-\eta^2)=mc^2\sqrt{1-\eta^2}.
\]
\Diagram{fig:m4relbohra}{相对论圆轨道的束缚能随$\eta$增加而加深}{\draw[axis](0,0)--(4.2,0)node[right]{$\eta$};\draw[axis](0,-2.6)--(0,.7)node[above]{$E_b/(mc^2)$};\draw[curve,domain=0:.98,samples=70]plot({3.6*\x},{2.3*(sqrt(1-\x*\x)-1)});\draw[guide,domain=0:1,samples=40]plot({3.6*\x},{-1.15*\x*\x});}
令$z=\eta^2$，对$f(z)=\sqrt{1-z}$求导有$f'(0)=-1/2$、$f''(0)=-1/4$。泰勒展开$f(z)=1+f'(0)z+f''(0)z^2/2+O(z^3)=1-z/2-z^2/8+O(z^3)$。减去静能后，图\ref{fig:m4relbohra}的束缚能为
\[
 E_b=-mc^2\left(\frac{\eta^2}2+\frac{\eta^4}8+O(\eta^6)\right).
\]
首项恢复玻尔能级。模型要求$\eta<1$即$n>Z\alpha$。它只修正半经典圆轨道，不含自旋、非圆轨道的角动量结构与完整量子动力学，因此不能把此式当作全部精细结构能级。
\end{solution}

\section{作用量量子化与转折边界}\label{th:m4action}
对一维束缚运动，定义一个完整周期的作用量$J=\oint p\dd x$。这是相平面闭合轨道所围面积。圆周运动则$J_\varphi=\int_0^{2\pi}L\dd\varphi=2\pi L$，旧量子条件$J_\varphi=nh$给$L=n\hbar$。

一维波函数在允许区的短波形式含相位$\int p\dd x/\hbar$。边界反射时还会附加相位：两端硬壁使$J=nh$，$n=1,2,\ldots$；两端平滑转折点使$J=(n+1/2)h$，$n=0,1,\ldots$；一硬壁、一平滑转折点使$J=(n+3/4)h$。后两式是半经典近似，一般在量子数较大时更可靠。谐振子恰能由该条件给出精确能级，但不能据此认定任意势的半经典最低能级也精确。

\begin{example}\label{ex:m4well}
质量$m$的粒子限制在$0<x<L$，外侧无限高势垒，内部势能为零。由驻波与作用量两种方法求能级，并写出归一化波函数。
\Diagram{fig:m4wellq}{无限深方阱的两个硬边界}{\draw[axis](-.5,0)--(4.8,0)node[right]{$x$};\draw[curve](0,3)--(0,0)--(3.5,0)--(3.5,3);\node[below]at(0,0){$0$};\node[below]at(3.5,0){$L$};\node[above]at(1.75,2.1){$V=0$};}
\end{example}
\begin{solution}
图\ref{fig:m4wellq}的阱内满足$-\hbar^2\psi''/(2m)=E\psi$。令$k=\sqrt{2mE}/\hbar$，通解为$A\sin kx+B\cos kx$。左壁$\psi(0)=0$给$B=0$；右壁$\psi(L)=0$要求$kL=n\pi$，$n=1,2,\ldots$。$n=0$只给处处为零的函数，不能归一化，故不属于允许态。
\[
 E_n=\frac{\hbar^2k_n^2}{2m}=\frac{n^2h^2}{8mL^2}.
\]
\Diagram{fig:m4wella}{方阱相轨道为两条水平动量支，面积为$2pL$}{\draw[axis](-.3,0)--(4.6,0)node[right]{$x$};\draw[axis](0,-1.8)--(0,2)node[above]{$p$};\draw[very thick,->](0,1.2)--(3.5,1.2);\draw[very thick,->](3.5,-1.2)--(0,-1.2);\draw[guide](0,1.2)--(0,-1.2);\draw[guide](3.5,1.2)--(3.5,-1.2);\node[right]at(3.5,1.2){$p_n$};}
归一化积分逐步为
\[
 1=|A|^2\int_0^L\frac{1-\cos(2n\pi x/L)}2\dd x
 =|A|^2\left[\frac x2-\frac{L}{4n\pi}\sin\frac{2n\pi x}{L}\right]_0^L
 =|A|^2L/2,
\]
故可取$\psi_n=\sqrt{2/L}\sin(n\pi x/L)$。图\ref{fig:m4wella}的相轨道上，向右支$p=p_n$、$\dd x>0$；向左支$p=-p_n$、$\dd x<0$，因此两支贡献都为正：
\[
 J=\int_0^L p_n\dd x+\int_L^0(-p_n)\dd x
 =[p_nx]_0^L+[-p_nx]_L^0=2p_nL.
\]
用两个硬壁条件$J=nh$同样得到$p_n=nh/(2L)$。竖直碰壁段满足$\dd x=0$，不增加作用量。这个定态是两个反向行波的叠加，$\langle p\rangle=0$，却有$\langle p^2\rangle=p_n^2$；不能将相轨道上的固定动量大小解释为定态粒子总沿某一个方向运动。
\end{solution}

\begin{example}\label{ex:m4osc}
一维谐振子势$V(x)=m\omega^2x^2/2$。用相平面面积求半经典能级，并比较直接使用$J=nh$会遗漏什么。
\Diagram{fig:m4oscq}{谐振子的两个平滑转折点}{\draw[axis](-2.3,0)--(2.5,0)node[right]{$x$};\draw[axis](0,0)--(0,3)node[above]{$V$};\draw[curve,domain=-1.9:1.9,samples=60]plot(\x,{.65*\x*\x});\draw[guide](-1.7,1.88)--(1.7,1.88);\node[right]at(1.7,1.88){$E$};}
\end{example}
\begin{solution}
图\ref{fig:m4oscq}的势为$V=m\omega^2x^2/2$。固定能量的相轨道满足
\[
 \frac{x^2}{2E/(m\omega^2)}+\frac{p^2}{2mE}=1,
\]
故两半轴$x_m=\sqrt{2E/(m\omega^2)}$、$p_m=\sqrt{2mE}$。作用量是顺运动方向的闭合积分，上下两支都给正面积贡献。
\Diagram{fig:m4osca}{椭圆所围面积$J=\pi x_mp_m$}{\draw[axis](-2.6,0)--(2.8,0)node[right]{$x$};\draw[axis](0,-1.9)--(0,2.2)node[above]{$p$};\draw[curve](0,0)ellipse(2 and 1.4);\node[below right]at(2,0){$x_m$};\node[above right]at(0,1.4){$p_m$};\draw[->](1.8,.61)--(1.95,.31);}
图\ref{fig:m4osca}的椭圆面积直接给$J=\pi x_mp_m=2\pi E/\omega$。也可令$x=x_m\sin\varphi$作积分：
\[
 J=4p_mx_m\int_0^{\pi/2}\cos^2\varphi\dd\varphi
 =4p_mx_m\left[\frac\varphi2+\frac{\sin2\varphi}{4}\right]_0^{\pi/2}
 =\pi p_mx_m.
\]
两个转点平滑，取$J=(n+1/2)h$，$n=0,1,\ldots$，于是$E_n=\hbar\omega(n+1/2)$，相邻能级间隔为$\hbar\omega$。对谐振子这一量子条件恰给出精确能级，但不能据此认定任意势的最低能级都由半经典条件精确给出。

若用$J=nh$并允许$n=0$，便给出$E=0$、$x=p=0$，与量子态不能同时具有确定位置和动量的要求不符。零点项来自转点相位，并非任意选取能量零点；给定势最低点为零后，基态相对最低点仍有$\hbar\omega/2$。由$\dd J/\dd E=2\pi/\omega$还恢复能量无关的经典周期。
\end{solution}

\begin{example}\label{ex:m4bounce}
粒子在$x\geq0$受重力势$V=mgx$，$x<0$为无限高硬壁。用一硬一软的量子条件求能级近似，并说明其量纲与随量子数的变化。
\Diagram{fig:m4bounceq}{底部硬壁与上方重力转折点}{\draw[axis](-.4,0)--(4.5,0)node[right]{$x$};\draw[axis](0,0)--(0,3.2)node[above]{$V$};\draw[very thick](0,0)--(0,3);\draw[curve](0,0)--(3.8,2.8);\draw[guide](0,2)--(2.714,2);\node[right]at(2.8,2){$x_t=E/(mg)$};}
\end{example}
\begin{solution}
图\ref{fig:m4bounceq}中的允许区是$0\leq x\leq E/(mg)$，动量大小$p(x)=\sqrt{2m(E-mgx)}$。完整往返作用量
\begin{align*}
 J&=2\int_0^{E/(mg)}\sqrt{2m(E-mgx)}\dd x\\
 &=-\frac{2\sqrt{2m}}{mg}\int_E^0 q^{1/2}\dd q
 =\frac{2\sqrt{2m}}{mg}\left[\frac23q^{3/2}\right]_0^E
 =\frac{4\sqrt2}{3g\sqrt m}E^{3/2}.
\end{align*}
这里取$q=E-mgx$，故$\dd x=-\dd q/(mg)$，空间上下限$0,E/(mg)$依次变为$E,0$。
\Diagram{fig:m4bouncea}{上下动量支在平滑转点相连，在硬壁处反射}{\draw[axis](-.3,0)--(4.5,0)node[right]{$x$};\draw[axis](0,-2)--(0,2.2)node[above]{$p$};\draw[curve,domain=0:3.5,samples=70]plot(\x,{1.6*sqrt(1-\x/3.5)});\draw[curve,domain=0:3.5,samples=70]plot(\x,{-1.6*sqrt(1-\x/3.5)});\draw[very thick](0,-1.6)--(0,1.6);\node[below]at(3.5,0){$x_t$};}
图\ref{fig:m4bouncea}中一端是硬壁，一端是平滑转点，故取$J=(n+3/4)h$，得到
\[
 E_n=\left[\frac{3hg\sqrt m}{4\sqrt2}\left(n+\frac34\right)\right]^{2/3},\quad n=0,1,\ldots.
\]
括号内量纲为能量的$3/2$次方。能量随$n^{2/3}$增长，所以高能级的相邻间隔减小。此式是半经典近似；精确解应以满足硬壁条件的艾里函数零点确定，不能把$3/4$条件称为最低几个能级的精确值。
\end{solution}
\section{从波长关系到束缚态方程}
\subsection{经典轨迹和量子态的不同描述}
玻尔模型把经典圆轨道与量子化条件相结合，能解释氢样原子的主要能量尺度，却没有给出一般波函数和测量概率。处理波函数节点、势垒穿透和一般角动量态，需要量子态的描述。更完整的一维描述使用复波函数$\psi(x,t)$，$|\psi|^2$为位置概率密度，满足归一化$\int|\psi|^2\dd x=1$。

自由粒子的德布罗意关系$p=\hbar k$、$E=\hbar\omega$与非相对论色散$E=p^2/(2m)$相容于方程$i\hbar\partial_t\psi=-\hbar^2\partial_x^2\psi/(2m)$。加入给定势能$V(x)$，非相对论量子理论采用
\[
 i\hbar\partial_t\psi=\left[-\frac{\hbar^2}{2m}\partial_x^2+V(x)\right]\psi.
\]
这是一条量子动力学的基本方程；前面的自由粒子关系只说明它与德布罗意关系相容。对不含时势，取$\psi=\varphi(x)\ee^{-iEt/\hbar}$，则$\partial_t\psi=(-iE/\hbar)\psi$、$\partial_x^2\psi=\varphi''\ee^{-iEt/\hbar}$。代入并约去共同的时间因子，得到$-\hbar^2\varphi''/(2m)+V\varphi=E\varphi$。允许能量还受边界条件约束，不能只由局部$p^2=2m(E-V)$确定。

\subsection{边界连续性与允许区}
有限势跳变处，积分定态方程跨过一个无穷小区间。只要势没有狄拉克型奇点，积分中的有限项趋零，得到导数连续；波函数本身也连续，避免二阶导数产生未被势抵消的奇异项。无限高硬壁是极限边界，波函数在壁外为零、壁处为零，壁内导数不需要与外侧零函数继续相等。

在$E>V$区，局部波数实，波函数可振荡；$E<V$区，局部波数纯虚，波函数可以指数衰减或增长。有限势垒内的非零尾部不表示经典粒子拥有负动能，而是波函数满足另一类局部微分方程。对半无限禁区排除增长解，对有限禁区则一般保留两支以匹配两端。

\Diagram{fig:m4allowed}{线性势的艾里解平滑连接允许区与禁区}{\draw[axis](-.4,0)--(6,0)node[right]{$x$};\draw[axis](0,-1.2)--(0,1.9)node[above]{$\varphi$};\draw[guide](3,-1.2)--(3,1.8);\draw[curve]plot coordinates {(0.00000,-0.9470357) (0.05046,-0.9038078) (0.10092,-0.8537957) (0.15138,-0.7974842) (0.20183,-0.7353909) (0.25229,-0.6680602) (0.30275,-0.5960570) (0.35321,-0.5199612) (0.40367,-0.4403617) (0.45413,-0.3578511) (0.50459,-0.2730203) (0.55505,-0.1864540) (0.60550,-0.0987258) (0.65596,-0.0103941) (0.70642,0.0780016) (0.75688,0.1659438) (0.80734,0.2529406) (0.85780,0.3385279) (0.90826,0.4222713) (0.95872,0.5037687) (1.00917,0.5826506) (1.05963,0.6585822) (1.11009,0.7312630) (1.16055,0.8004281) (1.21101,0.8658473) (1.26147,0.9273255) (1.31193,0.9847018) (1.36239,1.0378485) (1.41284,1.0866705) (1.46330,1.1311037) (1.51376,1.1711138) (1.56422,1.2066944) (1.61468,1.2378660) (1.66514,1.2646733) (1.71560,1.2871841) (1.76606,1.3054872) (1.81651,1.3196902) (1.86697,1.3299179) (1.91743,1.3363102) (1.96789,1.3390201) (2.01835,1.3382122) (2.06881,1.3340601) (2.11927,1.3267457) (2.16972,1.3164565) (2.22018,1.3033846) (2.27064,1.2877250) (2.32110,1.2696741) (2.37156,1.2494285) (2.42202,1.2271835) (2.47248,1.2031323) (2.52294,1.1774648) (2.57339,1.1503667) (2.62385,1.1220186) (2.67431,1.0925954) (2.72477,1.0622656) (2.77523,1.0311908) (2.82569,0.9995253) (2.87615,0.9674155) (2.92661,0.9349999) (2.97706,0.9024089) (3.02752,0.8697645) (3.07798,0.8371804) (3.12844,0.8047618) (3.17890,0.7726057) (3.22936,0.7408008) (3.27982,0.7094279) (3.33028,0.6785597) (3.38073,0.6482616) (3.43119,0.6185912) (3.48165,0.5895995) (3.53211,0.5613302) (3.58257,0.5338209) (3.63303,0.5071029) (3.68349,0.4812018) (3.73394,0.4561376) (3.78440,0.4319255) (3.83486,0.4085756) (3.88532,0.3860938) (3.93578,0.3644820) (3.98624,0.3437383) (4.03670,0.3238576) (4.08716,0.3048314) (4.13761,0.2866488) (4.18807,0.2692963) (4.23853,0.2527582) (4.28899,0.2370171) (4.33945,0.2220536) (4.38991,0.2078474) (4.44037,0.1943766) (4.49083,0.1816186) (4.54128,0.1695500) (4.59174,0.1581466) (4.64220,0.1473842) (4.69266,0.1372379) (4.74312,0.1276832) (4.79358,0.1186951) (4.84404,0.1102491) (4.89450,0.1023207) (4.94495,0.0948859) (4.99541,0.0879210) (5.04587,0.0814028) (5.09633,0.0753088) (5.14679,0.0696167) (5.19725,0.0643052) (5.24771,0.0593535) (5.29817,0.0547415) (5.34862,0.0504499) (5.39908,0.0464601) (5.44954,0.0427541) (5.50000,0.0393148)};\node[above]at(1.2,1.8){允许区};\node[above]at(4.5,1.8){禁区};}
图\ref{fig:m4allowed}画出线性势中的同一个艾里函数解：转点左侧振荡，右侧衰减，函数及其导数在转点处连续。两区振幅由这个连接条件确定，不能分别任取。有限方阱解中的$\arctan(k/\kappa)$边界相位同样来自函数和导数的同时匹配。

\section{WKB相位、振幅和量子偏移}
令定态波函数$\varphi=A\ee^{iS/\hbar}$。先求一次导数$\varphi'=(A'+iAS'/\hbar)\ee^{iS/\hbar}$，再对括号与指数分别求导，得到
\[
 \varphi''=\left[A''+\frac{2i}{\hbar}A'S'
 +\frac{i}{\hbar}AS''-\frac1{\hbar^2}A(S')^2\right]\ee^{iS/\hbar}.
\]
当$A,S$相对局部波长缓变，领先阶给$(S')^2=2m(E-V)=p^2(x)$；下一阶虚部给$2A'p+Ap'=0$。在$p>0$的一个允许区间内，以$x_0$为参考点分离变量，
\[
 \frac{A'}A=-\frac{p'}{2p},\qquad
 \int_{A(x_0)}^{A(x)}\frac{\dd a}{a}
 =-\frac12\int_{p(x_0)}^{p(x)}\frac{\dd q}{q},\qquad
 \ln\frac{A(x)}{A(x_0)}=-\frac12\ln\frac{p(x)}{p(x_0)}.
\]
因此$A=C/\sqrt p$，允许区的两支为
\[
 \varphi_\pm\simeq\frac{C_\pm}{\sqrt{p(x)}}
 \exp\left[\pm\frac i\hbar\int^x p(q)\dd q\right].
\]
检验忽略$A''$的条件，还须把振幅代回原方程。对$A=Cp^{-1/2}$连续求导，
\[
 A'=-\frac C2p^{-3/2}p',\qquad
 A''=\frac{3C}{4}p^{-5/2}(p')^2-\frac C2p^{-3/2}p'',
\]
故相对于领先项$p^2A$的余项为
\[
 \frac{\hbar^2A''}{p^2A}
 =\hbar^2\left[\frac{3(p')^2}{4p^4}-\frac{p''}{2p^3}\right].
\]
一组充分的局部缓变条件是$\hbar|p'|/p^2\ll1$及$\hbar^2|p''|/p^3\ll1$。长区间的累计误差还需结合具体势估计。在转点$p=0$附近条件失效，不能把上述振幅直接延伸穿过转点。

\subsection{平滑转点的四分之一周期相移}
取右侧简单转点$x_t$，即$V(x_t)=E$、$F_t=V'(x_t)>0$，允许区在$x<x_t$。局部线性化$V-E\simeq F_t(x-x_t)$，定态方程成为$\varphi''=(2mF_t/\hbar^2)(x-x_t)\varphi$。令$\ell_t=(\hbar^2/(2mF_t))^{1/3}$、$z=(x-x_t)/\ell_t$，则$\varphi''(x)=\ell_t^{-2}\varphi_{zz}$，代入得$\varphi_{zz}-z\varphi=0$。禁区$z>0$选衰减的$\operatorname{Ai}(z)$。利用艾里函数在$z<0$的渐近连接式，允许区的驻波为
\[
 \varphi(x)\propto\frac1{\sqrt{p(x)}}
 \sin\left[\frac1\hbar\int_x^{x_t}p(q)\dd q+\frac\pi4\right].
\]
这里的相移是局部线性势的连接结果，不能用于突然跳到无限高的硬壁。两个平滑转点各带相同类型修正，闭合相位条件成为$\int_{x_1}^{x_2}p\dd x=(n+1/2)\pi\hbar$，即$J=(n+1/2)h$。

若左端是硬壁$x=0$、右端是平滑转点，将上式连接解代入$\varphi(0)=0$，得到
\[
 \frac1\hbar\int_0^{x_t}p\dd x+\frac\pi4=(n+1)\pi,
\]
所以$J=2\int_0^{x_t}p\dd x=(n+3/4)h$。这给出重力反弹中的$3/4$。两硬壁之间则直接要求站波节点，$\int_0^Lp\dd x=n\pi\hbar$，$J=nh$，从$n=1$开始。三类边界的量子数起点和偏移不同，比较时须把编号一并说明。

\Diagram{fig:m4boundarytypes}{硬壁节点与平滑转点连接所对应的相位条件不同}{\begin{scope}[xshift=-2.4cm]\draw[axis](0,0)--(2.2,0);\draw[very thick](0,-.4)--(0,1.7);\draw[curve,domain=0:2,samples=60]plot(\x,{.9*sin(120*\x)});\node[below]at(1,-.7){硬壁节点};\end{scope}\begin{scope}[xshift=1.2cm]\draw[axis](0,0)--(2.5,0);\draw[curve](0,.7)..controls(.5,1.3)and(.7,1.1)..(1,.5)..controls(1.4,.1)and(1.8,.02)..(2.3,0);\draw[guide](1,-.4)--(1,1.7);\node[below]at(1.2,-.7){平滑转点};\end{scope}}
图\ref{fig:m4boundarytypes}仅区分边界形态。半经典能级公式通常对高量子数较准，因为大部分允许区包含许多波长，转点局部修正占比小。重力反弹可用艾里零点检验误差，谐振子则因二次势的特殊结构使领先WKB能级恰好精确。精确的一个例子不能消除近似条件。

\section{零点能和位置—动量展宽}
对归一化且边界项消失的波函数，定义$\Delta x^2=\langle(x-\langle x\rangle)^2\rangle$、$\Delta p^2=\langle(p-\langle p\rangle)^2\rangle$。取$f=(x-\langle x\rangle)\psi$、$g=(p-\langle p\rangle)\psi$，柯西不等式给$\Delta x^2\Delta p^2\geq|\langle f,g\rangle|^2$。由算符的厄米性及对易关系，
\[
 2i\,\operatorname{Im}\langle f,g\rangle
 =\langle f,g\rangle-\langle g,f\rangle
 =\langle\psi,[x,p]\psi\rangle=i\hbar.
\]
于是$|\langle f,g\rangle|^2\geq(\operatorname{Im}\langle f,g\rangle)^2=\hbar^2/4$，得到
\[
 \Delta x\Delta p\geq\frac\hbar2.
\]
这里的$\Delta x,\Delta p$是同一量子态的概率标准差。位置分布越窄，动量分布的最小展宽越大；方阱基态的非零动能和谐振子零点能都满足这一约束。

对以平衡点为中心、平均动量为零的谐振子态，能量平均满足
\[
 \langle E\rangle=\frac{\Delta p^2}{2m}+\frac{m\omega^2\Delta x^2}2
 \geq\frac{\hbar^2}{8m\Delta x^2}+\frac{m\omega^2\Delta x^2}2.
\]
记$z=\Delta x^2>0$，右侧函数为$f(z)=\hbar^2/(8mz)+m\omega^2z/2$。逐项求导，
\[
 f'(z)=-\frac{\hbar^2}{8mz^2}+\frac{m\omega^2}{2},\qquad
 f''(z)=\frac{\hbar^2}{4mz^3}>0.
\]
令$f'=0$，正根为$z=\hbar/(2m\omega)$，代回得$\langle E\rangle\geq\hbar\omega/2$。可取$\varphi=C\ee^{-b x^2/2}$核实等号态。两次求导为$\varphi'=-bx\varphi$、$\varphi''=(-b+b^2x^2)\varphi$，代入定态方程给
\[
 E\varphi=\left[\frac{\hbar^2b}{2m}
 +\left(\frac{m\omega^2}{2}-\frac{\hbar^2b^2}{2m}\right)x^2\right]\varphi.
\]
定态能量不依赖$x$，故$b=m\omega/\hbar$（取正根以使无穷远衰减），$E=\hbar\omega/2$。零点项同时由量子展宽下界和满足边界条件的波函数确定。

\section{黑体谱中经典近似的边界}
经典均分定理的前提是能量可在$\kb T$尺度上近似连续地取值。对低频模态，能量间隔$h\nu\ll\kb T$，相邻占据数许多，几何分布的离散和可近似为连续积分，于是每模平均$\kb T$。对高频模态，第一次激发就需远大于$\kb T$的能量，绝大多数模态停留在$n=0$，平均能量被玻尔兹曼因子压低。

紫外灾难的根源因此是把低频的连续能量近似用于全部无限多个高频模态。修正只改模态计数是不够的：真空电磁波的态密度$\nu^2$来自空间维数与色散；量子化改变的是各模态的热占据。计算时分别求出态密度和平均热占据，再将二者相乘。

黑体辐射的总能量按$T^4$、光子数按$T^3$变化，平均光子能量因而正比于$T$。把频率谱的横轴换为$h\nu/(\kb T)$、纵轴除以$T^3$，不同温度的曲线重合。实际物体若有选择性吸收、反射、有限模态或非平衡照明，其谱形可偏离这一关系。

\subsection{作用量中的反射相位与量子数编号}
一维束缚运动的作用量为$J(E)=\oint p\dd x=2\int_{x_-}^{x_+}|p(x)|\dd x$。返回支的$p$与$\dd x$同时反号，两支对作用量均给出正贡献。波的闭合相位还包括转点反射相位，因此不同边界不能统一写成$J=nh$。

两个平滑转点时，每个转点引入四分之一周的相位修正，量子条件为$J=(n+1/2)h$，$n=0,1,\ldots$。一硬一软边界给$J=(n+3/4)h$，仍从$n=0$编号；两个无限硬壁则给$J=nh$，应从$n=1$编号。三个式子的$n$从不同数值起算，分别对应各边界问题的节点编号。
\Diagram{fig:m4boundaryphase}{允许区的三类端点；硬壁与平滑转点采用不同反射条件}{\begin{scope}\draw[axis](0,0)--(1.8,0);\draw[curve]plot[smooth]coordinates{(0,1.4)(.3,.5)(.7,.1)(1.1,.5)(1.4,1.4)};\draw[guide](0,.8)--(1.5,.8);\node[below]at(.7,-.2){两软};\end{scope}\begin{scope}[xshift=2.2cm]\draw[axis](0,0)--(1.8,0);\draw[thick](.2,0)--(.2,1.5);\draw[curve](.2,.1)--(1.5,1.4);\draw[guide](.2,.8)--(1.6,.8);\node[below]at(.8,-.2){一硬一软};\end{scope}\begin{scope}[xshift=4.5cm]\draw[axis](0,0)--(1.8,0);\draw[thick](.2,1.5)--(.2,.1)--(1.4,.1)--(1.4,1.5);\draw[guide](.2,.8)--(1.4,.8);\node[below]at(.8,-.2){两硬};\end{scope}}
图\ref{fig:m4boundaryphase}的虚线表示同一类比较中的能量线，交点限定作用量的积分区间。有限高度的陡峭势垒还会产生随能量变化的反射相位，需由波函数与导数的连续条件求出；它通常介于软转点近似与无限硬壁极限之间。

还可把相位条件写成$J(E)=h[n+\eta(E)]$。对量子数求导，
\[
 \left[\frac{\dd J}{\dd E}-h\frac{\dd\eta}{\dd E}\right]\frac{\dd E}{\dd n}=h.
\]
若$\eta$为常数且半经典近似成立，便有$\Delta E\simeq h/T$，其中$T=\dd J/\dd E$是经典周期。若边界相位显著依赖能量，能级间隔还含$\dd\eta/\dd E$修正。无限深方阱可直接解边界方程，重力反弹低能级则应对照艾里函数零点；使用作用量时须区分精确边界解与渐近量子条件。

\section{光子统计与辐射能量尺度}
普朗克式中的$n$可以看作一个模态的光子占据数。黑体与腔壁可交换、产生和吸收光子，热平衡不限制光子总数，因而没有独立的光子数化学势。单模概率$p_n=(1-q)q^n$，其中$q=\ee^{-h\nu/(\kb T)}$。不同模态的占据独立时，总能量平均是各模能量平均之和；这不等于把每个频率区间强行塞入一个光子。

\begin{example}\label{ex:m4photonnumber}
求温度$T$的黑体辐射光子数密度$n_\gamma$和平均光子能量。可把$\zeta(3)=\sum_{j=1}^\infty j^{-3}$保留为常数。
\Diagram{fig:m4numberq}{能量谱除以单光子能量后才得到光子数谱}{\draw[curve](0,0)rectangle(2.3,1.1);\node at(1.15,.55){$u_\nu\dd\nu$};\draw[->](2.5,.55)--(3.5,.55)node[midway,above]{除$h\nu$};\node[right]at(3.7,.55){$\dd n_\gamma$};}
\end{example}
\begin{solution}
图\ref{fig:m4numberq}中，每个频率为$\nu$的光子能量是$h\nu$。把单位体积、单位频率能量密度除以它，再对频率积分：
\[
 n_\gamma=\int_0^\infty\frac{u_\nu}{h\nu}\dd\nu
 =\frac{8\pi\kb^3T^3}{h^3c^3}
 \int_0^\infty\frac{x^2}{\ee^x-1}\dd x,
 \quad x=\frac{h\nu}{\kb T}.
\]
展开$(\ee^x-1)^{-1}=\sum_{j\geq1}\ee^{-jx}$，各项非负，允许逐项积分。令$y=jx$、$\dd x=\dd y/j$，则
\[
 \int_0^\infty x^2\ee^{-jx}\dd x
 =j^{-3}\int_0^\infty y^2\ee^{-y}\dd y=\frac2{j^3}.
\]
最后一个积分的两次分部积分与端点计算已在例\ref{ex:m4stefan}给出。
\Diagram{fig:m4numbera}{光子数随$T^3$，总能量随$T^4$，平均能量随$T$}{\draw[axis](0,0)--(4.6,0)node[right]{$T$};\draw[axis](0,0)--(0,3)node[above]{$\langle E\rangle$};\draw[curve](0,0)--(4,2.5);\node[above]at(2.2,2.7){斜率$\pi^4\kb/[30\zeta(3)]$};}
积分因此为$2\zeta(3)$，得到
\[
 n_\gamma=\frac{16\pi\zeta(3)\kb^3T^3}{h^3c^3},\qquad
 \langle E\rangle=\frac{u}{n_\gamma}
 =\frac{\pi^4}{30\zeta(3)}\kb T\simeq2.701\kb T.
\]
图\ref{fig:m4numbera}给出按光子数加权的平均能量。相应概率密度为$p_N(\nu)=u_\nu/(h\nu n_\gamma)$，满足$\int p_N\dd\nu=1$，故$\langle E\rangle=\int h\nu p_N\dd\nu$。能量谱峰对应的$h\nu_{\max}\simeq2.821\kb T$则由谱函数的驻值确定。若先将能量谱归一化再平均$h\nu$，使用的已是能量权重，所得数值也会改变。
\end{solution}

\begin{example}\label{ex:m4fluctuation}
对单个频率为$\nu$的热平衡模态，求占据数方差和能量方差。讨论低频和高频极限。
\Diagram{fig:m4fluctq}{单模各占据数的几何概率分布}{\draw[axis](0,0)--(5.4,0)node[right]{$n$};\draw[axis](0,0)--(0,3)node[above]{$p_n$};\foreach\x in{0,1,2,3,4}{\draw[very thick](\x,0)--(\x,{2.5*.6^\x});}\node[below]at(0,0){$0$};}
\end{example}
\begin{solution}
对图\ref{fig:m4fluctq}的单模态，令$q=\ee^{-h\nu/(\kb T)}$，有$p_n=(1-q)q^n$。从几何级数$\sum_{n\geq0}q^n=(1-q)^{-1}$开始，分别施加算符$q\dd/\dd q$，
\[
 \sum nq^n=q\frac{\dd}{\dd q}(1-q)^{-1}
 =q[-(1-q)^{-2}(-1)]=\frac q{(1-q)^2},
\]
\[
 \sum n^2q^n=q\frac{\dd}{\dd q}[q(1-q)^{-2}]
 =q[(1-q)^{-2}+2q(1-q)^{-3}]
 =\frac{q(1+q)}{(1-q)^3}.
\]
乘以$(1-q)$即得$\bar n=q/(1-q)$、$\langle n^2\rangle=q(1+q)/(1-q)^2$。
\Diagram{fig:m4flucta}{方差由粒子项和二次涨落项组成}{\draw[axis](0,0)--(4.3,0)node[right]{$\bar n$};\draw[axis](0,0)--(0,3)node[above]{$\operatorname{Var}n$};\draw[curve,domain=0:1.4,samples=60]plot({2.3*\x},{.8*(\x+\x*\x)});\draw[guide](0,0)--(3.3,1.15);}
按方差定义相减，
\[
 \operatorname{Var}n=\frac{q(1+q)-q^2}{(1-q)^2}
 =\frac q{(1-q)^2}=\bar n(1+\bar n),
\]
\[
 \operatorname{Var}E=(h\nu)^2\bar n+(h\nu\bar n)^2.
\]
图\ref{fig:m4flucta}的两项在不同极限主导：高频$h\nu\gg\kb T$时$\bar n\ll1$，方差近于$\bar n$，但相对涨落$\sqrt{\operatorname{Var}n}/\bar n$很大；低频$\bar n\simeq\kb T/(h\nu)\gg1$，能量方差趋$(\kb T)^2$，单模态的相对能量涨落仍约为1。若测量包含$M$个独立、同分布模态，总平均能量乘$M$、方差乘$M$，相对涨落才减小为单模态值的$1/\sqrt M$。低频极限的平均模能为$\kb T$，单模测量仍有同量级的能量涨落。
\end{solution}

\begin{example}\label{ex:m4mixedcavity}
小腔中的某些方向由温度$T_1$的黑面照亮，其他方向由温度$T_2$的黑面照亮，立体角份额为$f,1-f$，腔内光不再充分热化。求角平均谱，判断它能否在所有频率上等价于单一温度的黑体。
\Diagram{fig:m4mixedq}{不同温度的方向混合一般不是单温度平衡谱}{\draw[curve](0,0)circle(1.6);\draw[curve](0,0)--(1.6,0);\draw[curve](0,0)--(-.8,1.386);\node at(.1,.8){$T_1$};\node at(0,-.7){$T_2$};\fill(0,0)circle(1.5pt);}
\end{example}
\begin{solution}
图\ref{fig:m4mixedq}中的两个方向集合互不重叠、各自的辐射互不相干，故角平均谱是强度密度之和：$\bar u_\nu=fu_\nu(T_1)+(1-f)u_\nu(T_2)$。按全频积分，
\[
 \bar u=a[fT_1^4+(1-f)T_2^4],\qquad
 T_E=[fT_1^4+(1-f)T_2^4]^{1/4},
\]
其中$a=4\sigma/c$。$T_E$只使某个单温黑体具有同一总能量，并未保证逐频相同。
\Diagram{fig:m4mixeda}{混合谱与等总能量单温谱；采用相同纵轴归一化}{
\draw[axis](0,0)--(4.8,0)node[right]{$h\nu/(\kb T_2)$};
\draw[axis](0,0)--(0,3)node[above]{谱密度};
\draw[curve]plot coordinates{(0.000000,0.0000000) (0.014400,0.0008289) (0.028800,0.0032713) (0.043200,0.0072621) (0.057600,0.0127375) (0.072000,0.0196350) (0.086400,0.0278933) (0.100800,0.0374527) (0.115200,0.0482547) (0.129600,0.0602420) (0.144000,0.0733589) (0.158400,0.0875507) (0.172800,0.1027643) (0.187200,0.1189476) (0.201600,0.1360501) (0.216000,0.1540223) (0.230400,0.1728162) (0.244800,0.1923848) (0.259200,0.2126826) (0.273600,0.2336651) (0.288000,0.2552893) (0.302400,0.2775131) (0.316800,0.3002959) (0.331200,0.3235979) (0.345600,0.3473808) (0.360000,0.3716074) (0.374400,0.3962414) (0.388800,0.4212479) (0.403200,0.4465930) (0.417600,0.4722437) (0.432000,0.4981683) (0.446400,0.5243361) (0.460800,0.5507174) (0.475200,0.5772835) (0.489600,0.6040067) (0.504000,0.6308602) (0.518400,0.6578183) (0.532800,0.6848562) (0.547200,0.7119499) (0.561600,0.7390765) (0.576000,0.7662137) (0.590400,0.7933404) (0.604800,0.8204360) (0.619200,0.8474809) (0.633600,0.8744564) (0.648000,0.9013443) (0.662400,0.9281274) (0.676800,0.9547892) (0.691200,0.9813139) (0.705600,1.0076863) (0.720000,1.0338921) (0.734400,1.0599175) (0.748800,1.0857493) (0.763200,1.1113752) (0.777600,1.1367833) (0.792000,1.1619623) (0.806400,1.1869015) (0.820800,1.2115909) (0.835200,1.2360208) (0.849600,1.2601823) (0.864000,1.2840667) (0.878400,1.3076660) (0.892800,1.3309728) (0.907200,1.3539798) (0.921600,1.3766804) (0.936000,1.3990685) (0.950400,1.4211381) (0.964800,1.4428840) (0.979200,1.4643010) (0.993600,1.4853847) (1.008000,1.5061307) (1.022400,1.5265350) (1.036800,1.5465942) (1.051200,1.5663049) (1.065600,1.5856643) (1.080000,1.6046696) (1.094400,1.6233185) (1.108800,1.6416090) (1.123200,1.6595392) (1.137600,1.6771076) (1.152000,1.6943129) (1.166400,1.7111540) (1.180800,1.7276300) (1.195200,1.7437404) (1.209600,1.7594848) (1.224000,1.7748628) (1.238400,1.7898746) (1.252800,1.8045203) (1.267200,1.8188001) (1.281600,1.8327147) (1.296000,1.8462646) (1.310400,1.8594508) (1.324800,1.8722742) (1.339200,1.8847358) (1.353600,1.8968370) (1.368000,1.9085792) (1.382400,1.9199637) (1.396800,1.9309923) (1.411200,1.9416666) (1.425600,1.9519886) (1.440000,1.9619600) (1.454400,1.9715831) (1.468800,1.9808598) (1.483200,1.9897924) (1.497600,1.9983832) (1.512000,2.0066345) (1.526400,2.0145489) (1.540800,2.0221287) (1.555200,2.0293767) (1.569600,2.0362954) (1.584000,2.0428876) (1.598400,2.0491560) (1.612800,2.0551034) (1.627200,2.0607328) (1.641600,2.0660469) (1.656000,2.0710489) (1.670400,2.0757416) (1.684800,2.0801282) (1.699200,2.0842116) (1.713600,2.0879951) (1.728000,2.0914818) (1.742400,2.0946748) (1.756800,2.0975774) (1.771200,2.1001928) (1.785600,2.1025243) (1.800000,2.1045750) (1.814400,2.1063484) (1.828800,2.1078478) (1.843200,2.1090764) (1.857600,2.1100377) (1.872000,2.1107350) (1.886400,2.1111716) (1.900800,2.1113510) (1.915200,2.1112766) (1.929600,2.1109517) (1.944000,2.1103799) (1.958400,2.1095644) (1.972800,2.1085087) (1.987200,2.1072163) (2.001600,2.1056906) (2.016000,2.1039350) (2.030400,2.1019529) (2.044800,2.0997477) (2.059200,2.0973230) (2.073600,2.0946820) (2.088000,2.0918283) (2.102400,2.0887651) (2.116800,2.0854961) (2.131200,2.0820244) (2.145600,2.0783536) (2.160000,2.0744870) (2.174400,2.0704280) (2.188800,2.0661799) (2.203200,2.0617461) (2.217600,2.0571299) (2.232000,2.0523347) (2.246400,2.0473638) (2.260800,2.0422205) (2.275200,2.0369080) (2.289600,2.0314297) (2.304000,2.0257888) (2.318400,2.0199885) (2.332800,2.0140321) (2.347200,2.0079227) (2.361600,2.0016636) (2.376000,1.9952580) (2.390400,1.9887089) (2.404800,1.9820195) (2.419200,1.9751930) (2.433600,1.9682324) (2.448000,1.9611408) (2.462400,1.9539212) (2.476800,1.9465767) (2.491200,1.9391103) (2.505600,1.9315250) (2.520000,1.9238237) (2.534400,1.9160093) (2.548800,1.9080849) (2.563200,1.9000533) (2.577600,1.8919172) (2.592000,1.8836797) (2.606400,1.8753436) (2.620800,1.8669115) (2.635200,1.8583863) (2.649600,1.8497707) (2.664000,1.8410675) (2.678400,1.8322793) (2.692800,1.8234088) (2.707200,1.8144586) (2.721600,1.8054314) (2.736000,1.7963297) (2.750400,1.7871562) (2.764800,1.7779132) (2.779200,1.7686034) (2.793600,1.7592292) (2.808000,1.7497931) (2.822400,1.7402974) (2.836800,1.7307446) (2.851200,1.7211371) (2.865600,1.7114771) (2.880000,1.7017670) (2.894400,1.6920091) (2.908800,1.6822056) (2.923200,1.6723587) (2.937600,1.6624707) (2.952000,1.6525437) (2.966400,1.6425798) (2.980800,1.6325812) (2.995200,1.6225499) (3.009600,1.6124880) (3.024000,1.6023975) (3.038400,1.5922803) (3.052800,1.5821386) (3.067200,1.5719742) (3.081600,1.5617889) (3.096000,1.5515848) (3.110400,1.5413636) (3.124800,1.5311272) (3.139200,1.5208773) (3.153600,1.5106158) (3.168000,1.5003444) (3.182400,1.4900647) (3.196800,1.4797785) (3.211200,1.4694874) (3.225600,1.4591931) (3.240000,1.4488971) (3.254400,1.4386010) (3.268800,1.4283063) (3.283200,1.4180147) (3.297600,1.4077275) (3.312000,1.3974462) (3.326400,1.3871723) (3.340800,1.3769072) (3.355200,1.3666522) (3.369600,1.3564088) (3.384000,1.3461781) (3.398400,1.3359617) (3.412800,1.3257606) (3.427200,1.3155762) (3.441600,1.3054097) (3.456000,1.2952623) (3.470400,1.2851352) (3.484800,1.2750295) (3.499200,1.2649463) (3.513600,1.2548868) (3.528000,1.2448520) (3.542400,1.2348430) (3.556800,1.2248608) (3.571200,1.2149064) (3.585600,1.2049808) (3.600000,1.1950850) (3.614400,1.1852198) (3.628800,1.1753862) (3.643200,1.1655851) (3.657600,1.1558174) (3.672000,1.1460838) (3.686400,1.1363853) (3.700800,1.1267226) (3.715200,1.1170965) (3.729600,1.1075077) (3.744000,1.0979571) (3.758400,1.0884452) (3.772800,1.0789729) (3.787200,1.0695407) (3.801600,1.0601494) (3.816000,1.0507995) (3.830400,1.0414917) (3.844800,1.0322266) (3.859200,1.0230048) (3.873600,1.0138267) (3.888000,1.0046931) (3.902400,0.9956043) (3.916800,0.9865609) (3.931200,0.9775634) (3.945600,0.9686122) (3.960000,0.9597078) (3.974400,0.9508506) (3.988800,0.9420410) (4.003200,0.9332795) (4.017600,0.9245664) (4.032000,0.9159021) (4.046400,0.9072870) (4.060800,0.8987213) (4.075200,0.8902054) (4.089600,0.8817396) (4.104000,0.8733242) (4.118400,0.8649595) (4.132800,0.8566457) (4.147200,0.8483831) (4.161600,0.8401719) (4.176000,0.8320124) (4.190400,0.8239046) (4.204800,0.8158489) (4.219200,0.8078454) (4.233600,0.7998943) (4.248000,0.7919957) (4.262400,0.7841498) (4.276800,0.7763567) (4.291200,0.7686165) (4.305600,0.7609294) (4.320000,0.7532953)};
\draw[guide]plot coordinates{(0.000000,0.0000000) (0.014400,0.0009450) (0.028800,0.0037358) (0.043200,0.0083068) (0.057600,0.0145934) (0.072000,0.0225320) (0.086400,0.0320602) (0.100800,0.0431165) (0.115200,0.0556406) (0.129600,0.0695730) (0.144000,0.0848556) (0.158400,0.1014309) (0.172800,0.1192429) (0.187200,0.1382362) (0.201600,0.1583567) (0.216000,0.1795513) (0.230400,0.2017678) (0.244800,0.2249551) (0.259200,0.2490631) (0.273600,0.2740428) (0.288000,0.2998459) (0.302400,0.3264256) (0.316800,0.3537356) (0.331200,0.3817309) (0.345600,0.4103673) (0.360000,0.4396016) (0.374400,0.4693918) (0.388800,0.4996966) (0.403200,0.5304757) (0.417600,0.5616899) (0.432000,0.5933006) (0.446400,0.6252706) (0.460800,0.6575634) (0.475200,0.6901432) (0.489600,0.7229755) (0.504000,0.7560264) (0.518400,0.7892631) (0.532800,0.8226536) (0.547200,0.8561668) (0.561600,0.8897724) (0.576000,0.9234410) (0.590400,0.9571441) (0.604800,0.9908540) (0.619200,1.0245439) (0.633600,1.0581875) (0.648000,1.0917598) (0.662400,1.1252363) (0.676800,1.1585933) (0.691200,1.1918080) (0.705600,1.2248582) (0.720000,1.2577227) (0.734400,1.2903808) (0.748800,1.3228128) (0.763200,1.3549994) (0.777600,1.3869223) (0.792000,1.4185638) (0.806400,1.4499069) (0.820800,1.4809353) (0.835200,1.5116333) (0.849600,1.5419860) (0.864000,1.5719790) (0.878400,1.6015987) (0.892800,1.6308321) (0.907200,1.6596666) (0.921600,1.6880904) (0.936000,1.7160924) (0.950400,1.7436618) (0.964800,1.7707887) (0.979200,1.7974635) (0.993600,1.8236772) (1.008000,1.8494215) (1.022400,1.8746884) (1.036800,1.8994705) (1.051200,1.9237611) (1.065600,1.9475538) (1.080000,1.9708426) (1.094400,1.9936222) (1.108800,2.0158876) (1.123200,2.0376344) (1.137600,2.0588585) (1.152000,2.0795562) (1.166400,2.0997244) (1.180800,2.1193603) (1.195200,2.1384615) (1.209600,2.1570259) (1.224000,2.1750521) (1.238400,2.1925387) (1.252800,2.2094848) (1.267200,2.2258900) (1.281600,2.2417538) (1.296000,2.2570766) (1.310400,2.2718587) (1.324800,2.2861009) (1.339200,2.2998042) (1.353600,2.3129699) (1.368000,2.3255998) (1.382400,2.3376955) (1.396800,2.3492594) (1.411200,2.3602937) (1.425600,2.3708012) (1.440000,2.3807847) (1.454400,2.3902473) (1.468800,2.3991923) (1.483200,2.4076233) (1.497600,2.4155440) (1.512000,2.4229584) (1.526400,2.4298704) (1.540800,2.4362845) (1.555200,2.4422051) (1.569600,2.4476368) (1.584000,2.4525843) (1.598400,2.4570527) (1.612800,2.4610469) (1.627200,2.4645721) (1.641600,2.4676338) (1.656000,2.4702372) (1.670400,2.4723880) (1.684800,2.4740919) (1.699200,2.4753545) (1.713600,2.4761817) (1.728000,2.4765795) (1.742400,2.4765539) (1.756800,2.4761110) (1.771200,2.4752570) (1.785600,2.4739980) (1.800000,2.4723404) (1.814400,2.4702905) (1.828800,2.4678547) (1.843200,2.4650395) (1.857600,2.4618512) (1.872000,2.4582965) (1.886400,2.4543818) (1.900800,2.4501138) (1.915200,2.4454989) (1.929600,2.4405438) (1.944000,2.4352551) (1.958400,2.4296395) (1.972800,2.4237034) (1.987200,2.4174536) (2.001600,2.4108967) (2.016000,2.4040393) (2.030400,2.3968879) (2.044800,2.3894491) (2.059200,2.3817296) (2.073600,2.3737358) (2.088000,2.3654744) (2.102400,2.3569517) (2.116800,2.3481742) (2.131200,2.3391484) (2.145600,2.3298807) (2.160000,2.3203775) (2.174400,2.3106450) (2.188800,2.3006896) (2.203200,2.2905175) (2.217600,2.2801349) (2.232000,2.2695479) (2.246400,2.2587626) (2.260800,2.2477851) (2.275200,2.2366214) (2.289600,2.2252774) (2.304000,2.2137590) (2.318400,2.2020720) (2.332800,2.1902221) (2.347200,2.1782152) (2.361600,2.1660568) (2.376000,2.1537524) (2.390400,2.1413077) (2.404800,2.1287281) (2.419200,2.1160189) (2.433600,2.1031855) (2.448000,2.0902331) (2.462400,2.0771668) (2.476800,2.0639919) (2.491200,2.0507134) (2.505600,2.0373361) (2.520000,2.0238651) (2.534400,2.0103051) (2.548800,1.9966609) (2.563200,1.9829372) (2.577600,1.9691386) (2.592000,1.9552697) (2.606400,1.9413348) (2.620800,1.9273385) (2.635200,1.9132849) (2.649600,1.8991785) (2.664000,1.8850232) (2.678400,1.8708233) (2.692800,1.8565828) (2.707200,1.8423055) (2.721600,1.8279955) (2.736000,1.8136565) (2.750400,1.7992922) (2.764800,1.7849063) (2.779200,1.7705024) (2.793600,1.7560840) (2.808000,1.7416546) (2.822400,1.7272175) (2.836800,1.7127761) (2.851200,1.6983335) (2.865600,1.6838929) (2.880000,1.6694575) (2.894400,1.6550302) (2.908800,1.6406141) (2.923200,1.6262119) (2.937600,1.6118266) (2.952000,1.5974609) (2.966400,1.5831174) (2.980800,1.5687989) (2.995200,1.5545078) (3.009600,1.5402466) (3.024000,1.5260178) (3.038400,1.5118238) (3.052800,1.4976668) (3.067200,1.4835491) (3.081600,1.4694729) (3.096000,1.4554402) (3.110400,1.4414532) (3.124800,1.4275139) (3.139200,1.4136242) (3.153600,1.3997860) (3.168000,1.3860011) (3.182400,1.3722713) (3.196800,1.3585982) (3.211200,1.3449837) (3.225600,1.3314293) (3.240000,1.3179365) (3.254400,1.3045068) (3.268800,1.2911418) (3.283200,1.2778427) (3.297600,1.2646110) (3.312000,1.2514480) (3.326400,1.2383548) (3.340800,1.2253328) (3.355200,1.2123831) (3.369600,1.1995068) (3.384000,1.1867050) (3.398400,1.1739786) (3.412800,1.1613287) (3.427200,1.1487563) (3.441600,1.1362621) (3.456000,1.1238471) (3.470400,1.1115121) (3.484800,1.0992578) (3.499200,1.0870851) (3.513600,1.0749945) (3.528000,1.0629867) (3.542400,1.0510624) (3.556800,1.0392222) (3.571200,1.0274665) (3.585600,1.0157960) (3.600000,1.0042110) (3.614400,0.9927121) (3.628800,0.9812997) (3.643200,0.9699741) (3.657600,0.9587356) (3.672000,0.9475847) (3.686400,0.9365216) (3.700800,0.9255466) (3.715200,0.9146598) (3.729600,0.9038615) (3.744000,0.8931520) (3.758400,0.8825312) (3.772800,0.8719994) (3.787200,0.8615567) (3.801600,0.8512031) (3.816000,0.8409386) (3.830400,0.8307634) (3.844800,0.8206774) (3.859200,0.8106806) (3.873600,0.8007729) (3.888000,0.7909543) (3.902400,0.7812247) (3.916800,0.7715841) (3.931200,0.7620322) (3.945600,0.7525690) (3.960000,0.7431942) (3.974400,0.7339078) (3.988800,0.7247095) (4.003200,0.7155991) (4.017600,0.7065763) (4.032000,0.6976410) (4.046400,0.6887929) (4.060800,0.6800317) (4.075200,0.6713571) (4.089600,0.6627688) (4.104000,0.6542665) (4.118400,0.6458498) (4.132800,0.6375185) (4.147200,0.6292721) (4.161600,0.6211104) (4.176000,0.6130329) (4.190400,0.6050392) (4.204800,0.5971289) (4.219200,0.5893017) (4.233600,0.5815571) (4.248000,0.5738947) (4.262400,0.5663140) (4.276800,0.5588146) (4.291200,0.5513961) (4.305600,0.5440579) (4.320000,0.5367997)};
\node[above]at(3.7,2.5){实线：混合谱};\node[above]at(3.7,2.1){虚线：$T_E$单温谱};
}
在低频端，普朗克谱$u_\nu\simeq8\pi\kb T\nu^2/c^3$，所以混合谱对应的低频温度为$T_L=fT_1+(1-f)T_2$。函数$T^4$在正温区严格凸，故对$0<f<1$、$T_1\ne T_2$有
\[
 T_E^4=fT_1^4+(1-f)T_2^4>[fT_1+(1-f)T_2]^4=T_L^4.
\]
因此同一温度不能同时匹配低频系数与总能量。图\ref{fig:m4mixeda}采用$T_1=2T_2$、$f=1/2$绘出两条普朗克组合曲线；$T_L=1.5T_2$、$T_E=8.5^{1/4}T_2\simeq1.7075T_2$。

高频端若$T_1>T_2$，较热分量的指数衰减更慢，也不会变成$T_E$的单一指数。只有两温相同或某份额为零等退化情形才得到全频单温黑体。若腔壁吸收、再发射足够充分，原来分方向保留的谱会被热化，本题不再适用。
\end{solution}

\section{能级间隔、周期与边界相位}
对一维哈密顿量$H=p^2/(2m)+V(x)$，在允许区$\partial p/\partial E=m/p=1/\dot x$。对闭合轨道作用量求能量导数，转点处$p=0$使端点贡献为零，得到
\begin{equation}\label{eq:m4period}
 \frac{\dd J}{\dd E}=\oint\frac{\dd x}{\dot x}=T(E).
\end{equation}
相邻半经典能级$\Delta J=h$，所以大量子数时$\Delta E\simeq h/T(E)$，即跃迁频率趋于经典周期频率。周期变化与能级疏密由同一个相空间面积导数决定。

\begin{example}\label{ex:m4period}
由作用量求重力反弹粒子的经典周期，并用它估计高能级间距，核对直接对能级公式求导的结果。
\Diagram{fig:m4periodq}{粒子由地面反弹到$x_t$后返回}{\draw[axis](0,0)--(0,3.7)node[above]{$x$};\draw[very thick](-1,0)--(1,0);\draw[->](-.4,.3)--(-.4,2.7);\draw[->](.4,2.7)--(.4,.3);\draw[guide](-1,2.7)--(1,2.7);\node[right]at(1,2.7){$x_t=E/(mg)$};}
\end{example}
\begin{solution}
图\ref{fig:m4periodq}的重力反弹轨道作用量为$J(E)=4\sqrt2E^{3/2}/(3g\sqrt m)$，对能量求导得到
\[
 T(E)=\frac{\dd J}{\dd E}
 =\frac{4\sqrt2}{3g\sqrt m}\frac32E^{1/2}
 =\frac{2\sqrt2}{g\sqrt m}\sqrt E
 =\frac{2v_0}g,\qquad v_0=\sqrt{2E/m}.
\]
这与粒子从硬壁起跳、到最高点、再回到硬壁的时间一致。硬壁瞬时反弹不增加有限飞行时间。
\Diagram{fig:m4perioda}{周期随能量增大，半经典相邻能级间距减小}{\draw[axis](0,0)--(4.5,0)node[right]{$E$};\draw[axis](0,0)--(0,3)node[above]{$T$};\draw[curve,domain=0:4,samples=70]plot(\x,{1.2*sqrt(\x)});}
图\ref{fig:m4perioda}中周期随能量增大。半经典条件$J(E_n)=h(n+3/4)$相邻相减，若$n$足够大，可将$J(E_{n+1})-J(E_n)$展开为$T(E_n)\Delta E+O(\Delta E^2)$，所以
\[
 \Delta E\simeq\frac h{T(E)}=\frac{hg\sqrt m}{2\sqrt{2E}}.
\]
直接从$E_n=C(n+3/4)^{2/3}$求导，
\[
 \frac{\dd E}{\dd n}=\frac23C(n+3/4)^{-1/3}
 =\frac{2C^{3/2}}{3\sqrt E}.
\]
再用$C^{3/2}=3hg\sqrt m/(4\sqrt2)$，恰得同一表达式。这里$h/T$估计相邻高能级间距。最低几个能级应保留有限差分，并与完整能级式或精确艾里函数根比较。经典周期越长，对应的相位随能量变化越快，达到下一个量子相位所需的能量增量越小。
\end{solution}

\begin{example}\label{ex:m4correspondence}
非相对论玻尔氢样原子中，比较$n\to n-1$跃迁频率和第$n$条圆轨道的转动频率，求大$n$极限下两者之比。
\Diagram{fig:m4correspondenceq}{相邻高量子数轨道间的跃迁}{\draw[curve](0,0)circle(1.6);\draw[curve](0,0)circle(1.2);\draw[->](1.6,0)--(1.2,0);\node[right]at(1.7,.3){$n$};\node[right]at(1.1,-.6){$n-1$};\fill(0,0)circle(2pt);}
\end{example}
\begin{solution}
令$C=mc^2(Z\alpha)^2/2$，玻尔能级为$E_n=-C/n^2$。图\ref{fig:m4correspondenceq}中$n\to n-1$释放的光子能量为
\[
 h\nu_{n,n-1}=C\left[\frac1{(n-1)^2}-\frac1{n^2}\right]
 =C\frac{2n-1}{n^2(n-1)^2},\quad n\geq2.
\]
第$n$条圆轨道的$v_n=Z\alpha c/n$、$r_n=\hbar n^2/(mZ\alpha c)$，所以转动频率为
\[
 \nu_{\rm orb}=\frac{v_n}{2\pi r_n}=\frac{2C}{hn^3}.
\]
\Diagram{fig:m4correspondencea}{高能级局部斜率给出跃迁频率的领先项}{\draw[axis](0,0)--(4.8,0)node[right]{$n$};\draw[axis](0,-2.7)--(0,.5)node[above]{$E_n$};\draw[curve,domain=1:4.2,samples=70]plot(\x,{-2.3/(\x*\x)});\foreach\x in{1,2,3,4}{\fill(\x,{-2.3/(\x*\x)})circle(1.5pt);}}
两者之比可不经近似直接写为
\[
 \frac{\nu_{n,n-1}}{\nu_{\rm orb}}
 =\frac{n(2n-1)}{2(n-1)^2}
 =1+\frac{3}{2n}+\frac{2}{n^2}+O(n^{-3}).
\]
图\ref{fig:m4correspondencea}中的两条频率在大$n$处相近，极限比值为1。展开时用$(1-1/n)^{-2}=1+2/n+3/n^2+\cdots$，再乘$1-1/(2n)$，才能保留正确的一阶修正。对固定整数$k\ll n$，$n\to n-k$的跃迁频率近于$k\nu_{\rm orb}$，对应经典周期运动的第$k$次谐波；这不等于有限$n$的每条光谱线都能视作电子沿单一圆轨道连续辐射。
\end{solution}

\begin{example}\label{ex:m4finitebarrier}
粒子在$0<x<L$势能为零，两侧势垒均为$V_0>E$。求束缚态的相位量子条件，并说明无限深方阱是怎样的极限。只推导允许能量的方程，不要求逐根求值。
\Diagram{fig:m4finiteq}{有限深方阱在壁外有指数尾部}{\draw[axis](-1.6,0)--(5,0)node[right]{$x$};\draw[curve](-1.3,2.4)--(0,2.4)--(0,0)--(3,0)--(3,2.4)--(4.5,2.4);\draw[guide](-1.3,1)--(4.5,1);\node[right]at(4.5,1){$E$};\node[left]at(-1.3,2.4){$V_0$};}
\end{example}
\begin{solution}
图\ref{fig:m4finiteq}中令$k=\sqrt{2mE}/\hbar$、$\kappa=\sqrt{2m(V_0-E)}/\hbar$。左侧选$A\ee^{\kappa x}$，排除$x\to-\infty$增长的支；阱内取$B\sin(kx+\phi)$。在$x=0$连续匹配波函数及导数：
\[
 A=B\sin\phi,\qquad \kappa A=kB\cos\phi.
\]
两式相除得$k\cot\phi=\kappa$，可选$0<\phi<\pi/2$，故$\phi=\arctan(k/\kappa)$。
\Diagram{fig:m4finitea}{有限势阱边界处波函数非零，壁外有连续衰减尾部}{\draw[axis](-1.3,0)--(4.6,0)node[right]{$x$};\draw[guide](0,-.4)--(0,2.3);\draw[guide](3,-.4)--(3,2.3);\draw[curve,domain=-1:0,samples=35]plot(\x,{1.8*cos(67.5)*exp(1.896119*\x)});\draw[curve,domain=0:3,samples=80]plot(\x,{1.8*cos(45*(\x-1.5))});\draw[curve,domain=3:4,samples=35]plot(\x,{1.8*cos(67.5)*exp(-1.896119*(\x-3))});}
图\ref{fig:m4finitea}的右尾取$D\ee^{-\kappa(x-L)}$。在$x=L$，$D=B\sin(kL+\phi)$、$-\kappa D=kB\cos(kL+\phi)$，所以
\[
 \cot(kL+\phi)=-\cot\phi,
 \quad kL+\phi=n\pi-\phi,
\]
得到允许能量方程
\[
 kL+2\arctan(k/\kappa)=n\pi,
 \qquad 0<E<V_0,\quad n=1,2,\ldots.
\]
求根须同时代入$k(E)$和$\kappa(E)$，不能把它们视为独立变量。左端随$E$从0单调增至$L\sqrt{2mV_0}/\hbar+\pi$的开端点；只有严格低于这个上限的$n\pi$才有束缚态根，因此根的数目有限。$V_0\to\infty$时$\kappa\to\infty$、附加相位趋零，恢复无限深阱。有限势垒允许波函数渗出，等效约束较弱，同一编号的能量低于无限深阱相应能量。
\end{solution}

\begin{example}\label{ex:m4triangular}
一维势$V(x)=F|x|$，$F>0$，无硬壁。求半经典能级；与只保留$x\geq0$并在原点设置硬壁的模型比较相位偏移和作用量。
\Diagram{fig:m4triq}{对称线性势具有两端软转点，原点不是硬壁}{\draw[axis](-2.7,0)--(2.9,0)node[right]{$x$};\draw[axis](0,0)--(0,3.2)node[above]{$V$};\draw[curve](-2.4,2.7)--(0,0)--(2.4,2.7);\draw[guide](-1.8,2.025)--(1.8,2.025);\node[right]at(1.8,2.025){$E$};}
\end{example}
\begin{solution}
图\ref{fig:m4triq}的转点为$\pm E/F$。完整作用量的四倍因子计入左右两半区及往返两支。取$q=E-Fx$，$\dd x=-\dd q/F$，有
\begin{align*}
 J&=4\int_0^{E/F}\sqrt{2m(E-Fx)}\dd x
 =-\frac{4\sqrt{2m}}F\int_E^0q^{1/2}\dd q\\
 &=\frac{4\sqrt{2m}}F\left[\frac23q^{3/2}\right]_0^E
 =\frac{8\sqrt{2m}}{3F}E^{3/2}.
\end{align*}
取两个软转点条件，
\[
 E_n\simeq\left[\frac{3Fh}{8\sqrt{2m}}(n+1/2)\right]^{2/3}.
\]
\Diagram{fig:m4tria}{对称势相轨道跨越原点；硬壁模型只保留右半边}{\draw[axis](-2.6,0)--(2.8,0)node[right]{$x$};\draw[axis](0,-2)--(0,2.2)node[above]{$p$};\draw[curve,domain=-2:2,samples=80]plot(\x,{1.6*sqrt(1-abs(\x)/2)});\draw[curve,domain=-2:2,samples=80]plot(\x,{-1.6*sqrt(1-abs(\x)/2)});}
若原点改硬壁，图\ref{fig:m4tria}只留右半区，相轨道面积减半，且条件变为$(n+3/4)h$，能级为$[3Fh(n+3/4)/(4\sqrt{2m})]^{2/3}$。这对应对称势的奇宇称能级子集的半经典极限，不能只更换偏移而不改积分区域。原点势的尖点不是反射硬壁，波函数可以穿过。
\end{solution}

\begin{example}\label{ex:m4photonrecoil}
自由原子初始静止，激发态与基态静质量分别为$M_i,M_f$。原子辐射一个光子退激，求光子能量和基态原子的反冲能，比较内部能差$\Delta=(M_i-M_f)c^2$。
\Diagram{fig:m4recoilq}{自发辐射必须连同原子反冲一起守恒}{\draw[curve](0,0)circle(.45);\node at(0,0){$M_i$};\draw[->](.7,0)--(2,0);\draw[->](2.6,.4)--(4.7,.4)node[above]{$E_\gamma$};\draw[->](2.4,-.4)--(1.2,-.4)node[below]{$M_f$反冲};}
\end{example}
\begin{solution}
在图\ref{fig:m4recoilq}的原子初始静止系，末光子沿$+x$，其能量为$E_\gamma$。基态原子反冲动量为$-E_\gamma/c$，总能量为$M_ic^2-E_\gamma$。对末原子写质壳条件，
\begin{align*}
 M_f^2c^4&=(M_ic^2-E_\gamma)^2-E_\gamma^2\\
 &=M_i^2c^4-2M_ic^2E_\gamma.
\end{align*}
解得$E_\gamma=(M_i^2-M_f^2)c^2/(2M_i)$。
\Diagram{fig:m4recoila}{内部能差包含光子能量和反冲动能}{\draw[curve](0,0)rectangle(5.4,1);\draw[curve](4.5,0)--(4.5,1);\node at(2.25,.5){$E_\gamma$};\node at(4.95,.5){$K_R$};\node[above]at(2.7,1){$\Delta$};}
令内部能差$\Delta=(M_i-M_f)c^2$，利用$M_f=M_i-\Delta/c^2$展开平方，得到
\[
 E_\gamma=\Delta\left(1-\frac{\Delta}{2M_ic^2}\right),\quad
 K_R=M_ic^2-E_\gamma-M_fc^2=\frac{\Delta^2}{2M_ic^2}.
\]
图\ref{fig:m4recoila}中的能量分配$E_\gamma+K_R=\Delta$是精确的两体运动学结果；$K_R=p_R^2/(2M_f)$仅是非相对论近似，不能先用它求取精确公式。末原子速率为
\[
 |V_R|=\frac{p_Rc^2}{E_R}
 =c\frac{M_i^2-M_f^2}{M_i^2+M_f^2}.
\]
当$\Delta\ll M_ic^2$时$|V_R|\simeq\Delta/(M_ic)$，反冲修正很小。若讨论反向过程，即静止基态原子吸收一个光子而整体升为激发态，所需能量为$(M_i^2-M_f^2)c^2/(2M_f)>\Delta$；发射线与自由吸收线因反冲而不重合。
\end{solution}

\section{有限束缚边界与势阱的时间改变}
\subsection{允许能量、归一化和边界外概率}
有限势阱的束缚态在禁阻区保留指数尾部。边界处波函数及其导数的连续性决定允许能级，完整波函数的归一化再决定阱内、阱外概率。这两步应使用同一套分段解。

取对称阱$|x|<a/2$内$V=0$，外侧$V=V_0$。束缚态满足$0<E<V_0$，定义$k=\sqrt{2mE}/\hbar$、$\kappa=\sqrt{2m(V_0-E)}/\hbar$。偶态写成
\[
 \psi(x)=\begin{cases}
 A\cos kx,&|x|<a/2,\\
 A\cos(ka/2)\ee^{-\kappa(|x|-a/2)},&|x|\geq a/2.
 \end{cases}
\]
在右边界令导数连续，得$k\tan(ka/2)=\kappa$。奇态相应给$-k\cot(ka/2)=\kappa$。令$z=ka/2$、$z_0=a\sqrt{2mV_0}/(2\hbar)$，则$\kappa a/2=\sqrt{z_0^2-z^2}$。应在每个正切或余切分支上寻找合法交点，排除$E=V_0$时不能归一化的阈值态。

偶态的归一化包含阱内余弦段和两个尾部。尾部取$s=|x|-a/2$，各自的上下限为$0,\infty$，于是
\begin{align*}
 A^{-2}&=\int_{-a/2}^{a/2}\cos^2kx\dd x
 +2\cos^2(ka/2)\int_0^\infty\ee^{-2\kappa s}\dd s\\
 &=\left[\frac x2+\frac{\sin2kx}{4k}\right]_{-a/2}^{a/2}
 +2\cos^2(ka/2)\left[-\frac{\ee^{-2\kappa s}}{2\kappa}\right]_0^\infty\\
 &=\frac a2+\frac{\sin ka}{2k}+\frac{\cos^2(ka/2)}\kappa.
\end{align*}
最后一项正是禁阻区概率的未归一化权重。$V_0$很大时$\kappa$很大，尾部减小；靠近上沿的束缚态$\kappa$很小，延伸长度$1/\kappa$可能很长，不能仅用阱内几何宽度估计波函数范围。

\begin{example}一个能精确计算尾部概率的有限阱
上述对称阱的基态满足$k=\kappa$。求所需$V_0$、基态能量、归一化常数及粒子处于阱外的概率。
\ExpansionDiagram{fig:dep-m4-well-q}{有限方势阱的宽度与基态能级，阱内势取零}{
\draw[->](-3,0)--(3.2,0)node[right]{$x$};\draw[thick](-3,2)--(-1,2)--(-1,0)--(1,0)--(1,2)--(3,2);\draw[dashed](-2.8,1)--(2.8,1)node[right]{$E$};\node[above]at(2,2){$V_0$};\node[below]at(-1,0){$-a/2$};\node[below]at(1,0){$a/2$};}
\end{example}
\begin{solution}
基态没有内部节点，$0<ka/2<\pi/2$。由$k\tan(ka/2)=\kappa=k$，得$ka/2=\pi/4$，所以$k=\kappa=\pi/(2a)$。于是
\[
 E=\frac{\pi^2\hbar^2}{8ma^2},\qquad
 V_0-E=E,\qquad V_0=\frac{\pi^2\hbar^2}{4ma^2}.
\]
图\ref{fig:dep-m4-well-q}中的能级恰在阱深一半，但不能由这个条件直接猜测阱外概率也为一半。代入归一化公式，$\sin ka=1$、$\cos^2(ka/2)=1/2$，得
\[
 A^{-2}=\frac a2+\frac a\pi+\frac a\pi
 =\frac{a(\pi+4)}{2\pi},\qquad
 A=\sqrt{\frac{2\pi}{a(\pi+4)}}.
\]
\ExpansionDiagram{fig:dep-m4-well-a}{波函数的余弦段与指数尾部在边界平滑接续}{
\draw[->](-3,0)--(3.2,0)node[right]{$x/(a/2)$};\draw[thick,domain=-1:1,samples=80]plot(\x,{1.8*cos(45*\x)});\draw[thick,domain=1:3,samples=60]plot(\x,{1.8/sqrt(2)*exp(-pi*(\x-1)/4)});\draw[thick,domain=-3:-1,samples=60]plot(\x,{1.8/sqrt(2)*exp(pi*(\x+1)/4)});\draw[dashed](-1,0)--(-1,2);\draw[dashed](1,0)--(1,2);\draw[->](-3,0)--(-3,2.4)node[above]{$\psi/A$};\draw(-3.06,1.8)--(-2.94,1.8);\node[left]at(-3,1.8){$1$};\foreach \x in{-1,0,1}{\draw(\x,-.04)--(\x,.04);\node[below]at(\x,0){$\x$};}}
由图\ref{fig:dep-m4-well-a}两侧尾部积分，
\[
 P_{\rm out}=A^2\frac{\cos^2(ka/2)}\kappa
 =\frac{2}{\pi+4}\simeq0.2800.
\]
所以约28\%概率位于经典禁阻区。这里“禁阻”指经典动能$E-V_0<0$，并不要求量子波函数严格为零。若只在阱内归一化，会把所有局部概率密度都高估。
\end{solution}

\subsection{突然改变与缓慢改变使用不同的接续条件}
势阱宽度缓慢改变时，满足绝热条件的非简并本征态沿同一能级分支演化；突然改变时，波函数来不及改变，应把原波函数投影到新本征态。两者都需要初态和改变方式，不能只因初末几何宽度相同便使用同一能量公式。

\begin{example}无限深阱突然加宽为两倍
粒子最初处于$0<x<L$无限深阱基态。右壁瞬间移到$2L$，左壁不动。求改变后各能级的概率、能量期望，并与缓慢加宽的结果比较。取瞬间改变为理想极限，外部势在新增区间中直接降为零。
\ExpansionDiagram{fig:dep-m4-sudden-q}{右壁移开的一瞬间，原波函数在新增区间仍为零}{
\draw[->](0,0)--(6,0)node[right]{$x$};\draw[very thick](.5,0)--(.5,2);\draw[dashed,very thick](3,0)--(3,2);\draw[very thick](5.5,0)--(5.5,2);\draw[thick,domain=.5:3,samples=70]plot(\x,{1.5*sin(180*(\x-.5)/2.5)});\draw[thick](3,0)--(5.5,0);\node[below]at(3,0){$L$};\node[below]at(5.5,0){$2L$};\draw[->](3,2.3)--(5.5,2.3);}
\end{example}
\begin{solution}
图\ref{fig:dep-m4-sudden-q}给定瞬间波函数：$\psi=\sqrt{2/L}\sin(\pi x/L)$于$0<x<L$，在$L<x<2L$为零。新阱归一化本征函数为$\phi_k=L^{-1/2}\sin(k\pi x/(2L))$，能级为$E_k'=k^2\pi^2\hbar^2/(8mL^2)$。系数
\[
 c_k=\frac{\sqrt2}{L}\int_0^L\sin\frac{\pi x}{L}
 \sin\frac{k\pi x}{2L}\dd x.
\]
对$k\ne2$，先使用$\sin a\sin b=[\cos(a-b)-\cos(a+b)]/2$，再分别积分：
\begin{align*}
 c_k&=\frac{\sqrt2}{2L}\int_0^L
 \left[\cos\frac{(2-k)\pi x}{2L}-\cos\frac{(2+k)\pi x}{2L}\right]\dd x\\
 &=\frac{\sqrt2}{\pi}
 \left[\frac{\sin((2-k)\pi x/(2L))}{2-k}
 -\frac{\sin((2+k)\pi x/(2L))}{2+k}\right]_0^L\\
 &=\frac{\sqrt2}{\pi}\sin\frac{k\pi}{2}
 \left(\frac1{2-k}+\frac1{2+k}\right).
\end{align*}
这里用了$\sin(\pi-k\pi/2)=\sin(k\pi/2)$、$\sin(\pi+k\pi/2)=-\sin(k\pi/2)$，下端点两项均为零。因此
\[
 c_k=\frac{4\sqrt2\sin(k\pi/2)}{\pi(4-k^2)}.
\]
$k=2$时应直接积分两个相同的正弦函数：
\[
 c_2=\frac{\sqrt2}{L}\int_0^L\sin^2\frac{\pi x}{L}\dd x
 =\frac{\sqrt2}{L}\left[\frac x2-\frac{L}{4\pi}\sin\frac{2\pi x}{L}\right]_0^L
 =\frac1{\sqrt2}.
\]
其余偶数$k$的系数为零，奇数$k$则
\[
 P_k=|c_k|^2=\frac{32}{\pi^2(4-k^2)^2},
 \qquad P_2=\frac12.
\]
\ExpansionDiagram{fig:dep-m4-sudden-a}{突然加宽后能量测量分布在多个新能级上}{
\draw[->](0,0)--(7,0)node[right]{$k$};\draw[->](0,0)--(0,3.5)node[above]{$P_k$};\draw[very thick](1,0)--(1,2.16);\draw[very thick](2,0)--(2,3);\draw[very thick](3,0)--(3,.78);\draw[very thick](5,0)--(5,.044);\draw[dashed](0,3)--(2,3);\node[left]at(0,3){$1/2$};\node[below]at(1,0){1};\node[below]at(2,0){2};\node[below]at(3,0){3};\node[below]at(5,0){5};}
图\ref{fig:dep-m4-sudden-a}以同一纵向比例画前几个概率，其中$P_1=32/(9\pi^2)\simeq0.3603$、$P_3=32/(25\pi^2)\simeq0.1297$。其总和为1，由新本征函数的完备性和初态归一化保证。

能量期望可直接用动能二次型，无需逐项硬求级数：
\begin{align*}
 \langle E'\rangle&=\frac{\hbar^2}{2m}\int_0^{2L}|\psi'(x)|^2\dd x
 =\frac{\pi^2\hbar^2}{mL^3}\int_0^L\cos^2\frac{\pi x}{L}\dd x\\
 &=\frac{\pi^2\hbar^2}{mL^3}
 \left[\frac x2+\frac{L}{4\pi}\sin\frac{2\pi x}{L}\right]_0^L\\
 &=\frac{\pi^2\hbar^2}{2mL^2}=E_{1,\rm old}.
\end{align*}
瞬间移去远侧势垒没有改变原波函数的动能期望；之后波包由多个新定态叠加演化。若缓慢加宽且保持基态分支，末态只有$k=1$，能量为原来的$1/4$，能量差由移动壁面的做功带走。突然加宽后的测量结果虽然包含低能级，也包含更高能级，不能只取最可能能级代替平均能量。
\end{solution}

\section{自由原子的吸收共振与反冲}
\subsection{内部能差不是所有过程中的同一个光子能量}
设自由原子基态、激发态静质量为$M,M_*$，内部能差$\Delta=(M_*-M)c^2>0$。激发态原子在静止时发射一个光子的能量已由两体守恒给出$E_{\rm em}=(M_*^2-M^2)c^2/(2M_*)$。反过来，基态静止原子吸收一个光子而成为同一激发态时，末态原子也要携带光子动量，吸收能量一般不同。

吸收过程初态总四动量平方为$M^2c^4+2Mc^2E_{\rm abs}$，末态为$M_*^2c^4$，所以
\[
 E_{\rm abs}=\frac{(M_*^2-M^2)c^2}{2M}
 =\Delta+\frac{\Delta^2}{2Mc^2}.
\]
吸收光子须同时供给内部激发能和反冲动能；发射时，内部能差则由光子与反冲原子分享。两题都把各自的初态原子设为静止，因此所选初末四动量没有互换，时间反演对称性并不要求这两个光子同频。

\begin{example}静止自由原子能否吸收同种静止原子发出的光
两种完全相同的自由原子最初分别处于激发态和基态，且都在实验室静止。第一原子辐射的单光子沿连线入射第二原子。只考虑同一窄能级跃迁，求发射能量与静止吸收共振能量的差，并估计为补偿该差所需的吸收原子速度。取$\Delta\ll Mc^2$，忽略自然线宽以外的其他展宽机制。
\ExpansionDiagram{fig:dep-m4-recoil-q}{发射原子与吸收原子各自需要分配反冲能}{
\fill(0,0)circle(3pt)node[below]{$M_*$};\draw[->,thick](.5,0)--(4.5,0)node[midway,above]{$E_{\rm em}$};\fill(5,0)circle(3pt)node[below]{$M$};\draw[->](0,-.6)--(-1,-.6)node[left]{发射反冲};\draw[->](5,-.6)--(6,-.6)node[right]{吸收反冲};}
\end{example}
\begin{solution}
图\ref{fig:dep-m4-recoil-q}中两原子的初态都静止，但两个末态都要反冲。能量差精确为
\begin{align*}
 E_{\rm abs}-E_{\rm em}
 &=\frac{(M_*^2-M^2)c^2}{2}
 \left(\frac1M-\frac1{M_*}\right)\\
 &=\frac{\Delta^2(M_*+M)}{2MM_*c^2}
 \simeq\frac{\Delta^2}{Mc^2}.
\end{align*}
这一正差约为两次单原子反冲能$E_R=\Delta^2/(2Mc^2)$之和。若线宽远小于$2E_R$，静止发射源的中心光子不足以达到静止吸收原子的共振能量，尽管两原子内部能级完全相同。
\ExpansionDiagram{fig:dep-m4-recoil-a}{发射与吸收的中心能量分处内部能差两侧}{
\draw[->](0,0)--(7,0)node[right]{光子能量};\draw[thick](2,-.4)--(2,1)node[above]{$E_{\rm em}$};\draw[dashed](3.5,-.4)--(3.5,1)node[above]{$\Delta$};\draw[thick](5,-.4)--(5,1)node[above]{$E_{\rm abs}$};\draw[<->](2,-.8)--(5,-.8)node[midway,below]{$\simeq2E_R$};}
若吸收原子以小速度$u$迎着入射光运动，在原子系测得光子能量约为$E_{\rm em}(1+u/c)$。共振要求$u/c\simeq(E_{\rm abs}-E_{\rm em})/\Delta\simeq\Delta/(Mc^2)$，即$u\simeq\Delta/(Mc)$。这次一阶多普勒蓝移恰好补偿图\ref{fig:dep-m4-recoil-a}中的能量缺口。

若气体处于温度$T$，沿光束方向的速度标准差为$\sqrt{\kb T/M}$，多普勒能量标准差约为$\sigma_E=\Delta\sqrt{\kb T/(Mc^2)}$。它与反冲缺口同量级时，$\kb T\sim\Delta^2/(Mc^2)$。计算完整吸收率还需要线形、速度分布和跃迁强度。原子若受外部结构约束，动量可由更大系统承担，能量和动量守恒也须对这个整体列式。
\end{solution}

\section{习题与解答}
\begin{exercise}\label{ex:m4tunnel}
质量$m$的粒子以能量$0<E<V_0$从左射向$0<x<a$的矩形势垒，势垒外$V=0$。求透射率，并给出厚势垒极限，解释它与束缚态边界的联系。
\Diagram{fig:m4tunnelq}{有限禁区两侧都是传播区}{\draw[axis](-1.5,0)--(5,0)node[right]{$x$};\draw[curve](-1.2,0)--(0,0)--(0,2.4)--(3,2.4)--(3,0)--(4.7,0);\draw[guide](-1.2,1)--(4.7,1);\node[right]at(4.7,1){$E$};\node[above]at(1.5,2.4){$V_0$};\node[below]at(3,0){$a$};}
\end{exercise}
\begin{solution}
图\ref{fig:m4tunnelq}取$k=\sqrt{2mE}/\hbar$、$\kappa=\sqrt{2m(V_0-E)}/\hbar$。三段波函数分别为$\ee^{ikx}+r\ee^{-ikx}$、$A\ee^{\kappa x}+B\ee^{-\kappa x}$、$t\ee^{ikx}$。右界面两连续条件给
\[
 A\ee^{\kappa a}=\frac{t\ee^{ika}}2(1+ik/\kappa),\quad
 B\ee^{-\kappa a}=\frac{t\ee^{ika}}2(1-ik/\kappa).
\]
代回左界面，先算出
\begin{align*}
 A+B&=t\ee^{ika}[\cosh\kappa a-i(k/\kappa)\sinh\kappa a],\\
 \kappa(A-B)&=t\ee^{ika}[-\kappa\sinh\kappa a+ik\cosh\kappa a].
\end{align*}
左界面条件为$1+r=A+B$、$ik(1-r)=\kappa(A-B)$。第二式除以$ik$，再与第一式相加，消去$r$，得
\[
 t\ee^{ika}=\left[\cosh\kappa a+
 i\frac{\kappa^2-k^2}{2k\kappa}\sinh\kappa a\right]^{-1}.
\]
\Diagram{fig:m4tunnela}{有限禁区需要两指数支，透射侧只保留向右传播波}{\draw[axis](-1.5,0)--(5,0)node[right]{$x$};\draw[guide](0,-.6)--(0,2.3);\draw[guide](3,-.6)--(3,2.3);\draw[curve,domain=0:3,samples=60]plot(\x,{1.8*exp(-.65*\x)});\draw[->](3.4,.6)--(4.7,.6)node[above]{透射};\node[above]at(1.5,2.4){振幅包络示意};}
两侧传播介质相同，$T=|t|^2$。分母的模平方为
\begin{align*}
 T^{-1}&=\cosh^2\kappa a+
 \frac{(\kappa^2-k^2)^2}{4k^2\kappa^2}\sinh^2\kappa a\\
 &=1+\frac{(\kappa^2+k^2)^2}{4k^2\kappa^2}\sinh^2\kappa a.
\end{align*}
使用$\kappa^2+k^2=2mV_0/\hbar^2$和$k^2\kappa^2=4m^2E(V_0-E)/\hbar^4$，得到
\[
 T=\left[1+\frac{V_0^2\sinh^2\kappa a}{4E(V_0-E)}\right]^{-1}.
\]
图\ref{fig:m4tunnela}的厚垒极限$\kappa a\gg1$给$T\simeq16E(V_0-E)\ee^{-2\kappa a}/V_0^2$。有限禁区内不能先删增长支；它与衰减支一起满足两端匹配，正如有限光隙中的受抑全反射。
\end{solution}

\begin{exercise}\label{ex:m4fieldemission}
金属表面外的电子势垒近似为$V(x)=W-eFx$，$0<x<W/(eF)$，其中$F>0$是外加电场大小。以表面内电子能量为零，求半经典穿透概率的主要指数，解释增大电场的作用。
\Diagram{fig:m4emissionq}{外场使表面势垒成为有限宽三角形}{\draw[axis](-.6,0)--(4.9,0)node[right]{$x$};\draw[axis](0,0)--(0,3.2)node[above]{$V$};\draw[curve](0,2.6)--(3.8,0);\node[left]at(0,2.6){$W$};\node[below]at(3.8,0){$W/(eF)$};\draw[->](-.5,.3)--(-.05,.3);}
\end{exercise}
\begin{solution}
图\ref{fig:m4emissionq}禁区局部虚动量大小$|p|=\sqrt{2m(W-eFx)}$。先计算禁区积分，取$q=W-eFx$，$\dd x=-\dd q/(eF)$，空间上下限对应$q=W,0$：
\begin{align*}
 \int_0^{W/(eF)}|p|\dd x
 &=-\frac{\sqrt{2m}}{eF}\int_W^0q^{1/2}\dd q\\
 &=\frac{\sqrt{2m}}{eF}\left[\frac23q^{3/2}\right]_0^W
 =\frac{2\sqrt{2m}W^{3/2}}{3eF}.
\end{align*}
WKB振幅的衰减因子为$\exp[-\int|p|\dd x/\hbar]$；概率需平方，故
\[
 \mathcal T\propto\exp\left[-\frac2\hbar\int_0^{W/(eF)}|p|\dd x\right]
 =\exp\left[-\frac{4\sqrt{2m}W^{3/2}}{3e\hbar F}\right].
\]
\Diagram{fig:m4emissiona}{电场增大使势垒变窄，指数抑制减弱}{\draw[axis](0,0)--(4.8,0)node[right]{$x$};\draw[axis](0,0)--(0,3.2)node[above]{$V$};\draw[curve](0,2.6)--(4,0);\draw[guide](0,2.6)--(2,0);\node[right]at(3,1.2){较弱场};\node[left]at(-.4,1.4){较强场};\draw[guide](-.3,1.4)--(.8,1.56);}
图\ref{fig:m4emissiona}中势垒高度仍为$W$，宽度却随$1/F$减小，所以穿透概率随电场迅速上升。若画$\ln\mathcal T$对$1/F$，领先指数给负斜率直线；实际电流还含电子供给与前因子，不能仅凭此式求出绝对电流。这个过程不要求入射光子能量超过功函数，和经典过垒光电逃逸属于不同模型。
\end{solution}

\begin{exercise}\label{ex:m4angularvariant}
光电角度模型改为只在向外半球产生电子，概率密度按$\cos\theta$分布，即$\dd P=C\cos\theta\dd\Omega$。法向势垒仍为$W$。求归一化常数和逃逸比例，与全空间各向同性模型比较。
\Diagram{fig:m4angularq}{方向权重偏向表面法线}{\draw[curve](-2,0)--(2,0);\draw[axis](0,0)--(0,2.6)node[above]{法线};\draw[->](0,0)--(.7,2.1);\draw[->](0,0)--(1.7,.8);\node[right]at(1.7,.8){$\cos\theta$权重};}
\end{exercise}
\begin{solution}
图\ref{fig:m4angularq}在向外半球归一化：
\[
 1=2\pi C\int_0^{\pi/2}\cos\theta\sin\theta\dd\theta
 =2\pi C\left[\frac{\sin^2\theta}{2}\right]_0^{\pi/2}
 =\pi C,
\]
故$C=1/\pi$。当$h\nu\geq W$时，可逸出角域为$0\leq\theta\leq\theta_c$，其中$\cos\theta_c=\sqrt{W/(h\nu)}$。在同一原函数中代入新的上限，
\[
 P_{\rm esc}=2\int_0^{\theta_c}\cos\theta\sin\theta\dd\theta
 =[\sin^2\theta]_0^{\theta_c}=1-\frac W{h\nu}.
\]
低于阈值时比例为零。
\Diagram{fig:m4angulara}{同一势垒，不同初始角分布产生不同逃逸曲线}{\draw[axis](0,0)--(4.6,0)node[right]{$h\nu/W$};\draw[axis](0,0)--(0,2.7)node[above]{$P_{\rm esc}$};\draw[curve,domain=1:4,samples=60]plot(\x,{2.3*(1-1/\x)});\draw[guide,domain=1:4,samples=60]plot(\x,{1.15*(1-sqrt(1/\x))});}
图\ref{fig:m4angulara}中此比例在高能趋1，而全空间各向同性模型趋$1/2$。势垒决定允许角域，产生机制决定各方向的权重。因此，相同能量阈值可以对应不同光电产额。
\end{solution}

\begin{exercise}\label{ex:m4radiativecool}
热容量$C$近似恒定的黑球表面积$A$，置于温度$T_0$的大黑腔。球温$T=T_0+\Delta T$，$|\Delta T|\ll T_0$。推导小温差弛豫时间，并说明面积和热容量怎样进入。
\Diagram{fig:m4coolq}{黑球与大腔壁双向辐射交换}{\draw[curve](0,0)circle(.8);\draw[curve](-2,-1.7)rectangle(2,1.7);\node at(0,0){$T$};\node[above]at(0,1.7){$T_0$};\draw[->](.9,.3)--(1.7,.3);\draw[->](-1.7,-.3)--(-.9,-.3);}
\end{exercise}
\begin{solution}
图\ref{fig:m4coolq}中，腔内平衡辐射各向同性，球吸收功率为$A\sigma T_0^4$，自身发射功率为$A\sigma T^4$。以整球热容量$C$列能量方程，得
\[
 C\frac{\dd T}{\dd t}=-A\sigma(T^4-T_0^4).
\]
写$T=T_0+\Delta T$，展开得$T^4-T_0^4=4T_0^3\Delta T+6T_0^2\Delta T^2+\cdots$。在$|\Delta T|/T_0\ll1$时保留线性项，
\[
 \frac{\dd\Delta T}{\dd t}=-\frac{\Delta T}{\tau},\qquad
 \tau=\frac{C}{4A\sigma T_0^3}.
\]
\Diagram{fig:m4coola}{小温差以指数形式衰减}{\draw[axis](0,0)--(4.5,0)node[right]{$t/\tau$};\draw[axis](0,0)--(0,2.6)node[above]{$\Delta T/\Delta T(0)$};\draw[curve,domain=0:4,samples=60]plot(\x,{2.2*exp(-\x)});}
当$\Delta T(0)\ne0$时分离变量，
\[
 \int_{\Delta T(0)}^{\Delta T(t)}\frac{\dd q}{q}
 =-\frac1\tau\int_0^t\dd s,
 \quad \ln\frac{|\Delta T(t)|}{|\Delta T(0)|}=-\frac t\tau.
\]
温差的符号不变，故$\Delta T(t)=\Delta T(0)\ee^{-t/\tau}$，得到图\ref{fig:m4coola}的曲线；初温恰等于腔温时，另有恒解$\Delta T=0$。初温低于腔温时，负温差也按同一式趋于零。量纲上$A\sigma T_0^3$为功率每温度，与$C$相除给时间。

对同材料均匀实心球，$C=\rho c_p(4\pi a^3/3)$、$A=4\pi a^2$，所以$\tau=\rho c_pa/(12\sigma T_0^3)$，半径越大弛豫越慢。该结论要求内部导热足够快，使整球可用单一温度描述。若环境温度可忽略，则原方程改为$\dot T=-(A\sigma/C)T^4$。分离变量，
\[
 \int_{T(0)}^{T(t)}T^{-4}\dd T=-\frac{A\sigma}{C}\int_0^t\dd s,
 \quad \left[-\frac1{3T^3}\right]_{T(0)}^{T(t)}=-\frac{A\sigma t}{C},
\]
得到$T^{-3}(t)=T^{-3}(0)+3A\sigma t/C$，为代数冷却。它采用不同的环境条件，不能沿用小温差指数律。
\end{solution}

\begin{exercise}\label{ex:m4spectralpeak}
求黑体“每对数频率间隔”的能量谱峰值方程，并说明它为什么不同于$u_\nu$的峰。
\Diagram{fig:m4logq}{相同对数宽度在高频处包含更大的线性频率区间}{\draw[axis](0,0)--(5.4,0)node[right]{$\nu$};\draw[very thick](.5,0)--(1,0);\draw[very thick](2,0)--(4,0);\node[above]at(.75,.5){$\nu\to2\nu$};\node[above]at(3,.5){$4\nu\to8\nu$};}
\end{exercise}
\begin{solution}
图\ref{fig:m4logq}比较等比例频率区间。因$\dd\nu=\nu\dd\ln\nu$，对数谱为$\nu u_\nu$，以$x=h\nu/(\kb T)$表示正比$f_4(x)=x^4/(\ee^x-1)$。在例\ref{ex:m4wien}已逐步求出的$f_j'$中取$j=4$，
\[
 f_4'(x)=\frac{x^3\ee^x}{(\ee^x-1)^2}
 [4(1-\ee^{-x})-x].
\]
除去边界根$0$，峰值方程为$4(1-\ee^{-x})=x$，非零正根$x\simeq3.92069$。
\Diagram{fig:m4loga}{对数谱多一个频率权重，峰向高频移动}{\draw[axis](0,0)--(5.3,0)node[right]{$h\nu/(\kb T)$};\draw[axis](0,0)--(0,2.8)node[above]{对数谱};\draw[curve,domain=.05:8,samples=80]plot({.6*\x},{.4*\x^4/(exp(\x)-1)});\draw[guide](2.3524,0)--(2.3524,2);}
图\ref{fig:m4loga}中的峰对应每相同倍频带的能量贡献，而$u_\nu$比较每相同赫兹的能量。它也等于每对数波长间隔能量谱的峰所对应的同一物理频率，因为$|\dd\ln\lambda|=|\dd\ln\nu|$；线性波长谱则还有不同雅可比。
\end{solution}

\begin{exercise}\label{ex:m4welltransition}
无限深方阱宽度缓慢从$L$变为$2L$，粒子始终留在同一量子数$n$的定态分支。求能量与典型动量的变化，并用作用量解释“缓慢”的意义。
\Diagram{fig:m4expandq}{缓慢移动硬壁使阱宽加倍}{\draw[curve](0,2)--(0,0)--(2,0)--(2,2);\draw[guide](4,0)--(4,2);\draw[->](2.2,1)--(3.8,1);\node[below]at(2,-.4){$L\to2L$};}
\end{exercise}
\begin{solution}
图\ref{fig:m4expandq}的定态能量$E_n=n^2h^2/(8mL^2)$，所以阱宽加倍，能量减为$1/4$，动量大小$p_n=nh/(2L)$减为$1/2$。定态内平均动量为0，但$\langle p^2\rangle^{1/2}=p_n$，不要把典型动量大小等同于平均动量。
\Diagram{fig:m4expanda}{相轨道长宽反比变化，面积$2pL=nh$保持}{\draw[axis](0,0)--(5,0)node[right]{$x$};\draw[axis](0,-1.8)--(0,2)node[above]{$p$};\draw[curve](0,-1.3)rectangle(2,1.3);\draw[guide](0,-.65)rectangle(4,.65);}
图\ref{fig:m4expanda}说明缓变过程中作用量$J=2pL$近似保持，量子上要求壁运动足够慢，避免显著跃迁到别的$n$。若壁突然移动，原波函数不再是新阱单一定态，需要展开为新本征态之和，本题能量比例就不能直接使用。
\end{solution}

\begin{exercise}\label{ex:m4oscforce}
谐振子势中加入恒定力$F$，势为$m\omega^2x^2/2-Fx$。求能级、平衡位置和能级间隔，说明零点项是否改变。
\Diagram{fig:m4forceq}{恒力使抛物势横移并下移}{\draw[axis](-2.2,0)--(3.4,0)node[right]{$x$};\draw[axis](0,-1)--(0,3)node[above]{$V$};\draw[guide,domain=-1.8:1.8,samples=50]plot(\x,{.7*\x*\x});\draw[curve,domain=-1:2.8,samples=50]plot(\x,{.7*(\x-.9)^2-.567});}
\end{exercise}
\begin{solution}
逐项求导得$V'(x)=m\omega^2x-F$，令其为零，平衡位置为$x_0=F/(m\omega^2)$；再求导得$V''=m\omega^2>0$，故这是稳定平衡。配方得$V=\tfrac12m\omega^2(x-x_0)^2-F^2/(2m\omega^2)$。图\ref{fig:m4forceq}中的势曲率不变，运动频率仍为$\omega$，能级为$E_n=\hbar\omega(n+1/2)-F^2/(2m\omega^2)$。
\Diagram{fig:m4forcea}{所有能级整体平移，间隔$\hbar\omega$保持}{\foreach\y in{0,.7,1.4,2.1}{\draw[curve](-2,\y)--(-.3,\y);\draw[curve](1,{\y-.5})--(2.7,{\y-.5});}\draw[->](0,1.4)--(.7,.9);\node[above]at(1.8,2.3){有恒力};}
图\ref{fig:m4forcea}中零点能相对势阱最低点仍为$\hbar\omega/2$，但相对原来的能量零点发生共同平移。边界仍是两个平滑转点，故作用量偏移$1/2$不变。只看总能量可能出现负值，并不表示振动激发能或频率为负。
\end{solution}

\begin{exercise}\label{ex:m4hydrogenscaling}
氢样离子核电荷由$Z$变为$2Z$，量子数保持$n$。分别讨论非相对论玻尔半径、速度、束缚能和相邻能级频率的缩放，指出何时这种缩放失效。
\Diagram{fig:m4Zq}{同一$n$下增大核电荷，轨道缩小}{\draw[curve](-1.5,0)circle(1.1);\fill(-1.5,0)circle(2pt);\draw[curve](2,0)circle(.55);\fill(2,0)circle(2pt);\node[below]at(-1.5,-1.3){$Z$};\node[below]at(2,-1.3){$2Z$};}
\end{exercise}
\begin{solution}
图\ref{fig:m4Zq}依$r_n\propto1/Z$、$v_n\propto Z$、$E_n\propto-Z^2$，半径减半、速率加倍、束缚能绝对值乘4。固定两量子数的谱线差同乘4，频率也乘4。
\Diagram{fig:m4Za}{速度尺度$Z\alpha/n$决定非相对论近似的适用性}{\draw[axis](0,0)--(4.5,0)node[right]{$Z\alpha/n$};\draw[axis](0,0)--(0,2.7)node[above]{$v/c$};\draw[curve](0,0)--(3.7,2.3);\node[right]at(3.7,2.3){近1时须改模型};}
图\ref{fig:m4Za}中$Z\alpha/n\ll1$是非相对论条件。增大$Z$会增强相对论修正；当参数逼近1，圆轨道半经典模型本身也失去常规适用性，有限核尺寸、自旋和完整量子动力学不可忽略。缩放律是在固定模型条件下的结果，不是任意核电荷都可无限外推。
\end{solution}

\begin{exercise}\label{ex:m4yield}
采用例\ref{ex:m4escape}的角度筛选模型，在固定吸收功率$P$下改变频率。求零电压光电流最大的光子能量，以及最大电流。
\Diagram{fig:m4yieldq}{光子能量增加时每秒光子数减少，逃逸锥却扩大}{\draw[axis](0,0)--(0,2.7)node[above]{法线};\draw[curve](0,0)--(.8,2);\draw[curve](0,0)--(1.8,2);\node[right]at(1.8,2){较高$h\nu$};\node[left]at(-.1,1){$P$固定};}
\end{exercise}
\begin{solution}
令$x=h\nu/W\geq1$，图\ref{fig:m4yieldq}的逃逸概率为$(1-x^{-1/2})/2$，所以$I=eP( x^{-1}-x^{-3/2})/(2W)$。求导得$\dd I/\dd x=(eP/2W)(-x^{-2}+\tfrac32x^{-5/2})$，零点为$\sqrt x=3/2$。
\Diagram{fig:m4yielda}{电流在$x=9/4$达到极大}{\draw[axis](0,0)--(5,0)node[right]{$x=h\nu/W$};\draw[axis](0,0)--(0,2.6)node[above]{相对电流};\draw[curve,domain=1:7,samples=90]plot({.6*\x},{14*(1/\x-1/(\x^1.5))});\draw[guide](1.35,0)--(1.35,2.075);\node[below]at(1.35,0){$9/4$};}
在$x=1$电流为0，$x\to\infty$也趋于0，故图\ref{fig:m4yielda}的驻点是全局极大。代入得$I_{\max}=2eP/(27W)$、$h\nu=9W/4$。若只看逃逸概率而忽略光子数$P/(h\nu)$，就会错误地认为频率越高电流必越大。
\end{solution}

\begin{exercise}\label{ex:m4cavitypressure}
各向同性平衡辐射能量密度为$u$，理想反射壁法线为$z$。由光子动量通量求辐射压强；再求一束垂直单向光被完全吸收、完全反射的压强。
\Diagram{fig:m4pressureq}{斜向光子反射时法向动量改变$2p\cos\theta$}{\draw[curve](-2,0)--(2,0);\draw[->](-1.6,1.8)--(0,0);\draw[->](0,0)--(1.6,1.8);\draw[guide](0,0)--(0,2.4);\node at(.45,1.3){$\theta$};}
\end{exercise}
\begin{solution}
图\ref{fig:m4pressureq}中一个光子的法向动量改变$2E\cos\theta/c$。单位时间单位面积入射光子的能量还带通量投影$c\cos\theta$，因此
\[
 p=\frac{2u}{4\pi}\int_0^{2\pi}\dd\varphi\int_0^{\pi/2}\cos^2\theta\sin\theta\dd\theta
 =u\left[-\frac{\cos^3\theta}{3}\right]_0^{\pi/2}=\frac u3.
\]
\Diagram{fig:m4pressurea}{单向垂直光束的吸收和反射动量通量}{\draw[curve](0,-.5)--(0,1.2);\draw[->](-2,.8)--(-.1,.8)node[midway,above]{$I/c$};\draw[->](-.1,.2)--(-2,.2)node[midway,below]{反射多一份};\node[right]at(.2,.6){壁};}
图\ref{fig:m4pressurea}的单向光束不需要半球积分，完全吸收时$p=I/c$，完全反射时$p=2I/c$。这里$I$是垂直于光束的功率面密度。腔辐射的$u/3$与孔口通量$cu/4$来自不同角度矩：前者有$\cos^2\theta$，后者有$\cos\theta$。
\end{solution}

\begin{exercise}\label{ex:m4equilibrium}
半径$R$的恒星近似黑体温度$T_s$，距离$r\gg R$处有半径$a$、表面吸收和发射均为黑体的小球。忽略其他热源，求球均匀温度。若只在受光半球均匀发热、背面不辐射，结果如何？
\Diagram{fig:m4equilibriumq}{恒星辐射照到小球的截面积$\pi a^2$}{\draw[curve](0,0)circle(.8);\node at(0,0){$T_s$};\draw[curve](5,0)circle(.35);\foreach\y in{-.2,.2}{\draw[->](1,\y)--(4.6,\y);}\node[below]at(2.8,-.5){$r$};\node[right]at(5.4,0){$a$};}
\end{exercise}
\begin{solution}
恒星总功率$4\pi R^2\sigma T_s^4$，在半径$r$球面上的通量$F=\sigma T_s^4R^2/r^2$。图\ref{fig:m4equilibriumq}的小球吸收面积是迎光截面$\pi a^2$，吸收功率$\pi a^2F$。
\Diagram{fig:m4equilibriuma}{吸收用投影面积，发射用实际表面积}{\draw[curve](0,0)circle(1);\draw[guide](0,-1)--(0,1);\draw[->](1,0)--(1.7,0);\draw[->](0,1)--(0,1.7);\draw[->](-1,0)--(-1.7,0);\draw[->](0,-1)--(0,-1.7);\node[right]at(2,0){全表面$4\pi a^2$};}
均匀温度时，图\ref{fig:m4equilibriuma}全表面发射$4\pi a^2\sigma T^4$。平衡给$T=T_s\sqrt{R/(2r)}$。若只有半球辐射，发射面积$2\pi a^2$，则$T=2^{1/4}T_s\sqrt{R/(2r)}$。该变式是规定的温度分布模型，真实表面是否均温取决于导热、自转等过程。
\end{solution}

\begin{exercise}\label{ex:m4reducedmass}
把玻尔模型的固定核改为质量$M$可动的原子核。求相对坐标的半径、能级及相对固定核结果的修正。
\Diagram{fig:m4reducedq}{电子和核绕共同质心转动}{\draw[curve](0,0)circle(1.7);\draw[curve](0,0)circle(.4);\fill(1.7,0)circle(2pt)node[right]{$m$};\fill(-.4,0)circle(3pt)node[above left]{$M$};\fill(0,0)circle(1pt);\draw[guide](0,0)--(.6,-.65)node[right]{质心};}
\end{exercise}
\begin{solution}
图\ref{fig:m4reducedq}中记电子、原子核位置为$\vct r_e,\vct r_N$，定义质心$\vct R=(m\vct r_e+M\vct r_N)/(m+M)$和相对坐标$\vct r=\vct r_e-\vct r_N$。反解为
\[
 \vct r_e=\vct R+\frac M{M+m}\vct r,\quad
 \vct r_N=\vct R-\frac m{M+m}\vct r.
\]
代入动能，交叉项相消，得到
\[
 K=\tfrac12(M+m)\dot{\vct R}^{\,2}
 +\tfrac12\mu\dot{\vct r}^{\,2},\qquad \mu=\frac{mM}{m+M}.
\]
库仑势仅依赖$r$，因此相对运动等价于质量$\mu$在固定库仑中心中运动。
\Diagram{fig:m4reduceda}{核、电子到质心距离按反质量比分配}{\draw[axis](-1.7,0)--(3.9,0);\fill(-1,0)circle(3pt)node[above]{$M$};\fill(0,0)circle(1.5pt)node[above]{$C$};\fill(3,0)circle(2pt)node[above]{$m$};\draw[<->](-1,-.5)--(0,-.5)node[midway,below]{$r_N$};\draw[<->](0,-.5)--(3,-.5)node[midway,below]{$r_e$};}
在质心系对相对角动量取$\mu vr=n\hbar$，向心方程为$\mu v^2/r=Ze^2/(4\pi\varepsilon_0r^2)$。两式相除先得$v=Z\alpha c/n$，再得
\[
 r_n=\frac{4\pi\varepsilon_0\hbar^2n^2}{\mu Ze^2},\qquad
 E_n=-\frac{\mu c^2(Z\alpha)^2}{2n^2}.
\]
图\ref{fig:m4reduceda}中的电子、核到质心距离分别为$Mr_n/(M+m)$、$mr_n/(M+m)$。相对固定核结果，能级绝对值乘$\mu/m=1/(1+m/M)$，核间距乘其倒数。若$M\gg m$，能级相对修正为$-m/M+O(m^2/M^2)$，相应跃迁频率也降低同一比例。核可动导致的约化质量修正改变内部能级；整个原子发光时的反冲则还需要在这些能级之上作总四动量守恒。
\end{solution}

\begin{exercise}\label{ex:m4powerwell}
粒子在一维势$V(x)=A|x|^s$中运动，$A>0,s>0$。用两平滑转点条件给出能级对量子数的幂次关系，并比较$s=2$与$s\to\infty$的极限。
\Diagram{fig:m4powerq}{幂次束缚势的对称转点}{\draw[axis](-2,0)--(2.3,0)node[right]{$x$};\draw[axis](0,0)--(0,3)node[above]{$V$};\draw[curve,domain=-1.7:1.7,samples=70]plot(\x,{.5*abs(\x)^3});\draw[guide](-1.5,1.6875)--(1.5,1.6875);\node[right]at(1.5,1.6875){$E$};}
\end{exercise}
\begin{solution}
图\ref{fig:m4powerq}中转点为$\pm x_t$，$x_t=(E/A)^{1/s}$。利用势的偶对称，作用量写为
\[
 J=4\int_0^{x_t}\sqrt{2m(E-Ax^s)}\dd x.
\]
令$y=x/x_t$，有$\dd x=x_t\dd y$、$Ax^s=Ey^s$，所以
\[
 J=4\sqrt{2m}A^{-1/s}E^{(s+2)/(2s)}C_s,\quad
 C_s=\int_0^1\sqrt{1-y^s}\dd y.
\]
\Diagram{fig:m4powera}{作用量积分换成固定区间的无量纲面积}{\begin{scope}[yscale=.45]\draw[axis](0,0)--(3.9,0)node[right]{$y$};\draw[axis](0,0)--(0,2.5)node[above]{$\sqrt{1-y^s}$};\draw[curve,domain=0:1,samples=70]plot({3*\x},{2*sqrt(1-\x^3)});\node[below]at(3,0){$1$};\end{scope}}
图\ref{fig:m4powera}的无量纲面积与$E$无关。两软转点条件给出
\[
 E_n\simeq\left[\frac{hA^{1/s}}{4\sqrt{2m}C_s}
 (n+1/2)\right]^{2s/(s+2)}.
\]
因此$s<2$时高能级间距减小，$s=2$时等间隔，$s>2$时增大。取$s=2$，$C_2=\pi/4$；再令$A=m\omega^2/2$，系数化为$\hbar\omega$，恢复谐振子结果。

讨论$s\to\infty$时应先固定长度尺度，写$V=V_*|x/L_*|^s$。对有限$E$，转点趋$L_*$、$C_s\to1$，作用量趋$4L_*\sqrt{2mE}$，与宽$2L_*$的方阱相同，能量的高量子数标度为$n^2$。但端点在极限中变为硬壁，量子数偏移须重新匹配，不能把软转点的$1/2$保留为硬壁最低能级的准确偏移。原点在$s\leq1$时不光滑，转点仍平滑；相关半经典式主要用于高能级，低能级须另解边界问题。
\end{solution}

\begin{exercise}\label{ex:m4gravityscale}
定义重力反弹模型的特征长度$\ell=(\hbar^2/(2m^2g))^{1/3}$。把其薛定谔方程化为无量纲形式，说明精确能量由何种边界条件确定，并比较半经典结果。
\Diagram{fig:m4gravityq}{硬壁处波函数为零，高处波函数衰减}{\draw[axis](-.3,0)--(4.7,0)node[right]{$x$};\node[above]at(0,1.6){$V$};\draw[very thick](0,0)--(0,1.4);\draw[curve](0,0)--(4,1.3);\draw[guide](0,.9)--(4.3,.9);\node[right]at(4.3,.9){$E$};}
\end{exercise}
\begin{solution}
图\ref{fig:m4gravityq}对应$-\hbar^2\psi''/(2m)+mgx\psi=E\psi$。令$\xi=x/\ell$、$\epsilon=E/(mg\ell)$，则链式法则给$\dd\psi/\dd x=\ell^{-1}\dd\psi/\dd\xi$、$\dd^2\psi/\dd x^2=\ell^{-2}\dd^2\psi/\dd\xi^2$。代入并除以$mg\ell$，用$\hbar^2/(2m\ell^2)=mg\ell$，化为
\[
 -\frac{\dd^2\psi}{\dd\xi^2}+\xi\psi=\epsilon\psi.
\]
\Diagram{fig:m4gravitya}{无量纲重力反弹基态在硬壁为零，在高处衰减}{\begin{scope}[yscale=.65]\draw[axis](0,0)--(7.5,0)node[right]{$\xi$};\draw[axis](0,-.4)--(0,2.1)node[above]{$\psi$};\draw[curve]plot coordinates {(0.00000,-0.0000000) (0.06422,0.1348816) (0.12844,0.2685052) (0.19266,0.3996896) (0.25688,0.5273458) (0.32110,0.6504842) (0.38532,0.7682206) (0.44954,0.8797792) (0.51376,0.9844945) (0.57798,1.0818115) (0.64220,1.1712844) (0.70642,1.2525734) (0.77064,1.3254415) (0.83486,1.3897488) (0.89908,1.4454471) (0.96330,1.4925734) (1.02752,1.5312426) (1.09174,1.5616405) (1.15596,1.5840157) (1.22018,1.5986721) (1.28440,1.6059610) (1.34862,1.6062736) (1.41284,1.6000332) (1.47706,1.5876887) (1.54128,1.5697070) (1.60550,1.5465674) (1.66972,1.5187550) (1.73394,1.4867556) (1.79817,1.4510510) (1.86239,1.4121141) (1.92661,1.3704052) (1.99083,1.3263691) (2.05505,1.2804315) (2.11927,1.2329971) (2.18349,1.1844479) (2.24771,1.1351413) (2.31193,1.0854092) (2.37615,1.0355578) (2.44037,0.9858666) (2.50459,0.9365891) (2.56881,0.8879525) (2.63303,0.8401586) (2.69725,0.7933847) (2.76147,0.7477840) (2.82569,0.7034872) (2.88991,0.6606034) (2.95413,0.6192213) (3.01835,0.5794110) (3.08257,0.5412249) (3.14679,0.5046992) (3.21101,0.4698558) (3.27523,0.4367028) (3.33945,0.4052368) (3.40367,0.3754436) (3.46789,0.3472996) (3.53211,0.3207732) (3.59633,0.2958260) (3.66055,0.2724135) (3.72477,0.2504864) (3.78899,0.2299915) (3.85321,0.2108727) (3.91743,0.1930715) (3.98165,0.1765277) (4.04587,0.1611805) (4.11009,0.1469685) (4.17431,0.1338306) (4.23853,0.1217063) (4.30275,0.1105360) (4.36697,0.1002615) (4.43119,0.0908260) (4.49541,0.0821748) (4.55963,0.0742549) (4.62385,0.0670155) (4.68807,0.0604080) (4.75229,0.0543862) (4.81651,0.0489061) (4.88073,0.0439261) (4.94495,0.0394069) (5.00917,0.0353115) (5.07339,0.0316054) (5.13761,0.0282559) (5.20183,0.0252329) (5.26606,0.0225080) (5.33028,0.0200552) (5.39450,0.0178499) (5.45872,0.0158698) (5.52294,0.0140941) (5.58716,0.0125036) (5.65138,0.0110808) (5.71560,0.0098096) (5.77982,0.0086750) (5.84404,0.0076638) (5.90826,0.0067634) (5.97248,0.0059627) (6.03670,0.0052514) (6.10092,0.0046203) (6.16514,0.0040610) (6.22936,0.0035658) (6.29358,0.0031280) (6.35780,0.0027412) (6.42202,0.0023999) (6.48624,0.0020990) (6.55046,0.0018341) (6.61468,0.0016011) (6.67890,0.0013964) (6.74312,0.0012167) (6.80734,0.0010592) (6.87156,0.0009211) (6.93578,0.0008004) (7.00000,0.0006948)};\node[below left]at(0,0){$\psi(0)=0$};\node[above]at(4.6,.6){$\operatorname{Ai}(\xi-\epsilon_0)$};\end{scope}}
令$z=\xi-\epsilon$，方程为$\psi''=z\psi$，通解为$A\operatorname{Ai}(z)+B\operatorname{Bi}(z)$。正无穷处$\operatorname{Bi}$增长，归一化要求$B=0$；$A\ne0$再由归一化确定。图\ref{fig:m4gravitya}的硬壁$x=0$对应$z=-\epsilon$，要求$\operatorname{Ai}(-\epsilon)=0$，所以精确能量由其负零点确定。半经典近似给$\epsilon_n\simeq[\tfrac{3\pi}2(n+\tfrac34)]^{2/3}$。无量纲化把质量和重力全部收进尺度$mg\ell$，可用于比较不同粒子而不重复解微分方程。
\end{solution}


\end{document}
